\chapter{Full Result Tables}
\label{app:fullresults}


% Figure and table captions in this appendix are set at footnotesize, single
% spaced, as they were when these figures lived in their own appendix; moving
% them must not change how they print.
\captionsetup{font={footnotesize,singlespacing}}
% ============================================================================
% Ground-up rebuild (W3, 2026-08). The June appendix (single-seed, observation-
% level HAC, 57-cell census, "D2 parity" narrative) is discarded in full: every
% table below is rebuilt from the frozen evidence files on the declared primary
% basis, and each carries a "% src:" comment naming the file it is read from.
% ============================================================================

% ----------------------------------------------------------------------------
% Float barriers.  From Section C.14 on, every section is short prose in front
% of one to five near-full-page tables, and the stock float queue could not
% clear them: with \floatpagefraction at 0.95 (figuresetup.tex:38) a float page
% is built only when floats nearly fill one, so tables were held back and pooled
% into a long block at the end of the appendix, the first of them falling ten
% pages after the section that introduces it.  Sixteen \clearpage commands
% survive from that first fix; six sit at a section opening and the other ten
% mid-section, where a run of tables would otherwise straddle a page.
% They are now backed by a \FloatBarrier before every \section in this file,
% added when the dissolved visualisation appendix put 43 figures in here and the
% same queueing began carrying FIGURES across section boundaries -- one landed
% eight pages late, inside a section whose evidence it does not draw.  placeins
% is loaded in config.tex WITHOUT the [section] option: the barriers are written
% out here and in the other appendices, so the page-capped chapters keep the
% stock queue and cannot be pushed past 60 pages by a barrier.
% ----------------------------------------------------------------------------

\FloatBarrier
\section{Scope, Basis and How to Read These Tables}
\label{app:res-basis}

This appendix carries, at full width, the result tables that
Chapter~\ref{ch:results} compresses to headline counts and
Chapter~\ref{ch:validation} draws on selectively. Nothing here is new
evidence: every cell is read mechanically from a committed evidence file,
almost always a table under \texttt{results/tables/} --- the two exceptions are
the HAR-RV reference rows of Tables~\ref{tab:app-standalone-ed}
and~\ref{tab:app-standalone-comb}, read from that run's own
\texttt{metrics.json}, as their comments say --- and each table names its own
evidence file in a provenance comment in the report source directly
beneath its caption. Most numbers in the main text can be traced to a row
here and from there to the artefact that produced it; where a chapter
quantity has no row here, its own provenance comment names the artefact
directly, and a minority of those comments name an internal working
document rather than a released file (Section~\ref{app:external:verify}).

All tables sit on the \emph{primary basis} declared at the opening of
Chapter~\ref{ch:results}, and repeated here because an appendix is read
out of order. The forecast object is the three-seed ensemble (seeds
2026--2028, per-observation mean). The two prompted large-language-model
arms, C6 and D4, are single-pass by construction, so their rows are
identical on either seed basis. Inference is the Diebold--Mariano test
\cite{diebold1995comparing} computed on \emph{day-clustered} loss
differentials: averaged within each trading day, with a
Newey--West correction at lag $h-1$ days and the small-sample adjustment
of \cite{harvey1997testing}. Holm's step-down correction
\cite{holm1979simple} is applied inside each pre-declared family. The
primary loss is QLIKE \cite{patton2011proxies}, evaluated in
\emph{volatility} units on the 69-cell combination grid and in
\emph{variance} units on the standalone leaderboard. Levels are never
compared across those two conventions, and no percentage in one
convention may be set beside a percentage in the other.

Two sign conventions recur and are opposite to one another, which is the
single most common way of misreading these tables. In the
\emph{standalone} tables a \emph{positive} DM statistic means the
challenger named in the row is \emph{worse} than the HAR-RV baseline
\cite{corsi2009har}. In the \emph{combination} tables a \emph{negative}
DM statistic means the text-augmented forecast $f_U$ is \emph{better}
than the reference $f_R$ named in the column group; the accompanying
relative improvement is positive in the same case. Each caption restates
the convention that applies to it.

Two denominators recur. The standalone census has 180 comparisons: 20
challengers carrying committed forecasts on all nine disclosure-by-horizon
panels (long-form, event-driven and combined at $h = 5$, 10 and 20
trading days). Three are price arms and 17 are text or fusion
arms; the latter form the 153-comparison sub-family that the
variance-unit statements are scoped to. The combination grid has 69
cells: 15 long-form and 8 event-driven challengers at three horizons.

Which 20 those are is worth naming, because the model matrix of
Chapter~\ref{ch:data} carries 22 identifiers and a reader cannot
otherwise reconstruct the denominator from it. Nine matrix rows sit
outside the standalone set (\textbf{A1, A6, A7, B1, B3, B4, C5x, C6
and D4}). A2 is the reference every comparison is taken against
rather than a challenger within it. % src: results/tables/dm_pairwise_clustered.csv (20 distinct challenger values x 9 panels = 180 rows; A1, B1, B3, B4 appear zero times)
The nine exclusions are of four kinds, and none of them is a suppressed
result. A1, the naive historical-volatility arm, is carried on the
leaderboard of Table~\ref{tab:app-leaderboard} but was never entered in
the pairwise census. A6 and A7 are scored on the return-matched
subsample and on the point-in-time VIX panel respectively, so their
losses are not computed on the rows the other arms share. Block~B enters
the census through one representative: the classical
representative is fixed as B2 TF--IDF ridge by the generating script's
declared best-of-block rule. That rule is the script's own, not a ranking:
on the best-QLIKE-over-three-horizons ordering this appendix prints, B3
L--M linear sits above B2. % src: scripts/analysis/dm_pairwise.py BASELINE\_MODELS ("Classical / baseline model set (best B = B2\_tfidf\_ridge)"); results/tables/cost_accuracy.csv (qlike_long_form_h20: B3_lm_linear 0.6918 against B2_tfidf_ridge 0.7083)
B1, B3 and B4 therefore carry no census row even
though all three appear in the 69-cell combination grid and, apart from
their absence here, are reported throughout.
And C5x, C6 and D4 are single-seed combination-track arms, excluded on
the same ground that keeps them off the leaderboard.
Expanding the twelve identifiers that remain by their strategy and
embedding variants (C1 at two truncation strategies, C2 at four, C5
and D3 at three embedders each) returns exactly the 20 challengers,
and $20 \times 9 = 180$ comparisons. Three of the 20 are price arms
(A3, A4, A5), leaving the 17 text and fusion arms whose $17 \times 9 =
153$ comparisons carry the variance-unit verdict.
Where a table departs from the primary basis (single seed, observation
level, variance unit, two-way clustered, or a different row support),
the departure is named in the caption and never left to the reader.

A note on row support, because three different supports appear below and
they do not track the reference. The full test panels carry
7{,}951/7{,}933/7{,}902 long-form and 25{,}109/25{,}001/24{,}732
event-driven observations at $h = 5/10/20$. % src: results/tables/dm_pairwise_clustered.md (n_obs)
The merged-grid rows on which the firm-identity and maximal-pool rungs
are evaluated are 7{,}550/7{,}167/7{,}097 and
23{,}855/22{,}785/22{,}318. % src: results/tables/maximal_pool_robustness.csv (n_test); TECHNICAL_SUPPLEMENT.md Table S12
The day counts behind the clustered tests are 809/803/794 long-form
and 996/991/981 event-driven on the full panels, 792/766/765 and
996/991/981 on the merged grid. % src: results/tables/dm_pairwise_clustered.csv (n_days); results/tables/maximal_pool_robustness.csv (n_days)
Percentages computed on one support are not restatements of percentages
computed on another.

Table~\ref{tab:app-denominators} reconciles those row sets in one place, so
that a reader need not hold them in memory. It runs in five groups, set apart by a space:
the standalone census, the combination grid and its ladder, the narrow families
around the 8-K residual, the economic and external panels, and last the two row
supports everything above sits on. It carries no column for \emph{why}
a row set is the size it is, and the paragraphs above therefore remain the only
account of the nine model-matrix rows that sit outside the standalone census.
Three of its entries print as 18 and are not restatements of one another --- the
loss panels of the superior-predictive-ability test, the portfolio cells and the
cells the SHAR substitution touches --- which is the concordance's point rather
than an exception to it. Two units are absent by design: the twenty-comparison
panel of the first two Holm families and the eight-cell arm family of the
fourteenth are correction units rather than denominators, and nothing is stated
on them.

\begin{table}[p]
\centering
\caption[The denominator concordance]{The denominator concordance: the row
sets out of which this report states a count, with what one of its rows is, the
observation support it sits on, and one claim stated on it.}
\label{tab:app-denominators}
\singlespacing\footnotesize
\setlength{\tabcolsep}{3pt}
% src: counts verified against results/tables/ (dm_pairwise_clustered.csv 180 rows;
% control_intersection_ensemble.csv 69 rows and five ladder verdict columns;
% deployable_combiner.csv 75; m1_multiseed.csv 144; m1_stratified.csv 675 rows with
% axis != 'all'; row13_spa_mcs.csv 18 distinct (disclosure, horizon, loss);
% vix_control.csv 40; portfolio_econ.csv 18; utility_value.csv 27 distinct cells;
% var_backtest.csv 36 distinct (disclosure, model, horizon) x 2 nominal levels;
% maec_audit.csv 24 rows at stage row5, alignment primary, over the three text arms).
% Supports read from the n_test column of each of those files.
\begin{tabular}{@{}>{\raggedright\arraybackslash}p{0.15\textwidth}>{\raggedright\arraybackslash}p{0.20\textwidth}>{\raggedright\arraybackslash}p{0.155\textwidth}>{\raggedright\arraybackslash}p{0.335\textwidth}r@{}}
\toprule
Count & One row is & Row support & A claim stated on it & p. \\
\midrule
\addlinespace
180 & one challenger against A2 on one panel & full panels & ``none favours the challenger'' & 28 \\ % src: results/tables/dm_pairwise_clustered.csv (180 rows)
153 & the same, over the 17 text and fusion arms & full panels & ``the 17 text and fusion arms remain 0 of 153 better even there'' & 29 \\ % src: results/tables/variance_unit_standalone180.csv
18 & one (channel, horizon, loss) panel & full panels & ``a consistent $p$ of 1.000 in all 18 panels'' & 50 \\ % src: results/tables/row13_spa_mcs.csv (18 distinct disclosure x horizon x loss)
\addlinespace
69 & one combination cell; 45 long-form, 24 event-driven & full panels; merged grid at the last two rungs & ``the primary basis leaves 38 of 69 cells genuine'' & 31 \\ % src: results/tables/m1_ensemble_primary.csv (n_test 7,951 / 25,109 ...) against firm_identity_ensemble.csv and maximal_pool_robustness.csv (n_test 7,550 / 23,855 ...)
345 & one cell's verdict at one of five rungs ($5 \times 69$) & mixed across the five rungs & ``all 345 cell-level decisions behind the ladder's five 69-cell columns'' & 33 \\ % src: results/tables/control_intersection_ensemble.csv (69 rows; primary_holm, maximal_holm, firm_holm, AND_maximal_firm_holm, AND_full_holm)
38, Holm in 40 & a cell an earlier single-seed screen passed & full panels, seed 2026 & ``leaves 19 of 38 standing'' & 32 \\ % src: results/tables/vix_control.csv (40 rows)
75 & the 69 plus a six-cell C5 extension & full panels & ``the same 75-cell family under expanding weights returns 7 genuine cells'' & 148 \\ % src: results/tables/deployable_combiner.csv (75 rows)
144 & one trained arm at one channel and horizon & full panels, single seed & ``37 of 144 model-by-panel cells flip the sign of their test statistic across seeds'' & 24 \\ % src: results/tables/m1_multiseed.csv (144 rows = 16 models x 3 disclosures x 3 horizons, the combined channel included)
675 & one stratum cell of the increment map & strata of the full panels & ``The placebo is larger in only 38 of the 675 cells'' & 141 \\ % src: results/tables/m1_stratified.csv (675 rows with axis != 'all'; panel (a) draws 567, being 513 of these plus 54 pooled cells)
\addlinespace
12 & C6, two channels, three horizons, two references & full panels (HAR), merged grid (identity) & ``a 12-cell family symmetric around C6 \dots\ it survives all three horizons in both'' & 36 \\ % src: results/tables/residual_symmetric_holm.csv (both n_test supports present)
3 & C6 event-driven, Item 2.02 reference & full event-driven panel & ``The residual retains $+0.21/+0.18/+0.20$ per cent'' & 35 \\ % src: results/tables/itemcode_control.csv
12, Holm in 24 & C6 at one channel, horizon and reference & full panels, split at 2024-07-01 & ``0 of 12 stratified cells'' & 36 \\ % src: results/tables/llm_contamination.md section 3 (24 rows = 12 cells x 2 strata)
6 & the 70B replication, one horizon and reference & full event-driven panel & ``After Holm in its own six-cell family \dots\ only the five-day cell survives'' & 36 \\ % src: results/tables/crossfamily_llama70_ens.csv
6, of 9 readouts & one executed anonymisation arm at one horizon & masked 8-K panel, seed 2026 & ``the pooled median over both executed arms' six vs-HAR cells, $0.51$'' & 38 \\ % src: results/tables/anon_arm.csv (7 rows, 6 of them scored); anon_share_ci.csv (6 per-cell shares, 2 arm medians, 1 pooled median)
\addlinespace
18 & one channel--model--horizon portfolio cell & 159/77/39 and 200/100/50 holding periods & ``a median of $+0.003$ over the 18 cells tested'' & 36 \\ % src: results/tables/portfolio_econ.csv (18 rows; n_periods, the file carrying no observation count)
72 & one (channel, model, horizon) row at one level & merged grid & ``text improves on $f_R$ in 55 of 72 cells'' & 43 \\ % src: results/tables/var_backtest.csv (36 distinct rows x alpha in .01, .05)
27 & one of nine channel--model pairs, one horizon & merged grid & ``the fee is positive in 14 of its own 27 cells'' & 43 \\ % src: results/tables/utility_value.csv (27 distinct disclosure x model x horizon cells)
24 & one MAEC arm, one horizon and reference & 672 rows, 143 call dates & ``0 of the 24 primary-alignment cells'' & 56 \\ % src: results/tables/maec_audit.csv (stage row5, alignment primary, three text arms)
18 & one text arm, SHAR substituted for HAR & return-matched sample & ``changes no verdict in the 18 incremental cells'' & 31 \\ % src: results/tables/stronger_baselines.csv (section m1_incremental)
\addlinespace
7{,}951/\allowbreak 7{,}933/\allowbreak 7{,}902 and 25{,}109/\allowbreak 25{,}001/\allowbreak 24{,}732 & one test observation; combined 33{,}060/\allowbreak 32{,}934/\allowbreak 32{,}634 & census, primary rung, deployable family & ``7{,}951/7{,}933/7{,}902 long-form, 25{,}109/25{,}001/24{,}732 event-driven'' & 28 \\ % src: results/tables/dm_pairwise_clustered.csv (n_obs); m1_ensemble_primary.csv and deployable_combiner.csv (n_test)
7{,}550/\allowbreak 7{,}167/\allowbreak 7{,}097 and 23{,}855/\allowbreak 22{,}785/\allowbreak 22{,}318 & the same, on the pooled references' join & identity and pool rungs, VaR and fee panels & ``on 23{,}855/22{,}785/22{,}318 merged-grid test rows'' & 34 \\ % src: results/tables/maximal_pool_robustness.csv and firm_identity_ensemble.csv (n_test); var_backtest.csv and utility_value.csv carry the same n
\bottomrule
\end{tabular}
\end{table}



The figures interleaved with these tables were drawn for the same evidence and
are placed beside it rather than gathered into a gallery of their own; four of
them, the model spectrum and the efficiency frontier, sit beside the
configuration fields of Appendix~\ref{app:hyperparams} instead. One rule holds
across all of them: each generator carries an evidence gate that aborts the
drawing if the file it reads has drifted from the values the figure was verified
against, so no figure in this report is transcribed by hand. The worked-filing
exhibit of Section~\ref{app:vis-worked} draws nothing and so has no generator to
gate; it names the committed file for every number on the line that states it.

\FloatBarrier
\section{The Fifteen Pre-Declared Holm Families}
\label{app:res-holmfams}

Section~\ref{sec:design-inference} declares that Holm's correction is
applied inside each of 15 pre-declared families and that families are
never pooled, and every survivor count in Chapters~\ref{ch:results}
and~\ref{ch:validation} rests on that declaration. A few corrections in those
chapters are run in blocks that are deliberately \emph{not} among the fifteen,
and each says so where it stands rather than here: the VIX control is corrected
inside its own 40-cell re-test, and the cross-family replication inside the
replication script's own six-cell block. Neither revises a survivor count.
% src: results/tables/holm_families.md lists F1-F15 and neither block is among them; 04_results.tex:181 and :209 carry the two statements
Because it is the
strongest claim this report makes about its own discipline, the
enumeration is carried here in full rather than left to the release
documentation. Table~\ref{tab:app-holm-families} gives, for each family,
what it contains, how large it is, the record that fixed its membership
before the statistics inside it were inspected, and the committed
evidence file that carries its corrected $p$-values.

Four kinds of record appear in the ``declared in'' column, and naming
them separately is more honest than describing all fifteen as
pre-registered. Six families were fixed in a registered analysis record
(the seven \texttt{configs/prereg\_*.md} files inventoried in
Appendix~\ref{app:external}). Five are cited by the version-control tag under
which that record was committed; the sixth, F7, is cited by the record's own
section number instead, its evidence row naming no tag and no bare
\texttt{prereg} tag existing.
% src: results/tables/holm_families.md (F7 declared-in reads "prereg (FACTS 13b)"); git tag -l lists prereg-ea/-rfa/-maec/-ef/-cd/-h/-kc and hpo-prereg, none of them bare
Five were fixed by the protocol itself,
stated in the numbered protocol sections that govern the whole model matrix. Two were fixed by a block
headed \emph{pre-declared Holm families} inside the committed combiner
table, placed above the results it governs and stating in its own heading
that it was written before any result below it was inspected. The
remaining two rest on weaker instruments and are marked as such: one on
the column header of the committed standalone table, and one on the
second-domain plan and its five sanity gates. % src: results/tables/holm_families.md (declared-in column: 5 tag-cited rows F8/F11/F12/F13/F14, plus F7 whose entry reads "prereg (FACTS 13b)" -- a section number, no tag, 5 protocol rows F1/F3-F6, 2 deployable_combiner.md rows F9/F10, 1 table-header row F2, 1 second-domain-plan row F15); results/tables/deployable_combiner.md ("PRE-DECLARED Holm families (declared before any result below was inspected)"); configs/prereg_*.md (7 files)
Every mechanism fixed its family in advance; they differ in how much
external evidence of that fact a reader can inspect, not in whether the
family was pre-declared.

Three conventions travel with the enumeration and are set out here rather
than in the caption. Placebo gates (label shuffle,
$|\mathrm{DM}| < 2$) apply \emph{after} Holm inside the same family, so a
genuine cell requires a negative DM, a Holm-adjusted $p < .05$ and a
clean placebo together. Arms that were never executed, and arms retired
at a reproduction gate, carry no tests and enter no family. And the
not-applicable rule of the anonymisation record governs share estimands
only, never family membership. % src: results/tables/holm_families.md notes (i)-(iii)

Every fixed size in the table is a denominator this report already
prints, and the cross-check is worth stating outright, because a
sixteenth denominator would be a defect. The 180 and 153 of F1 and F2 are
the standalone census and its text-and-fusion subset
(Section~\ref{app:res-basis}); the 69 of F3 to F6, and of each rung of
F12, is the combination grid; the 12 of F7 and the 3 of F8 are the
symmetric residual family and the earnings-window control of
Section~\ref{sec:res-residual}; the 75 of F9 and F10 is that same grid
widened by the six-cell C5 extension (Section~\ref{app:res-deploy}); and
the 24 of F14 is the earnings-call alignment of
Section~\ref{sec:val-external}. % src: results/tables/holm_families.md (size column)
Two rows fix no size of their own, and neither introduces one: F13
inherits the F1 and F3 structures on the public-variant panels, and F15
sets its family per gate.
The correction units named in the Members column of F1, F2 and F14 --- the
20-comparison panel and the eight cells of one MAEC arm --- are not
sixteenth and seventeenth denominators. No count is reported out of
either; they say only how wide the step-down runs, and the counts
themselves are still reported out of the 180, the 153 and the 24. The
distinction matters in one direction only, and it is the unfavourable one:
correcting inside a narrower unit is weaker than correcting across the
whole scope, so these three families are less conservative than their
sizes alone would suggest. The 20-comparison unit is the one every
standalone caption in this appendix already states.

One size does need a word of care, and it is F11. Its 6 is a family
\emph{per executed arm}, three horizons against two references, whereas
the six per-cell identity shares drawn in
Figure~\ref{fig:res-anon-price} panel~(a) are two arms at three horizons.
The two counts sit on the same event-driven rows and count different
objects, so neither is a restatement of the other. % src: results/tables/holm_families.md (F11 members: "3 horizons x 2 references, per executed arm (c6, c2)"); results/tables/anon_arm.csv (6 scored rows: c6 and c2 at h = 5/10/20; b2 carries no cells)

No family is declared and then quietly abandoned. F1 and F2 are reported
in Section~\ref{app:res-standalone}, F3 to F6 in
Sections~\ref{app:res-grid} and~\ref{app:res-ladder}, F7 and F8 in
Section~\ref{sec:res-residual}, F9 and F10 in
Section~\ref{app:res-deploy}, F11 in Section~\ref{sec:res-identity}, F12
in Section~\ref{app:res-rangebased}, F13 in
Section~\ref{app:res-public}, F14 in Section~\ref{app:res-maec} and F15
in Section~\ref{app:res-yelp}.

\begin{table}[htbp]
\centering
\caption[The 15 pre-declared Holm families]{The 15 pre-declared Holm
families, transcribed in full from the release enumeration. Holm's
step-down correction \cite{holm1979simple} is applied \emph{within} each
family and families are never pooled, so no count in this report is
corrected across two of these rows. ``Size'' is the reporting scope; the
``Members'' column names the narrower unit the step-down runs in where
there is one. ``Declared in'' names the record that
fixed the membership before the statistics inside it were inspected;
``carried by'' names the committed evidence file under
\texttt{results/tables/}.}
\label{tab:app-holm-families}
\singlespacing\footnotesize
\setlength{\tabcolsep}{3pt}
% src: results/tables/holm_families.md (all 15 rows verbatim: family, size, members, declared in, carried by; notes (i)-(iii))
\begin{tabular}{@{}l>{\raggedright\arraybackslash}p{0.14\textwidth}l>{\raggedright\arraybackslash}p{0.24\textwidth}>{\raggedright\arraybackslash}p{0.16\textwidth}>{\raggedright\arraybackslash}p{0.20\textwidth}@{}}
\toprule
& Family & Size & Members & Declared in & Carried by \\
\midrule
F1 & Standalone leaderboard & 180 & Challenger against A2 on squared-error DM: models $\times$ disclosures $\times$ horizons; Holm inside each 20-comparison panel & protocol & \texttt{dm\_\allowbreak pairwise\_\allowbreak clustered.csv} \\
F2 & Standalone variance-unit & 153 & The text and fusion subset of F1 under variance-unit QLIKE; a scope, not a unit --- Holm as in F1 & \texttt{variance\_\allowbreak unit\_\allowbreak standalone180.\allowbreak csv} header & \texttt{variance\_\allowbreak unit\_\allowbreak standalone180.\allowbreak csv} \\ % src: that file's holm_qlike_var_clu column carries F2's corrected p-values; F1's carrier, dm_pairwise_clustered.csv, has no variance-unit column at all
\addlinespace
F3 & M1 primary & 69 & Combination cells, volatility unit, seed ensemble & protocol & \texttt{m1\_\allowbreak ensemble\_\allowbreak primary.csv} \\
F4 & M1 variance-unit & 69 & F3 under variance-unit QLIKE & protocol & \texttt{m1\_\allowbreak variance\_\allowbreak unit.csv} \\
F5 & Firm-identity reference & 69 & F3 against the identity-augmented reference; the battery is four specifications, each its own 69-cell family & protocol & \texttt{firm\_\allowbreak identity\_\allowbreak control.csv}, \texttt{firm\_\allowbreak identity\_\allowbreak ensemble.csv} \\
F6 & Maximal pool & 69 & F3 against the five-model price-pool reference & protocol & \texttt{maximal\_\allowbreak reference\_\allowbreak ensemble.md} \\
\addlinespace
F7 & Residual symmetric family & 12 & The pre-declared two-sided 12-cell family for the 8-K residual & \texttt{prereg}, FACTS \S13b & \texttt{omnibus\_\allowbreak m1.md}, \texttt{residual\_\allowbreak symmetric\_\allowbreak holm.md} \\
F8 & Earnings-window control & 3 & C6 event-driven at $h \in \{5, 10, 20\}$ against a firm-identity-plus-Item-2.02 reference & \texttt{prereg-rfa} & \texttt{itemcode\_\allowbreak control.csv} \\
\addlinespace
F9 & Deployable, fixed weights & 75 & Pooled DM per cell, validation-frozen weight scheme & pre-declaration block in \texttt{deployable\_\allowbreak combiner.md} & \texttt{deployable\_\allowbreak combiner.csv} \\
F10 & Deployable, expanding weights & 75 & Pooled DM per cell, expanding weight scheme & as F9 & as F9 \\
\addlinespace
F11 & Anonymisation, per arm & 6 & Three horizons $\times$ two references, within each executed arm (C6, C2) & \texttt{prereg-ea} v1.0/v1.4 & \texttt{anon\_\allowbreak arm.csv} \\
F12 & Range-based Parkinson / Garman--Klass & 69 each & Cascade rungs under the Parkinson and Garman--Klass labels, one family per rung per proxy & \texttt{prereg-cd} v1.0 & \texttt{rangebased\_\allowbreak cascade.csv} \\
F13 & Public-variant panels & as F1/F3 & Panels B and C replicate the F1 and F3 family structures & \texttt{prereg-h} v1.0 & \texttt{public\_\allowbreak variant\_\allowbreak cascade.csv} \\
F14 & MAEC audit & 24 & Four horizons $\times$ three text arms $\times$ two references, clustered; Holm inside each arm's 8 & \texttt{prereg-maec} & \texttt{maec\_\allowbreak audit.md} \\
F15 & Yelp cascade & per gate & Month-clustered DM; label-shuffle placebo primary & second-domain plan, gates G1--G5 & \texttt{yelp\_\allowbreak cascade.md} \\
\bottomrule
\end{tabular}
\end{table}

\FloatBarrier
\section{The Standalone Leaderboard and the 180-Comparison Census}
\label{app:res-standalone}

Table~\ref{tab:app-leaderboard} gives the full long-form standalone
leaderboard: test QLIKE in variance units at each horizon for the 25
models carried on the efficiency-frontier artefact, ordered as that
artefact orders them. A6 sits outside it because it is scored on the
return-matched subsample and A7 on the point-in-time VIX panel, so neither
loss is computed on the rows the other arms share; C5x, C6 and D4 sit
outside it because they are single-seed combination-track arms.
Table~\ref{tab:res-leaderboard} in Chapter~\ref{ch:results} carries the
long-form levels of three of those five, A6, C6 and D4. A7 and C5x carry
no levelled leaderboard row anywhere in the report; their results are
reported in the form each was run for instead. A7 enters the price-pool
completeness audit of Section~\ref{app:res-frontier} (as the
\mbox{HAR-X(VIX)} row) and the joint tests of Tables~\ref{tab:app-spa}
and~\ref{tab:app-mcs}, where it is the beating arm in two of the nine
rejected panels; C5x enters the input-parity control of
Table~\ref{tab:app-c5x}. The GPU-hours are in this appendix's own table.
The ordering is by best QLIKE over the three horizons. Three readings matter. First, the
first five places belong to the five price arms, none of which consumes
a single GPU-hour, and the HAR-RV reference sits third
among them at 0.5564/0.4618/0.2970. Second, the fine-tuned
encoders occupy the bottom: C2 FinBERT at truncation reads
1.3428/1.2440/0.7262, roughly twice the reference's loss at the short
horizons, and no C-block arm reaches the reference at any horizon.
Third, the frozen 7--8B embedding arms (C5) and the price-plus-embedding
fusions (D3) close most of that gap at a fraction of the fine-tuning
cost: D3 with gte-Qwen2 reaches 0.3270 for 0.20 GPU-hours against
C4 Longformer's 0.6051 for 254.72. None of them beats the
reference either. Reading more of the document, and spending twenty
times the compute to do it, moves the text arm without moving it to the
reference.
The event-driven and combined panels carry the same ordering, tabulated in
Section~\ref{app:res-standalone-other}, and the leaderboard is tested jointly
against the objection that its best row was selected after the fact in
Section~\ref{app:res-spamcs}.

\begin{table}[htbp]
\centering
\caption[Full standalone leaderboard, long-form panel]{Full standalone
leaderboard: test QLIKE in \emph{variance units} on the long-form panel,
with total and per-run GPU-hours. C- and D-block rows are the mean over
the three training seeds of each run's test QLIKE (a seed-mean of
per-seed metrics, not the QLIKE of the seed-ensemble forecast: an
off-basis quantity, tagged). A- and B-block rows are seed-invariant and
were run on CPU, reported at $\sim$0 GPU-hours. Lower is better; the A2
HAR-RV reference is the row every DM test in
Table~\ref{tab:app-dm180} is taken against.}
\label{tab:app-leaderboard}
\singlespacing\small
\setlength{\tabcolsep}{4pt}
% src: results/tables/cost_accuracy.md (all QLIKE levels and GPU-hours; ranks as committed)
\begin{tabular}{@{}rllrrrrr@{}}
\toprule
Rank & Model & Block & $h{=}5$ & $h{=}10$ & $h{=}20$ & GPU-h & GPU-h/run \\
\midrule
1  & A3 GARCH            & A & 0.4829 & 0.3750 & 0.2686 & 0.00   & 0.000 \\
2  & A5 ARIMA            & A & 0.5317 & 0.4067 & 0.2732 & 0.00   & 0.000 \\
3  & A2 HAR-RV \emph{(reference)} & A & 0.5564 & 0.4618 & 0.2970 & 0.00 & 0.000 \\
4  & A4 EGARCH           & A & 0.4631 & 0.3541 & 0.3180 & 0.00   & 0.000 \\
5  & A1 Historical vol.  & A & 1.1458 & 0.5600 & 0.3219 & 0.00   & 0.000 \\
6  & D3 price+gte-Qwen2  & D & 0.6698 & 0.5321 & 0.3270 & 0.20   & 0.022 \\
7  & D3 price+E5-Mistral & D & 0.7127 & 0.5688 & 0.3400 & 0.21   & 0.023 \\
8  & D3 price+Qwen3      & D & 0.7486 & 0.6232 & 0.3594 & 0.20   & 0.023 \\
9  & C5 gte-Qwen2        & C & 0.6538 & 0.5714 & 0.3916 & 5.72   & 0.635 \\
10 & C5 Qwen3-emb.       & C & 0.6866 & 0.6045 & 0.3940 & 4.44   & 0.494 \\
11 & D1 Concat-MLP       & D & 0.6408 & 0.6445 & 0.4374 & 11.43  & 1.270 \\
12 & C5 E5-Mistral       & C & 0.8467 & 0.6480 & 0.4375 & 8.94   & 0.993 \\
13 & D2 Gated fusion     & D & 0.7964 & 0.6022 & 0.4410 & 11.60  & 1.288 \\
14 & C2 FinBERT (S2)     & C & 1.4870 & 1.1952 & 0.5720 & 71.27  & 7.919 \\
15 & C2 FinBERT (S4)     & C & 1.4320 & 1.0943 & 0.6032 & 62.43  & 6.936 \\
16 & C4 Longformer (S5)  & C & 1.2351 & 1.0637 & 0.6051 & 254.72 & 28.302 \\
17 & C2 FinBERT (S3)     & C & 1.2430 & 1.2546 & 0.6731 & 67.80  & 7.533 \\
18 & C1 BERT (S2)        & C & 1.4557 & 1.0372 & 0.6905 & 51.73  & 5.748 \\
19 & B3 L--M linear      & B & 1.2937 & 1.0649 & 0.6918 & 0.00   & 0.000 \\
20 & C1 BERT (S1)        & C & 1.6444 & 1.1968 & 0.6995 & 13.71  & 1.523 \\
21 & B2 TF--IDF ridge    & B & 1.2371 & 1.0522 & 0.7083 & 0.00   & 0.000 \\
22 & B4 L--M features    & B & 1.4166 & 1.1223 & 0.7231 & 0.00   & 0.000 \\
23 & C2 FinBERT (S1)     & C & 1.3428 & 1.2440 & 0.7262 & 12.65  & 1.406 \\
24 & C3 RoBERTa (S1)     & C & 1.6303 & 1.0863 & 0.7518 & 13.31  & 1.479 \\
25 & B1 BoW ridge        & B & 2.2844 & 1.5916 & 1.0668 & 0.00   & 0.000 \\
\bottomrule
\end{tabular}
\end{table}

Table~\ref{tab:app-dm180} is the census itself: all 180 day-clustered DM
statistics against HAR-RV as a standalone forecaster, one row per
challenger and one column per disclosure-by-horizon panel. Every one of
the 180 statistics is positive. That is a stronger statement than the
significance count, because it means no comparison anywhere in the
census carries even a favourable point estimate for a test to
adjudicate. In addition, 155 of the 180 are significantly worse after
Holm correction within each 20-challenger panel. The gradient across the
table is informative in its own right: the fusion arms that carry the
price signal inside them (D1--D3) produce the smallest statistics and
account for 21 of the 25 non-significant cells. The pure
encoders sit between $+2.86$ and $+9.57$ throughout. A fusion arm
failing to be separated from the reference is not a text result, and the
table is laid out so that the two cannot be confused.

\begin{table}[p]
\centering
\caption[The 180-comparison standalone census]{The full standalone
census: day-clustered Diebold--Mariano statistics on squared forecast
error against A2 HAR-RV as a standalone forecaster, seed-ensemble
forecasts, all 20 challengers on all nine disclosure-by-horizon panels.
\emph{Positive means the challenger is worse than the baseline.} An
asterisk marks Holm significance within that panel's 20-comparison
family. All 180 statistics are positive; 155 are Holm-significantly
worse. Panel supports: long-form 7{,}951/7{,}933/7{,}902 observations
over 809/803/794 days, event-driven 25{,}109/25{,}001/24{,}732 over
996/991/981, combined 33{,}060/32{,}934/32{,}634 over 996/991/981.}
\label{tab:app-dm180}
\singlespacing\footnotesize
\setlength{\tabcolsep}{3.5pt}
% src: results/tables/variance_unit_standalone180.csv (dm_se_clu, worse_se_holm); cross-checked against results/tables/dm_pairwise_clustered.csv
\begin{tabular}{@{}lrrrrrrrrr@{}}
\toprule
& \multicolumn{3}{c}{Long-form (10-K/Q)} & \multicolumn{3}{c}{Event-driven (8-K)} & \multicolumn{3}{c}{Combined} \\
\cmidrule(lr){2-4}\cmidrule(lr){5-7}\cmidrule(lr){8-10}
Challenger & $h{=}5$ & $h{=}10$ & $h{=}20$ & $h{=}5$ & $h{=}10$ & $h{=}20$ & $h{=}5$ & $h{=}10$ & $h{=}20$ \\
\midrule
A3 GARCH               & $+5.22$* & $+4.73$* & $+4.27$* & $+4.33$* & $+3.37$* & $+2.58$\phantom{*} & $+6.08$* & $+4.78$* & $+3.49$* \\
A4 EGARCH              & $+2.74$* & $+2.39$\phantom{*} & $+3.61$* & $+5.75$* & $+5.08$* & $+3.96$* & $+6.26$* & $+5.46$* & $+4.18$* \\
A5 ARIMA               & $+5.84$* & $+4.21$* & $+3.10$* & $+5.48$* & $+3.75$* & $+2.20$\phantom{*} & $+6.17$* & $+4.12$* & $+2.27$\phantom{*} \\
\addlinespace
B2 TF--IDF ridge       & $+6.35$* & $+4.92$* & $+3.62$* & $+9.54$* & $+8.18$* & $+7.01$* & $+9.47$* & $+8.06$* & $+6.64$* \\
C1 BERT (S1)           & $+7.64$* & $+4.66$* & $+3.66$* & $+8.13$* & $+6.83$* & $+5.28$* & $+9.19$* & $+6.05$* & $+6.12$* \\
C1 BERT (S2)           & $+7.21$* & $+4.42$* & $+3.46$* & $+8.71$* & $+7.31$* & $+5.82$* & $+9.56$* & $+6.85$* & $+5.67$* \\
C2 FinBERT (S1)        & $+6.68$* & $+5.53$* & $+3.32$* & $+7.18$* & $+7.07$* & $+5.34$* & $+8.19$* & $+5.97$* & $+5.25$* \\
C2 FinBERT (S2)        & $+7.63$* & $+5.11$* & $+3.12$* & $+8.54$* & $+7.29$* & $+5.14$* & $+8.73$* & $+6.89$* & $+5.50$* \\
C2 FinBERT (S3)        & $+6.68$* & $+5.50$* & $+3.41$* & $+7.43$* & $+6.00$* & $+5.04$* & $+8.89$* & $+6.50$* & $+5.35$* \\
C2 FinBERT (S4)        & $+7.07$* & $+4.89$* & $+2.86$* & $+8.08$* & $+6.29$* & $+6.06$* & $+9.49$* & $+6.65$* & $+4.98$* \\
C3 RoBERTa (S1)        & $+7.68$* & $+5.61$* & $+3.87$* & $+7.47$* & $+7.11$* & $+5.84$* & $+8.66$* & $+7.54$* & $+5.88$* \\
C4 Longformer (S5)     & $+6.06$* & $+4.56$* & $+2.92$* & $+7.87$* & $+6.51$* & $+5.73$* & $+9.01$* & $+6.80$* & $+5.36$* \\
C5 E5-Mistral          & $+5.24$* & $+6.96$* & $+6.76$* & $+8.38$* & $+7.44$* & $+6.62$* & $+9.07$* & $+7.70$* & $+6.52$* \\
C5 gte-Qwen2           & $+6.71$* & $+7.31$* & $+6.67$* & $+9.57$* & $+9.02$* & $+8.77$* & $+8.92$* & $+7.79$* & $+6.61$* \\
C5 Qwen3-emb.          & $+5.92$* & $+7.08$* & $+6.79$* & $+9.20$* & $+8.66$* & $+7.70$* & $+8.78$* & $+7.63$* & $+6.38$* \\
D1 Concat-MLP          & $+1.61$\phantom{*} & $+2.62$\phantom{*} & $+2.44$\phantom{*} & $+6.98$* & $+6.43$* & $+5.07$* & $+6.85$* & $+4.74$* & $+3.92$* \\
D2 Gated fusion        & $+3.53$* & $+1.86$\phantom{*} & $+2.03$\phantom{*} & $+6.47$* & $+5.73$* & $+4.30$* & $+7.22$* & $+5.34$* & $+5.65$* \\
D3 price+E5-Mistral    & $+3.19$* & $+1.96$\phantom{*} & $+0.63$\phantom{*} & $+3.92$* & $+2.51$* & $+1.83$\phantom{*} & $+4.07$* & $+2.14$\phantom{*} & $+2.42$\phantom{*} \\
D3 price+gte-Qwen2     & $+1.91$\phantom{*} & $+0.77$\phantom{*} & $+0.24$\phantom{*} & $+2.24$* & $+2.11$* & $+1.16$\phantom{*} & $+2.69$* & $+0.99$\phantom{*} & $+0.24$\phantom{*} \\
D3 price+Qwen3         & $+3.81$* & $+2.62$\phantom{*} & $+1.26$\phantom{*} & $+3.73$* & $+2.52$* & $+1.80$\phantom{*} & $+3.80$* & $+1.90$\phantom{*} & $+1.93$\phantom{*} \\
\bottomrule
\end{tabular}
\end{table}

Table~\ref{tab:app-census} re-adjudicates those same 180 comparisons
(the same forecasts, the same rows) under three loss conventions and
three error structures. It gives the 153-comparison text-and-fusion
sub-family alongside the full census in the four rows that carry it; the
superseded observation-level row and the two-way row print a dash there, that
decomposition never having been computed on either. The one convention
that produces winners is the variance-unit restatement, and all seven
winners are price arms: A3 GARCH in six cells and A4 EGARCH at long-form
$h{=}5$. Under that convention the text and fusion arms are not merely
0 of 153 better; they are 153 of 153 significantly \emph{worse}, the
harshest verdict any convention returns on them. Moving from
observation-level HAC to day clustering and then to two-way firm-by-day
clustering \cite{cameron2011twoway} costs the study evidence at every
step: the significantly-worse count falls 174, 155, 152 and the
median absolute statistic falls 15.44, 5.70, 4.86. All three are
therefore printed rather than the primary alone.

\begin{table}[htbp]
\centering
\caption[The 180-comparison census under three losses and three error
structures]{The standalone census re-adjudicated. Panel (a) holds the
error structure at the day-clustered primary and changes the loss; panel
(b) holds the loss at squared error and changes the error structure.
Identical forecasts, rows and Holm families throughout; only the named
axis moves. ``DM $<0$'' counts comparisons carrying a favourable point
estimate before any test; ``better'' and ``worse'' are Holm-corrected
within each 20-challenger panel. ``med.'' is the median absolute DM
statistic over the 180 comparisons. The text/fusion sub-family is the 17
non-price challengers over the nine panels; a dash means the source file
does not carry that decomposition.}
\label{tab:app-census}
\singlespacing\small
\setlength{\tabcolsep}{4pt}
% src: results/tables/variance_unit_standalone180.csv (all counts re-derived 2026-08-02); median |DM| from TECHNICAL_SUPPLEMENT.md Table S3(a); obs-level and two-way counts from results/tables/dm_pairwise_clustered.md and results/tables/twoway_cluster.md
\begin{tabular}{@{}lrrrrrr@{}}
\toprule
& \multicolumn{4}{c}{All 180 comparisons} & \multicolumn{2}{c}{Text/fusion (153)} \\
\cmidrule(lr){2-5}\cmidrule(lr){6-7}
Convention & DM $<0$ & better & worse & med. & better & worse \\
\midrule
\multicolumn{7}{@{}l}{\emph{(a) Day-clustered inference; loss convention varied}} \\
Squared error (primary) & 0  & 0 & 155 & 5.70  & 0 & 132 \\
QLIKE, volatility units & 1  & 0 & 161 & ---   & 0 & 146 \\
QLIKE, variance units   & 14 & 7 & 153 & ---   & 0 & 153 \\
\midrule
\multicolumn{7}{@{}l}{\emph{(b) Squared error; error structure varied}} \\
Observation-level HAC \emph{(superseded)} & 0 & 0 & 174 & 15.44 & --- & --- \\
Day-clustered, lag $h-1$ days             & 0 & 0 & 155 & 5.70  & 0 & 132 \\
Two-way firm $\times$ day                 & 0 & 0 & 152 & 4.86  & --- & --- \\
\bottomrule
\end{tabular}
\end{table}

\FloatBarrier
\section{The 69-Cell Combination Grid}
\label{app:res-grid}

Tables~\ref{tab:app-grid-ed} and~\ref{tab:app-grid-lf} are the
combination grid cell by cell: the 24 event-driven and 45 long-form
cells at the primary rung. Each carries its relative QLIKE improvement over
the recalibrated HAR reference, its day-clustered DM statistic, its
Holm-adjusted $p$ within the 69-cell family, its label-shuffle placebo
statistic and the genuine verdict that requires all three. The right-hand
column groups then carry the same cell's increment over the fitted
maximal price pool and over the firm-identity reference, with their own
Holm verdicts. Read together the two tables reproduce every count in
Chapter~\ref{ch:results}: 46 cells at raw $p<.05$, 38 genuine, 9
surviving the pool, 8 surviving firm identity, and no row anywhere
carrying a Y in both of the last two column pairs.
Every weight behind these cells is fitted once on the validation block and
frozen. Section~\ref{app:res-deploy} reruns the same cells under weights
refitted before each quarter on strictly earlier filings, which is what a
forecaster running the method would have had to do.

Three features of the grid are visible only at this resolution. The
surviving primary-rung increments are not uniformly negligible: they
reach $+5.92\%$ (long-form B2 TF--IDF ridge at $h{=}20$) and the grid
spans $-3.86\%$ to $+5.92\%$. The stronger references are therefore
adjudicating real quantities rather than rounding error. The placebo
column earns its place by not firing: of the 38 cells with a negative DM
statistic and a Holm-adjusted $p$ below .05, none has a label-shuffle
statistic outside the $|\mathrm{DM}|<2$ pass band. The placebo gate
removes no cell that the correction had not already removed: a
null result for that instrument on this grid, reported as such. The two
largest placebo statistics in either table, $-3.17$ at long-form B4 L--M
features $h{=}20$ and $+3.00$ at event-driven D2 gated fusion $h{=}5$,
sit on cells that are already excluded. The first has a negative
increment; the second fails Holm at $p = .988$. And the
seed column matters for interpretation rather than
for arithmetic: 33 of the 69 cells are three-seed ensembles and 36 are
single-seed by construction (the four B-block arms are deterministic
fits and the two prompted arms have no training seed). The ensemble basis is a
statement about the C1--C5 and D1--D3 rows and leaves the rest untouched: the
twelve C6 and D4 cells sit in the C and D blocks and are single-seed too. % src: results/tables/control_intersection_ensemble.csv (45 C/D rows: 33 with n_seeds 3, 12 with n_seeds 1, those twelve being C6_llmtext and D4_llmfused)

\begin{table}[htbp]
\centering
\caption[The 69-cell combination grid: event-driven cells]{The
combination grid, \textbf{event-driven (8-K) channel}, all 24 cells.
Basis: seed-ensemble forecasts, volatility-unit QLIKE, log-space
combiner fitted on validation and frozen on test, day-clustered DM with
Holm inside the 69-cell family. \emph{Negative DM and positive rel.\%
mean text improves on the reference.} $s$ = number of training seeds;
plac.\ = label-shuffle placebo DM (pass band $|$DM$|<2$); G = genuine at
the primary rung (DM $<0$, Holm $<.05$, placebo in band). The last two
column pairs are the same cell's increment over the fitted five-model
price pool and over the firm-identity-augmented reference, each with its
own Holm verdict on the merged-grid support.}
\label{tab:app-grid-ed}
\singlespacing\footnotesize
\setlength{\tabcolsep}{3.5pt}
% src: results/tables/control_intersection_ensemble.csv (all columns); rung counts cross-checked against m1_ensemble_primary.md, maximal_reference_ensemble.md, firm_identity_ensemble.md
\begin{tabular}{@{}lrrrrrrcrcrc@{}}
\toprule
& & & \multicolumn{5}{c}{Primary rung: recalibrated HAR} & \multicolumn{2}{c}{Maximal pool} & \multicolumn{2}{c}{Firm identity} \\
\cmidrule(lr){4-8}\cmidrule(lr){9-10}\cmidrule(lr){11-12}
Model & $h$ & $s$ & rel.\% & DM & Holm $p$ & plac. & G & rel.\% & Holm & rel.\% & Holm \\
\midrule
B1 BoW ridge       &  5 & 1 & $+1.331$ & $-3.35$ & .0252 & $+0.59$ & Y & $+0.312$ & -- & $+0.461$ & Y \\
B1 BoW ridge       & 10 & 1 & $+1.229$ & $-3.25$ & .0332 & $+0.05$ & Y & $-0.157$ & -- & $+0.250$ & Y \\
B1 BoW ridge       & 20 & 1 & $+1.531$ & $-3.10$ & .0489 & $+0.45$ & Y & $-0.267$ & -- & $-0.193$ & -- \\
B2 TF--IDF ridge   &  5 & 1 & $+1.207$ & $-2.76$ & .1237 & $+0.39$ & -- & $+0.147$ & -- & $+0.444$ & -- \\
B2 TF--IDF ridge   & 10 & 1 & $+1.355$ & $-2.97$ & .0681 & $+0.38$ & -- & $-0.058$ & -- & $+0.281$ & Y \\
B2 TF--IDF ridge   & 20 & 1 & $+1.840$ & $-3.11$ & .0489 & $+0.42$ & Y & $-0.127$ & -- & $-0.045$ & -- \\
B3 L--M linear     &  5 & 1 & $+0.248$ & $-2.43$ & .2730 & $+0.01$ & -- & $+0.161$ & -- & $+0.110$ & -- \\
B3 L--M linear     & 10 & 1 & $+0.200$ & $-1.23$ & 1.0000 & $+0.16$ & -- & $+0.087$ & -- & $+0.108$ & -- \\
B3 L--M linear     & 20 & 1 & $+0.260$ & $-1.12$ & 1.0000 & $+0.84$ & -- & $+0.094$ & -- & $+0.170$ & -- \\
B4 L--M feat.      &  5 & 1 & $+0.184$ & $-0.99$ & 1.0000 & $+0.46$ & -- & $+0.195$ & -- & $+0.218$ & -- \\
B4 L--M feat.      & 10 & 1 & $+0.082$ & $-2.10$ & .6041 & $-0.05$ & -- & $+0.084$ & -- & $+0.130$ & -- \\
B4 L--M feat.      & 20 & 1 & $+0.247$ & $-2.01$ & .6569 & $+0.47$ & -- & $+0.086$ & -- & $+0.110$ & -- \\
C2 FinBERT (S1)    &  5 & 3 & $+2.572$ & $-5.92$ & $<$.0001 & $-1.72$ & Y & $+0.768$ & -- & $-0.061$ & -- \\
C2 FinBERT (S1)    & 10 & 3 & $+2.416$ & $-5.68$ & $<$.0001 & $+0.96$ & Y & $+0.084$ & -- & $-0.665$ & -- \\
C2 FinBERT (S1)    & 20 & 3 & $+1.578$ & $-1.94$ & .6569 & $+0.33$ & -- & $-1.594$ & -- & $-0.613$ & -- \\
C6 prompted LLM    &  5 & 1 & $+1.214$ & $-5.04$ & $<$.0001 & $+0.10$ & Y & $+0.398$ & -- & $+0.518$ & Y \\
C6 prompted LLM    & 10 & 1 & $+0.998$ & $-3.76$ & .0065 & $+0.95$ & Y & $+0.121$ & -- & $+0.244$ & Y \\
C6 prompted LLM    & 20 & 1 & $+0.658$ & $-1.98$ & .6569 & $+1.11$ & -- & $+0.086$ & -- & $+0.215$ & Y \\
D2 Gated fusion    &  5 & 3 & $+0.559$ & $-1.70$ & .9875 & $+3.00$ & -- & $+0.297$ & -- & $+0.226$ & -- \\
D2 Gated fusion    & 10 & 3 & $+0.177$ & $-0.42$ & 1.0000 & $-0.20$ & -- & $-0.017$ & -- & $+0.050$ & -- \\
D2 Gated fusion    & 20 & 3 & $-2.764$ & $+3.55$ & .0139 & $-0.44$ & -- & $-2.486$ & -- & $-7.992$ & -- \\
D4 price+C6        &  5 & 1 & $-0.013$ & $+3.37$ & .0241 & $+1.58$ & -- & $+0.114$ & -- & $-0.072$ & -- \\
D4 price+C6        & 10 & 1 & $-0.008$ & $+0.66$ & 1.0000 & $-0.08$ & -- & $-0.030$ & -- & $-0.021$ & -- \\
D4 price+C6        & 20 & 1 & $-0.355$ & $+4.69$ & .0001 & $-0.19$ & -- & $-0.243$ & -- & $-0.251$ & -- \\
\midrule
\multicolumn{12}{@{}l}{\emph{Event-driven subtotals: 8 genuine of 24; 0 surviving the maximal pool; 6 surviving firm identity.}} \\
\bottomrule
\end{tabular}
\end{table}

\begin{table}[p]
\centering
\caption[The 69-cell combination grid: long-form cells]{The combination
grid, \textbf{long-form (10-K/Q) channel}, all 45 cells. Columns,
conventions and sign rules exactly as in Table~\ref{tab:app-grid-ed}.
One pool survivor carries an increment too small to print at three
decimals: long-form D1 Concat-MLP at $h{=}20$ reads $+0.001\%$ over the
pool, which is a Holm survivor on an increment of essentially zero and
is shown rather than rounded away.}
\label{tab:app-grid-lf}
\singlespacing\footnotesize
\setlength{\tabcolsep}{3.5pt}
% src: results/tables/control_intersection_ensemble.csv (all columns)
\begin{tabular}{@{}lrrrrrrcrcrc@{}}
\toprule
& & & \multicolumn{5}{c}{Primary rung: recalibrated HAR} & \multicolumn{2}{c}{Maximal pool} & \multicolumn{2}{c}{Firm identity} \\
\cmidrule(lr){4-8}\cmidrule(lr){9-10}\cmidrule(lr){11-12}
Model & $h$ & $s$ & rel.\% & DM & Holm $p$ & plac. & G & rel.\% & Holm & rel.\% & Holm \\
\midrule
B1 BoW ridge       &  5 & 1 & $+1.645$ & $-3.83$ & .0050 & $+0.43$ & Y & $+0.870$ & -- & $-0.138$ & -- \\
B1 BoW ridge       & 10 & 1 & $+1.443$ & $-4.15$ & .0015 & $+0.39$ & Y & $+0.506$ & -- & $-2.720$ & -- \\
B1 BoW ridge       & 20 & 1 & $+2.986$ & $-5.45$ & $<$.0001 & $-0.25$ & Y & $+1.608$ & -- & $-4.985$ & -- \\
B2 TF--IDF ridge   &  5 & 1 & $+3.331$ & $-5.39$ & $<$.0001 & $+0.51$ & Y & $+1.672$ & Y & $-0.615$ & -- \\
B2 TF--IDF ridge   & 10 & 1 & $+3.482$ & $-8.89$ & $<$.0001 & $+0.33$ & Y & $+1.952$ & Y & $-3.892$ & -- \\
B2 TF--IDF ridge   & 20 & 1 & $+5.920$ & $-9.04$ & $<$.0001 & $+0.30$ & Y & $+4.013$ & Y & $-8.089$ & -- \\
B3 L--M linear     &  5 & 1 & $+0.495$ & $-2.58$ & .1894 & $+0.59$ & -- & $-0.221$ & -- & $+0.207$ & -- \\
B3 L--M linear     & 10 & 1 & $+1.794$ & $-4.62$ & .0002 & $+0.27$ & Y & $+0.641$ & -- & $+0.261$ & Y \\
B3 L--M linear     & 20 & 1 & $+3.485$ & $-4.69$ & .0001 & $-1.26$ & Y & $+1.927$ & -- & $-0.017$ & -- \\
B4 L--M feat.      &  5 & 1 & $+0.113$ & $+1.02$ & 1.0000 & $+0.43$ & -- & $+0.023$ & -- & $-0.139$ & -- \\
B4 L--M feat.      & 10 & 1 & $-0.922$ & $+3.32$ & .0276 & $-1.69$ & -- & $-0.729$ & -- & $-2.767$ & -- \\
B4 L--M feat.      & 20 & 1 & $-1.920$ & $+3.38$ & .0241 & $-3.17$ & -- & $-1.554$ & -- & $-7.133$ & -- \\
C1 BERT (S1)       &  5 & 3 & $+3.420$ & $-3.25$ & .0332 & $-0.03$ & Y & $+1.088$ & -- & $+1.376$ & -- \\
C1 BERT (S1)       & 10 & 3 & $-0.849$ & $+3.06$ & .0533 & $-1.02$ & -- & $-0.112$ & -- & $-4.751$ & -- \\
C1 BERT (S1)       & 20 & 3 & $+2.702$ & $-6.96$ & $<$.0001 & $+0.39$ & Y & $+0.896$ & -- & $-6.488$ & -- \\
C2 FinBERT (S1)    &  5 & 3 & $+1.898$ & $-4.53$ & .0003 & $-0.72$ & Y & $+0.626$ & -- & $-0.268$ & -- \\
C2 FinBERT (S1)    & 10 & 3 & $+2.616$ & $-6.67$ & $<$.0001 & $+0.80$ & Y & $+0.931$ & -- & $-3.180$ & -- \\
C2 FinBERT (S1)    & 20 & 3 & $-0.590$ & $-0.40$ & 1.0000 & $+0.78$ & -- & $-5.513$ & -- & $-0.015$ & -- \\
C2 FinBERT (S2)    &  5 & 3 & $+1.207$ & $-5.57$ & $<$.0001 & $+0.21$ & Y & $+0.224$ & Y & $-2.196$ & -- \\
C2 FinBERT (S2)    & 10 & 3 & $+0.484$ & $-7.59$ & $<$.0001 & $+0.52$ & Y & $-0.140$ & -- & $-3.732$ & -- \\
C2 FinBERT (S2)    & 20 & 3 & $+1.678$ & $-4.89$ & $<$.0001 & $-0.39$ & Y & $-0.576$ & -- & $-3.815$ & -- \\
C2 FinBERT (S3)    &  5 & 3 & $+2.900$ & $-5.36$ & $<$.0001 & $+0.38$ & Y & $+1.253$ & -- & $-0.822$ & -- \\
C2 FinBERT (S3)    & 10 & 3 & $+2.313$ & $-5.26$ & $<$.0001 & $+0.55$ & Y & $-0.119$ & -- & $-1.207$ & -- \\
C2 FinBERT (S3)    & 20 & 3 & $-3.859$ & $+1.68$ & .9875 & $+1.06$ & -- & $-12.691$ & -- & $+0.799$ & -- \\
C2 FinBERT (S4)    &  5 & 3 & $+1.456$ & $-4.78$ & $<$.0001 & $+0.37$ & Y & $+0.505$ & -- & $-0.118$ & -- \\
C2 FinBERT (S4)    & 10 & 3 & $+0.364$ & $-6.82$ & $<$.0001 & $-0.24$ & Y & $-0.064$ & -- & $-4.073$ & -- \\
C2 FinBERT (S4)    & 20 & 3 & $+3.078$ & $-8.34$ & $<$.0001 & $+0.24$ & Y & $+1.663$ & Y & $-6.866$ & -- \\
C3 RoBERTa (S1)    &  5 & 3 & $+0.298$ & $-5.13$ & $<$.0001 & $+0.06$ & Y & $-0.347$ & -- & $-2.779$ & -- \\
C3 RoBERTa (S1)    & 10 & 3 & $+1.842$ & $-5.06$ & $<$.0001 & $-0.73$ & Y & $+0.424$ & -- & $-0.777$ & -- \\
C3 RoBERTa (S1)    & 20 & 3 & $+0.020$ & $-3.90$ & .0040 & $-0.44$ & Y & $-0.028$ & -- & $-3.563$ & -- \\
C4 Longformer      &  5 & 3 & $+1.473$ & $-6.18$ & $<$.0001 & $+0.09$ & Y & $+0.467$ & Y & $-2.206$ & -- \\
C4 Longformer      & 10 & 3 & $-2.953$ & $+10.32$ & $<$.0001 & $+0.23$ & -- & $-3.539$ & -- & $-11.383$ & -- \\
C4 Longformer      & 20 & 3 & $+0.915$ & $-4.38$ & .0005 & $+1.20$ & Y & $-0.289$ & -- & $-5.026$ & -- \\
C6 prompted LLM    &  5 & 1 & $+1.794$ & $-6.31$ & $<$.0001 & $+0.83$ & Y & $+1.046$ & Y & $-0.184$ & -- \\
C6 prompted LLM    & 10 & 1 & $+2.252$ & $-7.92$ & $<$.0001 & $-1.45$ & Y & $+1.066$ & Y & $-0.433$ & -- \\
C6 prompted LLM    & 20 & 1 & $+0.275$ & $-3.23$ & .0332 & $-1.44$ & Y & $+0.233$ & -- & $+0.079$ & Y \\
D1 Concat-MLP      &  5 & 3 & $-1.044$ & $+2.66$ & .1612 & $+0.79$ & -- & $-0.888$ & -- & $-1.380$ & -- \\
D1 Concat-MLP      & 10 & 3 & $-0.456$ & $+3.54$ & .0141 & $-0.76$ & -- & $-0.384$ & -- & $-1.701$ & -- \\
D1 Concat-MLP      & 20 & 3 & $+0.124$ & $-5.57$ & $<$.0001 & $-0.61$ & Y & $+0.001$ & Y & $-2.800$ & -- \\
D2 Gated fusion    &  5 & 3 & $+0.192$ & $-2.11$ & .6041 & $+0.76$ & -- & $+0.036$ & -- & $-0.158$ & -- \\
D2 Gated fusion    & 10 & 3 & $-0.023$ & $+2.02$ & .6569 & $-0.41$ & -- & $+0.141$ & -- & $+0.196$ & -- \\
D2 Gated fusion    & 20 & 3 & $-2.269$ & $+5.56$ & $<$.0001 & $-0.18$ & -- & $+0.781$ & -- & $-3.485$ & -- \\
D4 price+C6        &  5 & 1 & $+0.149$ & $-0.75$ & 1.0000 & $+1.00$ & -- & $-0.105$ & -- & $-0.001$ & -- \\
D4 price+C6        & 10 & 1 & $-0.148$ & $-0.43$ & 1.0000 & $+0.87$ & -- & $+0.288$ & -- & $-0.149$ & -- \\
D4 price+C6        & 20 & 1 & $+0.605$ & $-3.58$ & .0127 & $-0.49$ & Y & $+0.365$ & -- & $+0.083$ & -- \\
\midrule
\multicolumn{12}{@{}l}{\emph{Long-form subtotals: 30 genuine of 45; 9 surviving the maximal pool; 2 surviving firm identity.}} \\
\bottomrule
\end{tabular}
\end{table}

\FloatBarrier
\section{The Reference Ladder, Rung by Rung}
\label{app:res-ladder}

Table~\ref{tab:app-ladder} collects every rung and every off-basis
restatement of the ladder in one place, with each row's denominator
printed beside its counts, because the rungs do not all share one. The
VIX-augmented rung is the case that most needs the warning: its 19 is
19 \emph{of the 38 genuine cells}, not 19 of 69, because that control
was run on the cells an earlier single-seed, observation-level screen had
passed. That set is not the primary rung's own 38: the two overlap in 30
cells, and the control adds both twenty-day prompted-LLM cells, so its
Holm correction runs inside 40.
% src: results/tables/vix_control.csv (40 rows); results/tables/forecast_combination_summary.json (genuine_cells, 38); results/tables/m1_ensemble_primary.csv (genuine_ens_vol, 38; overlap with the 38 is 30, with the 40 is 31)
Two further
distinctions are easy to lose. The specification battery for the
firm-identity term sits on the single-seed 2026 basis, and its four
counts (8, 10, 14, 8) price the spread across four reasonable ways of
writing down ``the firm's typical volatility'', not four attempts at the
same quantity. The most permissive of them, the expanding
point-in-time mean, is also the weakest identity instrument, improving
the plain recalibrated HAR on its own in only two of six channel-horizon
comparisons against four to five for the fixed-window specifications.
And the maximal-pool row is not an artefact of fitting six weights: an
equal-weight pool that fits nothing beyond the two recalibration
coefficients already leaves 17. The
single validation-best price member leaves 34. Neither number is the 38 minus a
loss: the 38 is scored on the reference-plus-text join and every pooled
reference on the smaller five-price-model intersection, so the fall mixes a
change of reference with a change of row set --- which is the confound
Figure~\ref{fig:app-matched-rows} exists to separate.

Table~\ref{tab:app-survivors} names the survivor cells at the two
strengthened rungs, which is the only way to see that the empty
conjunction is earned rather than forced. All nine pool survivors are
long-form; the eight identity survivors are six event-driven and two
long-form. Eleven survivors therefore coexist inside the same 45
long-form rows without a single intersection. A reader who assumed the
two rungs were nested, and therefore that the conjunction had to be
$\min(9,8)$, would be wrong on this panel as a matter of fact and not
of arithmetic.
Two further price models were computed and left out of this pool, which is the
one place the ladder can be accused of stopping short of the strongest
available reference. Section~\ref{app:res-frontier} adds them and prices what
they do to the text verdict.

\begin{table}[htbp]
\centering
\caption[Reference-ladder survivor counts by rung]{Reference-ladder
survivor counts. Primary basis unless the row says otherwise; each row's
denominator is printed because they differ. The primary rung's Holm
count is additionally placebo-gated, and so is the VIX-augmented rung, whose
survivor flag requires a negative DM, a Holm-adjusted $p$ below .05 and a
placebo statistic inside the pass band; the remaining rungs are Holm only. % src: scripts/analysis/vix_control.py:190 (survives_vix = vix_dm_q < 0 and vix_holm < .05 and |vix_placebo_dm| < 2)
Indented rows are restatements of the row above them on a changed basis,
not additional rungs. The four-specification identity battery is written
on the single-seed 2026 basis, so its counts are comparable with each
other but not with the ensemble row above them. Its own reading of the
validation-window specification is 14 raw and 8 Holm, and that Holm
survivor set coincides exactly with the ensemble's eight.}
\label{tab:app-ladder}
\singlespacing\small
\setlength{\tabcolsep}{4pt}
\begin{tabular}{@{}llrrl@{}}
\toprule
Rung / restatement & Denom. & Raw & Holm & Basis and note \\
\midrule
Recalibrated HAR (primary rung)      & 69 & 46 & 38 & seed ensemble, placebo-gated \\ % src: results/tables/control_intersection_ensemble.md; m1_ensemble_primary.md
\quad single seed 2026               & 69 & 39 & 29 & single-seed restatement \\ % src: results/tables/m1_clustered.csv; control_intersection.md
\quad two-way firm $\times$ day      & 69 & --- & 27 & 24 on seed 2026 \\ % src: TECHNICAL_SUPPLEMENT.md Table S3(b); results/tables/twoway_cluster.md
\quad variance-unit QLIKE            & 69 & --- & 20 & unit restatement \\ % src: results/tables/m1_variance_unit.md
VIX-augmented reference              & 38 & --- & 19 & \textbf{of the 38}; seed 2026, Holm in 40 \\ % src: results/tables/vix_control.md
\midrule
Validation-best single price member  & 69 & --- & 34 & one fitted member \\ % src: results/tables/maximal_pool_robustness.csv (ens basis, valbest_single)
Equal-weight five-model pool         & 69 & --- & 17 & never fitted \\ % src: results/tables/maximal_pool_robustness.csv (ens basis, eqw_pool)
Fitted maximal five-model pool       & 69 & 26 & 9  & the reported pool rung \\ % src: results/tables/maximal_pool_robustness.csv (ens basis, fitted_pool); maximal_reference_ensemble.md
\midrule
Firm identity, validation-window mean & 69 & 15 & 8 & the reported identity rung \\ % src: results/tables/firm_identity_ensemble.md; control_intersection_ensemble.md
\quad train-period (2010--19) mean    & 69 & 17 & 10 & battery, seed 2026 \\ % src: results/tables/firm_fe_battery.md
\quad expanding point-in-time mean    & 69 & 27 & 14 & battery, seed 2026 \\ % src: results/tables/firm_fe_battery.md
\quad empirical-Bayes shrunk mean     & 69 & 14 & 8  & battery, seed 2026 \\ % src: results/tables/firm_fe_battery.md
\midrule
Pool \textbf{and} identity            & 69 & 4 & \textbf{0} & survivor sets disjoint \\ % src: results/tables/control_intersection_ensemble.md
\textbf{Full conjunction}             & 69 & 4 & \textbf{0} & primary, pool and identity \\ % src: results/tables/control_intersection_ensemble.md
\midrule
Injected signal at 0.3 / 0.5 / 1.0\%  & 69 & --- & 2 / 6 / 13 & conjunction recovery \\ % src: results/tables/signal_injection_power.csv
\bottomrule
\end{tabular}
\end{table}

\begin{table}[htbp]
\centering
\caption[Named survivor cells at the two strengthened rungs]{The
survivor cells at the two strengthened rungs, named. The two sets are
disjoint, and not trivially so: nine pool survivors and two identity
survivors coexist within the same 45 long-form rows.}
\label{tab:app-survivors}
\singlespacing\small
% src: results/tables/control_intersection_ensemble.csv (maximal_holm, firm_holm)
\begin{tabular}{@{}llp{0.52\textwidth}@{}}
\toprule
Rung & Count & Cells (channel / model / horizon) \\
\midrule
Maximal price pool & 9 of 69 &
long-form: B2 TF--IDF ridge $h{=}5,10,20$; C2 FinBERT-S2 $h{=}5$;
C2 FinBERT-S4 $h{=}20$; C4 Longformer $h{=}5$; C6 prompted
$h{=}5,10$; D1 Concat-MLP $h{=}20$. Event-driven: none. \\
\addlinespace
Firm identity & 8 of 69 &
event-driven: B1 BoW ridge $h{=}5,10$; B2 TF--IDF ridge $h{=}10$;
C6 prompted $h{=}5,10,20$. Long-form: B3 L--M linear $h{=}10$;
C6 prompted $h{=}20$. \\
\addlinespace
Intersection & 0 of 69 & --- \\
\bottomrule
\end{tabular}
\end{table}

\FloatBarrier
\section{Identity Instruments}
\label{app:res-identity}

Table~\ref{tab:app-identity} is the anonymisation arm in full, in two
panels because the two references are different denominators for the
same loss. The increments over them must not be read as one series.
Panel (a) measures the masked and unmasked increments against the
recalibrated HAR and forms the identity share from their ratio, with
day-block bootstrap intervals \cite{politis1994stationary}. Panel (b)
repeats the measurement against the firm-identity reference and adds the
two triangulating instruments the same arm table carries. The
instrument is reported at full strength where it contradicts the
headline. The fine-tuned FinBERT arm prices its own $h{=}5$ share at
$+0.02$, meaning masking removes essentially nothing. Its $h{=}20$
share is \emph{negative} at $-0.39$ with an interval running to $-4.94$,
meaning masking made that increment larger. The third registered arm,
the TF--IDF ridge, exited at its reproduction gate before scoring and is
shown as an empty row rather than omitted, so a reader counts two
executed arms out of three registered. Two firm-identity cells carry no
share at all because their unmasked increment is at or below zero, and
they are marked n/a and kept in the family carrying their DM test.

The seven summary statistics at the foot of the table are seven
different quantities that must never be substituted for one another, and three
of them circulate widely enough to name here. The
pooled median of $0.51$ $[0.44, 0.56]$ over the two executed arms' six
vs-HAR cells is the figure this report treats as the headline; the
prompted arm's own median is $0.56$. The $h{=}5$ cell share happens
also to round to $0.51$, which is a coincidence between different
estimands. All three lie inside the bracket the other two instruments
set (matched-swap retention $0.29$ below, per-cell reference-interval
bound $0.72$ above), which is what the interval logic predicts if it
is sound. None of them is an upper bound.

What the mask does to a document is legible only on a document. The
treatment is a named-entity pass over each filing (spaCy
\texttt{en\_core\_web\_lg}) together with a rule table of 2{,}004
registrant names, and it writes one of five bracketed placeholders,
\texttt{[FIRM]}, \texttt{[TICKER]}, \texttt{[PERSON]},
\texttt{[PRODUCT]} and \texttt{[CIK]}, over every span it finds. The
median 8-K carries 21 such spans and the mean 31. % src: results/anon/mask_stats.json (ner.model en_core_web_lg-3.8.0; n_name_entries 2004; median_spans_per_doc 21.0; mean_spans_per_doc 31.3)
The committed audit sample records the before and after on 100 documents,
of which the first opens as follows. % src: results/anon/mask_audit_sample.md (first entry: AAL, 8-K accession 0000006201-20-000014, spans 22, mask fraction 0.075)

\medskip
\noindent\fbox{\begin{minipage}{0.94\textwidth}
{\small
\noindent\textit{original}

\ttfamily\raggedright
Date of Report (Date of earliest event reported): January 29, 2020
AMERICAN AIRLINES GROUP INC. AMERICAN AIRLINES, INC. (Exact name of
registrant as specified in its charter) Delaware 1-8400 75-1825172
Delaware 1-2691 13-1502798 (State or other Jurisdiction of
Incorporation) (Commission File Number)

\rmfamily\noindent\textit{masked}

\ttfamily\raggedright
Date of Report (Date of earliest event reported): January 29, 2020
{}[FIRM]. [FIRM]. (Exact name of registrant as specified in its charter)
Delaware 1-8400 75-1825172 Delaware 1-2691 13-1502798 ([FIRM] or other
[FIRM]) (Commission File Number)
\par}
\end{minipage}}
\medskip

Two properties of the treatment are visible there and in no gate
statistic. The identifying material sits in the cover page, and the mask
reaches it. The pass also over-reaches on the boilerplate around it,
taking \emph{State} and \emph{Jurisdiction of Incorporation} with the
names, so the masked arm reads a document from which a little
non-identifying text has gone as well. Whether that removal rather than
the loss of identity is what moves the increment is asked and answered in
Section~\ref{sec:val-instruments}.


\begin{table}[htbp]
\centering
\caption[The anonymisation arm and the identity instruments]{Identity
instruments: the entity-anonymisation arm on the event-driven channel,
with the two triangulating instruments. Firm names, tickers,
executive names, product names and CIKs across all 112{,}528 8-K documents are
masked and the text forecasts regenerated in an identical environment. The
masking is not perfect and its gate says so: one document carries no mask at
all, and an own-firm-name token survives in 0.17 per cent of documents.
% src: results/anon/mask_stats.json (docs_with_any_mask_pct 99.99911 = 112,527 of 112,528; leak_own_name_token_pct 0.174)
Neither
executed arm reproduces its June original bit-identically: the prompted arm
matches 111{,}458 of 117{,}407 predictions exactly and the fine-tuned one 0 of
333{,}192, and both were retained through the pre-registered deviation route for
GPU nondeterminism. What the share depends on is the within-pair contrast, and
each pair's control and masked legs were produced in one environment.
The identity share is one minus the ratio of the masked to the unmasked
combination increment. Volatility-unit QLIKE, log-space combiner fitted
on validation and frozen on test, day-clustered DM, Holm within each
arm's declared six-cell family, seed 2026. Intervals are day-block
bootstrap 95\% percentile intervals, $B = 2{,}000$ draws, resampling
unit the effective trading day, combiner weights frozen inside every
draw. The estimate is scoped to the 8-K channel: the long-form masking
gate ran on 31{,}601 documents but that channel closed with zero
executed arms. In panels (a) and (b) ``unmask.''\ and ``masked'' are
relative QLIKE improvements in per cent over the reference named in the
panel heading; ``bound'' and ``swap'' appear in panel~(b) only, as the
per-cell reference-interval bound and the matched-firm swap retention, and
panel~(c) carries the bound as a median over the six cells and the swap as an arm median. Both are triangulating
instruments measured against the increment over the \emph{recalibrated
HAR}, not over panel~(b)'s own reference.
The DM and Holm columns test the \emph{masked} arm's
increment over the reference named in the panel heading. The unmasked
increments are the seed-2026 values of this arm; Table~\ref{tab:app-grid-ed} tests the same cells on the seed-ensemble basis and prints different figures ($+2.572/+2.416/+1.578$ against $+2.135/+2.104/+0.920$ here), so the two must not be read across. % src: results/tables/m1_clustered.csv (event_driven C2_finbert_s1 rel_impr_pct 2.135411/2.103684/0.919653); results/tables/control_intersection_ensemble.csv (the seed-ensemble counterparts)
Panel (c) collects the
aggregates, which are seven different quantities and not seven readings
of one.}
\label{tab:app-identity}
\singlespacing\small
\setlength{\tabcolsep}{4pt}
% src: results/tables/anon_arm.md (increments, DM, Holm, share, interval bound, swap retention); results/tables/anon_share_ci.md (all 95% intervals); results/tables/matched_firm_swap.md (grid-wide retention medians)

\emph{(a) Against the recalibrated HAR reference}\\[0.4ex]
\begin{tabular}{@{}lrrrrrrl@{}}
\toprule
Arm & $h$ & unmask.\ & masked & DM & Holm $p$ & share & 95\% CI \\
\midrule
C6 prompted   &  5 & $+1.214$ & $+0.597$ & $-3.91$ & .0006 & $+0.51$ & $[+0.40, +0.62]$ \\
C6 prompted   & 10 & $+0.998$ & $+0.438$ & $-2.39$ & .0677 & $+0.56$ & $[+0.43, +0.69]$ \\
C6 prompted   & 20 & $+0.658$ & $+0.192$ & $-1.11$ & .4810 & $+0.71$ & $[+0.57, +0.84]$ \\
C2 FinBERT    &  5 & $+2.135$ & $+2.083$ & $-3.81$ & .0005 & $+0.02$ & $[-0.14, +0.19]$ \\
C2 FinBERT    & 10 & $+2.104$ & $+1.031$ & $-5.90$ & $<$.0001 & $+0.51$ & $[+0.45, +0.57]$ \\
C2 FinBERT    & 20 & $+0.920$ & $+1.277$ & $-1.68$ & .0942 & $-0.39$ & $[-4.94, +0.13]$ \\
B2 TF--IDF    & \multicolumn{7}{l}{arm exited at reproduction gate G1; no share and no test} \\
\bottomrule
\end{tabular}

\vspace{1.2ex}
\emph{(b) Against the firm-identity reference, with the two triangulating instruments}\\[0.4ex]
\begin{tabular}{@{}lrrrrrrrr@{}}
\toprule
Arm & $h$ & unmask.\ & masked & DM & Holm $p$ & share & bound & swap \\
\midrule
C6 prompted   &  5 & $+0.448$ & $+0.334$ & $-3.43$ & .0032 & $+0.25$ & $+0.63$ & $+0.42$ \\
C6 prompted   & 10 & $+0.253$ & $+0.224$ & $-2.22$ & .0793 & $+0.12$ & $+0.75$ & $+0.18$ \\
C6 prompted   & 20 & $+0.196$ & $+0.080$ & $-1.17$ & .4810 & $+0.59$ & $+0.70$ & $+0.10$ \\
C2 FinBERT    &  5 & $-0.314$ & $+0.129$ & $-4.46$ & $<$.0001 & n/a & $+1.15$ & $+0.52$ \\
C2 FinBERT    & 10 & $-1.223$ & $-1.975$ & $+5.79$ & $<$.0001 & n/a & $+1.58$ & $+0.41$ \\
C2 FinBERT    & 20 & $+0.658$ & $-0.349$ & $+3.07$ & .0044 & $+1.53$ & $+0.28$ & $-2.69$ \\
\bottomrule
\end{tabular}

\vspace{1.2ex}
\emph{(c) Aggregates}\\[0.4ex]
\begin{tabular}{@{}lrl@{}}
\toprule
Statistic & Value & 95\% CI \\
\midrule
Pooled median share, six vs-HAR cells (\textbf{headline}) & $+0.51$ & $[+0.44, +0.56]$ \\
\quad C6 arm median                                       & $+0.56$ & $[+0.47, +0.69]$ \\
\quad C2 arm median                                       & $+0.02$ & $[-0.13, +0.26]$ \\
Median over the four defined firm-identity cells          & $+0.42$ & --- \\ % src: results/tables/anon_arm.csv (share_anon_firm: 0.2535, 0.1155, 0.5929, 1.5294; the two c2 cells at h=5 and h=10 are NaN because their unmasked increment is at or below zero)
Reference-interval bound, median over six cells            & $+0.72$ & --- \\
Matched-swap retention, arm median                        & $+0.29$ & --- \\
\quad grid-wide median over the 38 genuine cells          & $+0.31$ & --- \\
\bottomrule
\end{tabular}
\end{table}


\begin{figure}[tbp]
\centering
\includegraphics[width=\textwidth,height=0.83\textheight,keepaspectratio]{figures/AR3_matched_row_cascade__a.pdf}
% caption numbers src: results/tables/matched_row_cascade.csv (69 cells x 2 row sets x 2 rungs); results/tables/firm_identity_ensemble.csv (used only to verify that the matched arm's firm rung is set-identical to the committed one)
\caption[Separating the reference from the sample]{
Separating the reference
from the sample. The primary rung is scored on the reference-plus-text join
while the firm-identity rung is scored on the smaller five-price-model
intersection, so the reported fall from 38 survivors to 8 confounds a change
of reference with a change of sample. Here both rungs are re-scored on both
row sets and only the labelled thing changes. 
Panel~(a) reads either way round
and reaches the same eight cells. Holding the reference fixed and moving to
the matched rows costs 8 cells, after which the identity term costs a further
22 net. On the committed rows the identity term alone takes 38 to 11 and
the row change then costs 3. % src: results/tables/matched_row_cascade.csv (committed rows 38 and 11, matched rows 30 and 8; 8 lost and 0 gained to the row set; 26 lost and 4 gained to the control)
}
\label{fig:app-matched-rows}
\end{figure}

\begin{figure}[tbp]
\centering
\includegraphics[width=\textwidth,height=0.83\textheight,keepaspectratio]{figures/AR3_matched_row_cascade__b.pdf}
% caption numbers src: results/tables/matched_row_cascade.csv (69 cells x 2 row sets x 2 rungs); results/tables/firm_identity_ensemble.csv (used only to verify that the matched arm's firm rung is set-identical to the committed one)
\caption[Separating the reference from the sample (b, c)]{Continued from Figure~\ref{fig:app-matched-rows}. Separating the reference
from the sample. The primary rung is scored on the reference-plus-text join
while the firm-identity rung is scored on the smaller five-price-model
intersection, so the reported fall from 38 survivors to 8 confounds a change
of reference with a change of sample. Here both rungs are re-scored on both
row sets and only the labelled thing changes. 
Panel~(b) draws all 69 cells on the matched row set, the increment over
recalibrated HAR-RV against the increment over HAR-RV plus firm identity, with
the diagonal marking no change. Five cells fall off the bottom, the furthest at
$-11.4$ per cent. 
Panel~(c) prices the matched support: 90 to 95 per cent of test rows,
92.0 per cent pooled, and 95.4 to 100 per cent of trading days retained. % src: results/tables/matched_row_cascade.csv (n_test and n_days of the matched arm against the committed arm)
Basis: Holm survivors of 69 on the seed-ensemble basis, day-clustered
statistics, validation-fitted weights frozen on test; survivorship includes
the placebo gate, which binds only on the matched rows, where it removes 3 of
33. It also governs panel~(b)'s four colour classes: each cell carries its Holm
verdict under each reference, the 4/26/4/35 split the panel's own key names. % src: results/tables/matched_row_cascade.csv (matched arm: adds_primary and adds_firm both true 4, primary only 26, firm only 4, neither 35)
}
\label{fig:app-matched-rows-b}
\end{figure}

Figures~\ref{fig:app-matched-rows} and~\ref{fig:app-matched-rows-b} answer an objection that is easy to raise
and, without this decomposition, impossible to settle: that the ladder's
collapse is an accounting artefact of scoring successive rungs on successively
smaller samples. It is not. On one row set the reference alone takes 30
survivors to 8; on the other it takes 38 to 11; the sample change costs 8
cells on the first rung and 3 on the second. The control does most of the
work, and both paths through the diagram arrive at the same eight cells.

Two checks make the two arms comparable rather than two implementations of one
idea. The same code path run on the committed rows returns the committed 38 of
69, and the matched arm's identity rung reproduces the committed
firm-identity rung cell for cell: the same eight cells, not merely the same
count. % src: results/tables/matched_row_cascade.csv (committed_rows arm, adds_primary = 38); exact set comparison against results/tables/firm_identity_ensemble.csv (adds_holm)
Without the second of those checks the diagram would be an argument about two
tables; with it, it is an argument about one.

What the figure does not licence is a claim that the row change is
inconsequential. It costs eight cells on the primary rung, six long-form and
two event-driven, and three on the identity rung, all three event-driven. The
five-price-model intersection thins by 5.0 to 10.2 per cent of test rows on
the long-form channel and by 5.0 to 9.8 per cent on the event-driven one, so
neither panel bears the loss alone. % src: results/tables/matched_row_cascade.csv (adds_primary and adds_firm on the committed_rows arm against the matched arm; n_test by disc and h on the two arms)
A reader comparing a long-form count in Chapter~\ref{ch:results} against one
in Chapter~\ref{ch:validation} should check which row basis each sits on.
The second step is also not
monotone: 26 cells fall and 4 rise, so 22 is a net, and a figure reporting
only the net would hide four cells that the identity term \emph{sharpens}. The
figure shows neither the maximal-pool rung, nor effect sizes on the committed
rows, nor any cell outside the 69-cell grid. Finally, the firm-mean term is
itself estimated on validation and is defined for 91.5 to 91.9 per cent of
test observations, firms without validation filings receiving the global
validation mean. % src: results/tables/matched_row_cascade.csv (firm_val_coverage_test_obs on the matched arm)
That weakens the identity term rather than strengthening
it and is therefore conservative for the text side.

\clearpage

\FloatBarrier
\section{Cross-Family Probes and the Forecaster-Health Screen}
\label{app:res-crossfamily}

Table~\ref{tab:app-crossfamily} is the per-probe record behind the
cross-family replication of Section~\ref{sec:res-residual} and the
health screen of Section~\ref{sec:val-instruments}. Panel (a) is the
screen: a run passes when, at every horizon, its standalone
variance-unit QLIKE stays below 4 \emph{and} the share of its forecasts
taking the single most common value after rounding to two decimals stays
below 60 per cent. Neither term reads the contrast under test --- one scores the run's own
accuracy, the other its dispersion --- which is why the screen can
discard the block's largest positive increment (Gemma-3-27B at
$+2.97\%$) without that being selection on the outcome. The verdict
attaches to a whole run, so a marker below the modal-share line at one
horizon can still belong to a failing run. The ranges printed are per
run for exactly that reason. Counted per family, four of the five probes
fail.

Panel (b) gives the increments and, critically, names the multiplicity
convention each probe's committed $p$-values carry: Holm within three
per reference for one family, Holm within six for another, raw $p$ only
for two more. A single ``Holm $p$'' column across these rows would not
be a well-defined quantity, which is why the convention is a column
rather than a footnote. The panel is complete rather than 8-K-scoped.
Yi-1.5-34B is the one probe family that also ran on long-form, and its
long-form row is carried even though it is the worst result in the block:
$-0.64/-2.71/-9.86$ per cent against the recalibrated HAR, worsening
monotonically with horizon. It is kept because the run truncates at 4{,}096
tokens, about 14 per cent of the median long-form filing
(Section~\ref{sec:data-corpus}), and a reader is entitled to see the arm
fail where it was starved as well as where it was fed. The rows are
flagged \textsc{4k-truncated} in the evidence file and the flag is what
makes them interpretable. % src: results/tables/crossfamily_llama70.md (long_form yi\_34b rows, flagged 4K-TRUNCATED)
Llama-3.1-70B and Phi-4-14B have no long-form row at all: the 70B's was
not run, on GPU budget, and 8-K is the channel the replication was
declared on. % src: results/tables/crossfamily_llama70.md disclosures ("long\_form was NOT run for llama70")
The one instrument-healthy probe outside the
primary family is also the largest of the five, at 70B against the
anchor's 32B and the failures' 14B to 34B. That probe, Llama-3.1-70B,
reproduces the sign of the 8-K residual
at all three horizons with a larger point estimate and weaker
significance. Its forecasts are 2.7 to 3.0 times more dispersed than
the anchor's on identical rows, which inflates the loss-differential
variance rather than shrinking the effect. % src: results/tables/crossfamily_llama70.csv (pred sd vs anchor)
That is a directional
replication and a bound. It is not an independent confirmation, and
Table~\ref{tab:app-crossfamily} is laid out so that it cannot be read as
one.

Panel (c) closes the objection that the replication could itself be firm-name
recall rather than reading. The 70B is put through the \emph{same}
zero-content prompt the anchor's contamination arm uses (form type and
item codes, filing date and ticker, the line stating that no filing text
is provided, and nothing else) on the identical event-driven rows, in the identical
quantisation and decoding stack. The probe and the full-text run differ
in the document and in nothing else. % src: results/tables/crossfamily_llama70_probe.md Disclosures block
The probe carries a positive increment over the recalibrated HAR at all
three horizons, which is the expected signature of identity and era memory
in a model that has seen the ticker before. The increment is small, and it
is the ratio that matters. At the two horizons where the denominator is
well identified the probe reproduces \textbf{5 and 7 per cent} of the
70B's own full-text edge, and 0.1 and 0.5 per cent of its edge over the
firm-identity reference. The $h{=}20$ share of 31 per cent is printed
struck from use rather than omitted: its denominator, a full-text
increment of $+0.69\%$ that is not itself significant, falls below the
stable-denominator rule the probe tables apply. Small denominators
inflate shares mechanically. Running the argument the other way round is
the stronger form, and panel~(d) gives it. With the 70B's own zero-content
forecast placed \emph{inside} the reference, the full text still adds in
two of three horizons after Holm correction within that three-cell block,
retaining 94 to 98 per cent of the uncontrolled increment. The
replication survives its own identity control; what it does not do is
become an independent confirmation by surviving it.

\begin{table}[htbp]
\centering
\caption[Cross-family probes, health screen and increments]{The five
cross-family probes and the primary-family anchor. Panel (a): the
registered health screen (variance-unit QLIKE below 4 and modal share of
the two-decimal forecast below 60\% at every horizon); ranges are per
run, and the verdict attaches to a (family, channel) run. The
Gemma-3-27B pilot row records the same family on the registered
2{,}000-filing validation pilot, which it cleared on both terms before
breaching the modal-share ceiling on the full 8-K panel drawn in the row
above it. Panel (b):
volatility-unit relative QLIKE improvements against two references, with
the multiplicity convention each family's committed $p$-values carry;
the three conventions are never pooled. Panel (a) uses the variance-unit
QLIKE convention and panels (b)--(d) the volatility-unit one; no number
moves between them. Panel~(b)'s 8-K rows are scored
on the full event-driven test panel (25{,}109/25{,}001/24{,}732
observations over 996/991/981 trading days), a wider support than the
merged grid of Tables~\ref{tab:app-grid-ed} and~\ref{tab:app-grid-lf}; its
10-K/Q rows are scored on $7{,}951/7{,}933/7{,}902$ observations over
$809/803/794$ days. % src: results/tables/crossfamily_llama70.csv (n_test and n_days by disc)
The Qwen3-32B rows read $+0.45/+0.25/+0.20$ here against
$+0.52/+0.24/+0.21$ there: the same arm and reference on different rows,
not a discrepancy. The Qwen3-32B rows
are the primary family and are not probes.
Panel (c) is the replication family's own zero-content date{+}ticker
probe: the identical prompt with the filing text removed, the identical
rows, quantisation and decoding stack, single seed 2026. A share is the
probe's increment divided by the full-text increment in the same cell,
and is quotable only where the denominator is well identified (full-text
increment at least 1\% \emph{and} raw-significant). The $h{=}20$ share is
printed in parentheses because it is not, and it is shown solely so that
the row is not silently dropped. Panel (d) reverses the direction: the
probe forecast enters the reference as $[1, \log f_{\mathrm{HAR}},
\log f_{\mathrm{probe}}]$ and the full text is tested on top of it, with
Holm inside that three-cell block; ``retained'' is the residual increment
as a share of the uncontrolled full-text-versus-HAR increment in the same
cell. Panels (c) and (d) are descriptive readouts and open no new
pre-registered branch.}
\label{tab:app-crossfamily}
\singlespacing\footnotesize
\setlength{\tabcolsep}{3.5pt}
% src: results/tables/crossfamily_llama70.csv, crossfamily_mistral24.csv, crossfamily_gemma27.csv, crossfamily_llm.csv, crossfamily_llama70_ens.csv; consolidated in TECHNICAL_SUPPLEMENT.md Table S11

\emph{(a) Instrument, scope and the registered health screen}\\[0.4ex]
\begin{tabular}{@{}llrrl@{}}
\toprule
Model (precision) & Channel & var-QLIKE & modal \% & Health verdict \\
\midrule
Qwen3-32B bf16 \emph{(anchor)} & 8-K    & 1.18--1.32 & 49.2--50.3 & passes \\
                               & 10-K/Q & 0.93--1.23 & 45.5--76.8 & fails modal share \\
\quad \emph{by horizon, $h=5/10/20$} & 10-K/Q & 1.23/1.09/0.93 & 64.4/76.8/45.5 & \emph{breached at 5 and 10} \\ % src: results/tables/crossfamily_standalone.csv (long_form qwen3_32b rows: qlike_var 1.225874/1.094998/0.925773, mode_share_pct 64.406993/76.818354/45.545432)
Llama-3.1-70B AWQ-INT4         & 8-K    & 1.12--2.10 & 40.4--51.2 & passes \\
Yi-1.5-34B bf16                & 8-K    & 7.60--8.19 & 55.9--73.6 & fails both terms \\
                               & 10-K/Q & 5.25--6.20 & 26.1--56.6 & fails QLIKE \\
Phi-4-14B bf16                 & 8-K    & 0.53--0.68 & 61.3--75.1 & fails modal share \\
Mistral-Small-24B bf16         & 8-K    & 0.42--0.69 & 70.6--88.6 & fails modal share \\
Gemma-3-27B bf16               & 8-K    & 0.38--0.90 & 71.3--71.5 & fails modal share \\
\quad \emph{pilot, 2{,}000 filings} & 8-K & 3.66 & 45.2 & \emph{passed the pilot} \\
\bottomrule
\end{tabular}

\vspace{1.2ex}
\emph{(b) Increments (\%, volatility-unit QLIKE, $h = 5/10/20$) and the multiplicity convention}\\[0.4ex]
\begin{tabular}{@{}llrrl@{}}
\toprule
Model & Channel & vs.\ recal.\ HAR & vs.\ firm identity & Multiplicity \\
\midrule
Qwen3-32B           & 8-K    & $+1.21/+1.00/+0.66$ & $+0.45/+0.25/+0.20$ & Holm(6) inherited \\
                    & 10-K/Q & $+1.79/+2.25/+0.27$ & $-0.14/-0.17/+0.02$ & \\
Llama-70B (single)  & 8-K    & $+1.39/+1.17/+0.70$ & $+0.83/+0.64/+0.39$ & Holm(6) \\
Llama-70B (ens.\ 3) & 8-K    & $+1.39/+1.16/+0.69$ & $+0.84/+0.64/+0.38$ & Holm(6), ensemble \\
Yi-1.5-34B          & 8-K    & $+0.37/+0.07/-0.62$ & $+0.22/+0.07/-0.07$ & raw $p$ only \\
                    & 10-K/Q & $-0.64/-2.71/-9.86$ & $-0.21/-1.61/-6.05$ & raw $p$ only; 4K-trunc. \\
Phi-4-14B           & 8-K    & $+0.38/+0.18/-0.12$ & $+0.20/+0.14/+0.00$ & raw $p$ only \\
Mistral-24B (ens.\ 3) & 8-K  & $+0.31/+0.15/+0.13$ & $-0.10/-0.45/-0.38$ & own-family Holm(6) \\
Gemma-3-27B (ens.\ 3) & 8-K  & $+1.84/+2.59/+2.97$ & $+1.08/+1.50/+1.26$ & Holm(3) per reference \\
\bottomrule
\end{tabular}

\vspace{1.2ex}
\emph{(c) The replication family's own zero-content date{+}ticker probe, 8-K panel}\\[0.4ex]
% src: results/tables/crossfamily_llama70_probe.csv (probe_rel_har, probe_rel_firm, share_har_pct, share_firm_pct, denominator_well_identified)
\begin{tabular}{@{}lrrrrrl@{}}
\toprule
& \multicolumn{2}{c}{Probe increment (\%)} & \multicolumn{2}{c}{70B full text (\%)} & \multicolumn{2}{c}{Probe as a share} \\
\cmidrule(lr){2-3}\cmidrule(lr){4-5}\cmidrule(lr){6-7}
$h$ & vs.\ HAR & vs.\ ident. & vs.\ HAR & vs.\ ident. & vs.\ HAR & denominator \\
\midrule
5  & $+0.073$ & $+0.001$ & $+1.395$ & $+0.835$ & $5.2\%$    & well identified \\
10 & $+0.085$ & $+0.003$ & $+1.163$ & $+0.637$ & $7.3\%$    & well identified \\
20 & $+0.218$ & $+0.022$ & $+0.692$ & $+0.383$ & \emph{(31.5\%)} & \textbf{not}: do not quote \\
\bottomrule
\end{tabular}

\vspace{1.2ex}
\emph{(d) Full text beyond the same model's own identity forecast, 8-K panel}\\[0.4ex]
% src: results/tables/crossfamily_llama70_probe.csv (beyond_identity block: rel_pct, dm_clu, p_holm) and .md section 2 (retained share)
\begin{tabular}{@{}lrrrrl@{}}
\toprule
$h$ & rel.\ \% & DM & Holm(3) $p$ & retained & verdict \\
\midrule
5  & $+1.374$ & $-3.17$ & .0046 & 98\% & adds \\
10 & $+1.136$ & $-2.27$ & .0465 & 98\% & adds \\
20 & $+0.653$ & $-0.75$ & .4512 & 94\% & not significant \\
\bottomrule
\end{tabular}
\end{table}


\begin{figure}[tbp]
\centering
\includegraphics[width=\textwidth,height=0.83\textheight,keepaspectratio]{figures/FP1_prompted_family_panel__a.pdf}
% caption numbers src: results/tables/crossfamily_llama70.csv (panel a health ranges; the Qwen3, Yi, Phi-4 and single-seed Llama-70B increments); results/tables/crossfamily_llama70_ens.csv (the three-seed Llama-70B row); results/tables/crossfamily_gemma27.csv and results/tables/crossfamily_mistral24.csv (the Gemma-3-27B and Mistral-Small-24B rows with their own Holm columns); results/tables/crossfamily_llama70_probe.csv (panels c and d)
% Holm-block itemisation src: the holm_family column of each source. crossfamily_gemma27.csv: qwen3_32b "Holm(6): 3 horizons x {HAR, HAR+firmID}" and gemma27_ens3 "Holm(3) per reference (prereg-rfa v1.3 B2); p_*_holm6 = pooled Holm(6), B1-comparable info only" -- the figure rings gemma27_ens3 from the committed Holm(3) column p_har_holm (0.00654, 0.01616), NOT p_har_holm6. crossfamily_llama70_ens.csv: llama70_awq_ens3 "ens Holm(6) ... (prereg B0)". crossfamily_mistral24.csv: mistral24_ens3 "Holm(6) ... (prereg B1, own family)". crossfamily_llama70.csv: p_har_holm and p_firm_holm are empty for yi_34b and phi4_14b, so those two rows carry raw p only and are tagged "(raw p)" in panel (b).
% (c)/(d) figure gloss src: results/tables/crossfamily_llama70_probe.csv -- probe_share.share_har_pct (5.23 / 7.31 / 31.48 per cent of fulltext_ens_rel_har) is the figure printed in panel (c); 100 x beyond_identity.rel_pct / fulltext_anchor_committed.rel_har for llama70_awq_ens3 (1.374/1.395, 1.136/1.163, 0.653/0.692 = 98 / 98 / 94 per cent) is the figure printed in panel (d)
\caption[The prompted arms, family by family]{
The prompted arms, family by
family. 
Panel~(a) is the registered forecaster-health screen. A run passes
only if, at \emph{every} horizon, its standalone variance-unit QLIKE stays
below 4 \emph{and} the share of its forecasts taking the single commonest
value after rounding to two decimals stays below 60 per cent. Neither term reads
the contrast under test (one scores the run's own accuracy, the other its
dispersion),
which is why the screen may discard a large positive increment without that
being selection on the outcome. The verdict attaches to a whole
family-and-channel run. Four of the five probe
families fail, and only the Llama-3.1-70B run joins the Qwen3-32B 8-K anchor
in passing. % src: results/tables/crossfamily_llama70.csv and crossfamily_gemma27.csv (qlike_var and mode_share_pct by family, channel and horizon against the two ceilings)
Panel~(a) runs no test, so it carries no family.
The identity audit of panels~(c) and~(d) is
Figure~\ref{fig:app-prompted-families-c}.
}
\label{fig:app-prompted-families}
\end{figure}

\begin{figure}[tbp]
\centering
\includegraphics[width=\textwidth,height=0.83\textheight,keepaspectratio]{figures/FP1_prompted_family_panel__b.pdf}
% caption numbers src: results/tables/crossfamily_llama70.csv (panel a health ranges; the Qwen3, Yi, Phi-4 and single-seed Llama-70B increments); results/tables/crossfamily_llama70_ens.csv (the three-seed Llama-70B row); results/tables/crossfamily_gemma27.csv and results/tables/crossfamily_mistral24.csv (the Gemma-3-27B and Mistral-Small-24B rows with their own Holm columns); results/tables/crossfamily_llama70_probe.csv (panels c and d)
% Holm-block itemisation src: the holm_family column of each source. crossfamily_gemma27.csv: qwen3_32b "Holm(6): 3 horizons x {HAR, HAR+firmID}" and gemma27_ens3 "Holm(3) per reference (prereg-rfa v1.3 B2); p_*_holm6 = pooled Holm(6), B1-comparable info only" -- the figure rings gemma27_ens3 from the committed Holm(3) column p_har_holm (0.00654, 0.01616), NOT p_har_holm6. crossfamily_llama70_ens.csv: llama70_awq_ens3 "ens Holm(6) ... (prereg B0)". crossfamily_mistral24.csv: mistral24_ens3 "Holm(6) ... (prereg B1, own family)". crossfamily_llama70.csv: p_har_holm and p_firm_holm are empty for yi_34b and phi4_14b, so those two rows carry raw p only and are tagged "(raw p)" in panel (b).
% (c)/(d) figure gloss src: results/tables/crossfamily_llama70_probe.csv -- probe_share.share_har_pct (5.23 / 7.31 / 31.48 per cent of fulltext_ens_rel_har) is the figure printed in panel (c); 100 x beyond_identity.rel_pct / fulltext_anchor_committed.rel_har for llama70_awq_ens3 (1.374/1.395, 1.136/1.163, 0.653/0.692 = 98 / 98 / 94 per cent) is the figure printed in panel (d)
\caption[The prompted arms, family by family (b)]{Continued from Figure~\ref{fig:app-prompted-families}. The prompted arms, family by
family. 
Panel~(b) gives each family's relative QLIKE improvement over the recalibrated
HAR and over the firm-identity-augmented reference on the 8-K channel, scored
on $25{,}109/25{,}001/24{,}732$ observations
over $996/991/981$ trading days. A marker is filled at raw $p<.05$ and ringed
where that family's own committed Holm block also falls below .05; blocks are
six cells, three per reference for Gemma-3-27B, none for the two raw-$p$ rows.
Panels~(c) and~(d) audit identity inside the same replication family and are
Figure~\ref{fig:app-prompted-families-c}. Basis:
volatility-unit QLIKE, day-clustered inference; the anchor's firm-referenced
triple is $+0.45/+0.25/+0.20$ per cent and the replication's three-seed
ensemble $+0.84/+0.64/+0.38$. 
}
\label{fig:app-prompted-families-b}
\end{figure}

\begin{figure}[tbp]
\centering
\includegraphics[width=\textwidth,height=0.83\textheight,keepaspectratio]{figures/FP1_prompted_family_panel__c.pdf}
% caption numbers src: results/tables/crossfamily_llama70.csv (panel a health ranges; the Qwen3, Yi, Phi-4 and single-seed Llama-70B increments); results/tables/crossfamily_llama70_ens.csv (the three-seed Llama-70B row); results/tables/crossfamily_gemma27.csv and results/tables/crossfamily_mistral24.csv (the Gemma-3-27B and Mistral-Small-24B rows with their own Holm columns); results/tables/crossfamily_llama70_probe.csv (panels c and d)
% Holm-block itemisation src: the holm_family column of each source. crossfamily_gemma27.csv: qwen3_32b "Holm(6): 3 horizons x {HAR, HAR+firmID}" and gemma27_ens3 "Holm(3) per reference (prereg-rfa v1.3 B2); p_*_holm6 = pooled Holm(6), B1-comparable info only" -- the figure rings gemma27_ens3 from the committed Holm(3) column p_har_holm (0.00654, 0.01616), NOT p_har_holm6. crossfamily_llama70_ens.csv: llama70_awq_ens3 "ens Holm(6) ... (prereg B0)". crossfamily_mistral24.csv: mistral24_ens3 "Holm(6) ... (prereg B1, own family)". crossfamily_llama70.csv: p_har_holm and p_firm_holm are empty for yi_34b and phi4_14b, so those two rows carry raw p only and are tagged "(raw p)" in panel (b).
% (c)/(d) figure gloss src: results/tables/crossfamily_llama70_probe.csv -- probe_share.share_har_pct (5.23 / 7.31 / 31.48 per cent of fulltext_ens_rel_har) is the figure printed in panel (c); 100 x beyond_identity.rel_pct / fulltext_anchor_committed.rel_har for llama70_awq_ens3 (1.374/1.395, 1.136/1.163, 0.653/0.692 = 98 / 98 / 94 per cent) is the figure printed in panel (d)
\caption[The prompted arms, family by family (c, d)]{Continued from Figure~\ref{fig:app-prompted-families}. The prompted arms, family by family.
Panels~(c) and~(d) audit identity inside the replication family: figures by the
bars are the probe's share of the full-text bar in~(c), and of the uncontrolled
increment in~(d). Basis: volatility-unit QLIKE, day-clustered inference; the
anchor's firm-referenced triple is $+0.45/+0.25/+0.20$ per cent and the
replication's three-seed ensemble $+0.84/+0.64/+0.38$. Panels~(c) and~(d) are
corrected inside the replication script's own six-cell block, which carries no
pre-registered branch and is descriptive; their denominator is those six cells,
not the 69 of the primary grid.
}
\label{fig:app-prompted-families-c}
\end{figure}

The health screen of Figure~\ref{fig:app-prompted-families} is the part of this appendix most likely to be read as convenient. Both ceilings are computed from the
forecast distribution alone: a run fails if its standalone loss is
catastrophic, or if it collapses onto a single modal value (symptoms of a
model that is not forecasting so much as emitting a constant). Neither
test can see the label or the increment, so a failing run cannot have been
discarded because its result was inconvenient. The screen's honesty is
demonstrated by what it discards. The largest single point estimate anywhere
in the cross-family block is the Gemma-3-27B run's $+2.97$ per cent over the
recalibrated HAR at $h=20$, and that run fails the modal-share ceiling at all
three horizons, at roughly 71 per cent. % src: results/tables/crossfamily_gemma27.csv (gemma27_ens3 rel_har 1.838 / 2.594 / 2.974 with mode_share_pct 71.4 / 71.2 / 71.2 against the 60 per cent ceiling)
Had the screen been drawn to flatter the study's positive finding, it would
have been drawn to keep that number rather than to lose it.

Panel~(b) is the reason the replication is reported as a replication rather
than as a second confirmation. Of six families run on the 8-K channel, one
anchor and one replication pass the health screen, and of the four that fail,
one (Gemma-3-27B) carries point estimates larger than the anchor's at
every horizon. The other three sit below it at every horizon on both
references. % src: results/tables/crossfamily_llama70.csv and crossfamily_gemma27.csv and crossfamily_mistral24.csv (event_driven rel_har and rel_firm: anchor qwen3_32b +1.214/+0.998/+0.658 and +0.448/+0.253/+0.196; gemma27_ens3 +1.838/+2.594/+2.974 and +1.076/+1.497/+1.265; yi_34b, phi4_14b and mistral24_ens3 all strictly below the anchor at each horizon)
The report's position,
carried in Chapter~\ref{ch:validation}, is that the cross-family evidence
bounds rather than establishes: it shows the anchor's residual is not unique
to one model lineage, and it does not show that prompted models generally
produce it. The panel also refuses to pool Holm conventions across families:
the anchor and the 70B are corrected within six tests each, the Mistral
family within its own six, Gemma within three per reference, and the Yi and
Phi-4 families carry raw $p$ only. Collapsing incompatible
correction families into one significance column is precisely the convenience
this audit exists to detect.

Panel~(b), Figure~\ref{fig:app-prompted-families-b}, draws no long-form cells
for any family, the anchor included, so it is an 8-K comparison throughout;
the health screen of panel~(a) does carry them, for the anchor and for Yi alike. The Yi family's long-form cells
($-0.64/-2.71/-9.86$ per cent) are excluded from panel~(b) because that
model's 4,096-token context truncates the long-form excerpts, and they are
reported here rather than silently dropped. % src: results/tables/crossfamily_llama70.csv (yi_34b long_form rel_har, flag 4K-TRUNCATED)
It shows no single-seed restatement of the Llama-70B arm, whose ensemble row
is the one drawn, and no economic quantity of any kind: every percentage in
the figure is a QLIKE improvement and none may be set beside a return or a
variance-unit percentage.

\clearpage

\FloatBarrier
\section{Economic Adjudication}
\label{app:res-econ}

Table~\ref{tab:app-sharpe} gives all 18 cells of the portfolio
adjudication, and Table~\ref{tab:app-econ-summary} the downside-risk and
utility tallies that accompany them. None of these quantities is a
QLIKE quantity and none may be placed on an axis with a relative-QLIKE
percentage: panel Sharpe differences are differences of two annualised
ratios, violation rates are proportions, and fees are basis points a
year.

The portfolio table repays a careful reading in three places. The
prompted arm, the identity rung's only three-horizon 8-K survivor, is priced at
zero of its six cells, with every interval covering zero. The single
cell whose interval excludes zero, event-driven TF--IDF ridge at
$h{=}10$ at $+0.026$, is about what fifteen nominal-95\% intervals
produce by chance. It is a configuration the ladder had already
rejected: it fails the primary rung's Holm correction and its increment
over the maximal price pool is negative, so it cannot be attached to the
surviving residual. And the three long-form $h{=}20$ cells have no
interval at all, because 39 holding periods cannot support a bootstrap
over 20-period blocks. Their reference book also loses money at a
Sharpe of $-0.104$, so a positive difference there improves a losing
position rather than evidencing a profitable signal. The book is
long-only and frictionless throughout, so every portfolio figure --- the
Sharpe panel and the utility-based fee --- is an upper bound on what an
implementable strategy would deliver. The value-at-risk panel carries no book
at all: its figures are violation-rate counts and tick-loss tallies computed
cell by cell, and the caveat does not reach them.

\begin{table}[htbp]
\centering
\caption[Portfolio economics: all 18 cells]{Portfolio adjudication, all
18 cells. On each rebalancing day every firm with a live filing signal
is held long at a weight inversely proportional to its forecast
variance. The book is held for non-overlapping $h$-day periods with
combination weights fitted on validation and frozen on test. $\Delta$Sh
is the change in annualised Sharpe ratio \cite{sharpe1966mutual} when
the text-augmented forecast $f_U$ replaces the recalibrated reference
$f_R$. Intervals are day-block bootstrap 95\% intervals over blocks of
$h$ periods. Grid median $\Delta$Sh $= +0.003$; intervals cover zero in
14 of the 15 cells that have one.}
\label{tab:app-sharpe}
\singlespacing\small
\setlength{\tabcolsep}{3pt}
% src: results/tables/portfolio_econ.md and portfolio_econ.csv (all 18 rows verbatim)
\begin{tabular}{@{}llrrrrrrrl@{}}
\toprule
Channel & Model & $h$ & periods & names & Sharpe $f_R$ & Sharpe $f_U$ & $\Delta$Sh & 95\% CI & sig. \\
\midrule
Event-driven & B2 TF--IDF   &  5 & 200 & 101.8 & $+0.329$ & $+0.336$ & $+0.006$ & $[-0.022,+0.038]$ & no \\
Event-driven & B2 TF--IDF   & 10 & 100 & 174.6 & $+0.125$ & $+0.151$ & $+0.026$ & $[+0.004,+0.047]$ & \textbf{yes} \\
Event-driven & B2 TF--IDF   & 20 &  50 & 282.8 & $+0.340$ & $+0.352$ & $+0.012$ & $[-0.012,+0.045]$ & no \\
Event-driven & C2 FinBERT   &  5 & 200 & 101.8 & $+0.329$ & $+0.331$ & $+0.002$ & $[-0.028,+0.031]$ & no \\
Event-driven & C2 FinBERT   & 10 & 100 & 174.6 & $+0.125$ & $+0.121$ & $-0.004$ & $[-0.028,+0.021]$ & no \\
Event-driven & C2 FinBERT   & 20 &  50 & 282.8 & $+0.340$ & $+0.335$ & $-0.005$ & $[-0.043,+0.060]$ & no \\
Event-driven & C6 prompted  &  5 & 200 & 101.8 & $+0.329$ & $+0.324$ & $-0.005$ & $[-0.024,+0.012]$ & no \\
Event-driven & C6 prompted  & 10 & 100 & 174.6 & $+0.125$ & $+0.118$ & $-0.007$ & $[-0.020,+0.006]$ & no \\
Event-driven & C6 prompted  & 20 &  50 & 282.8 & $+0.340$ & $+0.344$ & $+0.004$ & $[-0.003,+0.014]$ & no \\
Long-form    & B2 TF--IDF   &  5 & 159 &  47.4 & $+0.945$ & $+0.966$ & $+0.021$ & $[-0.036,+0.080]$ & no \\
Long-form    & B2 TF--IDF   & 10 &  77 &  92.9 & $+0.207$ & $+0.223$ & $+0.016$ & $[-0.011,+0.048]$ & no \\
Long-form    & B2 TF--IDF   & 20 &  39 & 181.4 & $-0.104$ & $-0.089$ & $+0.015$ & undefined & --- \\
Long-form    & C2 FinBERT   &  5 & 159 &  47.4 & $+0.945$ & $+0.941$ & $-0.004$ & $[-0.020,+0.012]$ & no \\
Long-form    & C2 FinBERT   & 10 &  77 &  92.9 & $+0.207$ & $+0.235$ & $+0.028$ & $[-0.020,+0.077]$ & no \\
Long-form    & C2 FinBERT   & 20 &  39 & 181.4 & $-0.104$ & $-0.104$ & $+0.001$ & undefined & --- \\
Long-form    & C6 prompted  &  5 & 159 &  47.4 & $+0.945$ & $+0.962$ & $+0.017$ & $[-0.018,+0.053]$ & no \\
Long-form    & C6 prompted  & 10 &  77 &  92.9 & $+0.207$ & $+0.207$ & $-0.001$ & $[-0.011,+0.016]$ & no \\
Long-form    & C6 prompted  & 20 &  39 & 181.4 & $-0.104$ & $-0.104$ & $-0.000$ & undefined & --- \\
\bottomrule
\end{tabular}
\end{table}

\begin{table}[htbp]
\centering
\caption[Downside-risk and utility adjudication]{Downside-risk and
utility tallies. The value-at-risk panel audits three forecasts (raw
HAR, the recalibrated reference $f_R$ and the text-augmented $f_U$)
on 36 (channel, model, horizon) rows at each of two nominal levels, 72
cells in all, under the pooled-drift convention. The tick loss is the
asymmetric score of \cite{gonzalezrivera2004forecasting} and its DM
statistics are observation-level with HAC lag $h-1$ (off the primary
day-clustered basis, tagged). The utility panel is the annualised
performance fee a mean--variance investor at risk aversion $\gamma = 2$
and a 6\% expected-return assumption would pay to swap $f_R$ for $f_U$
over 27 cells. Neither panel contains any prompted-arm row, so both
corroborate the grid without adjudicating the surviving 8-K residual.}
\label{tab:app-econ-summary}
\singlespacing\small
\begin{tabular}{@{}lp{0.52\textwidth}@{}}
\toprule
Quantity & Outcome \\
\midrule
\multicolumn{2}{@{}l}{\emph{Value-at-risk backtest (72 cells; \cite{christoffersen1998evaluating,kupiec1995techniques})}} \\
Recalibration vs.\ text        & the raw-HAR-to-$f_R$ move shifts the violation rate further than the $f_R$-to-$f_U$ move in \textbf{69 of 72} cells; median moves 0.95 against 0.05 percentage points \\ % src: results/tables/var_backtest.csv (re-derived 2026-08-02)
Tick loss, text vs.\ $f_R$     & lower in 55 of 72; significantly lower in 29; significantly \emph{worse} in 12 \\ % src: results/tables/var_backtest.md
Tick loss, recalibration alone & improves 60 of 72, 54 of them significantly \\ % src: results/tables/var_backtest.md
Violation rate closer to nominal & 35 of 72: no better than a coin \\ % src: results/tables/var_backtest.md
Residual miscalibration        & at the 1\% level 107 of 108 forecast-cells over-violate; the one exception is long-form C5 Qwen3 at $h{=}20$ under $f_U$, at 0.99\% \\ % src: results/tables/var_backtest.csv (re-derived 2026-08-02)
\midrule
\multicolumn{2}{@{}l}{\emph{Utility-based fee, 27 cells, $\gamma = 2$, $\mu = 6\%$ \cite{fleming2001timing}}} \\
Fee, unconstrained             & positive in 14 of 27; median $+0.03$ basis points a year; cells span $-26.7$ to $+96.3$ \\ % src: results/tables/utility_value.csv (re-derived 2026-08-02)
Fee, weights capped at the 99th percentile & positive in 14 of 27; median $+0.01$ basis points \\ % src: results/tables/utility_value.csv (re-derived 2026-08-02)
Target-volatility Sharpe (leverage-controlled) & $f_U$ ahead in 12 of 27; median gain $-0.00003$ \\ % src: results/tables/utility_value.csv (re-derived 2026-08-02)
Agreement with the statistical verdict & the fee agrees in sign with the QLIKE increment in only 6 of 27 cells; the largest fees occur where QLIKE judges $f_U$ worse \\ % src: results/tables/utility_value.md
\bottomrule
\end{tabular}
\end{table}

\FloatBarrier
\section{The Label-Proxy Relabelling in Full}
\label{app:res-rangebased}

Section~\ref{sec:val-perturb} reports the range-based relabelling as a
count-weakening perturbation. It is more than that, and this section gives
it whole, because the perturbation is one of only two places in the study
where the central negative (no cell clears the full conjunction) acquires an
exception; the other is the licence-free panel reported in the same section,
where the conjunction moves from 0 to 1 on the covered rows under either
label source. An exception reported only as a count reads as buried.

The construction is fixed in the registered analysis record and is
asymmetric against text. Realised volatility is recomputed from the
Parkinson high--low range estimator (Garman--Klass restated beside it).
The two references built on realised-volatility features are refitted on the
new labels, along with the recalibration and the combiner weights; GARCH,
EGARCH and ARIMA are label-free return-based forecasters with no range-based
analogue, so they stay frozen and are merely recalibrated. The text arms, trained against
close-to-close targets, are \emph{not} retrained. Text is therefore judged
on forecasts built for a different target than the one it is now scored
against, which biases the exercise against text rather than for it.
Figure~\ref{fig:app-rangebased} draws the whole perturbation and
Table~\ref{tab:app-rangebased} collects the committed blocks behind it.
Section~\ref{app:res-rangebased-cells} gives the same perturbation cell by
cell, so that which cells move is visible rather than only how many.

\begin{figure}[tbp]
\centering
\includegraphics[width=\textwidth,height=0.83\textheight,keepaspectratio]{figures/AP1_label_proxy_cascade.pdf}
% caption numbers src: results/tables/rangebased_cascade.csv (all four panels); rangebased_cascade.md (branch-(d) rank check)
\caption[The label-proxy relabelling, whole]{The label-proxy relabelling,
whole. (a) The four rungs under the committed close-to-close labels and
the two range-based proxies, on the same 69 cells and the same frozen
text forecasts: the primary rung falls under both while the
maximal-pool rung \emph{rises} from 9 to 15, and the firm-identity rung barely
moves, 8 to 7 under Parkinson and 8 to 8 under Garman--Klass.
% src: results/tables/rangebased_cascade.csv (primary bars are the placebo-gated old/pk/gk_genuine columns, 38/19/21; firm_detect 8/7/8 and pool_detect 9/15/15 are Holm-only, those rungs carrying no placebo gate. The Holm-only primary column, primary_detect, reads 38/21/23 and is not what the panel draws)
The relabelling
re-orders the ladder rather than weakening it uniformly. (b) Every cell's primary
increment under both labellings, ordered by the committed value, with the
one cell that clears all three rungs under Parkinson ringed. One Parkinson
value, $-8.19\%$, falls below the panel floor and is marked there by a
hollow caret carrying its value. (c) That
survivor's descent down the rungs against its own minimum detectable
effect: what reaches the firm-identity reference is $+0.016\%$. (d) The
reference-ordering check: recalibrated HAR's mean rank among the five
single recalibrated price references, which \emph{improves} under
Parkinson from 3.67 to 1.50. Primary-rung counts are placebo-gated; the
strengthened rungs are Holm-only within their own 69-cell families. That
applies to panels~(a) to~(c). Panel~(d) counts no cells: its denominator is
the five single recalibrated price references it ranks the HAR against, over
the six disclosure-by-horizon panels, and a mean rank carries no test and no
family.}
\label{fig:app-rangebased}
\end{figure}

\begin{table}[htbp]
\centering
\caption[The range-based relabelling: rungs, power and reference
ordering]{The range-based relabelling. Panel (a): Holm-detected cells per
rung under the committed close-to-close labels and under the two
range-based proxies, all on the same 69-cell family, the same frozen text
forecasts and the same machinery. Panel (b): per-cell minimum detectable
effects at 80\% power and 5\% two-sided size, which is why part of the fall
in panel (a) is power rather than signal. Panel (c): recovery of an
injected firm-orthogonal signal at three calibrated sizes, Holm-detected
within each pre-declared (stage, level) family. Panel (d): the
reference-ordering check, which asks whether the relabelling unseats HAR-RV
as the price reference the ladder is built on: the rank of a
recalibrated HAR against the five single recalibrated price references, one
ranking per disclosure-by-horizon panel. A lower mean rank is a stronger
HAR.}
\label{tab:app-rangebased}
\singlespacing\small
\setlength{\tabcolsep}{4pt}
% src: results/tables/rangebased_cascade.md (all four panels transcribed from the committed summary blocks)

\emph{(a) Holm-detected cells per rung, out of 69}\\[0.4ex]
\begin{tabular}{@{}lrrr@{}}
\toprule
Rung & Committed & Parkinson & Garman--Klass \\
\midrule
Primary: text over recalibrated HAR & 38 & 21 & 23 \\
\quad after the label-shuffle placebo gate & \textbf{38} & \textbf{19} & \textbf{21} \\
Firm-identity-augmented reference & 8 & 7 & 8 \\
Maximal five-model price pool & 9 & 15 & 15 \\
\textbf{Full conjunction} & \textbf{0} & \textbf{1} & \textbf{2} \\
\bottomrule
\end{tabular}

\vspace{1.2ex}
\emph{(b) Per-cell minimum detectable effect, relative QLIKE \%}\\[0.4ex]
\begin{tabular}{@{}lrrr@{}}
\toprule
& Committed & Parkinson & Garman--Klass \\
\midrule
Median MDE & 0.823 & 0.913 & 1.018 \\
Interquartile range & [0.37, 1.27] & [0.60, 1.44] & [0.77, 1.83] \\
Median shrink against committed & --- & $-10.9\%$ & $-23.6\%$ \\
\bottomrule
\end{tabular}

\vspace{1.2ex}
\emph{(c) Injected-signal recovery, cells detected of 69 (primary / identity / pool / conjunction)}\\[0.4ex]
\begin{tabular}{@{}llll@{}}
\toprule
Injected size & Committed & Parkinson & Garman--Klass \\
\midrule
0.3\% & 12 / 7 / 6 / 2    & 7 / 1 / 4 / 1    & 4 / 2 / 3 / 2 \\
0.5\% & 20 / 11 / 12 / 6  & 11 / 4 / 9 / 3   & 8 / 5 / 6 / 4 \\
1.0\% & 41 / 20 / 19 / 13 & 26 / 16 / 20 / 11 & 19 / 14 / 15 / 6 \\
\bottomrule
\end{tabular}

\vspace{1.2ex}
\emph{(d) Reference-ordering check: rank of recalibrated HAR among the five single price references}\\[0.4ex]
\begin{tabular}{@{}lrr@{}}
\toprule
Labels & Mean rank over the six panels & Panels where HAR ranks first \\
\midrule
Committed close-to-close & 3.67 & 0 of 6 \\
Parkinson                & \textbf{1.50} & \textbf{3 of 6} \\
\bottomrule
\end{tabular}
\end{table}

Panel (d) is the leg of this perturbation that supports the design rather
than straining it, and it is the reason the relabelling was registered with
a reference-ordering branch in the first place. The choice of HAR-RV as the
ladder's price reference is open to the objection that it was chosen for
familiarity rather than for strength. Section~\ref{sec:res-standalone}
concedes that HAR is not the QLIKE-strongest price arm on the committed
labels: ranked against the five single recalibrated price references, it
averages 3.67 of 5 and is first in none of the six panels. Under the
Parkinson labels that ordering reverses (mean rank 1.50, first in three
of six), so the relabelling that weakens the counts simultaneously
\emph{strengthens} the price side of the comparison. The registered
branch that would have fired had HAR degraded materially does not fire.
The two readings belong together: the proxy that costs text the most cells
is also the proxy on which text is being asked to beat the strongest
version of its opponent. % src: results/tables/rangebased_cascade.md branch-(d) block

The exception itself is one cell, and it is named here because a count
cannot be interrogated. Under Parkinson labels the sole cell clearing the
primary rung, the maximal price pool and the firm-identity reference
together is \textbf{FinBERT-S1 on the event-driven channel at $h{=}5$}
(25{,}109 observations over 996 trading days). Its profile down the ladder
is $+2.511\%$ over the recalibrated HAR (DM $-4.73$, Holm
$1.5\times10^{-4}$), $+1.130\%$ over the fitted five-model price pool (DM
$-3.46$, Holm .0257) and $+0.016\%$ over the firm-identity reference (DM
$-5.90$, Holm $3.1\times10^{-7}$), with a label-shuffle placebo of $+0.29$,
comfortably inside the pass band. % src: results/tables/rangebased_cascade.csv (event\_driven / C2\_finbert\_s1 / h5; pk\_conj = True)
On the committed labels the same cell reads $+2.572\%$, $+0.768\%$ and
$-0.061\%$ down the same three rungs, so what the relabelling changes in
sign is the last of them alone. What it changes in verdict is more: under
Parkinson the firm-identity and pool rungs both turn from undetected to
detected, the primary rung being detected either way. % src: results/tables/rangebased_cascade.csv (old\_rel, old\_pool\_rel, old\_firm\_rel on the same row)

Two things must be said about that $+0.016\%$ at once, and the second is
the one that matters. It is statistically robust: a DM statistic of
$-5.90$ and a Holm-adjusted $p$ of three parts in ten million is not a
borderline detection, and nothing here is being dismissed on a
technicality. And it is, in magnitude, two orders below the increment that
survived the previous rung, and roughly one sixty-fourth of the cell's own
minimum detectable effect of $1.04\%$: a precisely estimated near-zero
rather than a gain the price side failed to see. Under one label proxy,
therefore, the conjunction is not empty; what occupies it is an increment
over the identity reference of sixteen thousandths of one per cent. That
is the honest form of the exception, and it is entirely consistent with the
character of the residual reported in
Section~\ref{sec:res-residual}: real, repeatedly detectable, and far too
small to matter.

The Garman--Klass restatement admits two cells rather than one, and they
are a different arm: the Loughran--McDonald engineered-feature model on the
event-driven channel at $h{=}5$ and $h{=}10$, at $+0.360\%$ and $+0.111\%$
over the recalibrated HAR and $+0.412\%$ and $+0.154\%$ over the identity
reference (Holm .0027 and .0466) on cells whose own
minimum detectable effects are $0.183\%$ and $0.085\%$, and neither of
which is genuine on the committed labels at all. % src: results/tables/rangebased_cascade.csv (event\_driven / B4\_lm\_features / h5, h10; gk\_conj = True; primary-rung gk\_primary\_rel 0.3605 and 0.1108, identity-rung gk\_firm\_rel 0.4121 and 0.1543 with gk\_firm\_holm 0.0027 and 0.0466; gk\_mde 0.1828 and 0.0850; old\_genuine False on both)
The three conjunction survivor sets (empty, one FinBERT cell, two
dictionary cells) are therefore pairwise disjoint, which is the finding
Section~\ref{sec:val-perturb} reports. No cell anywhere clears the
conjunction under every label proxy, and the cross-proxy conjunction is 0
of 69.

\FloatBarrier
\section{The Tuned Challenger Arm in Full}
\label{app:res-hpo}

The ASHA arm specified in Appendix~\ref{app:hyperparams} exists to answer
one objection: that the fixed training recipe, not the text, produces the
null. Its evaluation was run once under a single-shot rule, the
script refusing a second execution. Chapter~\ref{ch:results} reports the
verdict in two clauses; Table~\ref{tab:app-hpo} is the whole evaluation,
both halves of it, because a tuning defence that is quoted rather than
shown invites the reply that the favourable half was quoted.

Two properties of this table are easy to misread and are stated before it.
Its ensemble rule is the per-observation \emph{log-space} mean of the three
seeds, not the arithmetic mean used everywhere else in this report: the
object tested is the object the search selected, which is the only way the
test answers the question asked. No level in this table may therefore be
set beside a level elsewhere. And its Holm family is the six cells here and
nothing else, the wider tuning programme having been descoped before
execution and the descoping recorded in the evidence file rather than in
retrospect.

\begin{table}[htbp]
\centering
\caption[The tuned challenger arm: standalone and combination]{The
ASHA-tuned FinBERT arm, the complete one-shot test evaluation. Panel (a)
is the standalone comparison against the validation-recalibrated HAR
reference. Panel (b) is the same tuned forecast entering the combination
grid's log-space combiner, scored as an increment over that reference.
Decision loss is volatility-unit QLIKE in both panels; day-clustered DM,
Holm within the six-cell family. \emph{Positive rel.\% means the tuned arm
lowers loss}; in panel (a) a positive DM therefore means the tuned arm is
worse. The placebo gate runs only on cells that reach DM $<0$ with Holm
$p<.05$, which is why most cells have no placebo entry. Here ``ls'' is the
label-shuffle placebo and ``wd'' the within-date one, and a cell must pass
both. The arm uses a log-space seed ensemble, not the arithmetic ensemble
of the rest of this appendix, so its levels are not comparable with any
other table here.}
\label{tab:app-hpo}
\singlespacing\small
\setlength{\tabcolsep}{4pt}
% src: results/tables/hpo_test_eval.md panels (a) tuned_standalone and (b) tuned_cascade, transcribed in full

\emph{(a) Standalone: tuned ensemble against the recalibrated HAR}\\[0.4ex]
\begin{tabular}{@{}lrrrrrrl@{}}
\toprule
Channel & $h$ & $n$ & days & QLIKE $f_R$ & QLIKE tuned & rel.\% & DM / Holm \\
\midrule
Long-form    &  5 & 7{,}951  & 809 & 0.1209 & 0.1444 & $-19.46$ & $+5.36$ / $<$.0001 \\
Long-form    & 10 & 7{,}933  & 803 & 0.0873 & 0.1204 & $-37.93$ & $+3.20$ / .0028 \\
Long-form    & 20 & 7{,}902  & 794 & 0.0701 & 0.0875 & $-24.79$ & $+1.01$ / .3115 \\
Event-driven &  5 & 25{,}109 & 996 & 0.1265 & 0.1451 & $-14.72$ & $+8.27$ / $<$.0001 \\
Event-driven & 10 & 25{,}001 & 991 & 0.0883 & 0.1106 & $-25.33$ & $+5.89$ / $<$.0001 \\
Event-driven & 20 & 24{,}732 & 981 & 0.0645 & 0.0875 & $-35.58$ & $+5.09$ / $<$.0001 \\
\bottomrule
\end{tabular}

\vspace{1.2ex}
\emph{(b) Combination: the tuned arm's increment over the same reference}\\[0.4ex]
\begin{tabular}{@{}lrrrrrrl@{}}
\toprule
Channel & $h$ & rel.\% & DM & Holm $p$ & plac.\ ls & plac.\ wd & Gate \\
\midrule
Long-form    &  5 & $+1.57$ & $-0.88$ & 1.0000 & --- & --- & --- \\
Long-form    & 10 & $+0.10$ & $+0.53$ & 1.0000 & --- & --- & --- \\
Long-form    & 20 & $-3.54$ & $+3.14$ & .0088  & --- & --- & --- \\
Event-driven &  5 & $+2.88$ & $-4.89$ & $<$.0001 & $-2.40$ & $-2.46$ & \textbf{artefact} \\
Event-driven & 10 & $+2.13$ & $-2.54$ & .0446  & $+0.40$ & $+1.38$ & \textbf{genuine} \\
Event-driven & 20 & $-0.33$ & $+0.03$ & 1.0000 & --- & --- & --- \\
\bottomrule
\end{tabular}
\end{table}

Panel (a) settles the standalone half and settles it heavily. Tuning
recovers about three quarters of the long-form standalone \emph{validation}
gap and none of the test one: the tuned arm is worse than the recalibrated
HAR in all six test cells, by 14.7 to 37.9 per cent, five of the six
Holm-significantly. A search over eight hyperparameters, 32 and 24 Sobol
trials, three seeds and three promotion rungs does not move a fine-tuned
encoder to the price baseline on this panel.

Panel (b) settles the half that Chapter~\ref{ch:results} states in a clause
and that a reader is entitled to see, because the combination question is
the one the report is actually about. Of the six tuned combination cells,
one clears Holm and both placebos: event-driven $h{=}10$, at $+2.13\%$.
That cell was \emph{already} genuine under the fixed recipe, at $+2.42\%$
(Table~\ref{tab:app-grid-ed}). The two figures are not comparable and are not
being compared: the tuned arm is averaged in log space, matching the rule that
selected it, and the grid arithmetically, and the source that carries the tuned
row says the two conventions are never mixed. What the cell shows is that
tuning buys nothing the fixed recipe did not already have. % src: results/tables/hpo_test_eval.md, ensemble-rule disclosure ("The two conventions ... are NOT mixed anywhere; do not compare rows across the two conventions")
The arm's larger increment, $+2.88\%$ at event-driven $h{=}5$,
fails the gate outright, both placebos returning statistics beyond the
band on a cell the fixed recipe passes --- though not cleanly: at $-1.72$
against the $\pm2$ band that cell carries the largest-magnitude placebo
statistic of all 38 genuine cells, and the fixed recipe is never asked to
clear the within-date placebo the tuned arm fails. % src: results/tables/m1_ensemble_primary.csv (vol_placebo_dm_clu over the 38 genuine cells: event_driven C2_finbert_s1 h=5 at -1.7185, next largest -1.4455)
At long-form $h{=}20$ the
tuned arm is Holm-significantly \emph{worse} than the untuned reference.
The tuned arm therefore adds no combination cell the fixed recipe did not
already have, and the ``you simply did not tune'' objection is closed on
both the standalone and the incremental sides rather than on one of them.

\FloatBarrier
\section{The Cutoff Stratification in Full}
\label{app:res-cutoff}

The prompted arm's residual invites the objection that an instruction-tuned
model may have seen the outcomes it is forecasting. Chapter~\ref{ch:validation}
answers it with the direction of the point estimates and with the Holm verdicts
in five of six cells. Table~\ref{tab:app-cutoff} gives
the whole stratification so that the answer can be checked cell by cell.

The test years are split at 2024-07-01, an approximation of the prompted
model's training-data boundary, and the combination weights are the full
validation fit, frozen. Only the test evaluation is stratified, so
nothing is refitted to either stratum. Two arms are tracked in every cell:
the full-text increment over the recalibrated HAR, and the same increment
over a reference that already contains the same model's own date-and-ticker
forecast.

\begin{table}[htbp]
\centering
\caption[Increments before and after the model's training cutoff]{The
prompted arm's increment on filings before and after 2024-07-01, all 12
cells at both references. \emph{Post-cutoff outcomes postdate the model's
approximate training-data boundary.}
Volatility-unit QLIKE, day-clustered DM, Holm within the 24-cell family;
combination weights are the full-validation fit, frozen, and only the test
evaluation is stratified. Bold marks Holm significance. The boundary is an
approximation, and it is a conservative one: placing it too early moves
memorisable filings into the post stratum, which would bias the post-cutoff
increment upward. The observed pattern (post above pre in all six
cells) is the opposite of what memorised outcomes would produce.}
\label{tab:app-cutoff}
\singlespacing\footnotesize
\setlength{\tabcolsep}{4pt}
% src: results/tables/llm_contamination.md section 3 (cutoff-date stratification), all 24 rows
\begin{tabular}{@{}llrrlrrl@{}}
\toprule
& & \multicolumn{3}{c}{Pre-cutoff} & \multicolumn{3}{c}{Post-cutoff} \\
\cmidrule(lr){3-5}\cmidrule(lr){6-8}
Channel & $h$ & $n$ & rel.\% & Holm $p$ & $n$ & rel.\% & Holm $p$ \\
\midrule
\multicolumn{8}{@{}l}{\emph{(a) Full text against the recalibrated HAR reference}} \\
Long-form    &  5 & 4{,}957  & $+1.25$ & \textbf{$<$.0001} & 2{,}994 & $+2.56$ & \textbf{.0013} \\
Long-form    & 10 & 4{,}956  & $+1.71$ & \textbf{$<$.0001} & 2{,}977 & $+3.09$ & \textbf{.0001} \\
Long-form    & 20 & 4{,}949  & $+0.11$ & .5217            & 2{,}953 & $+0.53$ & \textbf{.0355} \\
Event-driven &  5 & 15{,}804 & $+0.87$ & .0669            & 9{,}305 & $+1.71$ & \textbf{.0002} \\
Event-driven & 10 & 15{,}786 & $+0.70$ & .1837            & 9{,}215 & $+1.40$ & \textbf{.0355} \\
Event-driven & 20 & 15{,}751 & $+0.51$ & .1280            & 8{,}981 & $+0.86$ & .9822 \\
\addlinespace
\multicolumn{8}{@{}l}{\emph{(b) Full text beyond the same model's date{+}ticker forecast}} \\
Long-form    &  5 & 4{,}957  & $+1.03$ & \textbf{.0001}   & 2{,}994 & $+2.16$ & \textbf{.0055} \\
Long-form    & 10 & 4{,}956  & $+1.44$ & \textbf{$<$.0001} & 2{,}977 & $+2.73$ & \textbf{.0004} \\
Long-form    & 20 & 4{,}949  & $+0.05$ & .9822            & 2{,}953 & $+0.35$ & .1280 \\
Event-driven &  5 & 15{,}804 & $+0.55$ & .1837            & 9{,}305 & $+1.07$ & \textbf{.0324} \\
Event-driven & 10 & 15{,}786 & $+0.47$ & .1996            & 9{,}215 & $+0.98$ & .2381 \\
Event-driven & 20 & 15{,}751 & $+0.46$ & .1042            & 8{,}981 & $+0.86$ & .9822 \\
\midrule
\multicolumn{8}{@{}l}{\emph{Post-cutoff Holm counts: 5 of 6 against HAR, 3 of 6 beyond identity (raw 5 of 6).}} \\
\multicolumn{8}{@{}l}{\emph{Post-cutoff point estimate exceeds pre-cutoff in 6 of 6 cells at both references.}} \\
\bottomrule
\end{tabular}
\end{table}

The counts are what convert an impression into a bounded claim. Against the
recalibrated HAR the post-cutoff increment is Holm-significant in five of
the six cells. The single failure, event-driven $h{=}20$, is the cell that
is already non-significant on the undivided test set, so nothing is lost
there that the full panel had. Beyond the same model's own identity
forecast the post-cutoff increment is Holm-significant in three of six and
raw-significant in five of six, with positive point estimates in all six.
The weaker of those two counts is the honest headline: on unmemorisable
filings, with the model's own name-recall inside the reference, the
residual is demonstrable in half the cells and directionally present in all
of them. That is enough to exclude memorisation as the explanation of the
residual, and not enough to claim the residual is uniform across the grid;
Chapter~\ref{ch:validation} claims the former only. % src: results/tables/llm_contamination.md bottom-line block

\FloatBarrier
\section{Reference Robustness and the Representation Ladder}
\label{app:res-refrobust}

Two questions about the reference survive everything reported above, and
both are answered by evidence the main text has room only to allude to.

The first is whether the ladder's verdict depends on HAR-RV specifically.
It does not. Table~\ref{tab:app-shar} refits the whole incremental
apparatus with realised-semivariance HAR in place of plain HAR as the price
reference (a strictly better price model on these labels, beating HAR
under QLIKE in all nine panels). The text increments are
indistinguishable from the ones the report quotes. Across 18 text-by-channel-by-horizon
cells no verdict changes: the ten cells where text adds against HAR add
against SHAR, the seven where it loses lose, and the one weak cell stays
weak. The largest disagreement between the two reference columns anywhere
in the table is 0.181 percentage points, 0.18 at the precision printed, at
long-form C5 Qwen3 at $h{=}10$; differencing the two rounded columns as printed
gives 0.19, which is the arithmetic of the display and not the gap. % src: results/tables/stronger\_baselines.csv (section m1\_incremental over the 18 cells: 10 adds, 7 loses, 1 weak, matching the Verdict column below; largest absolute gap between A6\_shar\_rel\_impr\_pct and A2\_rel\_impr\_pct is 0.181 pp at long\_form / C5\_qwen3 / h10, i.e. 0.18 at two decimals; the columns themselves print as -2.89 and -3.08, whose difference is 0.19)

\begin{table}[htbp]
\centering
\caption[Text's increment against a semivariance-HAR reference]{The
incremental-value test with realised-semivariance HAR
\cite{pattonsheppard2015shar} substituted for HAR-RV as the price
reference. Panel (a) establishes that the substitution is an upgrade:
variance-unit test QLIKE, DM against A2 on the common kept sample, negative
meaning SHAR is better. Panel (b) is the incremental test itself, with the
two reference columns side by side on the identical rows.
\textbf{Inference in this table is observation-level, not day-clustered
(off the primary basis, tagged)}, because that is the convention the
committed comparison was run under. Magnitudes are therefore roughly two to
three times those of the clustered tables and are not comparable with them.
The row support is the return-matched sample: rows whose ticker-recycled
return joins are unreliable are dropped, retaining 90.4 to 91.5 per cent.
That is why A2's levels differ slightly from
Table~\ref{tab:app-leaderboard}'s. What the table is evidence for is the
\emph{invariance of the verdict} to the reference, not the magnitudes.}
\label{tab:app-shar}
\singlespacing\small
\setlength{\tabcolsep}{4pt}
% src: results/tables/stronger_baselines.md (both blocks; A6_shar columns)

\emph{(a) Is SHAR the better price model? Variance-unit test QLIKE}\\[0.4ex]
\begin{tabular}{@{}lrrrrl@{}}
\toprule
Channel & $h$ & $n$ & A2 (common) & SHAR & $\Delta$QLIKE \% / DM \\
\midrule
Long-form    &  5 & 7{,}550  & 0.5635 & 0.5568 & $+1.18$ / $-3.74$ \\
Long-form    & 10 & 7{,}167  & 0.4226 & 0.4106 & $+2.84$ / $-6.49$ \\
Long-form    & 20 & 7{,}097  & 0.3076 & 0.2948 & $+4.17$ / $-8.02$ \\
Event-driven &  5 & 23{,}855 & 0.5992 & 0.5927 & $+1.09$ / $-5.76$ \\
Event-driven & 10 & 22{,}785 & 0.4576 & 0.4454 & $+2.65$ / $-8.37$ \\
Event-driven & 20 & 22{,}318 & 0.3327 & 0.3192 & $+4.08$ / $-11.73$ \\
Combined     &  5 & 31{,}405 & 0.5908 & 0.5841 & $+1.14$ / $-6.77$ \\
Combined     & 10 & 29{,}952 & 0.4497 & 0.4373 & $+2.74$ / $-9.79$ \\
Combined     & 20 & 29{,}415 & 0.3286 & 0.3151 & $+4.12$ / $-13.50$ \\
\bottomrule
\end{tabular}

\vspace{1.2ex}
\emph{(b) Text's increment over each reference, volatility-unit QLIKE, identical rows}\\[0.4ex]
\begin{tabular}{@{}llrrrrrl@{}}
\toprule
Channel & Text arm & $h$ & $n$ & vs.\ A2 \% & DM & vs.\ SHAR \% & Verdict \\
\midrule
Long-form & B2 TF--IDF   &  5 & 7{,}550  & $+3.39$ & $-11.18$ & $+3.53$ & adds \\
Long-form & B2 TF--IDF   & 10 & 7{,}167  & $+3.63$ & $-15.21$ & $+3.78$ & adds \\
Long-form & B2 TF--IDF   & 20 & 7{,}097  & $+5.27$ & $-19.10$ & $+5.29$ & adds \\
Long-form & C2 FinBERT   &  5 & 7{,}550  & $+0.49$ & $-5.51$  & $+0.55$ & adds \\
Long-form & C2 FinBERT   & 10 & 7{,}167  & $+4.41$ & $-12.46$ & $+4.49$ & adds \\
Long-form & C2 FinBERT   & 20 & 7{,}097  & $-0.33$ & $+4.96$  & $-0.30$ & loses \\
Long-form & C5 Qwen3     &  5 & 7{,}550  & $-1.03$ & $+3.91$  & $-0.94$ & loses \\
Long-form & C5 Qwen3     & 10 & 7{,}167  & $-3.08$ & $+4.24$  & $-2.89$ & loses \\
Long-form & C5 Qwen3     & 20 & 7{,}097  & $-6.14$ & $+4.19$  & $-5.98$ & loses \\
Event-dr. & B2 TF--IDF   &  5 & 23{,}855 & $+1.17$ & $-6.76$  & $+1.18$ & adds \\
Event-dr. & B2 TF--IDF   & 10 & 22{,}785 & $+1.27$ & $-6.78$  & $+1.29$ & adds \\
Event-dr. & B2 TF--IDF   & 20 & 22{,}318 & $+1.68$ & $-7.93$  & $+1.69$ & adds \\
Event-dr. & C2 FinBERT   &  5 & 23{,}855 & $+2.01$ & $-14.24$ & $+2.02$ & adds \\
Event-dr. & C2 FinBERT   & 10 & 22{,}785 & $+1.92$ & $-13.37$ & $+1.92$ & adds \\
Event-dr. & C2 FinBERT   & 20 & 22{,}318 & $+0.75$ & $-1.68$  & $+0.78$ & weak \\
Event-dr. & C5 Qwen3     &  5 & 23{,}855 & $-0.14$ & $+2.67$  & $-0.15$ & loses \\
Event-dr. & C5 Qwen3     & 10 & 22{,}785 & $-0.15$ & $+3.41$  & $-0.16$ & loses \\
Event-dr. & C5 Qwen3     & 20 & 22{,}318 & $-0.52$ & $+6.51$  & $-0.53$ & loses \\
\bottomrule
\end{tabular}
\end{table}

The second question is the one Chapter~\ref{ch:background} raises about
domain adaptation, which Section~\ref{sec:res-standalone} answers in a
single sentence against this table. Does
pretraining a language model on financial prose, rather than merely
making it larger, lift text over the simple pipelines and over price?
Table~\ref{tab:app-domain} settles it in the negative, and settles it in
the direction the evidence actually supports.

\begin{table}[htbp]
\centering
\caption[Domain adaptation against the simple pipeline and the price
reference]{Does domain pretraining lift text? Seed-averaged test QLIKE in
\emph{variance units} on the long-form panel, at each horizon, lower being
better. C-block rows are the mean over the three seeds of each run's test
QLIKE. The comparison the table is built to support is the vertical one
(every text row against the last two), not a FinBERT-versus-BERT
ordering, which the numbers do not support in either direction.}
\label{tab:app-domain}
\singlespacing\small
\setlength{\tabcolsep}{5pt}
% src: results/tables/seed_aggregate.csv (long_form, qlike_mean by model x horizon); A2 row from results/tables/cost_accuracy.csv
\begin{tabular}{@{}llrrr@{}}
\toprule
Model & Pretraining & $h{=}5$ & $h{=}10$ & $h{=}20$ \\
\midrule
C2 FinBERT (S1) & financial prose & 1.3428 & 1.2440 & 0.7262 \\
C1 BERT (S1)    & generic         & 1.6444 & 1.1968 & 0.6995 \\
C1 BERT (S2)    & generic         & 1.4557 & 1.0372 & 0.6905 \\
C3 RoBERTa (S1) & generic         & 1.6303 & 1.0863 & 0.7518 \\
\midrule
B2 TF--IDF ridge & none (counts)  & \textbf{1.2371} & \textbf{1.0522} & \textbf{0.7083} \\
A2 HAR-RV \emph{(reference)} & none (price) & \textbf{0.5564} & \textbf{0.4618} & \textbf{0.2970} \\
\bottomrule
\end{tabular}
\end{table}

The defensible reading is the one against the pipelines and the price arm,
and it is unambiguous. Domain-pretrained FinBERT does not lift text above
the TF--IDF pipeline at any horizon (it is worse at all three), let
alone above the price reference, whose loss it exceeds by factors of 2.4,
2.7 and 2.4. What the table does \emph{not} license is the neater claim
that domain pretraining fails to beat generic pretraining. Against BERT the
two are overlapping rather than ordered: FinBERT is clearly better at
$h{=}5$ (1.343 against 1.644 and 1.456) and slightly worse at $h{=}10$ and
$h{=}20$. The ordering also moves with the truncation strategy. The
finding is that neither scale nor domain adaptation lifts text above simple
representations or price, not that domain adaptation is worthless
relative to generic pretraining, which this panel cannot show.


\begin{figure}[tbp]
\centering
\includegraphics[width=\textwidth,height=0.83\textheight,keepaspectratio]{figures/AR1_reference_spec_cell_matrix__a.pdf}
% caption numbers src: results/tables/maximal_pool_robustness.csv (414 rows = 69 cells x 3 reference specs x 2 seed bases); results/tables/m1_ensemble_primary.csv (the recalibrated-HAR column on the primary and archived single-seed bases)
\caption[The reference ladder by specification, cell by cell]{
The reference
ladder by specification, cell by cell. Panel~(a) draws each of the 69 combination cells
(event-driven panel above, long-form below; one row per model, one column per
horizon) against four price references, ordered left to right by how
much they absorb: the recalibrated HAR-RV of the primary rung, the
validation-best single price model, a never-fitted equal-weight pool of five
price models, and the validation-fitted pool of the same five. A filled square
marks a cell where text adds at Holm, an open square with a dot one
significant at the raw level only, and a vermillion square a cell where text
significantly \emph{hurts}. 
}
\label{fig:app-reference-spec-matrix}
\end{figure}

\begin{figure}[tbp]
\centering
\includegraphics[width=\textwidth,height=0.83\textheight,keepaspectratio]{figures/AR1_reference_spec_cell_matrix__b.pdf}
% caption numbers src: results/tables/maximal_pool_robustness.csv (414 rows = 69 cells x 3 reference specs x 2 seed bases); results/tables/m1_ensemble_primary.csv (the recalibrated-HAR column on the primary and archived single-seed bases)
\caption[The reference ladder by specification, cell by cell (b, c)]{
The reference
ladder by specification, cell by cell, continued from the cell matrix of
Figure~\ref{fig:app-reference-spec-matrix}.
Panel~(b) counts how many of the three stronger
references each cell adds under (34 cells under none, 16 under one, 13
under two, 6 under all three). 
Panel~(c) repeats the four survivor
counts on the archived single-seed basis (29, 28, 19, 8) beside the primary
(38, 34, 17, 9). % src: results/tables/maximal_pool_robustness.csv (adds_holm by reference and basis); results/tables/m1_ensemble_primary.csv (genuine_ens_vol 38, genuine_s26_clu 29)
Basis: the seed-ensemble forecast object, volatility-unit QLIKE, day-clustered
Diebold--Mariano statistics, Holm inside each 69-cell family; the denominator
is 69 in every column, so the four counts are directly comparable.
}
\label{fig:app-reference-spec-matrix-b}
\end{figure}

Figures~\ref{fig:app-reference-spec-matrix} and~\ref{fig:app-reference-spec-matrix-b} take panel~(a) of
Figure~\ref{fig:app-pool-audit} down to the cell, and in doing so establish
something counts alone cannot. The two views report the same four numbers;
only the cell-level view can say whether the sets those numbers count are
nested. They are. The 35 cells that add under at least one of the three pool
specifications form a strict subset of the 38 that add against the
recalibrated HAR-RV alone: not one cell anywhere in the grid is
\emph{rescued} by a stronger reference. % src: results/tables/maximal_pool_robustness.csv (union of adds_holm over the three ensemble reference specs = 35 cells, all contained in the 38 of m1_ensemble_primary.csv)
Only six cells add under all three pool specifications (TF-IDF ridge at all
three horizons, the Longformer at $h=5$, and the prompted arm at $h=5$ and
$h=10$), and every one of the six is long-form. % src: results/tables/maximal_pool_robustness.csv (cells with adds_holm under fitted_pool, eqw_pool and valbest_single on the ensemble basis)

Two readings follow. First, the survivor count is specification-dependent,
running from 34 against the strongest single member to 9 against the fitted
pool, so no reader may quote one of those numbers as ``the'' absorption result
without naming the reference it belongs to. The report accordingly reports all
four. Second, and not specification-dependent, the nesting means the ladder
subtracts rather than reshuffles. That matters, because a reshuffling ladder
(one in which a stronger reference lets new cells through) would suggest
the rungs were measuring different things, whereas a subtracting one is
consistent with a single shrinking residual, which is the reading
Chapter~\ref{ch:results} adopts.

The figure does not show the firm-identity reference or the conjunction rungs,
which belong to Chapter~\ref{ch:results}, so it is a picture of one arm of the
ladder and not of the ladder. It does not show magnitudes: a filled square is
a verdict, and the primary-rung magnitudes are the margin bars of
Figure~\ref{fig:res-ladder-matrix}. And the placebo term is absent from the
pool specifications' source table. The absence is immaterial at the primary rung,
where 38 cells survive with the placebo gate and 38 without it, but it means the
three pool columns are Holm-only columns and should not be described as
placebo-confirmed. % src: results/tables/m1_ensemble_primary.csv (vol_dmq_holm_clu and genuine_ens_vol agree at 38)

\FloatBarrier
\section{The Yelp Portability Cascade in Full}
\label{app:res-yelp}

Section~\ref{sec:val-external} reports the Yelp panel as the cross-domain
case where the identity control changes sign and becomes costly.
Table~\ref{tab:app-yelp} carries the whole cascade for both fitted text arms,
and it carries one
quantity the main text does not reach: what the identity term recovers of
the \emph{field-standard} design's apparent gain.

\begin{table}[htbp]
\centering
\caption[The Yelp cascade, both horizons]{The Yelp business-month cascade
in full: 8{,}474 businesses, 407{,}385 business-months, train to 2016,
validation 2017, test 2018--21. Loss is squared error on the star rating,
inference is month-clustered DM with a business-by-month two-way
robustness check. \emph{Nothing in this table shares a loss, a
denominator or a panel with any QLIKE percentage elsewhere in this report}.
Each row names its own reference; rows 3 and 5 are increments over
different references and are not a series. Rows 3b and 5b repeat those two
rungs for the frozen-embedding text arm against the identical references, and
carry no mean squared error because the multi-arm record holds relative
increments only. The recovery share in the note
is the zero-text entity-mean predictor's improvement expressed as a
fraction of row 1's apparent gain, both measured under row 1's own
field-standard random split.}
\label{tab:app-yelp}
\singlespacing\small
\setlength{\tabcolsep}{4pt}
% src: results/tables/yelp_cascade.csv and yelp_cascade.md (all five arms, both horizons, row notes and table note)
\begin{tabular}{@{}rlp{0.20\textwidth}rrrr@{}}
\toprule
& Arm & Reference & \multicolumn{2}{c}{$h = 1$ month} & \multicolumn{2}{c}{$h = 3$ months} \\
\cmidrule(lr){4-5}\cmidrule(lr){6-7}
& & & MSE & $\Delta$\% & MSE & $\Delta$\% \\
\midrule
1 & Text, pooled random 80/20 \emph{(field design)} & pooled mean & 0.4259 & $+28.43$ & 0.2981 & $+36.64$ \\
2 & Text alone, chronological & recal.\ AR $f_R$ & 0.4924 & $-24.70$ & 0.3537 & $-35.29$ \\
3 & AR $+$ text combiner $f_U$ & recal.\ AR $f_R$ & 0.3935 & $+0.35$ & 0.2592 & $+0.87$ \\
4 & AR $+$ entity mean \emph{(zero text)} & recal.\ AR $f_R$ & 0.4367 & $-10.58$ & 0.2863 & $-9.48$ \\
5 & AR $+$ entity mean $+$ text \emph{(residual)} & AR $+$ entity mean & 0.4351 & $+0.37$ & 0.2845 & $+0.61$ \\
\midrule
\multicolumn{7}{@{}l}{\emph{The frozen-embedding arm (Qwen3-Embedding-8B), same protocol and references}} \\
3b & AR $+$ text combiner $f_U$ & recal.\ AR $f_R$ & --- & $+0.68$ & --- & $+1.51$ \\ % src: results/tables/yelp_multiarm.csv (arm Qwen3-Emb-8B, combiner 0.677 at h=1 and 1.513 at h=3)
5b & AR $+$ entity mean $+$ text \emph{(residual)} & AR $+$ entity mean & --- & $+0.38$ & --- & $+0.79$ \\ % src: results/tables/yelp_multiarm.csv (arm Qwen3-Emb-8B, residual 0.376 and 0.786; res_p2way 4.43e-15 and 1.11e-14; placebo_dm +0.76 p=0.367 and +0.84 p=0.470)
\midrule
\multicolumn{7}{@{}p{0.95\textwidth}@{}}{\emph{Recovery of row 1's apparent gain by a zero-text entity mean under row 1's own split:
\textbf{139\%} at one month, \textbf{148\%} at three.}} \\
\multicolumn{7}{@{}p{0.95\textwidth}@{}}{\emph{Minimum detectable effect at 80\% power: entity stage 0.08\% at both horizons; AR stage 0.18\% at one month and 0.40\% at three.}} \\
\bottomrule
\end{tabular}
\end{table}

Row 1 is the field-standard design: pool the business-months, split at
random, and text appears to cut mean squared error by 28.43 per cent at one
month and 36.64 at three. Under that same split a predictor with no text in
it at all (each business's own mean rating) recovers 139 and 148 per
cent of that gain. % src: results/tables/yelp\_cascade.csv row-1 note; yelp\_cascade.md HONEST HEADLINE
The apparent gain is therefore not merely contaminated by entity identity;
it is more than fully accounted for by it, with the text arm delivering
less than the identity term alone. This is the exact cross-domain mirror of
the SEC result, obtained on a corpus where consumer review text is
uncontroversially informative. It is the strongest single piece of
evidence in the report that the identity artefact is a property of
entity-panel designs rather than of financial disclosure.

What the chronological protocol then shows is the second half of the story
and the reason the identity control cannot be sold as free hygiene. Adding
the entity mean to the reference \emph{costs} 10.58 and 9.48 per cent here,
where on the SEC panel it removed roughly half the increment. A rating
history already encodes the business, so the control's sign and size are
properties of the panel. And a small content residual nonetheless survives
the full reference: $+0.37$ per cent at one month, against the entity-stage
minimum detectable effect of 0.08 per cent that governs this row (0.18 at the
AR stage above it), placebo-clean. The three-month residual of $+0.61$ per cent
carries the same 0.08 per cent entity-stage threshold (0.40 at the AR stage)
and a borderline within-month swap diagnostic, and is shown rather than claimed
for that reason. % src: results/tables/yelp\_cascade.md table note (MDE by stage and horizon)
Two distinct probe shares circulate for this panel and must not be
substituted for one another: the 139 and 148 per cent above are the
zero-text entity mean against \emph{row 1's field-design gain}. The
149 and 103 per cent quoted in Section~\ref{sec:val-external} are a
prompted model's zero-content probe against \emph{that model's own
chronological combination increment}. Different numerators, different
denominators, different arms. % src: results/tables/yelp\_multiarm.md Reading block (149\%/103\%) vs yelp\_cascade.csv row-1 notes (139\%/148\%)

\clearpage

\begin{figure}[tbp]
\centering
\includegraphics[width=\textwidth,height=0.83\textheight,keepaspectratio]{figures/FP2_yelp_ladder__a.pdf}
% caption numbers src: results/tables/yelp_cascade.csv (all five cascade rows at both horizons; the recovery ratios are the committed row-1 note field); results/tables/yelp_cascade.md (the pre-registered minimum detectable effects and the reproduction gates); results/tables/yelp_entity_disjoint.csv (the split-rule experiment, its month-clustered p-values, entity coverage and row counts)
\caption[The second-domain ladder, in full]{
The second-domain ladder, in full:
the same audit run on a domain in which text is supposed to work. 
Panel~(a)
holds the three large moves. Under the field-standard pooled random split the
text model appears to cut mean squared error by 28.43 per cent at one month
and 36.64 at three. The outline bars show what a \emph{zero-text} business
mean recovers under that same split: 139 and 148 per cent of the apparent
gain, so the design's own headline is more than fully explained by entity
identity before any text is read. % src: results/tables/yelp_cascade.csv (row 1 delta_rel_pct 28.43 and 36.64; the 139 and 148 per cent ratios are the committed row-1 note field)
Under the chronological protocol text alone loses to the recalibrated
autoregressive reference ($-24.70$ and $-35.29$ per cent), and inserting the
zero-text business mean into that reference costs 10.58 and 9.48 per cent.
}
\label{fig:app-yelp-ladder}
\end{figure}

\begin{figure}[tbp]
\centering
\includegraphics[width=\textwidth,height=0.83\textheight,keepaspectratio]{figures/FP2_yelp_ladder__b.pdf}
% caption numbers src: results/tables/yelp_cascade.csv (all five cascade rows at both horizons; the recovery ratios are the committed row-1 note field); results/tables/yelp_cascade.md (the pre-registered minimum detectable effects and the reproduction gates); results/tables/yelp_entity_disjoint.csv (the split-rule experiment, its month-clustered p-values, entity coverage and row counts)
\caption[The second-domain ladder, in full (b)]{Continued from Figure~\ref{fig:app-yelp-ladder}. The second-domain ladder, in full:
the same audit run on a domain in which text is supposed to work. 
Panel~(b) is the same evidence on an axis a hundred times finer: the two
surviving increments, each against its own pre-registered minimum detectable
effect at 80 per cent power (0.18 and 0.08 per cent at one month, 0.40 and
0.08 at three). % src: results/tables/yelp_cascade.csv (rows 3 and 5: +0.346 and +0.365 at h = 1 m, +0.871 and +0.612 at h = 3 m); results/tables/yelp_cascade.md (table note: the four minimum detectable effects)
}
\label{fig:app-yelp-ladder-b}
\end{figure}

\begin{figure}[tbp]
\centering
\includegraphics[width=\textwidth,height=0.83\textheight,keepaspectratio]{figures/FP2_yelp_ladder__c.pdf}
% caption numbers src: results/tables/yelp_cascade.csv (all five cascade rows at both horizons; the recovery ratios are the committed row-1 note field); results/tables/yelp_cascade.md (the pre-registered minimum detectable effects and the reproduction gates); results/tables/yelp_entity_disjoint.csv (the split-rule experiment, its month-clustered p-values, entity coverage and row counts)
\caption[The second-domain ladder, in full (c)]{Continued from Figure~\ref{fig:app-yelp-ladder}. The second-domain ladder, in full:
the same audit run on a domain in which text is supposed to work. 
Panel~(c) changes one thing only, the split rule, over the same universe and
audit. Basis: squared error on stars; month-clustered Diebold--Mariano with a
Newey--West correction at lag $h-1$ months and the small-sample adjustment,
business-by-month two-way clustering as robustness; combiner weights fitted on
validation and frozen on test, business means from training and validation
observations only. Family F15 is declared per gate and governs panels~(a)
and~(b); 
panel~(c) changes the split rule and is scored cell by cell with no
Holm correction across its rows, a diagnostic of the design rather than a
survivor claim.
}
\label{fig:app-yelp-ladder-c}
\end{figure}

Figures~\ref{fig:app-yelp-ladder}, \ref{fig:app-yelp-ladder-b} and~\ref{fig:app-yelp-ladder-c} carry the one number the second-domain
section of Chapter~\ref{ch:validation} states in neither of its halves: that a
zero-text business mean, under the field-standard split, recovers 139 and 148
per cent of the apparent gain that split produces. That is the exact
cross-domain mirror of the disclosure finding. The chapter shows the identity
control as a \emph{cost} under the chronological protocol, which is the
conservative half of the story. The recovery ratio is the aggressive half, and
it says that the published-style headline in this second domain is more than
fully accounted for by knowing which business is being reviewed, before a word
of review text is read.

Panel~(b) is included because without it the surviving residual would be
unreadable. The two increments that clear the audit are of the order of a
third of one per cent: roughly a hundredth of the field design's headline.
The only way to see that they are nonetheless above what the design
could have detected is to draw each against its own pre-registered minimum
detectable effect. Both clear it at one month; at three months the
autoregressive-stage residual clears a threshold four times larger, which the
panel prints rather than leaves implicit.

Panel~(c) is the transferable methodological point and the reason the
experiment was run. The field's standard remedy for entity leakage, an
entity-disjoint split, is a statement about the researcher's own data. It
works exactly as intended on the fitted pipeline, whose entity memory came
from the training split and disappears with it: at three months that arm falls
from $+0.534$ per cent ($p=.012$) to $+0.152$ per cent ($p=.394$). It does
nothing at all to a prompted model whose entity prior was acquired from a
pretraining corpus no split of this panel can touch. A probe carrying the
business name, city, categories and month, and no review text whatever,
remains significant at both horizons after the fix, while the share of test
entities for which the entity-mean control has a mean falls from 83.4 per cent
to zero. % src: results/tables/yelp_entity_disjoint.csv (TF-IDF fitted h=3m under both split rules; the zero-content 70B probe at p = .0075 and .0011 under split B; entity_coverage_test 0.834 and 0.000)
The detector the report recommends in its place is cheap and is used
throughout this appendix (an arm-matched partial-input probe), and its
results for the disclosure panel are
Figures~\ref{fig:app-elicitation-b} and~\ref{fig:app-elicitation-c},
and panels~(c) and~(d) of Figure~\ref{fig:app-prompted-families-c}.

Three scope limits. Because the split rule reassigns rows, the test panels of
panel~(c) differ in size: 28,134 rows against 7,183 at one month and 38,399
against 9,843 at three. % src: results/tables/yelp_entity_disjoint.csv (n_test by split and horizon)
The claim there is the contrast between the two
splits and not the magnitude of either; every absolute effect in that
panel is below 0.6 per cent.
The figure shows no quantity from the main disclosure panel, no QLIKE and no
standalone loss for the entity-mean control, and the outline bars of panel~(a)
are the committed recovery ratios applied to the committed apparent gain
rather than separately reported losses. And the scope is one domain and one
prompted family: the claim is existence and mechanism, not a law about how
large the blind spot is.

\clearpage

\FloatBarrier
\section{Standalone Accuracy on the Event-Driven and Combined Panels}
\label{app:res-standalone-other}

Table~\ref{tab:app-leaderboard} is a long-form table, because the
efficiency-frontier artefact it is read from carries long-form levels and
GPU-hours only. The census of Table~\ref{tab:app-dm180} tests all three
panels but reports test statistics rather than levels. A reader who
wants to know \emph{how badly} a text arm loses on 8-K filings, or on the
two channels pooled, has nowhere in the report to look. Tables
\ref{tab:app-standalone-ed} and~\ref{tab:app-standalone-comb} close that
gap: every arm carrying a committed seed-aggregate row, at every horizon,
with the A2 HAR-RV reference printed at the head of each table so the
comparison is direct.

Three readings survive the change of panel. First, the ordering is the
ordering of each panel's own best QLIKE, with the reference row pulled to the
top: three price arms lead, the price-plus-embedding D3
fusions next, the frozen 7--8B embedding arms after them, and the
fine-tuned encoders and dictionary pipelines at the bottom. Second, no text
arm reaches the reference on either panel: the best pure-text figure is
0.4732 on the event-driven panel (C5 gte-Qwen2 at $h{=}20$, against the
reference's 0.3185) and 0.6104 on the combined panel (the same arm at the
same horizon, against 0.3151). Third, and most directly, the
coefficient-of-determination columns are negative in all 45 B- and C-block
cells on \emph{each} panel, exactly as they are on long-form. A text-only
forecast is worse than predicting the test-period mean, on every channel,
at every horizon, for every representation tried. The only non-price arms
with a positive test $R^2$ anywhere are the fusions carrying the price
signal inside them, and even there one cell is negative (D2 gated fusion,
combined, $h{=}20$, at $-0.005$).

\begin{table}[p]
\centering
\caption[Standalone accuracy, event-driven panel]{Standalone accuracy on
the \textbf{event-driven (8-K) panel}: test QLIKE in \emph{variance units}
and test $R^2$ for all 24 arms carrying a committed seed-aggregate row,
with the A2 HAR-RV reference at the head of the table. C- and
D-block figures are the mean over the three training seeds of each run's
own test metric (a seed-mean of per-seed metrics, not the metric of the
seed-ensemble forecast: an off-basis quantity, tagged). A- and B-block
rows are seed-invariant, and $s$ is the number of seeds averaged. Ordered
by best QLIKE over the three horizons, the ordering of
Table~\ref{tab:app-leaderboard}. Support 25{,}109/25{,}001/24{,}732
observations at $h = 5/10/20$. A6, A7, C5x, C6 and D4 do not appear because
the seed-aggregate artefact does not carry them; their event-driven
increments appear instead in Tables~\ref{tab:app-grid-ed}
and~\ref{tab:app-crossfamily}. Lower QLIKE is better; a negative $R^2$
means the arm is beaten by the test-period mean.}
\label{tab:app-standalone-ed}
\singlespacing\small
\setlength{\tabcolsep}{4pt}
% src: results/tables/seed_aggregate.csv (qlike_mean, r2_mean, n_seeds); results/runs/A2_har_rv_full_event_driven_seed2026/metrics.json (reference row)
\begin{tabular}{@{}llrrrrrrr@{}}
\toprule
& & & \multicolumn{3}{c}{Test QLIKE (variance units)} & \multicolumn{3}{c}{Test $R^2$} \\
\cmidrule(lr){4-6}\cmidrule(lr){7-9}
Model & Bl. & $s$ & $h{=}5$ & $h{=}10$ & $h{=}20$ & $h{=}5$ & $h{=}10$ & $h{=}20$ \\
\midrule
A2 HAR-RV \emph{(reference)} & A & --- & $0.5902$ & $0.4408$ & $0.3185$ & $+0.189$ & $+0.231$ & $+0.277$ \\
\addlinespace
A3 GARCH & A & 1 & $0.5095$ & $0.3735$ & $0.2660$ & $+0.098$ & $+0.160$ & $+0.193$ \\
A5 ARIMA & A & 1 & $0.6112$ & $0.4528$ & $0.3269$ & $+0.026$ & $+0.115$ & $+0.190$ \\
D3 price+gte-Qwen2 & D & 3 & $0.6800$ & $0.4952$ & $0.3455$ & $+0.168$ & $+0.210$ & $+0.266$ \\
D3 price+Qwen3 & D & 3 & $0.7326$ & $0.5072$ & $0.3518$ & $+0.141$ & $+0.203$ & $+0.256$ \\
D3 price+E5-Mistral & D & 3 & $0.7441$ & $0.5116$ & $0.3534$ & $+0.137$ & $+0.203$ & $+0.258$ \\
A1 Historical vol. & A & 1 & $1.3619$ & $0.6108$ & $0.3781$ & $-0.339$ & $-0.124$ & $-0.021$ \\
A4 EGARCH & A & 1 & $0.7161$ & $0.6118$ & $0.4595$ & $-0.135$ & $-0.083$ & $-0.012$ \\
C5 gte-Qwen2 & C & 3 & $0.8081$ & $0.6364$ & $0.4732$ & $-0.023$ & $-0.021$ & $-0.014$ \\
C5 Qwen3-emb. & C & 3 & $0.8133$ & $0.6284$ & $0.4835$ & $-0.021$ & $-0.012$ & $-0.011$ \\
D2 Gated fusion & D & 3 & $0.9571$ & $0.6791$ & $0.4974$ & $+0.050$ & $+0.104$ & $+0.126$ \\
D1 Concat-MLP & D & 3 & $0.9892$ & $0.7781$ & $0.5298$ & $+0.042$ & $+0.041$ & $+0.092$ \\
C5 E5-Mistral & C & 3 & $0.9963$ & $0.7649$ & $0.5885$ & $-0.076$ & $-0.065$ & $-0.076$ \\
C2 FinBERT (S1) & C & 3 & $0.9763$ & $1.1358$ & $0.7024$ & $-0.032$ & $-0.156$ & $-0.084$ \\
C2 FinBERT (S2) & C & 3 & $1.4556$ & $1.1665$ & $0.7287$ & $-0.171$ & $-0.161$ & $-0.115$ \\
C1 BERT (S1) & C & 3 & $1.2640$ & $1.0992$ & $0.7859$ & $-0.123$ & $-0.135$ & $-0.127$ \\
B4 L--M features & B & 1 & $1.4165$ & $1.0774$ & $0.8171$ & $-0.209$ & $-0.203$ & $-0.230$ \\
C2 FinBERT (S3) & C & 3 & $1.1119$ & $0.8271$ & $0.9199$ & $-0.059$ & $-0.033$ & $-0.227$ \\
C4 Longformer (S5) & C & 3 & $1.2275$ & $0.9733$ & $0.8356$ & $-0.088$ & $-0.098$ & $-0.153$ \\
B3 L--M linear & B & 1 & $1.4582$ & $1.1248$ & $0.8569$ & $-0.219$ & $-0.221$ & $-0.253$ \\
C2 FinBERT (S4) & C & 3 & $1.3306$ & $0.9871$ & $0.8596$ & $-0.121$ & $-0.094$ & $-0.176$ \\
C1 BERT (S2) & C & 3 & $1.5037$ & $1.2292$ & $0.9014$ & $-0.158$ & $-0.177$ & $-0.177$ \\
C3 RoBERTa (S1) & C & 3 & $1.2291$ & $1.1132$ & $0.9290$ & $-0.086$ & $-0.149$ & $-0.211$ \\
B2 TF--IDF ridge & B & 1 & $1.6632$ & $1.3636$ & $1.1125$ & $-0.222$ & $-0.252$ & $-0.325$ \\
B1 BoW ridge & B & 1 & $1.7707$ & $1.4092$ & $1.1743$ & $-0.570$ & $-0.624$ & $-0.538$ \\
\bottomrule
\end{tabular}
\end{table}

\begin{table}[p]
\centering
\caption[Standalone accuracy, combined panel]{Standalone accuracy on the
\textbf{combined panel} (long-form and event-driven filings pooled):
columns, conventions, seed basis and ordering exactly as in
Table~\ref{tab:app-standalone-ed}. Support 33{,}060/32{,}934/32{,}634
observations at $h = 5/10/20$. The combined panel is not a third
disclosure channel but the union of the other two, so its levels are a
row-weighted blend of theirs and must not be read as independent evidence.
It is tabulated because the census of Table~\ref{tab:app-dm180} tests on it
and the report otherwise never states what is being tested.}
\label{tab:app-standalone-comb}
\singlespacing\small
\setlength{\tabcolsep}{4pt}
% src: results/tables/seed_aggregate.csv (qlike_mean, r2_mean, n_seeds); results/runs/A2_har_rv_full_combined_seed2026/metrics.json (reference row)
\begin{tabular}{@{}llrrrrrrr@{}}
\toprule
& & & \multicolumn{3}{c}{Test QLIKE (variance units)} & \multicolumn{3}{c}{Test $R^2$} \\
\cmidrule(lr){4-6}\cmidrule(lr){7-9}
Model & Bl. & $s$ & $h{=}5$ & $h{=}10$ & $h{=}20$ & $h{=}5$ & $h{=}10$ & $h{=}20$ \\
\midrule
A2 HAR-RV \emph{(reference)} & A & --- & $0.5823$ & $0.4448$ & $0.3151$ & $+0.183$ & $+0.217$ & $+0.269$ \\
\addlinespace
A3 GARCH & A & 1 & $0.5031$ & $0.3739$ & $0.2667$ & $+0.025$ & $+0.084$ & $+0.103$ \\
A5 ARIMA & A & 1 & $0.5921$ & $0.4417$ & $0.3139$ & $+0.002$ & $+0.091$ & $+0.164$ \\
D3 price+gte-Qwen2 & D & 3 & $0.6831$ & $0.4893$ & $0.3365$ & $+0.155$ & $+0.207$ & $+0.267$ \\
D3 price+Qwen3 & D & 3 & $0.7167$ & $0.5126$ & $0.3478$ & $+0.136$ & $+0.193$ & $+0.244$ \\
D3 price+E5-Mistral & D & 3 & $0.7365$ & $0.5204$ & $0.3562$ & $+0.129$ & $+0.190$ & $+0.240$ \\
A1 Historical vol. & A & 1 & $1.3099$ & $0.5986$ & $0.3645$ & $-0.566$ & $-0.284$ & $-0.085$ \\
A4 EGARCH & A & 1 & $0.6493$ & $0.5472$ & $0.4281$ & $-0.170$ & $-0.135$ & $-0.059$ \\
D1 Concat-MLP & D & 3 & $0.9651$ & $0.6770$ & $0.5388$ & $+0.039$ & $+0.085$ & $+0.071$ \\
C5 gte-Qwen2 & C & 3 & $1.0738$ & $0.8463$ & $0.6104$ & $-0.112$ & $-0.104$ & $-0.106$ \\
C5 Qwen3-emb. & C & 3 & $1.0730$ & $0.8438$ & $0.6213$ & $-0.110$ & $-0.097$ & $-0.109$ \\
D2 Gated fusion & D & 3 & $1.0080$ & $0.7004$ & $0.6445$ & $+0.021$ & $+0.076$ & $-0.005$ \\
C5 E5-Mistral & C & 3 & $1.2412$ & $0.9687$ & $0.7167$ & $-0.165$ & $-0.151$ & $-0.175$ \\
C2 FinBERT (S4) & C & 3 & $1.6165$ & $1.0603$ & $0.7423$ & $-0.222$ & $-0.123$ & $-0.111$ \\
C4 Longformer (S5) & C & 3 & $1.5939$ & $1.0328$ & $0.7616$ & $-0.209$ & $-0.116$ & $-0.110$ \\
C1 BERT (S1) & C & 3 & $1.5902$ & $0.7889$ & $0.9343$ & $-0.189$ & $-0.037$ & $-0.213$ \\
B4 L--M features & B & 1 & $1.3986$ & $1.0961$ & $0.7920$ & $-0.215$ & $-0.204$ & $-0.227$ \\
C2 FinBERT (S1) & C & 3 & $1.3786$ & $0.8203$ & $0.8455$ & $-0.134$ & $-0.043$ & $-0.151$ \\
B3 L--M linear & B & 1 & $1.4403$ & $1.1367$ & $0.8351$ & $-0.227$ & $-0.220$ & $-0.254$ \\
C3 RoBERTa (S1) & C & 3 & $1.4075$ & $1.2695$ & $0.8631$ & $-0.160$ & $-0.213$ & $-0.191$ \\
C1 BERT (S2) & C & 3 & $1.6921$ & $1.1525$ & $0.8803$ & $-0.217$ & $-0.158$ & $-0.166$ \\
C2 FinBERT (S2) & C & 3 & $1.5138$ & $1.1488$ & $0.8834$ & $-0.157$ & $-0.155$ & $-0.187$ \\
B1 BoW ridge & B & 1 & $1.5852$ & $1.2233$ & $0.9022$ & $-0.322$ & $-0.297$ & $-0.355$ \\
C2 FinBERT (S3) & C & 3 & $1.4810$ & $1.2171$ & $0.9043$ & $-0.161$ & $-0.158$ & $-0.175$ \\
B2 TF--IDF ridge & B & 1 & $1.5213$ & $1.2359$ & $0.9631$ & $-0.202$ & $-0.211$ & $-0.268$ \\
\bottomrule
\end{tabular}
\end{table}

Table~\ref{tab:app-census} records that the variance-unit restatement is
the one convention under which any challenger beats the reference, and that
all seven winners are price arms. Table~\ref{tab:app-varwinners} names
them, and names the seven further comparisons carrying a favourable point
estimate without reaching significance, so that the ``DM $<0$'' count of 14
in that table can be audited cell by cell. The pattern is that the entire
variance-unit reversal is a GARCH-family result concentrated on the panels
containing 8-K filings: A3 GARCH takes all six event-driven and combined
winners and A4 EGARCH the single long-form one. The seven
non-significant favourable estimates are those two arms plus A5 ARIMA. No
text or fusion arm appears in the table at all, which is the 0-of-153 count
of Table~\ref{tab:app-census} written out.

\begin{table}[htbp]
\centering
\caption[The fourteen favourable variance-unit comparisons]{Every
comparison in the 180-cell standalone census whose day-clustered DM
statistic is \emph{negative} under the variance-unit QLIKE restatement;
that is, every comparison in the census whose point estimate beats the
A2 HAR-RV reference under any of the three loss conventions.
Fourteen of 180, of which seven are Holm-significant within their
20-comparison panel family. \emph{Negative means the challenger is better}:
the census sign convention is preserved and only the loss changes, so the
reading is the opposite of Table~\ref{tab:app-dm180}'s. The last two
columns give the same cell's statistic under the two conventions that
produce no winners at all, so the reader can see how far the reversal
travels. No text or fusion arm appears.}
\label{tab:app-varwinners}
\singlespacing\small
\setlength{\tabcolsep}{4pt}
% src: results/tables/variance_unit_standalone180.csv (dm_qlike_var_clu, holm_qlike_var_clu, better_qlike_var_holm, dm_se_clu, dm_qlike_vol_clu)
\begin{tabular}{@{}llrrrcrr@{}}
\toprule
& & & \multicolumn{3}{c}{QLIKE, variance units} & \multicolumn{2}{c}{Same cell, other conventions} \\
\cmidrule(lr){4-6}\cmidrule(lr){7-8}
Panel & $h$ & Challenger & DM & Holm $p$ & Wins & DM (sq.\ err.) & DM (QLIKE, vol.) \\
\midrule
Long-form & 5 & A4 EGARCH & $-2.92$ & .0107 & \textbf{Y} & $+2.74$ & $+2.46$ \\
Long-form & 5 & A3 GARCH & $-1.69$ & .1846 & -- & $+5.22$ & $+4.50$ \\
Long-form & 5 & A5 ARIMA & $-0.10$ & .9216 & -- & $+5.84$ & $+4.74$ \\
Long-form & 10 & A4 EGARCH & $-1.72$ & .2593 & -- & $+2.39$ & $+2.15$ \\
Long-form & 10 & A3 GARCH & $-0.79$ & .8577 & -- & $+4.73$ & $+3.95$ \\
Long-form & 10 & A5 ARIMA & $-0.46$ & .8577 & -- & $+4.21$ & $+2.75$ \\
Long-form & 20 & A5 ARIMA & $-0.51$ & 1.0000 & -- & $+3.10$ & $+1.65$ \\
Event-driven & 5 & A3 GARCH & $-4.71$ & $<$.0001 & \textbf{Y} & $+4.33$ & $+2.15$ \\
Event-driven & 10 & A3 GARCH & $-3.91$ & .0003 & \textbf{Y} & $+3.37$ & $+0.79$ \\
Event-driven & 20 & A3 GARCH & $-3.91$ & .0004 & \textbf{Y} & $+2.58$ & $-0.30$ \\
Combined & 5 & A3 GARCH & $-4.34$ & $<$.0001 & \textbf{Y} & $+6.08$ & $+3.11$ \\
Combined & 10 & A3 GARCH & $-3.59$ & .0010 & \textbf{Y} & $+4.78$ & $+1.50$ \\
Combined & 20 & A3 GARCH & $-3.61$ & .0013 & \textbf{Y} & $+3.49$ & $+0.11$ \\
Combined & 20 & A5 ARIMA & $-0.02$ & 1.0000 & -- & $+2.27$ & $+1.88$ \\
\bottomrule
\end{tabular}
\end{table}

\clearpage
\FloatBarrier
\section{Seed Dispersion Behind the Ensemble Basis}
\label{app:res-seeds}

The primary basis is a three-seed ensemble, and the report justifies that
choice with a single count: 37 of 144 trained cells flip the sign of their
DM statistic across seeds. Table~\ref{tab:app-seeds} is the evidence that
count is read from, restricted to the 33 cells where the choice bites:
the three-seed C- and D-block rows of the 69-cell combination grid. The
other 36 grid cells are single-seed by construction, the deterministic B-block and
prompted arms having no training seed to vary, and cannot disperse. The
remaining rows of the source file are the same trained arms scored on the
combined panel and at grid positions outside the primary family, and they
are summarised in the caption rather than printed.

Two things are visible here that a count cannot convey. The first is the
size of the dispersion relative to the effect: long-form C1 BERT at
$h{=}20$ reports $+2.95$, $+2.09$ and $-0.83$ per cent on the three seeds
(a spread wider than most of the increments the grid credits). Its
DM statistic runs $-15.11$, $-4.40$, $+14.88$. A single-seed study would
have reported that cell as a strong positive or a strong negative depending
only on which integer it fixed. The second is that dispersion is not
uniform: it concentrates in the fine-tuned encoders, whose every weight is
trained, and is smaller in the fusion arms whose price channel is
deterministic. Of the 14 cells that seed 2026 alone would have declared
genuine, 5 keep a significant same-sign increment in all three seeds and 12
in at least two. Eight of the 14 contain a seed that disagrees on the
sign of the test statistic. Averaging the forecasts before testing, rather
than the test statistics after, is the response the study adopts, and it is
a conservative one: it removes the favourable draws along with the
unfavourable.

\begin{table}[p]
\centering
\caption[Seed dispersion over the three-seed grid cells]{Seed dispersion
over the 33 three-seed C- and D-block cells of the 69-cell combination
grid. Each row re-runs the whole combination protocol from scratch on one
training seed: relative QLIKE improvement over the recalibrated HAR
reference and its observation-level HAC DM statistic --- the superseded
error structure of Table~\ref{tab:app-census} panel (b), under which the
seed campaign was scored --- per seed, with Holm applied
inside each seed's own C/D family. \emph{Negative DM and positive rel.\%
mean text improves}, the convention of Table~\ref{tab:app-grid-lf}.
``sig.'' counts the seeds in which the cell is Holm-significant with a
favourable sign; ``flip'' marks cells whose DM statistic changes sign
across the three seeds; ``gen.'' marks the cells the single-seed 2026 grid
would have declared genuine. This is a \emph{single-seed} table by
construction: it is the object the ensemble basis was chosen against,
not a restatement of it. None of its columns may be set beside the
ensemble figures of Table~\ref{tab:app-grid-lf}, on basis or on
inference: the DM values here are observation-level, those there are
day-clustered. Over all 144 trained cells
in the source file: 37 flip sign, 62 are Holm-significant with a favourable
sign in all three seeds, and of the 14 genuine on seed 2026 alone, 5 hold
in three seeds and 12 in at least two.}
\label{tab:app-seeds}
\singlespacing\footnotesize
\setlength{\tabcolsep}{1.5pt}
% src: results/tables/m1_multiseed.csv (per-seed rel_impr_pct and dm_q, rel_mean, rel_std, n_sig_holm, sign_disagree, grid2026_genuine); cell membership from results/tables/control_intersection_ensemble.csv (n_seeds = 3)
\begin{tabular}{@{}llrrrrrrrrrrcc@{}}
\toprule
& & & \multicolumn{5}{c}{Relative improvement (\%) by seed} & \multicolumn{3}{c}{Obs.-level DM by seed} & & & \\
\cmidrule(lr){4-8}\cmidrule(lr){9-11}
Panel & Model & $h$ & 2026 & 2027 & 2028 & mean & s.d. & 2026 & 2027 & 2028 & sig. & flip & gen. \\
\midrule
Long-form & C1 BERT (S1) & 5 & $+0.12$ & $+3.71$ & $+2.64$ & $+2.16$ & $1.85$ & $-0.96$ & $-6.91$ & $-7.55$ & 2/3 & -- & -- \\
Long-form & C1 BERT (S1) & 10 & $-1.28$ & $+0.67$ & $-0.97$ & $-0.53$ & $1.05$ & $+4.37$ & $-6.77$ & $+4.85$ & 1/3 & \textbf{Y} & -- \\
Long-form & C1 BERT (S1) & 20 & $+2.95$ & $+2.09$ & $-0.83$ & $+1.40$ & $1.98$ & $-15.11$ & $-4.40$ & $+14.88$ & 2/3 & \textbf{Y} & Y \\
Long-form & C2 FinBERT (S1) & 5 & $+0.56$ & $+0.85$ & $+2.37$ & $+1.26$ & $0.97$ & $-5.57$ & $-6.23$ & $-8.35$ & 3/3 & -- & Y \\
Long-form & C2 FinBERT (S1) & 10 & $+4.56$ & $-1.09$ & $+0.94$ & $+1.47$ & $2.86$ & $-12.77$ & $+15.54$ & $-1.93$ & 1/3 & \textbf{Y} & Y \\
Long-form & C2 FinBERT (S1) & 20 & $-0.08$ & $-5.36$ & $+1.17$ & $-1.43$ & $3.47$ & $+5.98$ & $+6.82$ & $-7.44$ & 1/3 & \textbf{Y} & -- \\
Long-form & C2 FinBERT (S2) & 5 & $+1.73$ & $+0.90$ & $-0.90$ & $+0.58$ & $1.34$ & $-7.25$ & $-10.76$ & $+5.75$ & 2/3 & \textbf{Y} & Y \\
Long-form & C2 FinBERT (S2) & 10 & $+0.24$ & $-3.88$ & $+1.90$ & $-0.58$ & $2.98$ & $-7.92$ & $+7.58$ & $-10.11$ & 2/3 & \textbf{Y} & Y \\
Long-form & C2 FinBERT (S2) & 20 & $-5.19$ & $+3.64$ & $+0.79$ & $-0.25$ & $4.50$ & $+7.93$ & $-11.42$ & $-2.71$ & 1/3 & \textbf{Y} & -- \\
Long-form & C2 FinBERT (S3) & 5 & $+2.39$ & $+1.94$ & $+2.48$ & $+2.27$ & $0.29$ & $-9.85$ & $-7.50$ & $-6.35$ & 3/3 & -- & Y \\
Long-form & C2 FinBERT (S3) & 10 & $+3.03$ & $+0.05$ & $+4.99$ & $+2.69$ & $2.49$ & $-9.38$ & $-0.35$ & $-10.34$ & 2/3 & -- & Y \\
Long-form & C2 FinBERT (S3) & 20 & $-0.76$ & $-0.40$ & $-11.10$ & $-4.09$ & $6.08$ & $+2.07$ & $+0.71$ & $+11.76$ & 0/3 & -- & -- \\
Long-form & C2 FinBERT (S4) & 5 & $-1.09$ & $+2.43$ & $-0.74$ & $+0.20$ & $1.94$ & $+5.94$ & $-4.38$ & $+6.00$ & 1/3 & \textbf{Y} & -- \\
Long-form & C2 FinBERT (S4) & 10 & $+2.28$ & $-2.09$ & $-0.56$ & $-0.13$ & $2.22$ & $-10.36$ & $+5.98$ & $+1.28$ & 1/3 & \textbf{Y} & Y \\
Long-form & C2 FinBERT (S4) & 20 & $+3.45$ & $+5.73$ & $-0.13$ & $+3.02$ & $2.95$ & $-16.89$ & $-12.76$ & $+14.18$ & 2/3 & \textbf{Y} & Y \\
Long-form & C3 RoBERTa (S1) & 5 & $+0.67$ & $+0.31$ & $-0.26$ & $+0.24$ & $0.47$ & $-10.47$ & $-6.80$ & $+9.69$ & 2/3 & \textbf{Y} & Y \\
Long-form & C3 RoBERTa (S1) & 10 & $+1.98$ & $+1.47$ & $+0.42$ & $+1.29$ & $0.79$ & $-6.37$ & $-7.83$ & $-9.80$ & 3/3 & -- & Y \\
Long-form & C3 RoBERTa (S1) & 20 & $+0.24$ & $-0.14$ & $+0.20$ & $+0.10$ & $0.21$ & $-6.89$ & $+4.41$ & $-7.22$ & 2/3 & \textbf{Y} & Y \\
Long-form & C4 Longformer (S5) & 5 & $+1.53$ & $+0.91$ & $+1.45$ & $+1.30$ & $0.33$ & $-8.73$ & $-8.30$ & $-9.91$ & 3/3 & -- & Y \\
Long-form & C4 Longformer (S5) & 10 & $-2.61$ & $-5.38$ & $+0.25$ & $-2.58$ & $2.81$ & $+16.84$ & $+16.64$ & $-9.41$ & 1/3 & \textbf{Y} & -- \\
Long-form & C4 Longformer (S5) & 20 & $-9.34$ & $-1.43$ & $-1.05$ & $-3.94$ & $4.68$ & $+8.43$ & $+2.21$ & $+16.20$ & 0/3 & -- & -- \\
Long-form & D1 Concat-MLP & 5 & $-0.01$ & $-2.60$ & $-0.12$ & $-0.91$ & $1.47$ & $+0.63$ & $+6.75$ & $+1.70$ & 0/3 & -- & -- \\
Long-form & D1 Concat-MLP & 10 & $+0.18$ & $-0.36$ & $-0.16$ & $-0.11$ & $0.27$ & $-1.68$ & $+5.87$ & $+1.07$ & 0/3 & \textbf{Y} & -- \\
Long-form & D1 Concat-MLP & 20 & $-0.08$ & $+0.04$ & $-0.83$ & $-0.29$ & $0.47$ & $+0.33$ & $-9.94$ & $+4.83$ & 1/3 & \textbf{Y} & -- \\
Long-form & D2 Gated fusion & 5 & $-0.12$ & $+0.90$ & $-0.29$ & $+0.16$ & $0.64$ & $+1.12$ & $-5.74$ & $+4.00$ & 1/3 & \textbf{Y} & -- \\
Long-form & D2 Gated fusion & 10 & $+0.13$ & $+0.11$ & $-0.76$ & $-0.17$ & $0.51$ & $-0.56$ & $-1.47$ & $+5.39$ & 0/3 & \textbf{Y} & -- \\
Long-form & D2 Gated fusion & 20 & $-6.00$ & $+1.70$ & $-1.47$ & $-1.92$ & $3.87$ & $+4.28$ & $-9.77$ & $+2.71$ & 1/3 & \textbf{Y} & -- \\
\addlinespace
Event-driven & C2 FinBERT (S1) & 5 & $+2.14$ & $+2.00$ & $+2.06$ & $+2.06$ & $0.07$ & $-14.36$ & $-16.09$ & $-10.82$ & 3/3 & -- & -- \\
Event-driven & C2 FinBERT (S1) & 10 & $+2.10$ & $+2.10$ & $+1.94$ & $+2.05$ & $0.10$ & $-14.19$ & $-11.84$ & $-9.38$ & 3/3 & -- & Y \\
Event-driven & C2 FinBERT (S1) & 20 & $+0.92$ & $+0.52$ & $+1.36$ & $+0.93$ & $0.42$ & $-2.03$ & $-2.02$ & $-9.69$ & 1/3 & -- & -- \\
Event-driven & D2 Gated fusion & 5 & $+0.51$ & $-0.00$ & $+0.87$ & $+0.46$ & $0.44$ & $-5.58$ & $+2.02$ & $-10.09$ & 2/3 & \textbf{Y} & -- \\
Event-driven & D2 Gated fusion & 10 & $+0.23$ & $+0.03$ & $+0.14$ & $+0.13$ & $0.10$ & $-2.22$ & $-2.69$ & $-2.87$ & 0/3 & -- & -- \\
Event-driven & D2 Gated fusion & 20 & $-2.87$ & $-0.59$ & $-2.32$ & $-1.93$ & $1.19$ & $+8.58$ & $+5.49$ & $+10.56$ & 0/3 & -- & -- \\
\bottomrule
\end{tabular}
\end{table}

\clearpage

\begin{figure}[tbp]
\centering
\includegraphics[width=\textwidth,height=0.83\textheight,keepaspectratio]{figures/stability_seed_dispersion__a.pdf}
% caption numbers src: results/tables/m1_multiseed.csv (144 disclosure x model x horizon cells, seeds 2026/2027/2028; per-seed rel%, DM and Holm p; sign_disagree; grid2026_genuine)
\caption[How much of a single-seed verdict is the seed]{
How much of a
single-seed verdict is the seed. Every fine-tuned text and fusion arm was
trained three times, at seeds 2026, 2027 and 2028, the combination grid
recomputed for each. Panels~(a) and~(b) cover all 144
disclosure-by-model-by-horizon cells (48 in each of the long-form,
event-driven and combined panels), the price and classical arms being
seed-invariant and excluded. 
Panel~(a) plots the across-seed
standard deviation of the text increment against its mean. Inside
the shaded wedge the spread exceeds the increment being reported,
true of 40 of the 144 cells, and the 37 cells whose Diebold--Mariano statistic
changes sign across seeds are marked separately. % src: results/tables/m1_multiseed.csv (rel_std > |rel_mean| in 40 of 144; sign_disagree in 37 of 144)
Panel~(b) localises the instability by block: the three frozen-embedding text
arms flip in 0 of 27 cells and the three frozen-embedding fusions in 3 of 27,
whereas each fully fine-tuned fusion arm flips in 6 of its 9. % src: results/tables/m1_multiseed.csv (sign_disagree grouped by model: C5 trio 0/27, D3 trio 3/27, D1 6/9, D2 6/9)
}
\label{fig:app-seed-dispersion}
\end{figure}

\begin{figure}[tbp]
\centering
\includegraphics[width=\textwidth,height=0.83\textheight,keepaspectratio]{figures/stability_seed_dispersion__b.pdf}
% caption numbers src: results/tables/m1_multiseed.csv (144 disclosure x model x horizon cells, seeds 2026/2027/2028; per-seed rel%, DM and Holm p; sign_disagree; grid2026_genuine)
\caption[How much of a single-seed verdict is the seed (c)]{Continued from Figure~\ref{fig:app-seed-dispersion}. How much of a
single-seed verdict is the seed. Every fine-tuned text and fusion arm was
trained three times, at seeds 2026, 2027 and 2028, the combination grid
recomputed for each. Panels~(a) and~(b) cover all 144
disclosure-by-model-by-horizon cells (48 in each of the long-form,
event-driven and combined panels), the price and classical arms being
seed-invariant and excluded. 
Panel~(c) zooms on the fourteen cells seed 2026 alone would have declared
genuine and shows each seed's own increment, filled where that seed is
Holm-significant \emph{and} favours text. Eight of the eleven hollow markers are
Holm-significant with the sign reversed; only three are genuinely
non-significant. % src: results/tables/m1_multiseed.csv (of the 11 marker-cells failing p_holm<0.05 AND dm_q<0: 8 have p_holm<0.05 with dm_q>0, min p 1.38e-51; 3 have p_holm>=0.05, at p 1, 1 and 0.702)
Of the 14, 5 hold in all three
seeds, 12 in at least two, two in only one. Basis: single-seed combination
grids, volatility-unit QLIKE, obs.-level HAC (superseded); denominators 144
in~(a) and~(b) and 14 in~(c); (a) and~(b) mark no significance.
}
\label{fig:app-seed-dispersion-b}
\end{figure}

Figures~\ref{fig:app-seed-dispersion} and~\ref{fig:app-seed-dispersion-b} are the evidence behind a number the
report states three times and never draws. Panel~(a) is the blunt version. In
40 of 144 cells the across-seed spread of the increment exceeds the increment
itself, so a single-seed study of this panel would be reporting, in more than
a quarter of its cells, a quantity smaller than its own run-to-run noise.
Panel~(b) says where that noise lives, and the ordering is intuitive. Arms
whose only trained component is a small regression head over frozen embeddings
are almost perfectly stable, and arms in which a full encoder and a fusion
head are trained together are the least stable of all. Panel~(c) gives the
consequence for verdicts rather than for estimates: of the fourteen cells a
seed-2026 study would have called genuine, five hold in all three seeds.

The figure is deliberately not a criticism of any particular single-seed study
and should not be read as one. It is the justification for a design choice
this project made and paid for: three seeds per trained arm across three
disclosure channels, and a per-observation ensemble as the primary forecast
object. That choice is what allows every count in
Chapters~\ref{ch:results} and~\ref{ch:validation} to be quoted without a seed
caveat, and its cost (roughly three times the training fleet) is the
larger part of the compute budget of Figure~\ref{fig:app-compute-accuracy}.

What the figure does not show is the primary basis itself. Every cell drawn
here is a single-seed combination grid, and no count in it appears anywhere
else in this report; the study's answer to the dispersion is the ensemble, and
the ensemble is elsewhere. Holm is applied within each seed's own family, so a
cell can move between seeds either because its own statistic moved or because
its neighbours' did, and panel~(c) therefore shows verdict stability rather
than estimate stability. The 144 cells also include the combined channel,
whose rows overlap both others, so the three panels of 48 are not independent
samples of the same quantity.

% The second of the two removals; see the note above Section E.2. Removing this
% one clears the 80 per cent white page that used to close Section E.4 and,
% because everything below moves up a page, the 48 per cent one that used to
% close Section E.5 as well.

\FloatBarrier
\section{Joint Tests Across the Whole Leaderboard}
\label{app:res-spamcs}

Every test in Tables~\ref{tab:app-dm180} and~\ref{tab:app-grid-lf} is
pairwise, and Holm controls error only across the pairs actually run. The
standing objection to a leaderboard study is aggregate data-snooping: with
two dozen challengers, the best-looking one is selected after the fact.
Hansen's test for superior predictive ability \cite{hansen2005spa} and the
model confidence set \cite{hansen2011mcs} answer that objection directly,
and Tables~\ref{tab:app-spa} and~\ref{tab:app-mcs} report them panel by
panel rather than as the two sentences the main text has room for.

Three readings matter, and the second is the one an examiner should press
on. First, the text-and-fusion SPA $p$-value is 1.000 in all 18 panels:
against the null that no text or fusion challenger beats HAR-RV, the data
provide no evidence whatever, on either loss and every channel. Second,
the \emph{full-set} SPA does reject HAR-RV in nine of the 18 panels, so
HAR-RV is not the best forecast available. In all nine the beating
challenger is a price model: HAR-X augmented with the VIX, semivariance
HAR, or HARQ. The leaderboard has genuine winners; none of them reads a
document. Third, no pure-text model enters the 90 per cent model confidence
set in any of the 18 panels, which is the aggregate form of the 0-of-180
pairwise count. Price-carrying fusion models do enter, in 24 slots across
12 panels, and that is a tie rather than a win. Membership means the model
cannot be separated from the best, and the SPA never rejects in a fusion
arm's favour anywhere.

\begin{table}[htbp]
\centering
\caption[Superior-predictive-ability tests, all 18 panels]{Hansen's test
for superior predictive ability \cite{hansen2005spa} and 90\% model
confidence set \cite{hansen2011mcs}, one row per loss-by-channel-by-horizon
panel. Losses are block-aggregated to one value per model per trading day,
matching the day-clustered primary; stationary bootstrap
\cite{politis1994stationary} over days with expected block length equal to
the horizon, $B = 2{,}000$. ``SPA $p$'' is the consistent $p$-value for the
null that HAR-RV is not inferior to the best of \emph{all} alternatives
(a \emph{small} value means some model beats HAR-RV), and ``text/fus.\
$p$'' is the same test restricted to the text-and-fusion block. The MCS
columns count survivors by block out of the models present. Membership in
the confidence set means ``cannot be statistically separated from the
best'', not ``beats HAR-RV''. The common inner-joined sample is used
throughout, so the day counts are the merged-grid ones and the levels
behind these tests are not the own-sample levels of
Table~\ref{tab:app-leaderboard}.}
\label{tab:app-spa}
\singlespacing\footnotesize
\setlength{\tabcolsep}{2.5pt}
% src: results/tables/row13_spa_mcs_panels.csv (spa_p_consistent, spa_best_challenger, spa_best_tstat, spa_textfusion_p_consistent, mcs90_* and n_* counts)
\begin{tabular}{@{}llrrrrcrlrl@{}}
\toprule
& & & & & \multicolumn{3}{c}{Full alternative set} & Text/fus. & \multicolumn{2}{c}{MCS90 price/text/fusion} \\
\cmidrule(lr){6-8}\cmidrule(lr){9-9}\cmidrule(lr){10-11}
Channel & $h$ & Loss & days & $K$ & SPA $p$ & Rej. & Best challenger ($t$) & $p$ & in set & of \\
\midrule
Long-form & 5 & Sq.\ error & 792 & 26 & 0.462 & -- & A7 HAR-X ($+1.05$) & 1.000 & 4/0/2 & 8/12/6 \\
Long-form & 5 & QLIKE (vol) & 792 & 26 & 0.385 & -- & A7 HAR-X ($+1.00$) & 1.000 & 5/0/1 & 8/12/6 \\
Long-form & 10 & Sq.\ error & 766 & 26 & 0.343 & -- & A7 HAR-X ($+1.22$) & 1.000 & 4/0/1 & 8/12/6 \\
Long-form & 10 & QLIKE (vol) & 766 & 26 & 0.350 & -- & A7 HAR-X ($+1.08$) & 1.000 & 7/0/4 & 8/12/6 \\
Long-form & 20 & Sq.\ error & 765 & 26 & 0.116 & -- & A7 HAR-X ($+1.78$) & 1.000 & 4/0/3 & 8/12/6 \\
Long-form & 20 & QLIKE (vol) & 765 & 26 & 0.150 & -- & A7 HAR-X ($+1.61$) & 1.000 & 7/0/4 & 8/12/6 \\
\addlinespace
Event-driven & 5 & Sq.\ error & 996 & 27 & 0.074 & -- & A7 HAR-X ($+1.94$) & 1.000 & 2/0/0 & 8/13/6 \\
Event-driven & 5 & QLIKE (vol) & 996 & 27 & 0.021 & \textbf{Y} & A7 HAR-X ($+2.58$) & 1.000 & 1/0/0 & 8/13/6 \\
Event-driven & 10 & Sq.\ error & 991 & 27 & 0.014 & \textbf{Y} & A6 HARQ ($+2.84$) & 1.000 & 2/0/0 & 8/13/6 \\
Event-driven & 10 & QLIKE (vol) & 991 & 27 & 0.021 & \textbf{Y} & A6 SHAR ($+2.60$) & 1.000 & 6/0/1 & 8/13/6 \\
Event-driven & 20 & Sq.\ error & 981 & 27 & 0.005 & \textbf{Y} & A6 HARQ ($+3.22$) & 1.000 & 3/0/0 & 8/13/6 \\
Event-driven & 20 & QLIKE (vol) & 981 & 27 & 0.002 & \textbf{Y} & A6 SHAR ($+3.58$) & 1.000 & 6/0/2 & 8/13/6 \\
\addlinespace
Combined & 5 & Sq.\ error & 996 & 26 & 0.080 & -- & A7 HAR-X ($+1.90$) & 1.000 & 3/0/0 & 8/12/6 \\
Combined & 5 & QLIKE (vol) & 996 & 26 & 0.030 & \textbf{Y} & A7 HAR-X ($+2.33$) & 1.000 & 1/0/0 & 8/12/6 \\
Combined & 10 & Sq.\ error & 991 & 26 & 0.052 & -- & A6 HARQ ($+2.18$) & 1.000 & 3/0/1 & 8/12/6 \\
Combined & 10 & QLIKE (vol) & 991 & 26 & 0.027 & \textbf{Y} & A6 SHAR ($+2.42$) & 1.000 & 6/0/2 & 8/12/6 \\
Combined & 20 & Sq.\ error & 981 & 26 & 0.015 & \textbf{Y} & A6 HARQ ($+2.69$) & 1.000 & 3/0/1 & 8/12/6 \\
Combined & 20 & QLIKE (vol) & 981 & 26 & 0.003 & \textbf{Y} & A6 SHAR ($+3.52$) & 1.000 & 6/0/2 & 8/12/6 \\
\bottomrule
\end{tabular}
\end{table}

\begin{table}[htbp]
\centering
\caption[Model-confidence-set membership, all 18 panels]{The 90\% model
confidence set \cite{hansen2011mcs} itself: which models survive in each of
the 18 panels. Every surviving model is either a price arm (A-block) or a
price-carrying fusion (D-block). The B- and C-blocks (bag-of-words,
dictionaries, every fine-tuned encoder, every frozen 7--8B embedding arm
and both prompted large-language-model arms) are absent from all 18
sets. This is the joint statement the pairwise census can only make one
pair at a time.}
\label{tab:app-mcs}
\singlespacing\footnotesize
\setlength{\tabcolsep}{4pt}
% src: results/tables/row13_spa_mcs.csv (in_mcs90 by model, disclosure, horizon and loss)
\begin{tabular}{@{}lllp{0.56\textwidth}@{}}
\toprule
Channel & $h$ & Loss & Members of the 90\% model confidence set \\
\midrule
Long-form & 5 & Sq.\ error & A2 HAR-RV; A6 HARQ; A6 SHAR; A7 HAR-X(VIX); D1 Concat-MLP; D3 price+gte-Qwen2 \\
Long-form & 5 & QLIKE (vol) & A2 HAR-RV; A4 EGARCH; A6 HARQ; A6 SHAR; A7 HAR-X(VIX); D1 Concat-MLP \\
Long-form & 10 & Sq.\ error & A2 HAR-RV; A6 HARQ; A6 SHAR; A7 HAR-X(VIX); D3 price+gte-Qwen2 \\
Long-form & 10 & QLIKE (vol) & A2 HAR-RV; A3 GARCH; A4 EGARCH; A5 ARIMA; A6 HARQ; A6 SHAR; A7 HAR-X(VIX); D2 Gated fusion; D3 price+E5-Mistral; D3 price+Qwen3; D3 price+gte-Qwen2 \\
Long-form & 20 & Sq.\ error & A2 HAR-RV; A6 HARQ; A6 SHAR; A7 HAR-X(VIX); D3 price+E5-Mistral; D3 price+Qwen3; D3 price+gte-Qwen2 \\
Long-form & 20 & QLIKE (vol) & A2 HAR-RV; A3 GARCH; A4 EGARCH; A5 ARIMA; A6 HARQ; A6 SHAR; A7 HAR-X(VIX); D2 Gated fusion; D3 price+E5-Mistral; D3 price+Qwen3; D3 price+gte-Qwen2 \\
\addlinespace
Event-driven & 5 & Sq.\ error & A6 HARQ; A7 HAR-X(VIX) \\
Event-driven & 5 & QLIKE (vol) & A7 HAR-X(VIX) \\
Event-driven & 10 & Sq.\ error & A6 HARQ; A7 HAR-X(VIX) \\
Event-driven & 10 & QLIKE (vol) & A2 HAR-RV; A3 GARCH; A4 EGARCH; A6 HARQ; A6 SHAR; A7 HAR-X(VIX); D3 price+gte-Qwen2 \\
Event-driven & 20 & Sq.\ error & A6 HARQ; A6 SHAR; A7 HAR-X(VIX) \\
Event-driven & 20 & QLIKE (vol) & A2 HAR-RV; A3 GARCH; A4 EGARCH; A6 HARQ; A6 SHAR; A7 HAR-X(VIX); D3 price+Qwen3; D3 price+gte-Qwen2 \\
\addlinespace
Combined & 5 & Sq.\ error & A6 HARQ; A6 SHAR; A7 HAR-X(VIX) \\
Combined & 5 & QLIKE (vol) & A7 HAR-X(VIX) \\
Combined & 10 & Sq.\ error & A6 HARQ; A6 SHAR; A7 HAR-X(VIX); D3 price+gte-Qwen2 \\
Combined & 10 & QLIKE (vol) & A2 HAR-RV; A3 GARCH; A4 EGARCH; A6 HARQ; A6 SHAR; A7 HAR-X(VIX); D3 price+Qwen3; D3 price+gte-Qwen2 \\
Combined & 20 & Sq.\ error & A6 HARQ; A6 SHAR; A7 HAR-X(VIX); D3 price+gte-Qwen2 \\
Combined & 20 & QLIKE (vol) & A2 HAR-RV; A3 GARCH; A4 EGARCH; A6 HARQ; A6 SHAR; A7 HAR-X(VIX); D3 price+Qwen3; D3 price+gte-Qwen2 \\
\bottomrule
\end{tabular}
\end{table}

Figures~\ref{fig:app-spa-mcs} and~\ref{fig:app-spa-mcs-b} draw the same
evidence: the membership matrix panel by panel, and the two
superior-predictive-ability $p$-value series beside it.

\clearpage

\begin{figure}[tbp]
\centering
\includegraphics[width=\textwidth,height=0.83\textheight,keepaspectratio]{figures/FP3_joint_spa_mcs__a.pdf}
% caption numbers src: results/tables/row13_spa_mcs.csv (474 rows: per-model, per-panel in_mcs90, block labels, benchmark flag); results/tables/row13_spa_mcs_panels.csv (18 rows: SPA p-values, the text-and-fusion-only p-value, per-block counts); results/tables/row13_spa_mcs.md (specification and the identity of the beating models)
\caption[Joint tests, panel by panel]{
Joint tests, panel by panel. Pairwise
Diebold--Mariano testing controls error only across the pairs actually tested;
the test for superior predictive ability \cite{hansen2005spa} and the model
confidence set \cite{hansen2011mcs} answer the aggregate data-snooping
objection instead. 
Panel~(a) is 90 per cent model-confidence-set membership
for every model on each of the 18 loss-by-disclosure-by-horizon panels. A
filled square means the model is in that panel's set, a dot that it was run
and excluded, a cross that it was not run on that disclosure, the 70B prompted
arm being 8-K only. The whole text band is empty (0 of its 222 cells),
while price arms occupy 73 of 144 and price-carrying fusion arms 24 of 108,
spread over 12 of the 18 panels. % src: results/tables/row13_spa_mcs.csv (in_mcs90 by block); results/tables/row13_spa_mcs_panels.csv (mcs90_fusion > 0 in 12 panels)
}
\label{fig:app-spa-mcs}
\end{figure}

\begin{figure}[tbp]
\centering
\includegraphics[width=\textwidth,height=0.83\textheight,keepaspectratio]{figures/FP3_joint_spa_mcs__b.pdf}
% caption numbers src: results/tables/row13_spa_mcs.csv (474 rows: per-model, per-panel in_mcs90, block labels, benchmark flag); results/tables/row13_spa_mcs_panels.csv (18 rows: SPA p-values, the text-and-fusion-only p-value, per-block counts); results/tables/row13_spa_mcs.md (specification and the identity of the beating models)
\caption[Joint tests, panel by panel (b)]{Continued from Figure~\ref{fig:app-spa-mcs}. Joint tests, panel by panel. Pairwise
Diebold--Mariano testing controls error only across the pairs actually tested;
the test for superior predictive ability \cite{hansen2005spa} and the model
confidence set \cite{hansen2011mcs} answer the aggregate data-snooping
objection instead. 
Panel~(b) plots each panel's consistent $p$-value on a logarithmic axis.
Against the entire alternative set the null that HAR-RV is not inferior is
rejected in 9 of 18 panels, and in all nine the beating model is a price arm:
semivariance HAR in four, HARQ in three, the VIX-augmented HAR-X in two.
Against the text-and-fusion block alone the $p$-value is 1.000 in all 18. % src: results/tables/row13_spa_mcs_panels.csv (spa_p_consistent < .05 in 9 of 18 with spa_best_challenger for those nine; spa_textfusion_p_consistent = 1.000 in 18 of 18)
Basis: a common inner-joined sample per panel, three-seed ensembles for the
neural arms, losses aggregated to one value per model per effective trading
day, a stationary bootstrap over days with expected block length $h$,
$B = 2{,}000$, model-confidence-set size 0.10. Neither panel carries a Holm family, and
neither needs one: the bootstrap controls error across each panel's whole
alternative set. The 18 panels are not corrected across.
}
\label{fig:app-spa-mcs-b}
\end{figure}

Figures~\ref{fig:app-spa-mcs} and~\ref{fig:app-spa-mcs-b} give the one view in this report in which the
strongest joint claim can be verified by eye rather than by counting a table.
The text band (thirteen arms, 222 cells --- twelve arms across all
eighteen panels and the 8-K-only 70B replication across six) contains
not one filled square. Chapter~\ref{ch:results} states that fact in a clause;
here it is a visible absence, and the absence is what an examiner can check in
a second. The fusion band beneath it is the necessary qualification: 24 of its
108 cells are filled, in 12 of the 18 panels, and every one of those arms
carries the price signal inside it. Membership means ``cannot be statistically
separated from the best model'', not ``beats HAR-RV'', so a fusion arm
entering the set is evidence that price plus text cannot be separated from the
best model and not that text contributes to the tie. Panel~(b) settles
which reading is right, the test for superior predictive ability never
rejecting in favour of the text-and-fusion block on any panel and returning a
$p$-value of exactly 1.000 in all eighteen.

What the figure cannot do is quantify. Confidence-set membership is a binary
and a $p$-value is not an effect size, so nothing here says how far behind the
text arms are; the leaderboard of Chapter~\ref{ch:results} and
Appendix~\ref{app:fullresults} carries that. The panel losses also sit on a
common inner-joined sample so that all models on a panel are compared on
identical rows, which is a smaller row set than each model's own evaluation
sample, so no level in the underlying file is interchangeable with a
leaderboard level. And the figure carries no combination-grid quantity at all:
it is a standalone-forecast object throughout, and its 0 of 222 is not the
same statement as the ladder's counts out of 69.

\FloatBarrier
\section{Detection Without Attribution}
\label{app:res-detect}

Chapter~\ref{ch:results} closes on an ordering (detectable, attributable,
realisable) and supports its first term with two numbers: 48 of 48
encompassing tests return a positive text coefficient, and a pooled test
over the 69-cell grid returns $t = +8.29$ with a minimum detectable effect
of 0.375 per cent. Neither is tabulated anywhere in the report. Tables
\ref{tab:app-encompass} to~\ref{tab:app-power} give both in full, because
``detectable'' is the one favourable verdict the study reaches and it
should be the most, not the least, auditable.

The encompassing tables are unanimous, and unanimity is itself a reason for caution: \cite{harvey1998encompassing} show that the least-squares form of this test over-rejects a true encompassing null when forecast errors are non-normal, which volatility forecast errors are. The 48 of 48 is therefore an upper bound, and the day-clustered Diebold--Mariano test remains the test of record. The tables should be read with their scope
attached. The record is split at the channel boundary: the long-form half is
Table~\ref{tab:app-encompass}, 33 regressions, and the event-driven half
Table~\ref{tab:app-encompass-ed}, the remaining 15. % src: results/tables/encompassing_regression.md (33 rows in the long_form block, 15 in the event_driven block, 48 in total)
The fitted text coefficient is positive and significantly so in
all 48 regressions, the smallest $t$ statistic anywhere being $+2.28$, and
the coefficients span $+0.052$ to $+2.882$. But the regression is fitted on
the test sample itself, so it measures a conditional correlation in sample
rather than an out-of-sample loss reduction, and it is scored with
observation-level HAC standard errors rather than the day-clustered primary.
Those are two departures from the basis, both tagged here and both working in
text's favour. Conditional on the HAR forecast the text forecasts carry
information; that is what detection means, and it is compatible with every
one of those arms losing to HAR-RV standalone and with the ladder
dissolving their increments.

The pooled omnibus is the stronger of the two instruments because it is an
out-of-sample loss test, and its power calibration is what turns its
rejection into a bounded claim. Averaging the daily loss differentials
across all 69 cells and testing the single resulting series gives $t =
+8.29$ over 996 days, with both subfamilies rejecting on their own.
Table~\ref{tab:app-omnibus} gives the three rows. The
implied cross-cell increment is about $+0.84$ per cent. Table~\ref{tab:app-power}
then injects a firm-orthogonal signal calibrated to a known size in every
cell simultaneously and records how often the same test finds it. The
recovery rate is 0.10 at an injected 0.1 per cent, 0.68 at 0.3 per cent and
1.00 from 0.5 per cent up, giving an 80 per cent-power minimum detectable
effect of 0.375 per cent. The pooled test detects; what it detects is a
systematic micro-increment: small in level, but not near the edge of what
the test could have seen. At 0.84 per cent it is about 2.2 times the
80 per cent-power minimum detectable effect of 0.375 per cent, and the
calibration recovers an injected 0.5 per cent at a rate of 1.00. % src: results/tables/omnibus_m1.md (all_69 approx rel 0.8428 per cent; empirical 80%-power MDE 0.375 per cent; MBB rejection rate 1.00 at the 0.5 per cent injection)

\begin{table}[p]
\centering
\caption[The 48 encompassing regressions, long-form half]{The complete
forecast-encompassing record \cite{harvey1998encompassing}, \textbf{long-form
channel}: 33 of the 48 regressions, the remaining 15 event-driven ones
following in Table~\ref{tab:app-encompass-ed}. Realised
volatility regressed on the HAR-RV forecast and the text forecast jointly,
$\mathrm{RV} = a + b f_{\mathrm{HAR}} + g f_{\mathrm{text}} + e$, one
regression per text-arm-by-channel-by-horizon cell. $g$ is the coefficient
on the text forecast conditional on the price forecast; a significantly
positive $g$ means HAR-RV does not encompass the text arm.
\textbf{Off-basis and tagged in two ways}: $a$, $b$ and $g$ are fitted on
the \emph{test} sample, so this is an in-sample conditional-correlation
statistic and not a validation-frozen out-of-sample loss reduction. The
$t$ statistics are observation-level HAC at lag $h-1$ observations, not the
day-clustered primary, so they are inflated relative to every clustered
table in this appendix. Both departures favour text, which is why the
unanimous result is reported as detection and never as attribution. Across
both halves all 48 coefficients are positive and significant at 5\%; the
smallest statistic anywhere in the 48 is the $+2.28$ in this table.}
\label{tab:app-encompass}
\singlespacing\footnotesize
\setlength{\tabcolsep}{4pt}
% src: results/tables/encompassing_regression.md (the long_form block, all 33 rows; the source's editorialising verdict column is not reproduced)
\begin{tabular}{@{}llrrrr@{}}
\toprule
Channel & Text arm & $h$ & $g$ & HAC $t$ & $p$ \\
\midrule
Long-form & B1 BoW ridge & 5 & $+0.134$ & $+4.76$ & $<$.0001 \\
Long-form & B1 BoW ridge & 10 & $+0.179$ & $+4.18$ & $<$.0001 \\
Long-form & B1 BoW ridge & 20 & $+0.199$ & $+4.25$ & $<$.0001 \\
\addlinespace
Long-form & B2 TF--IDF ridge & 5 & $+0.698$ & $+16.23$ & $<$.0001 \\
Long-form & B2 TF--IDF ridge & 10 & $+0.668$ & $+17.26$ & $<$.0001 \\
Long-form & B2 TF--IDF ridge & 20 & $+0.690$ & $+21.12$ & $<$.0001 \\
\addlinespace
Long-form & B3 L--M linear & 5 & $+0.652$ & $+6.82$ & $<$.0001 \\
Long-form & B3 L--M linear & 10 & $+0.722$ & $+7.31$ & $<$.0001 \\
Long-form & B3 L--M linear & 20 & $+0.707$ & $+9.99$ & $<$.0001 \\
\addlinespace
Long-form & B4 L--M features & 5 & $+1.123$ & $+5.93$ & $<$.0001 \\
Long-form & B4 L--M features & 10 & $+0.573$ & $+5.99$ & $<$.0001 \\
Long-form & B4 L--M features & 20 & $+0.566$ & $+8.60$ & $<$.0001 \\
\addlinespace
Long-form & C1 BERT (S1) & 5 & $+0.626$ & $+12.96$ & $<$.0001 \\
Long-form & C1 BERT (S1) & 10 & $+0.706$ & $+13.21$ & $<$.0001 \\
Long-form & C1 BERT (S1) & 20 & $+0.611$ & $+16.01$ & $<$.0001 \\
\addlinespace
Long-form & C2 FinBERT (S1) & 5 & $+0.601$ & $+12.69$ & $<$.0001 \\
Long-form & C2 FinBERT (S1) & 10 & $+0.444$ & $+12.80$ & $<$.0001 \\
Long-form & C2 FinBERT (S1) & 20 & $+0.554$ & $+16.14$ & $<$.0001 \\
\addlinespace
Long-form & C2 FinBERT (S2) & 5 & $+0.052$ & $+2.28$ & .0225 \\
Long-form & C2 FinBERT (S2) & 10 & $+0.560$ & $+14.11$ & $<$.0001 \\
Long-form & C2 FinBERT (S2) & 20 & $+0.350$ & $+11.23$ & $<$.0001 \\
\addlinespace
Long-form & C2 FinBERT (S3) & 5 & $+0.181$ & $+6.77$ & $<$.0001 \\
Long-form & C2 FinBERT (S3) & 10 & $+0.336$ & $+12.79$ & $<$.0001 \\
Long-form & C2 FinBERT (S3) & 20 & $+0.364$ & $+13.82$ & $<$.0001 \\
\addlinespace
Long-form & C2 FinBERT (S4) & 5 & $+0.294$ & $+7.00$ & $<$.0001 \\
Long-form & C2 FinBERT (S4) & 10 & $+0.112$ & $+4.52$ & $<$.0001 \\
Long-form & C2 FinBERT (S4) & 20 & $+0.397$ & $+13.10$ & $<$.0001 \\
\addlinespace
Long-form & C3 RoBERTa (S1) & 5 & $+0.569$ & $+11.09$ & $<$.0001 \\
Long-form & C3 RoBERTa (S1) & 10 & $+0.417$ & $+11.04$ & $<$.0001 \\
Long-form & C3 RoBERTa (S1) & 20 & $+0.334$ & $+10.49$ & $<$.0001 \\
\addlinespace
Long-form & C4 Longformer (S5) & 5 & $+0.484$ & $+12.82$ & $<$.0001 \\
Long-form & C4 Longformer (S5) & 10 & $+0.476$ & $+12.59$ & $<$.0001 \\
Long-form & C4 Longformer (S5) & 20 & $+0.602$ & $+17.37$ & $<$.0001 \\
\bottomrule
\end{tabular}
\end{table}

\begin{table}[p]
\centering
\caption[The 48 encompassing regressions, event-driven half]{The complete
forecast-encompassing record, \textbf{event-driven channel}: the remaining
15 of the 48 regressions begun in
Table~\ref{tab:app-encompass}. Specification, estimator, basis tags and
reading are identical to that table and are not restated: $g$ is the
coefficient on the text forecast conditional on the HAR-RV forecast, and it
is fitted on the \emph{test} sample with observation-level HAC standard
errors at lag $h-1$. This is therefore an in-sample conditional-correlation
statistic and not an out-of-sample loss reduction, and its $t$ statistics
are inflated relative to the day-clustered primary. The encompassing
artefact covers a fixed arm list rather than the full panel: the eleven
long-form and five event-driven arms for which a text-only forecast series
was regressed on the HAR-RV forecast. The frozen-embedding (C5), prompted
(C6) and fusion (D) arms carry no encompassing regression in either
channel, and the event-driven C1, C3, C4 and S2--S4 FinBERT runs, which do
exist on this panel and are scored standalone in
Table~\ref{tab:app-standalone-ed}, were not passed through it. That is a
scope limit of the encompassing script and not a suppressed result: the
event-driven C6 increments on which the chapter's residual claim rests are
reported in full in Table~\ref{tab:app-grid-ed}. All 15
coefficients are positive and significant at 5\%, which is what makes the
combined verdict 48 of 48.}
\label{tab:app-encompass-ed}
\singlespacing\footnotesize
\setlength{\tabcolsep}{4pt}
% src: results/tables/encompassing_regression.md (the event_driven block, all 15 rows; the source's editorialising verdict column is not reproduced)
\begin{tabular}{@{}llrrrr@{}}
\toprule
Channel & Text arm & $h$ & $g$ & HAC $t$ & $p$ \\
\midrule
Event-driven & B1 BoW ridge & 5 & $+0.054$ & $+3.85$ & .0001 \\
Event-driven & B1 BoW ridge & 10 & $+0.053$ & $+4.20$ & $<$.0001 \\
Event-driven & B1 BoW ridge & 20 & $+0.089$ & $+3.59$ & .0003 \\
\addlinespace
Event-driven & B2 TF--IDF ridge & 5 & $+0.372$ & $+13.95$ & $<$.0001 \\
Event-driven & B2 TF--IDF ridge & 10 & $+0.366$ & $+14.77$ & $<$.0001 \\
Event-driven & B2 TF--IDF ridge & 20 & $+0.376$ & $+17.22$ & $<$.0001 \\
\addlinespace
Event-driven & B3 L--M linear & 5 & $+1.184$ & $+5.71$ & $<$.0001 \\
Event-driven & B3 L--M linear & 10 & $+1.024$ & $+5.73$ & $<$.0001 \\
Event-driven & B3 L--M linear & 20 & $+1.023$ & $+6.73$ & $<$.0001 \\
\addlinespace
Event-driven & B4 L--M features & 5 & $+2.882$ & $+6.81$ & $<$.0001 \\
Event-driven & B4 L--M features & 10 & $+1.953$ & $+6.10$ & $<$.0001 \\
Event-driven & B4 L--M features & 20 & $+1.236$ & $+6.63$ & $<$.0001 \\
\addlinespace
Event-driven & C2 FinBERT (S1) & 5 & $+0.194$ & $+12.74$ & $<$.0001 \\
Event-driven & C2 FinBERT (S1) & 10 & $+0.365$ & $+19.66$ & $<$.0001 \\
Event-driven & C2 FinBERT (S1) & 20 & $+0.303$ & $+20.37$ & $<$.0001 \\
\bottomrule
\end{tabular}
\end{table}

\begin{table}[htbp]
\centering
\caption[The pooled omnibus test over the 69-cell grid]{The pooled omnibus
test. For each cell the day-clustered loss differential is averaged within
each trading day. The daily series are then averaged across the cells
present on that day, and the single pooled series is tested with the same
Newey--West and small-sample machinery as the per-cell tests. \emph{A
positive $t$ means text helps}: the sign convention is the negation of
the per-cell DM tables. Holm is applied across the three
pre-declared subfamilies. ``approx.\ rel.\%'' is the mean daily
differential expressed against the member cells' mean reference loss and is
descriptive scale only, not a QLIKE improvement comparable with the
per-cell percentages. Days enter with equal weight, so the pooled test is a
statement about trading days rather than about filings.}
\label{tab:app-omnibus}
\singlespacing\footnotesize
\setlength{\tabcolsep}{3pt}
% src: results/tables/omnibus_m1.csv (section = omnibus)
\begin{tabular}{@{}lrrrrrlll@{}}
\toprule
Subfamily & Cells & Days & mean $D(t)$ & approx.\ rel.\% & $t$ & $p$ & Holm $p$ & Rejects \\
\midrule
Long-form subfamily & 45 & 809 & $+0.001100$ & $+1.19$\% & $+6.90$ & $1.06\times10^{-11}$ & $2.12\times10^{-11}$ & \textbf{Y} \\
Event-driven subfamily & 24 & 996 & $+0.000508$ & $+0.55$\% & $+4.57$ & $5.61\times10^{-6}$ & $5.61\times10^{-6}$ & \textbf{Y} \\
\textbf{All 69 cells} & 69 & 996 & $+0.000783$ & $+0.84$\% & $+8.29$ & $3.76\times10^{-16}$ & $1.13\times10^{-15}$ & \textbf{Y} \\
\bottomrule
\end{tabular}
\end{table}

\begin{table}[htbp]
\centering
\caption[Power calibration of the pooled omnibus test]{Power calibration of
the pooled test. A firm-orthogonal signal (the within-firm demeaned test
log-residual of the reference forecast) is injected into every one of
the 69 cells and scaled, cell by cell, until each cell's realised relative
QLIKE improvement hits the target level. The pooled test is then re-run.
\textbf{This is an oracle construction and the single declared exception to
the no-look-ahead rule in the whole study: it calibrates power and is never
citable as forecasting performance.} ``$\kappa<0$'' counts the cells whose
existing increment exceeded the target and therefore had signal
\emph{removed} to reach it. The rejection rate is over 100 day-block
moving-bootstrap replications of the injected pooled series. Interpolating
the rate gives an 80\% -power minimum detectable effect of \textbf{0.375\%}
(analytic counterpart 0.302\%, a normal approximation reported because the
registered grid may saturate at its smallest level).}
\label{tab:app-power}
\singlespacing\small
\setlength{\tabcolsep}{4pt}
% src: results/tables/omnibus_m1.csv (sections power and mde)
\begin{tabular}{@{}lrrrrrll@{}}
\toprule
Injected level & Converged & $\kappa<0$ & mean $D(t)$ & $t$ & $p$ & Rejects & Rate \\
\midrule
0.1\% & 69/69 & 51 & +0.000089 & $+0.84$ & .4001 & no & 0.10 \\
0.2\% & 69/69 & 45 & +0.000178 & $+1.72$ & .0859 & no & 0.46 \\
0.3\% & 69/69 & 40 & +0.000266 & $+2.63$ & .0088 & yes & 0.68 \\
0.5\% & 69/69 & 37 & +0.000442 & $+4.51$ & $7.36\times10^{-6}$ & yes & 1.00 \\
1.0\% & 69/69 & 32 & +0.000881 & $+9.41$ & $3.20\times10^{-20}$ & yes & 1.00 \\
\bottomrule
\end{tabular}
\end{table}

Figure~\ref{fig:app-encompassing} draws both
instruments (the 48 fitted coefficients and the pooled power curve) with the
in-sample scope of the encompassing panel carried on its face.


\clearpage

\begin{figure}[p]
\centering
\includegraphics[width=\textwidth,height=0.83\textheight,keepaspectratio]{figures/F10_encompassing_pooled_power.pdf}
% caption numbers src: results/tables/encompassing_regression.md (all 48 tests: g, HAC t, raw p; no CSV companion exists); results/tables/omnibus_m1.csv and .md (pooled statistics, injection recovery rates, the two minimum detectable effects)
\caption[Detectable, not attributable]{Detectable, not attributable.
Panel~(a) plots the fitted text coefficient in all 48 forecast-encompassing
regressions: realised volatility on a constant, the HAR forecast and the
text forecast, all three in levels, with a Newey--West correction at lag $h-1$.
They cover 11 long-form and 5 event-driven challenger arms at three horizons
each, with the coefficient and its HAC $t$ printed beside every bar. All 48
coefficients are positive and significant at the 5 per cent level, running
from $+0.052$ to $+2.882$, with $t$ from $+2.28$ to $+21.12$ and a largest
$p$-value of 0.0225. % src: results/tables/encompassing_regression.md (48 rows; minimum at long-form FinBERT-S2 h=5, maximum at event-driven Loughran--McDonald features h=5)
Panel~(b) is the pooled test: per-cell daily loss differentials averaged
across whichever of the 69 cells are present each trading day and the single
resulting series tested day-clustered, giving $t = +8.29$
($p = 3.8\times10^{-16}$ over 996 days). Its curve is what that test is powered
to see, yielding an 80 per cent-power minimum detectable effect of 0.375 per
cent. % src: results/tables/omnibus_m1.md (omnibus table and power-calibration table)
Basis: panel~(a) is an \emph{in-sample} statistic fitted on the test sample
with raw uncorrected $p$-values; panel~(b) is the seed-ensemble primary family
with a denominator of 69 cells.}
\label{fig:app-encompassing}
\end{figure}

Figure~\ref{fig:app-encompassing} carries the evidence that most favours text,
at full strength, and it is placed here rather than suppressed because a
report whose headline is a near-null owes its reader the strongest contrary
evidence it possesses. Forty-eight of forty-eight encompassing tests return a
positive and significant text coefficient, and a pooled test over all 69 grid
cells detects a systematic micro-increment comfortably above what that test is
powered to see. Chapter~\ref{ch:results} quotes both, and the ordering it
draws from them (detectable, then attributable, then realisable) is the
project's principal methodological contribution. The power curve is measured rather than assumed: with an oracle firm-orthogonal signal injected at 0.1 to 1.0 per cent, the rejection rate over 100 day-block bootstrap replications runs 0.10, 0.46, 0.68, 1.00 and 1.00.
% src: results/tables/omnibus_m1.md (power-calibration table)
The pre-declared subfamilies of that pooled test agree with it, at $t=+6.90$ over the 45 long-form cells and $+4.57$ over the 24 event-driven ones, and the analytic counterpart of the interpolated minimum detectable effect is 0.302 per cent.
% src: results/tables/omnibus_m1.md (subfamily rows; analytic MDE)

Neither panel is in tension with the empty conjunction, and the reasons are on
the figure's face rather than left to the reader. Detection at the pooled
level and attribution at the cell level are different claims. Averaging 69
dependent cells buys power against noise, it does not identify which cells
carry an effect, and the identity instruments of Chapter~\ref{ch:results} price
roughly half of the 8-K increment as which firm filed. That share is measured
on six event-driven cells only; no identity instrument was executed on the
long-form arms, which supply most of the cells this pooled test averages. % src: results/tables/anon_arm.csv (every executed row carries channel = ed; six cells, C6 and C2 at h = 5/10/20); results/tables/anon_arm_lf.md (long-form channel closed with zero executed arms)

What the figure does not show is the more important half of its own reading.
Panel~(a) contains no out-of-sample loss reduction, no Holm correction across
its 48 tests and no economic magnitude. The estimating equation is in levels
and is fitted on the test sample itself, so a coefficient of $+2.882$ is a
conditional-correlation loading and not a forecasting gain, it cannot be
converted into one from this chart, and it entitles no cell to a survivor
count. It is the only in-sample quantity in this report, and the 16 arms
tested are a subset of the 23 in the 69-cell family, so its 48 is not
comparable cell for cell with any ladder count. Panel~(b) contains no per-cell
information and no placebo gating, and its injection is an oracle construction
built from test labels, declared as such and usable only as a power
calibration. At its five injected levels the procedure \emph{subtracts} signal
in 51, 45, 40, 37 and 32 of the 69 cells to equalise the realised effect, so
each point is a recovery rate at a fixed effect size and not an addition to
the observed data. These are the omnibus calibration's own counts, and at
the three levels it shares with Section~\ref{sec:val-power} they are the ones
that section prints: both count the cells receiving a negative \emph{effective}
injection $\kappa = g_{1}\delta$. The larger 50, 47 and 42 that the per-cell
calibration table reports instead count a negative $\delta$, the coefficient
applied to the raw text forecast; the two sets part company in exactly the
thirty cell-levels where the text leg's own combination weight is negative and
the two signs cancel. % src: results/tables/omnibus_m1.md (power-calibration table, kappa<0 cells:
% 51/45/40/37/32 at 0.1/0.2/0.3/0.5/1.0 per cent);
% results/tables/signal_injection_power.csv (same file, two columns: kappa1<0 gives
% 40/37/32 and delta<0 gives 50/47/42 at the three shared levels; the 30 cell-levels
% that differ all have g1<0)

\begin{figure}[tbp]
\centering
\includegraphics[width=\textwidth,height=0.83\textheight,keepaspectratio]{figures/AF3_stratified_increment_map__a.pdf}
% caption numbers src: results/tables/m1_stratified.csv (756 rows; 567 drawn, 675 partition cells used in panel b); results/tables/m1_stratified.md (method header)
\caption[The stratum map of the apparent increment]{
The stratum map of the
apparent increment over a recalibrated HAR. 
Panel~(a) draws 567 cells: nine
text and fusion arms at three horizons (rows) against the pooled cell and
three stratum axes (volatility tercile, filing period and disclosure form)
for the long-form and event-driven channels (columns). Colour is the
relative QLIKE improvement clipped at $\pm5$ per cent, with carets on the 32
cells beyond the scale and grey dots on the 148 of 567 that are not
significant; 388 are positive. % src: results/tables/m1_stratified.csv (567 drawn cells: 148 with dm_q_p >= .05, 388 with rel_impr_pct > 0, 32 with |rel_impr_pct| > 5)
The pattern is organised by row rather than by column: the sign and size of an
increment are a property of the arm and horizon and not of the stratum. The
deep negative band is the long-document arms at $h=20$, reaching $-43.14$ per
cent on annual reports for the frozen 8B embedding arm. 
}
\label{fig:app-stratum-map}
\end{figure}

\begin{figure}[tbp]
\centering
\includegraphics[width=\textwidth,height=0.83\textheight,keepaspectratio]{figures/AF3_stratified_increment_map__b.pdf}
% caption numbers src: results/tables/m1_stratified.csv (756 rows; 567 drawn, 675 partition cells used in panel b); results/tables/m1_stratified.md (method header)
\caption[The stratum map of the apparent increment (b)]{Continued from Figure~\ref{fig:app-stratum-map}. The stratum map of the
apparent increment over a recalibrated HAR. 
Panel~(b) plots all
675 partition cells (the 513 drawn above plus 162 combined-channel cells
panel~(a) omits) against their placebos, inside $\pm0.78$ per cent against
real values spanning $-43.14$ to $+7.17$. % src: results/tables/m1_stratified.csv (placebo_rel_impr_pct and rel_impr_pct ranges over the 675 partition cells; 675 = 513 drawn partition cells + 162 combined-channel cells)
Basis: single seed 2026; observation-level Diebold--Mariano rather than the
day-clustered primary, and no Holm family: an off-basis map. Supports are
$7{,}951/7{,}933/7{,}902$ long-form and $25{,}109/25{,}001/24{,}732$
event-driven rows at $h=5/10/20$, which the strata subdivide.
}
\label{fig:app-stratum-map-b}
\end{figure}

The stratum map of Figures~\ref{fig:app-stratum-map} and~\ref{fig:app-stratum-map-b} is the densest artefact in this report, and
it earns the density by settling a reading the project once held and has since
retired. An earlier stage of this work reported that the text signal
concentrated in distressed or high-volatility firms; were that true, the
high-volatility columns of panel~(a) would be visibly warmer than their
neighbours across many rows. They are not. Heat is organised by row (by
which arm at which horizon), and the stratum axes are close to uninformative
about it. The retired reading is therefore not merely unsupported by the
identity controls of Chapter~\ref{ch:results}; it is absent from the raw
stratum map at the rung most favourable to it.

Panel~(b) forestalls the opposite over-correction. Because the apparent
increments are small and the reference ladder removes them, a reader might
suspect they were noise all along. They were not: the label-shuffle placebo,
which permutes the text forecast and leaves everything else in place, returns
values inside $\pm0.78$ per cent, two orders of magnitude below the typical
real increment (median $0.025$ against $1.10$). The placebo is larger in only
38 of the 675 cells, each with a real increment below $0.12$ per cent. % src: results/tables/m1_stratified.csv (over the 675 axis != 'all' cells: max |placebo_rel_impr_pct| = 0.777; median |placebo| = 0.025, median |rel_impr_pct| = 1.103; |placebo| >= |real| in 38 cells, whose largest |rel_impr_pct| is 0.119)
The apparent increment is real signal; the audit's
finding concerns what that signal turns out to be, not whether it exists.

What the map does not show is attributable value. Its reference is the
recalibrated HAR alone, so a warm cell is an apparent increment of the kind
the identity and pool rungs dissolve, and the figure says nothing about
survival at either. It is off-basis in three respects simultaneously (one
seed, observation-level inference, no multiplicity family), so no cell in it
is a verdict, and the grey dots mark uncorrected significance only.
Combination weights are fitted once on validation and frozen, and the strata
partition the test residuals with nothing refitted inside a stratum, so a
stratum with few rows can carry a very large percentage. The extreme negatives
are concentrated in the long-document arms at the longest horizon. The
combined channel is omitted because its rows are the union of the two drawn.
Finally, the event-driven form axis is the four most frequent literal
item-code combinations plus a residual class, a different partition from the
earnings-conceding grouping of Figure~\ref{fig:app-item-geography}, and the
two must not be read across.

\clearpage

\FloatBarrier
\section{Completing the Price Frontier}
\label{app:res-frontier}

The pool rung of Table~\ref{tab:app-ladder} fits five price models, and two
further price models were computed and left out of it: HARQ, which corrects
the measurement error in a realised-variance regressor
\cite{bollerslev2016harq}, and HAR-X augmented with the VIX. That omission
is the one place in the ladder where the study can be accused of stopping
short of the strongest available reference. The accusation bites
precisely because the member most obviously worth adding, HARQ,
addresses the weakness of a daily close-to-close proxy. It cannot be
settled by assertion, because a pool member can fail in two opposite ways:
an unstable member makes a validation-fitted pool \emph{worse} out of
sample, which would credit text for the pool's own breakage. Tables
\ref{tab:app-pool-quality} and~\ref{tab:app-pool-cascade} measure both
directions.

Table~\ref{tab:app-pool-quality} answers the stability question first, and
finds no instability anywhere: the largest absolute fitted log weight is
about $1.8$ in every pool, including both pools containing HARQ. The
characterisation of HARQ as log-space-unstable is not evidenced in the pool
context. On value the two candidates separate. HAR-X(VIX) belongs in the
pool: it lowers mean test QLIKE from 0.0916 to 0.0887 and is
clustered-significantly better than the five-member pool in four of the six
panels. HARQ is mixed, better in three panels and a hair worse on
the mean. Table~\ref{tab:app-pool-cascade} prices the consequence for the
text verdict, and it runs in the conservative direction: completing the
frontier cuts the Holm survivor count from 8 of 69 to 3 of 69. All
three survivors are long-form cells at $h{=}20$. The two survivor sets are
not nested: B2 TF--IDF ridge and C2 FinBERT-S4 at $h{=}20$ survive both
pools, but B1 BoW ridge at $h{=}20$ survives only the seven-member pool,
having failed the five-member one. The change therefore reshuffles a little as
well as subtracting. Strengthening the reference the study was accused of
weakening nonetheless deepens the null rather than relieving it.

The basis warning on these two tables matters more than usual. The frontier
recomputation exists only on the archived single-seed 2026 basis, on which
the five-member pool's own survivor count is 8 rather than the 9 the
seed-ensemble primary reports in Table~\ref{tab:app-ladder}. The internally
consistent claim is the difference (8 to 3, same basis, same code path),
and the 3 must never be set beside the ensemble-basis 9.

\begin{table}[htbp]
\centering
\caption[Price-pool quality as the frontier is completed]{Price-pool
quality as the two omitted members are put back. Each pool is a log-space
combination of its members fitted on validation and frozen on test. The
entries are test QLIKE in \emph{volatility units} on each
disclosure-by-horizon panel, so they are not comparable with the
variance-unit levels of Tables~\ref{tab:app-leaderboard},
\ref{tab:app-standalone-ed} and~\ref{tab:app-standalone-comb}. ``Better''
counts the panels in which the enlarged pool is clustered-significantly
better than the five-member pool at raw $p<.05$, with no correction across the
eighteen pool-by-panel tests; Holm over those eighteen would leave 8 of the 10
wins standing.
% src: results/tables/pool_frontier_audit.csv (p_vs_pool5 over 18 tests, better_than_pool5 true in 10; step-down Holm over the 18 leaves 8)
``max $|w|$'' is the largest absolute
fitted log weight over the six panels (the instability diagnostic),
and no pool exceeds 1.82, including both pools containing HARQ.
\textbf{Basis: archived single seed 2026}, the only basis on which the
frontier recomputation exists.}
\label{tab:app-pool-quality}
\singlespacing\small
\setlength{\tabcolsep}{3.5pt}
% src: results/tables/pool_frontier_audit.csv (qlike_test, max_abs_weight, better_than_pool5, n_members)
\begin{tabular}{@{}lrrrrrrrrlr@{}}
\toprule
& & \multicolumn{3}{c}{Long-form} & \multicolumn{3}{c}{Event-driven} & & & \\
\cmidrule(lr){3-5}\cmidrule(lr){6-8}
Price pool & $m$ & $h{=}5$ & $h{=}10$ & $h{=}20$ & $h{=}5$ & $h{=}10$ & $h{=}20$ & mean & Better & max $|w|$ \\
\midrule
Five-model pool \emph{(reported)} & 5 & $0.1174$ & $0.0856$ & $0.0731$ & $0.1233$ & $0.0869$ & $0.0634$ & $0.0916$ & --- & $1.82$ \\
\quad $+$ HARQ & 6 & $0.1182$ & $0.0865$ & $0.0735$ & $0.1230$ & $0.0866$ & $0.0631$ & $0.0918$ & 3 of 6 & $1.82$ \\
\quad $+$ HAR-X(VIX) & 6 & $0.1125$ & $0.0769$ & $0.0725$ & $0.1226$ & $0.0845$ & $0.0632$ & $0.0887$ & 4 of 6 & $1.81$ \\
\quad $+$ both (seven members) & 7 & $0.1142$ & $0.0785$ & $0.0729$ & $0.1224$ & $0.0843$ & $0.0630$ & $0.0892$ & 3 of 6 & $1.81$ \\
\bottomrule
\end{tabular}
\end{table}

\begin{table}[htbp]
\centering
\caption[The text cascade against the completed price frontier]{What
completing the price frontier does to the 69-cell text cascade. Both rows
are computed on the same rows, with the same combiner, the same
day-clustered DM and the same 69-cell Holm family; only the membership of
the price pool forming the reference changes. \textbf{Basis: archived
single seed 2026}, on which the five-member pool's own count is 8 rather
than the 9 Table~\ref{tab:app-ladder} reports on the seed-ensemble primary.
The valid claim is the difference between the two rows, not either row
against an ensemble-basis figure.}
\label{tab:app-pool-cascade}
\singlespacing\footnotesize
% src: results/tables/pool_frontier_cascade.csv (genuine_raw, genuine_holm and the per-cell survivor identities); pool_frontier_audit.md
\begin{tabular}{@{}lrrp{0.52\textwidth}@{}}
\toprule
Reference pool & Raw & Holm & Surviving cells (all long-form) \\
\midrule
Five-model pool & 23 of 69 & 8 of 69 &
B2 TF--IDF ridge $h{=}5,10,20$; C2 FinBERT-S4 $h{=}10,20$; C4 Longformer
$h{=}5$; C6 prompted $h{=}5,10$ \\
\addlinespace
Seven-model pool & 17 of 69 & \textbf{3 of 69} &
B1 BoW ridge $h{=}20$; B2 TF--IDF ridge $h{=}20$; C2 FinBERT-S4 $h{=}20$ \\
\bottomrule
\end{tabular}
\end{table}

Figures~\ref{fig:app-pool-frontier} and~\ref{fig:app-pool-frontier-b} draw the
frontier audit: the change in pool quality, the fitted weight on HAR-RV itself, the
cascade under both references and the VIX-augmented reference alone, all on the
same archived single-seed basis as the two tables above.

\clearpage

\begin{figure}[tbp]
\centering
\includegraphics[width=\textwidth,height=0.83\textheight,keepaspectratio]{figures/AR2_price_frontier_completeness__a.pdf}
% caption numbers src: results/tables/pool_frontier_audit.csv (4 pools x 6 panels: qlike_test, fitted weights, max_abs_weight, dm_vs_pool5); results/tables/pool_frontier_cascade.csv (69 cells x 2 pools); results/tables/vix_control_harx.csv (9 panels)
\caption[Is the maximal price pool actually maximal?]{
Is the maximal price
pool actually maximal? Two price models were computed for this project and
then left out of the five-member pool; this figure puts them back. 
Panel~(a)
gives each enlarged pool's test QLIKE relative to the reported five, one point
per disclosure-by-horizon panel, filled where the improvement is significant
under day-clustered Diebold--Mariano. Adding the VIX-augmented HAR-X is better
in four of six panels and lowers mean test QLIKE from $0.0916$ to $0.0887$,
while HARQ is of mixed value at three of six. % src: results/tables/pool_frontier_audit.csv (better_than_pool5 by pool; mean qlike_test 0.0916 for the five-member pool, 0.0887 with HAR-X)
Panel~(b) plots the fitted log-space weight the pool places on HAR-RV itself,
already negative in four of six panels at five members and reaching $-1.60$ at
seven. 
}
% src: results/tables/pool_frontier_cascade.csv (all nine union cells disc=long_form; family 45 long_form + 24 event_driven)
\label{fig:app-pool-frontier}
\end{figure}

\begin{figure}[tbp]
\centering
\includegraphics[width=\textwidth,height=0.83\textheight,keepaspectratio]{figures/AR2_price_frontier_completeness__b.pdf}
% caption numbers src: results/tables/pool_frontier_audit.csv (4 pools x 6 panels: qlike_test, fitted weights, max_abs_weight, dm_vs_pool5); results/tables/pool_frontier_cascade.csv (69 cells x 2 pools); results/tables/vix_control_harx.csv (9 panels)
\caption[Is the maximal price pool actually maximal? (c, d)]{Continued from Figure~\ref{fig:app-pool-frontier}. Is the maximal price
pool actually maximal? Two price models were computed for this project and
then left out of the five-member pool; this figure puts them back. 
Panel~(c) carries the consequence for the text cascade, every cell
genuine under either pool being drawn: Holm survivors fall from 8 of 69 to 3
of 69, six lost, two kept, one gained. 
Panel~(c) is a Holm-corrected cascade, run inside its own 69-cell family under
each enlarged pool; all nine cells it draws are long-form, no event-driven
cell surviving under either pool.
Panel~(d) isolates the VIX-augmented
reference: 3.1 to 9.6 per cent better than HAR-RV, raw-significant in six of
nine panels, surviving Holm in none. Basis: the archived single-seed basis
only, where the reported pool carries 8 of 69 rather than the primary's 9.
Panels~(a) and~(b) carry no Holm family and (a)'s fills are uncorrected;
panel~(d)'s ``surviving Holm in none'' is a nine-panel Holm verdict.
% src: results/tables/vix_control_harx.csv (holm = p x 9: event_driven h=5 carries p=0.008135, holm=0.07322)
}
% src: results/tables/pool_frontier_cascade.csv (all nine union cells disc=long_form; family 45 long_form + 24 event_driven)
\label{fig:app-pool-frontier-b}
\end{figure}

The frontier audit of Figures~\ref{fig:app-pool-frontier} and~\ref{fig:app-pool-frontier-b} is a self-criticism, and it is in the
record because the direction it points is the conservative one. Completing the
price frontier makes the price side harder to beat: the pool this study
reports is not the best pool available from its own computed models, and the
better pool leaves fewer text survivors rather than more. Panel~(b) carries a
second observation that is easy to miss and hard to unsee. Inside a fitted
pool the reference model does not behave as the forecast: its weight is
already negative in most panels at five members and strongly negative at
seven, so HAR-RV is acting as a correction term applied to the other members.
A reader who imagines the pool as ``HAR-RV plus some help'' has the wrong
picture of what the strongest rung of the ladder is.

Two guards against over-reading. The enlargement is accompanied by no weight
explosion: the largest absolute fitted log-space weight anywhere is 1.82
and no pool exceeds it, including both pools containing HARQ. % src: results/tables/pool_frontier_audit.csv (max_abs_weight; per-pool maxima over the six panels are 1.8209 for pool5, 1.8162 for pool6+HARQ, 1.8134 for pool6+HARX and 1.8127 for pool7, so 1.82 is the maximum across pools rather than a per-pool equality)
The common
dismissal that HARQ is unusable in log space because it destabilises the fit is
not what this evidence shows.
And a better single reference is not automatically a defensible one. HAR-X
reads an exogenous implied-volatility series the rest of the price block does
not, its advantage over HAR-RV survives correction in none of the nine panels,
and admitting it would change what the benchmark's reference \emph{means} as
well as how strong it is.

What the figure does not show should be stated with equal firmness. The
recomputation exists only on the archived single-seed basis, so what it
licenses is the five-to-seven \emph{change}, both halves of which are computed
by one code path on one basis. It is not a restatement of the primary rung's
9 of 69, and the two counts must never be netted. There is no seed-ensemble
recomputation of the seven-member cascade, no effect size for any cell of
panel~(c), and no firm-identity reference anywhere in the figure. That is a
separate rung, and Figures~\ref{fig:app-firm-identity-rung} and~\ref{fig:app-firm-identity-rung-b} carry it. The panel's axis is symmetric-logarithmic so the extremes are drawn at full length: the Longformer at $h=10$ loses 11.38 per cent and TF-IDF ridge at $h=20$ loses 8.09 per cent.
% src: results/tables/firm_identity_ensemble.csv (extreme cells)
The firm-mean term covers 0.629 and 0.630 of long-form and event-driven firms but 0.915 and 0.919 of test observations, so the thin coverage falls on firms that file rarely.
% src: results/tables/firm_identity_ensemble.csv (firm and row coverage)

\begin{figure}[tbp]
\centering
\includegraphics[width=\textwidth,height=0.83\textheight,keepaspectratio]{figures/F5_maximal_pool_audit__a.pdf}
% caption numbers src: results/tables/maximal_pool_robustness.csv (panel a counts, both bases); results/tables/control_intersection_ensemble.csv (the 38-cell primary rung drawn as the dashed rule); results/tables/maximal_pool_robustness_panels.csv (panel b); results/tables/maximal_reference_single_refs.csv (panel c); results/tables/stronger_baselines.csv, section m1_incremental (panel d)
\caption[Is the maximal price pool a shopped reference?]{
Is the maximal price
pool a shopped reference? 
Panel~(a) counts Holm survivors out of 69 under each
reference specification: 38 against the single recalibrated HAR of the primary
rung, 34 against the validation-best single price member, 17 against a
never-fitted equal-weight pool of five price models and 9 against the
validation-fitted pool of the same five, with the archived single-seed counts
as ghost bars. % src: results/tables/maximal_pool_robustness.csv (basis ens: 34, 17, 9; basis s26: 28, 19, 8); results/tables/m1_ensemble_primary.csv (genuine_ens_vol = 38)
Panel~(b) gives test QLIKE for four price references, one axis per
disclosure-by-horizon panel. 
}
\label{fig:app-pool-audit}
\end{figure}

\begin{figure}[tbp]
\centering
\includegraphics[width=\textwidth,height=0.83\textheight,keepaspectratio]{figures/F5_maximal_pool_audit__b.pdf}
% caption numbers src: results/tables/maximal_pool_robustness.csv (panel a counts, both bases); results/tables/control_intersection_ensemble.csv (the 38-cell primary rung drawn as the dashed rule); results/tables/maximal_pool_robustness_panels.csv (panel b); results/tables/maximal_reference_single_refs.csv (panel c); results/tables/stronger_baselines.csv, section m1_incremental (panel d)
\caption[Is the maximal price pool a shopped reference? (c)]{Continued from Figure~\ref{fig:app-pool-audit}. Is the maximal price
pool a shopped reference? 
Panel~(c) is the single-reference sweep: fifteen
cells scored against each of the five single price models and against the
fitted pool. 
}
\label{fig:app-pool-audit-b}
\end{figure}

\begin{figure}[tbp]
\centering
\includegraphics[width=\textwidth,height=0.83\textheight,keepaspectratio]{figures/F5_maximal_pool_audit__c.pdf}
% caption numbers src: results/tables/maximal_pool_robustness.csv (panel a counts, both bases); results/tables/control_intersection_ensemble.csv (the 38-cell primary rung drawn as the dashed rule); results/tables/maximal_pool_robustness_panels.csv (panel b); results/tables/maximal_reference_single_refs.csv (panel c); results/tables/stronger_baselines.csv, section m1_incremental (panel d)
\caption[Is the maximal price pool a shopped reference? (d)]{Continued from Figure~\ref{fig:app-pool-audit}. Is the maximal price
pool a shopped reference? 
Panel~(d) swaps the recalibrated HAR for a recalibrated
semivariance HAR on eighteen identical cells. Basis: seed-ensemble forecasts
with day-clustered inference and Holm inside each 69-cell family in panel~(a);
the archived single-seed basis throughout panel~(c); observation-level
statistics with neither Holm nor a placebo in panel~(d), which is a magnitude
comparison only. Denominators: 69 cells in panel~(a), 6 panels in~(b), 15
cells in~(c), 18 cells in~(d).
}
\label{fig:app-pool-audit-c}
\end{figure}

The charge Figures~\ref{fig:app-pool-audit}, \ref{fig:app-pool-audit-b} and~\ref{fig:app-pool-audit-c} answer has two halves: that the
reference was chosen to minimise the credit text receives, and that the pool
wins only because six weights were fitted on validation. Panel~(c) answers the
first half, and answers it unfavourably to any single-reference study, this
one included. On one cell (long-form TF-IDF ridge at $h=5$, marked in the
row labels) the reported increment is $+3.39$ per cent against HAR,
$+3.53$ against semivariance HAR, $+1.66$ against GARCH, $+1.06$ against
EGARCH and $+3.24$ against ARIMA. % src: results/tables/maximal_reference_single_refs.csv (long_form / B2_tfidf_ridge / h=5, one row per single price reference)
The same forecasts, the same rows and the same loss produce a range of two and
a half percentage points, purely from the choice of price model. HAR is
not the reference that flatters text least, being the best of the five single
price references in none of the six disclosure-by-horizon panels. % src: results/tables/maximal_pool_robustness_panels.csv (test_qlike_A2_har_rv against the four other single-reference columns, per panel)
Panels~(a) and~(b) answer the second half in both directions. Because the
equal-weight pool fits nothing beyond the two recalibration coefficients every
reference receives, most of the absorption cannot be weight fitting: an
unfitted pool already leaves 17. That 17 is scored on the five-price-model
rows, where the primary rung itself leaves 30 rather than 38, so the fall must
be read against the 30 and never netted against the 38. % src: results/tables/matched_row_cascade.csv (arm = matched: adds_primary 30); results/tables/maximal_pool_robustness.csv (equal-weight pool 17, same intersection rows)

The panel that cuts against the reported rung is drawn at full height rather
than described. Across all six panels the fitted pool loses on test QLIKE to
the never-fitted equal-weight pool, so the ordering in panel~(a), by how
much a reference absorbs, is not an ordering by forecast accuracy. The
panel title says so. At long-form $h=20$ the fitted pool is significantly
worse than the validation-best single member (day-clustered DM $+8.49$,
$p \approx 10^{-16}$) and at $0.0731$ is the worst of the four references
drawn, behind the validation-best member at $0.0714$, the hindsight test-best
member at $0.0689$ and the equal-weight pool at $0.0680$. In the other five
panels the fitted pool beats the member a forecaster could have selected on
validation, at day-clustered statistics of $-7.24$ to $-2.05$. % src: results/tables/maximal_pool_robustness_panels.csv (pool_worse_than_valbest_test, dm_pool_vs_valbest, and the four test-QLIKE columns at long_form h=20)
A reference that is a worse forecaster and yet absorbs more of the text
increment is an uncomfortable object, and the report's response is to report
the pool rung as one rung among several rather than as the verdict.

What this figure does not show is effect size. Panel~(a) counts survivors, so
a specification leaving fewer survivors is not thereby shown to leave smaller
increments; and, as panel~(b) establishes, it is not thereby the better
forecaster either. Panel~(b) reports pooled panel-level QLIKE and says nothing
about which cells move, its within-row positions being normalised to that
row's own range so distances are readable within a row and not between rows.
Panel~(c) covers only the fifteen cells of the committed single-reference
sweep and is drawn entirely on the archived single-seed basis. Panel~(d)
carries neither Holm correction nor a placebo and its statistics are
observation-level. Its finding is a statement about magnitudes and not about
verdicts: swapping HAR for semivariance HAR moves no cell by more than 0.18
percentage points, the largest shift being long-form Qwen3-embedding at
$h=10$. % src: results/tables/stronger_baselines.csv (section m1_incremental: |A6_shar_rel_impr_pct - A2_rel_impr_pct| over the 18 cells)

\FloatBarrier
\section{The Combiner Under Deployment}
\label{app:res-deploy}

The primary grid fits the combination weights once, on the 2020--21
validation block, and freezes them across the whole test span. That is feasible
in real time --- the block closes before the test span opens, and the generator
fits nothing on test rows --- but it is not what a forecaster would have done:
the weights would have been refitted each quarter as filings accrued. % src: scripts/analysis/deployable_combiner.py:28 ("NO LOOK-AHEAD: fixed weights are fit on validation only")
Section~\ref{sec:res-residual} reports the consequence in one sentence.
Tables~\ref{tab:app-deploy-ed} and~\ref{tab:app-deploy-lf} give it cell by
cell, on both schemes at once, so that a reader can see \emph{which} cells
the deployable scheme destroys rather than only how many.

The accounting has to be read carefully in two places, and both are stated
on the tables. The Holm family here is 75 cells rather than 69, because the
deployable rerun carries a six-cell C5 extension outside the primary grid.
That widening alone moves the frozen-weight count from 38 to 36, and it is
an accounting change, not a finding. The finding is what happens next: the
same 75-cell family under expanding weights returns 7 genuine cells, 6 of
them inside the primary grid. Thirty-one cells are lost on deployment, two
are gained, and five survive. Almost every casualty is a long-form cell,
and the average pooled increment falls by 3.2 percentage points. The
bag-of-words and fine-tuned encoder arms collapse furthest: long-form B1 BoW
ridge at $h{=}5$ goes from $+1.65$ per cent frozen to $-17.44$ per cent
deployable, and FinBERT-S4 falls furthest of any arm on the mean. The
dictionary arms are among the steadiest, one of them improving, but they are
not the steadiest: the prompted arm and its fused variant move less than
either, on the per-cell average and on its absolute counterpart alike.
% src: results/tables/deployable_combiner.csv (mean exp_pooled_rel minus fixed_pooled_rel over the primary-grid cells: C2_finbert_s4 -12.44 and B1_bow_ridge -11.84 at one end; D4_llmfused +0.06 and C6_llmtext -0.08 at the other, then B3_lm_linear -0.50 and B4_lm_features +0.45. Mean ABSOLUTE change orders the same four: D4 0.09, C6 0.16, B3 0.51, B4 0.74. D1_concat_mlp's mean is +0.05 only because its per-cell moves cancel; its mean absolute change is 4.78)
What survives is the prompted arm on both channels at the two
shorter horizons, plus long-form FinBERT-S1 at $h{=}5$.

Table~\ref{tab:app-deploy-quarters} then opens the four surviving prompted
cells quarter by quarter, which is the resolution at which the pooled
statistic can be checked for dependence on a single period. The
frozen-weight and deployable paths track each other closely in those
cells: the pooled gap is $-0.02$ percentage points at both event-driven
horizons. The increment is positive in 13 to 15 of the 16 quarters
under either scheme, so neither the survival nor its size rests on one
quarter. The per-quarter statistics are descriptive: at $h{=}20$ a single
quarter carries roughly 60 filing days against a Newey--West lag of 19, so
only the pooled path is treated as evidence.

\begin{table}[p]
\centering
\caption[The deployable combiner, event-driven cells]{The deployable
combination grid, \textbf{event-driven cells}, all 27. ``Fixed'' is the
committed scheme (weights fitted once on validation and frozen),
reproduced here to machine precision; ``expanding'' refits the weights
before each test quarter on the full pre-quarter filing history and
concatenates the sixteen quarterly blocks into one pseudo-out-of-sample
path. \emph{Negative DM and positive rel.\% mean text improves.} ``sig.\
$q$'' counts the quarters in which that cell's own per-quarter clustered
test is significant, a descriptive column. Holm is applied within a
pre-declared \emph{75}-cell family spanning both channels and a six-cell
C5 extension marked $^{\dagger}$, so these Holm values are slightly more
conservative than the 69-cell family used in Table~\ref{tab:app-grid-ed}.
The two must not be compared cell for cell. The pooled expanding
comparison is read in the Giacomini--White sense
\cite{giacominiwhite2006conditional}, as a comparison of
forecasting \emph{methods including their estimation scheme}.}
\label{tab:app-deploy-ed}
\singlespacing\footnotesize
\setlength{\tabcolsep}{3pt}
% src: results/tables/deployable_combiner.csv (fixed_* and exp_* columns, genuine flags, deploy_status)
\begin{tabular}{@{}lrrrrrcrrrrrcl@{}}
\toprule
& & \multicolumn{5}{c}{Validation-frozen weights} & \multicolumn{6}{c}{Expanding deployable weights} & \\
\cmidrule(lr){3-7}\cmidrule(lr){8-13}
Model & $h$ & sig.\ $q$ & rel.\% & DM & Holm & G & sig.\ $q$ & rel.\% & DM & Holm & plac. & G & Status \\
\midrule
B1 BoW ridge&5&0/16&$+1.33$&$-3.35$& .0277 &Y&0/16&$-3.91$&$+4.67$& .0002 &$+1.15$&--&lost\\
B1 BoW ridge&10&2/16&$+1.23$&$-3.25$& .0367 &Y&0/16&$-4.84$&$+3.92$& .0057 &$+0.76$&--&lost\\
B1 BoW ridge&20&0/16&$+1.53$&$-3.10$& .0548 &--&0/16&$-6.20$&$+3.20$& .0746 &$+1.12$&--&---\\
\addlinespace
B2 TF--IDF ridge&5&0/16&$+1.21$&$-2.76$& .1296 &--&0/16&$-4.48$&$+4.85$& .0001 &$+0.78$&--&---\\
B2 TF--IDF ridge&10&2/16&$+1.35$&$-2.97$& .0743 &--&0/16&$-5.28$&$+3.74$& .0112 &$+0.26$&--&---\\
B2 TF--IDF ridge&20&0/16&$+1.84$&$-3.11$& .0548 &--&0/16&$-6.47$&$+2.87$& .2059 &$+1.25$&--&---\\
\addlinespace
B3 L--M linear&5&0/16&$+0.25$&$-2.43$& .2730 &--&0/16&$-0.16$&$+5.30$& $<$.0001 &$+1.44$&--&---\\
B3 L--M linear&10&0/16&$+0.20$&$-1.23$& 1.0000 &--&0/16&$+0.01$&$-0.93$& 1.0000 &$+0.52$&--&---\\
B3 L--M linear&20&0/16&$+0.26$&$-1.12$& 1.0000 &--&2/16&$+0.01$&$-0.57$& 1.0000 &$-0.16$&--&---\\
\addlinespace
B4 L--M features&5&0/16&$+0.18$&$-0.99$& 1.0000 &--&0/16&$-0.11$&$+3.50$& .0266 &$+0.05$&--&---\\
B4 L--M features&10&0/16&$+0.08$&$-2.10$& .6041 &--&0/16&$+0.18$&$-0.50$& 1.0000 &$+0.10$&--&---\\
B4 L--M features&20&1/16&$+0.25$&$-2.01$& .6569 &--&3/16&$+0.10$&$-1.16$& 1.0000 &$+0.22$&--&---\\
\addlinespace
C2 FinBERT (S1)&5&4/16&$+2.57$&$-5.92$& $<$.0001 &Y&1/16&$+1.08$&$+0.70$& 1.0000 &$-0.13$&--&lost\\
C2 FinBERT (S1)&10&6/16&$+2.42$&$-5.68$& $<$.0001 &Y&3/16&$-5.44$&$+2.38$& .7315 &$+0.34$&--&lost\\
C2 FinBERT (S1)&20&3/16&$+1.58$&$-1.94$& .6569 &--&2/16&$+1.56$&$-1.03$& 1.0000 &$-0.76$&--&---\\
\addlinespace
C5 Qwen3-emb.$^{\dagger}$&5&0/16&$-0.13$&$+2.83$& .1087 &--&0/16&$-0.45$&$+7.03$& $<$.0001 &$-1.28$&--&---\\
C5 Qwen3-emb.$^{\dagger}$&10&0/16&$-0.13$&$+3.04$& .0602 &--&0/16&$-0.20$&$+6.61$& $<$.0001 &$+0.59$&--&---\\
C5 Qwen3-emb.$^{\dagger}$&20&0/16&$-0.45$&$+3.40$& .0253 &--&2/16&$+0.07$&$-4.18$& .0020 &$+0.42$&Y&gained\\
\addlinespace
C6 prompted LLM&5&3/16&$+1.21$&$-5.04$& $<$.0001 &Y&3/16&$+1.19$&$-4.64$& .0003 &$-0.39$&Y&survives\\
C6 prompted LLM&10&3/16&$+1.00$&$-3.76$& .0074 &Y&3/16&$+0.98$&$-3.44$& .0328 &$-0.60$&Y&survives\\
C6 prompted LLM&20&1/16&$+0.66$&$-1.98$& .6569 &--&1/16&$+0.50$&$-1.56$& 1.0000 &$+1.32$&--&---\\
\addlinespace
D2 Gated fusion&5&0/16&$+0.56$&$-1.70$& .9875 &--&0/16&$-0.22$&$+2.65$& .3771 &$+0.79$&--&---\\
D2 Gated fusion&10&0/16&$+0.18$&$-0.42$& 1.0000 &--&0/16&$-2.43$&$+4.81$& .0001 &$-0.55$&--&---\\
D2 Gated fusion&20&0/16&$-2.76$&$+3.55$& .0155 &--&2/16&$-1.83$&$+0.50$& 1.0000 &$+0.59$&--&---\\
\addlinespace
D4 price+C6&5&0/16&$-0.01$&$+3.37$& .0264 &--&1/16&$+0.03$&$-0.92$& 1.0000 &$+0.61$&--&---\\
D4 price+C6&10&0/16&$-0.01$&$+0.66$& 1.0000 &--&0/16&$-0.03$&$+1.07$& 1.0000 &$-0.09$&--&---\\
D4 price+C6&20&0/16&$-0.35$&$+4.69$& .0002 &--&1/16&$+0.00$&$-0.27$& 1.0000 &$-0.25$&--&---\\
\bottomrule
\end{tabular}
\end{table}

\begin{table}[p]
\centering
\caption[The deployable combiner, long-form cells]{The deployable
combination grid, \textbf{long-form cells}, all 48. Columns, schemes, sign
rules and the 75-cell Holm family exactly as in
Table~\ref{tab:app-deploy-ed}; $^{\dagger}$ again marks the C5 extension
cells outside the 69-cell primary grid. This is where deployment does its
damage: 27 of the 31 cells lost on deployment are here, and the
bag-of-words and TF--IDF arms (the largest frozen-weight increments in
the whole grid) turn sharply negative once their weights must be
estimated without seeing the test span.}
\label{tab:app-deploy-lf}
\singlespacing\footnotesize
\setlength{\tabcolsep}{3pt}
% 48 rows plus a fourteen-line caption stood 11.6 pt over the text block at the
% stock 0.5 em booktabs group separation; the group rules still read at 0.25 em.
\setlength{\defaultaddspace}{0.25em}
% src: results/tables/deployable_combiner.csv (fixed_* and exp_* columns, genuine flags, deploy_status)
\begin{tabular}{@{}lrrrrrcrrrrrcl@{}}
\toprule
& & \multicolumn{5}{c}{Validation-frozen weights} & \multicolumn{6}{c}{Expanding deployable weights} & \\
\cmidrule(lr){3-7}\cmidrule(lr){8-13}
Model & $h$ & sig.\ $q$ & rel.\% & DM & Holm & G & sig.\ $q$ & rel.\% & DM & Holm & plac. & G & Status \\
\midrule
B1 BoW ridge&5&6/16&$+1.65$&$-3.83$& .0057 &Y&0/16&$-17.44$&$+4.23$& .0017 &$-0.52$&--&lost\\
B1 BoW ridge&10&4/16&$+1.44$&$-4.15$& .0016 &Y&0/16&$-15.06$&$+3.02$& .1310 &$+0.77$&--&lost\\
B1 BoW ridge&20&5/16&$+2.99$&$-5.45$& $<$.0001 &Y&0/16&$-13.46$&$+2.29$& .8635 &$-0.00$&--&lost\\
\addlinespace
B2 TF--IDF ridge&5&5/16&$+3.33$&$-5.39$& $<$.0001 &Y&1/16&$-4.47$&$+3.02$& .1310 &$-0.22$&--&lost\\
B2 TF--IDF ridge&10&10/16&$+3.48$&$-8.89$& $<$.0001 &Y&3/16&$-4.12$&$+1.11$& 1.0000 &$+0.42$&--&lost\\
B2 TF--IDF ridge&20&10/16&$+5.92$&$-9.04$& $<$.0001 &Y&1/16&$-3.06$&$+0.68$& 1.0000 &$-0.84$&--&lost\\
\addlinespace
B3 L--M linear&5&3/16&$+0.49$&$-2.58$& .1894 &--&1/16&$+0.54$&$-2.30$& .8586 &$-0.27$&--&---\\
B3 L--M linear&10&3/16&$+1.79$&$-4.62$& .0002 &Y&2/16&$+1.16$&$-3.19$& .0775 &$+0.19$&--&lost\\
B3 L--M linear&20&2/16&$+3.48$&$-4.69$& .0002 &Y&1/16&$+1.92$&$-2.60$& .4193 &$-0.77$&--&lost\\
\addlinespace
B4 L--M features&5&1/16&$+0.11$&$+1.02$& 1.0000 &--&0/16&$-0.30$&$+2.37$& .7413 &$+0.48$&--&---\\
B4 L--M features&10&0/16&$-0.92$&$+3.32$& .0304 &--&1/16&$+0.10$&$-1.54$& 1.0000 &$+0.14$&--&---\\
B4 L--M features&20&0/16&$-1.92$&$+3.38$& .0264 &--&0/16&$+0.54$&$-1.20$& 1.0000 &$-0.59$&--&---\\
\addlinespace
C1 BERT (S1)&5&4/16&$+3.42$&$-3.25$& .0367 &Y&1/16&$-8.45$&$+4.08$& .0030 &$+0.14$&--&lost\\
C1 BERT (S1)&10&1/16&$-0.85$&$+3.06$& .0602 &--&1/16&$+0.96$&$-1.05$& 1.0000 &$-0.03$&--&---\\
C1 BERT (S1)&20&5/16&$+2.70$&$-6.96$& $<$.0001 &Y&1/16&$-4.42$&$+1.02$& 1.0000 &$-0.75$&--&lost\\
\addlinespace
C2 FinBERT (S1)&5&2/16&$+1.90$&$-4.53$& .0003 &Y&3/16&$+8.58$&$-5.50$& $<$.0001 &$-0.75$&Y&survives\\
C2 FinBERT (S1)&10&6/16&$+2.62$&$-6.67$& $<$.0001 &Y&1/16&$-10.05$&$+2.18$& 1.0000 &$+0.32$&--&lost\\
C2 FinBERT (S1)&20&3/16&$-0.59$&$-0.40$& 1.0000 &--&1/16&$-3.04$&$+1.09$& 1.0000 &$+0.08$&--&---\\
\addlinespace
C2 FinBERT (S2)&5&4/16&$+1.21$&$-5.57$& $<$.0001 &Y&1/16&$-4.82$&$+3.75$& .0109 &$-0.43$&--&lost\\
C2 FinBERT (S2)&10&12/16&$+0.48$&$-7.59$& $<$.0001 &Y&2/16&$-2.36$&$-0.50$& 1.0000 &$-0.34$&--&lost\\
C2 FinBERT (S2)&20&4/16&$+1.68$&$-4.89$& $<$.0001 &Y&0/16&$-0.16$&$-0.40$& 1.0000 &$-0.26$&--&lost\\
\addlinespace
C2 FinBERT (S3)&5&5/16&$+2.90$&$-5.36$& $<$.0001 &Y&1/16&$-4.35$&$+3.56$& .0220 &$-0.50$&--&lost\\
C2 FinBERT (S3)&10&7/16&$+2.31$&$-5.26$& $<$.0001 &Y&0/16&$-19.82$&$+3.82$& .0083 &$+0.09$&--&lost\\
C2 FinBERT (S3)&20&2/16&$-3.86$&$+1.68$& .9875 &--&1/16&$-0.25$&$-0.12$& 1.0000 &$+0.09$&--&---\\
\addlinespace
C2 FinBERT (S4)&5&6/16&$+1.46$&$-4.78$& .0001 &Y&0/16&$-14.21$&$+5.94$& $<$.0001 &$+0.33$&--&lost\\
C2 FinBERT (S4)&10&10/16&$+0.36$&$-6.82$& $<$.0001 &Y&1/16&$-13.94$&$+2.09$& 1.0000 &$+0.21$&--&lost\\
C2 FinBERT (S4)&20&9/16&$+3.08$&$-8.34$& $<$.0001 &Y&1/16&$-4.28$&$+0.75$& 1.0000 &$-0.05$&--&lost\\
\addlinespace
C3 RoBERTa (S1)&5&5/16&$+0.30$&$-5.13$& $<$.0001 &Y&1/16&$-3.93$&$+2.65$& .3771 &$-0.96$&--&lost\\
C3 RoBERTa (S1)&10&5/16&$+1.84$&$-5.06$& $<$.0001 &Y&3/16&$+3.32$&$-2.80$& .2515 &$+0.01$&--&lost\\
C3 RoBERTa (S1)&20&6/16&$+0.02$&$-3.90$& .0045 &Y&0/16&$-17.26$&$+2.25$& .9372 &$-0.10$&--&lost\\
\addlinespace
C4 Longformer (S5)&5&6/16&$+1.47$&$-6.18$& $<$.0001 &Y&4/16&$+2.00$&$-1.48$& 1.0000 &$-0.39$&--&lost\\
C4 Longformer (S5)&10&0/16&$-2.95$&$+10.32$& $<$.0001 &--&3/16&$-3.01$&$-0.54$& 1.0000 &$-0.14$&--&---\\
C4 Longformer (S5)&20&4/16&$+0.91$&$-4.38$& .0006 &Y&2/16&$-3.16$&$-0.01$& 1.0000 &$+0.26$&--&lost\\
\addlinespace
C5 Qwen3-emb.$^{\dagger}$&5&0/16&$-1.02$&$+2.64$& .1698 &--&0/16&$-0.27$&$+2.40$& .7201 &$+0.33$&--&---\\
C5 Qwen3-emb.$^{\dagger}$&10&0/16&$-3.10$&$+3.62$& .0123 &--&0/16&$-0.14$&$+0.51$& 1.0000 &$+0.66$&--&---\\
C5 Qwen3-emb.$^{\dagger}$&20&0/16&$-6.39$&$+5.38$& $<$.0001 &--&2/16&$+0.69$&$-0.99$& 1.0000 &$-0.86$&--&---\\
\addlinespace
C6 prompted LLM&5&6/16&$+1.79$&$-6.31$& $<$.0001 &Y&4/16&$+2.04$&$-4.71$& .0002 &$+0.84$&Y&survives\\
C6 prompted LLM&10&11/16&$+2.25$&$-7.92$& $<$.0001 &Y&9/16&$+2.03$&$-5.64$& $<$.0001 &$+0.53$&Y&survives\\
C6 prompted LLM&20&2/16&$+0.27$&$-3.23$& .0371 &Y&1/16&$-0.04$&$-0.61$& 1.0000 &$+0.11$&--&lost\\
\addlinespace
D1 Concat-MLP&5&0/16&$-1.04$&$+2.66$& .1693 &--&1/16&$+6.20$&$-4.03$& .0037 &$+0.81$&Y&gained\\
D1 Concat-MLP&10&0/16&$-0.46$&$+3.54$& .0158 &--&1/16&$-1.05$&$+0.14$& 1.0000 &$-0.35$&--&---\\
D1 Concat-MLP&20&5/16&$+0.12$&$-5.57$& $<$.0001 &Y&0/16&$-6.37$&$+1.65$& 1.0000 &$+0.37$&--&lost\\
\addlinespace
D2 Gated fusion&5&3/16&$+0.19$&$-2.11$& .6041 &--&1/16&$+0.79$&$+1.61$& 1.0000 &$+0.67$&--&---\\
D2 Gated fusion&10&0/16&$-0.02$&$+2.02$& .6569 &--&0/16&$+1.57$&$-1.03$& 1.0000 &$-0.31$&--&---\\
D2 Gated fusion&20&0/16&$-2.27$&$+5.56$& $<$.0001 &--&1/16&$-5.36$&$+0.93$& 1.0000 &$+0.19$&--&---\\
\addlinespace
D4 price+C6&5&3/16&$+0.15$&$-0.75$& 1.0000 &--&0/16&$+0.14$&$-0.31$& 1.0000 &$+0.52$&--&---\\
D4 price+C6&10&0/16&$-0.15$&$-0.43$& 1.0000 &--&0/16&$-0.19$&$-0.49$& 1.0000 &$+1.52$&--&---\\
D4 price+C6&20&1/16&$+0.60$&$-3.58$& .0142 &Y&1/16&$+0.67$&$-2.79$& .2515 &$-0.65$&--&lost\\
\bottomrule
\end{tabular}
\end{table}

\begin{table}[htbp]
\centering
\caption[The prompted arm quarter by quarter, both weight schemes]{The four
prompted-arm cells that survive deployment, quarter by quarter under both
weight schemes: relative QLIKE improvement over the recalibrated HAR
reference in each of the sixteen test quarters. \emph{These sixteen figures
are descriptive and the pooled path is the evidence}: a single quarter
carries roughly 500 long-form and 1{,}400--1{,}800 event-driven
observations. At the longer horizons the per-quarter clustered test
sits on too few filing days to be reliable. The columns are shown so that a
reader can confirm that neither the survival nor its size depends on one
period. The increment is positive in 14, 15, 14 and 13 of the sixteen
quarters in the four frozen-weight columns respectively.}
\label{tab:app-deploy-quarters}
\singlespacing\small
\setlength{\tabcolsep}{4pt}
% src: results/tables/deployable_combiner_quarters.csv (fixed_rel, exp_rel for the four C6_llmtext cells)
\begin{tabular}{@{}lrrrrrrrr@{}}
\toprule
& \multicolumn{4}{c}{Event-driven, C6 prompted} & \multicolumn{4}{c}{Long-form, C6 prompted} \\
\cmidrule(lr){2-5}\cmidrule(lr){6-9}
& \multicolumn{2}{c}{$h{=}5$} & \multicolumn{2}{c}{$h{=}10$} & \multicolumn{2}{c}{$h{=}5$} & \multicolumn{2}{c}{$h{=}10$} \\
\cmidrule(lr){2-3}\cmidrule(lr){4-5}\cmidrule(lr){6-7}\cmidrule(lr){8-9}
Quarter & frozen & deploy. & frozen & deploy. & frozen & deploy. & frozen & deploy. \\
\midrule
2022Q1 & $+1.00$ & $+1.00$ & $+0.85$ & $+0.85$ & $-1.08$ & $-1.08$ & $-0.43$ & $-0.43$ \\
2022Q2 & $-0.10$ & $+0.00$ & $+0.19$ & $+0.29$ & $-1.14$ & $-0.68$ & $-1.69$ & $-1.25$ \\
2022Q3 & $+1.12$ & $+1.18$ & $+1.06$ & $+1.10$ & $+3.55$ & $+3.11$ & $+3.63$ & $+2.94$ \\
2022Q4 & $+1.33$ & $+1.42$ & $+0.97$ & $+1.05$ & $+0.43$ & $+0.99$ & $-0.15$ & $+0.27$ \\
2023Q1 & $-0.43$ & $-0.34$ & $-0.08$ & $-0.03$ & $+2.36$ & $+2.12$ & $+4.27$ & $+4.34$ \\
2023Q2 & $+1.19$ & $+1.18$ & $+1.55$ & $+1.48$ & $+1.60$ & $+1.61$ & $+1.80$ & $+1.63$ \\
2023Q3 & $+1.32$ & $+1.18$ & $+0.74$ & $+0.64$ & $+1.96$ & $+1.76$ & $+2.32$ & $+2.02$ \\
2023Q4 & $+1.59$ & $+1.60$ & $+0.82$ & $+0.81$ & $+1.32$ & $+1.30$ & $+2.26$ & $+2.00$ \\
2024Q1 & $+1.15$ & $+1.06$ & $+0.71$ & $+0.67$ & $+2.30$ & $+2.66$ & $+1.39$ & $+0.93$ \\
2024Q2 & $+0.87$ & $+0.84$ & $+0.47$ & $+0.53$ & $+2.30$ & $+2.66$ & $+3.05$ & $+3.63$ \\
2024Q3 & $+2.72$ & $+2.72$ & $+1.71$ & $+1.74$ & $+1.87$ & $+1.93$ & $+2.75$ & $+2.40$ \\
2024Q4 & $+2.29$ & $+2.19$ & $+2.20$ & $+2.20$ & $+2.37$ & $+3.11$ & $+2.58$ & $+2.75$ \\
2025Q1 & $+1.10$ & $+0.94$ & $+0.29$ & $+0.16$ & $+1.54$ & $+2.25$ & $+0.21$ & $-1.23$ \\
2025Q2 & $+1.24$ & $+1.10$ & $+1.69$ & $+1.45$ & $+3.50$ & $+3.23$ & $+5.18$ & $+4.68$ \\
2025Q3 & $+1.25$ & $+1.20$ & $+0.81$ & $+0.76$ & $+3.38$ & $+4.13$ & $+3.70$ & $+2.98$ \\
2025Q4 & $+1.80$ & $+1.87$ & $+1.82$ & $+1.91$ & $+2.59$ & $+3.25$ & $+4.06$ & $+4.38$ \\
\bottomrule
\end{tabular}
\end{table}

The figures below draw this deployment evidence rather than
this appendix repeating it. Figures~\ref{fig:app-quarterly} and~\ref{fig:app-quarterly-b} give the
quarter-by-quarter path of Table~\ref{tab:app-deploy-quarters} as a heat map
over the 69 primary cells under both weighting schemes, with the cross-cell
median and interquartile band beneath it. Figures~\ref{fig:app-freeze-point} and~\ref{fig:app-freeze-point-b}
plot the frozen against the deployable increment for all 75 cells of this
family, marking the six extension cells separately. The third panel
separates the accounting part of the fall in the headline count (the same
evidence read inside a Holm family widened from 69 to 75) from the part the
refit itself causes. Neither figure adds evidence to
Tables~\ref{tab:app-deploy-ed}, \ref{tab:app-deploy-lf}
and~\ref{tab:app-deploy-quarters}; they are the same cells drawn, and the
counts they carry are this family's and not the 38 of
Table~\ref{tab:app-ladder}.

\clearpage

\begin{figure}[tbp]
\centering
\includegraphics[width=\textwidth,height=0.83\textheight,keepaspectratio]{figures/stability_quarter_by_quarter__a.pdf}
% caption numbers src: results/tables/deployable_combiner_quarters.csv (per-quarter fixed_rel and exp_rel, 75 cells x 16 quarters); results/tables/deployable_combiner.csv (membership of the 69-cell primary grid, genuine_fixed and genuine_exp)
\caption[Quarter-by-quarter stability of the text increment]{
Quarter-by-quarter
stability of the text increment, 2022Q1--2025Q4. Each row is one of the 69
primary combination cells (disclosure channel by text model by horizon;
event-driven above the rule, long-form below, and within each model band the
horizons run $h=5$, 10, 20 downwards) and each column one test quarter. Each
heat map holds 1,104 quarter-cells. Colour is the within-quarter relative
QLIKE improvement from adding the text forecast to the recalibrated price
reference, on a symmetric-logarithmic scale; blue means text helps. 
Panel~(a)
is the committed recipe, weights fitted once on the 2020--21 validation block
and frozen; 
panel~(b) refits them before each quarter on strictly earlier
filings only. The right-hand strip records each cell's pooled verdict: genuine
under frozen weights in 36 of 69 cells, under expanding weights in 6. % src: results/tables/deployable_combiner.csv (genuine_fixed and genuine_exp restricted to in_primary_grid)
White dots mark the 50 quarter-cells outside $\pm30$ percentage points, drawn
at the colour limit; none in panel~(a), whose extremes are $-18.9$/$+14.1$. % src: results/tables/deployable_combiner_quarters.csv (min and max of fixed_rel over the 69 primary cells)
Basis: seed-ensemble forecasts,
volatility-unit QLIKE, day-clustered inference. The right-hand strip is the
one part of this figure carrying a survivor verdict, and its two columns come
from the deployable families F9 and F10, whose Holm correction runs over all
75 deployable cells; the 36 and the 6 are then counted over the 69 primary
cells shown here, so the denominator printed is not the denominator corrected
over. The heat maps themselves mark no significance.
}
\label{fig:app-quarterly}
\end{figure}

\begin{figure}[tbp]
\centering
\includegraphics[width=\textwidth,height=0.83\textheight,keepaspectratio]{figures/stability_quarter_by_quarter__b.pdf}
% caption numbers src: results/tables/deployable_combiner_quarters.csv (per-quarter fixed_rel and exp_rel, 75 cells x 16 quarters); results/tables/deployable_combiner.csv (membership of the 69-cell primary grid, genuine_fixed and genuine_exp)
\caption[Quarter-by-quarter stability of the text increment (c, d)]{
Quarter-by-quarter
stability of the text increment, 2022Q1--2025Q4, continued.
Panels~(c) and~(d) give the cross-cell median with the interquartile band over
the sixteen test quarters. The 69 primary combination cells behind that median
are the ones drawn as heat maps in Figure~\ref{fig:app-quarterly}. Under
frozen weights that median is positive in all sixteen quarters and spans only
$+0.1$ to $+1.2$ percentage points. Under the deployable path it falls to
$-14.5$ in 2022Q1, recovers through 2023--24, and breaks again to $-5.5$ in
2025Q1. % src: results/tables/deployable_combiner_quarters.csv (cross-cell median of fixed_rel and exp_rel by quarter over the 69 primary cells)
The colour bar beneath the panels is the shared scale for those heat maps:
within-quarter relative QLIKE improvement from adding the text forecast to the
recalibrated price reference, on a symmetric-logarithmic scale, blue meaning
text helps. Basis: seed-ensemble forecasts, volatility-unit QLIKE,
day-clustered inference. Neither panel here marks significance; the survivor
verdict is the right-hand strip of Figure~\ref{fig:app-quarterly}.
}
\label{fig:app-quarterly-b}
\end{figure}

The point of drawing 1,104 quarter-cells twice is that the two panels answer
two different objections with one picture.
Figures~\ref{fig:app-quarterly} and~\ref{fig:app-quarterly-b} answer the charge that the committed
verdict is an artefact of one favourable window. Under frozen weights the
cross-cell median is positive in all sixteen test quarters and never exceeds
1.2 percentage points, so the increment is small, persistent and not
window-selected. Panel~(b) refuses the comfort that would follow. Under the
weight scheme a forecaster could actually have run, the same 69 cells collapse
in 2022Q1 (when the refitting window is at its shortest), recover through
the middle of the test span, and break again in 2025Q1. The sensitivity is
displayed rather than smoothed; the report treats the frozen recipe as primary
and the expanding one as the deployability check, which
Chapter~\ref{ch:results} reports and Chapter~\ref{ch:conclusions} carries into
its verdict.

What the figure cannot support is any per-quarter claim, and it carries no
per-quarter significance at all. A quarter-cell holds 465 to 504 filings in
the long-form channel and 1,159 to 1,822 in the event-driven channel. % src: results/tables/deployable_combiner_quarters.csv (per-quarter n, by channel, over the 69 primary cells)
A per-quarter day-clustered test at $h=20$ rests on roughly sixty filing days,
so a single unusual day moves a whole square.
That is why the colour scale is symmetric-logarithmic, why the fifty extreme
quarter-cells are marked rather than allowed to dominate the range, and why
the pooled statistics of Figures~\ref{fig:app-freeze-point} and~\ref{fig:app-freeze-point-b} rather than
anything here are the inferential evidence. The figure also does not compare
text models with one another: every panel compares the same text forecast
against the same recalibrated price reference, so a warm row means that cell's
text arm helped, not that it beat its neighbours.

\begin{figure}[tbp]
\centering
\includegraphics[width=\textwidth,height=0.83\textheight,keepaspectratio]{figures/stability_freeze_point_deployability__a.pdf}
% caption numbers src: results/tables/deployable_combiner.csv (fixed_pooled_rel, exp_pooled_rel, genuine_fixed, genuine_exp, deploy_status, m1_genuine, in_primary_grid)
\caption[What moves when the combiner's freeze point moves]{
What moves when
the combiner's freeze point moves. 
Panel~(a) plots, for each of the 75 cells
of the deployable family, the pooled relative QLIKE improvement from adding
text under frozen validation weights against the same quantity under expanding
deployable weights. Filled markers are the 69 primary-grid cells and open
markers the six frozen-embedding extension cells outside it. Fifty-four of the
75 fall below the line of no change. The collapse concentrates in the
classical pipelines: long-form bag-of-words moves from $+1.4$ to $+3.0$ per
cent frozen to $-13.5$ to $-17.4$ per cent refitted. The prompted
cells sit almost exactly on the line; three markedly improve: two primary-grid
and the third the unlabelled open marker at far left ($-6.4$ to $+0.7$), null
under both schemes. % src: results/tables/deployable_combiner.csv (B1_bow_ridge long_form fixed 1.65/1.44/2.99 against expanding -17.44/-15.06/-13.46; over all 75 cells the three improvements are long_form D1_concat_mlp h=5 from -1.04 to +6.20, long_form C5_qwen3 h=20 from -6.39 to +0.69 - an extension cell - and long_form C2_finbert_s1 h=5 from +1.90 to +8.58; the fourth largest move is +3.6 pp)
}
\label{fig:app-freeze-point}
\end{figure}

\begin{figure}[tbp]
\centering
\includegraphics[width=\textwidth,height=0.83\textheight,keepaspectratio]{figures/stability_freeze_point_deployability__b.pdf}
% caption numbers src: results/tables/deployable_combiner.csv (fixed_pooled_rel, exp_pooled_rel, genuine_fixed, genuine_exp, deploy_status, m1_genuine, in_primary_grid)
\caption[What moves when the combiner's freeze point moves (b, c)]{Continued from Figure~\ref{fig:app-freeze-point}. What moves when
the combiner's freeze point moves. 
Panel~(b) traces the 69 primary cells: of the 36 the frozen recipe calls
genuine, 5 survive the refit and 31 are lost, 1 further cell is gained and 32
are null under both. % src: results/tables/deployable_combiner.csv (deploy_status over the 69 in_primary_grid cells)
Panel~(c) separates the two reasons the headline count falls: 38 to 36 is an
accounting change, the same evidence read inside a Holm family widened from 69
to 75 cells, whereas 36 to 6 is the refit itself. Genuine means a day-clustered
statistic below zero, a Holm-adjusted $p$ below .05 within the scheme's own
pre-declared family, and a permuted-text placebo statistic below 2 in
magnitude. Basis: seed-ensemble forecasts where multi-seed runs exist,
volatility-unit QLIKE. The scheme's own family is F9 for frozen weights and
F10 for expanding ones, each corrected over all 75 deployable cells; the
transitions counted in panel~(b) are then read over the 69 primary cells, the
two denominators panel~(c) exists to separate.
}
\label{fig:app-freeze-point-b}
\end{figure}

Panel~(c) of Figure~\ref{fig:app-freeze-point-b} is the panel a marker should
look at first, because it disaggregates a fall that is easy to overstate. The
report's deployability sentence moves a count from 38 to 6, and only one of
those two steps is about deployment. The step from 38 to 36 is an accounting
change: the same cells and the same statistics, read inside a Holm family
widened from the 69-cell primary grid to the 75-cell deployable family that
adds six frozen-embedding extension cells. The step from 36 to 6 is the
substantive one. Reporting the two together as a single collapse would
overstate what refitting costs; reporting the second without the first would
leave a reader unable to reconcile the counts across chapters. Both are drawn;
the two denominators (69 and 75) must not be swapped.

Panel~(a) names cells rather than summarising them, because the collapse is
not uniform. The classical bag-of-words pipelines lose most, their increments
inverting from small positives to double-digit negatives, which is what one
expects of weights fitted on a short and unrepresentative early window. The
prompted arms are almost unmoved. Three cells are markedly better under
the refit than under the freeze: two of the 69 primary-grid cells, long-form
concatenation fusion at $h=5$ and long-form FinBERT at $h=5$, and one of the
six frozen-embedding extension cells, long-form Qwen3-embedding at
$h=20$. % src: results/tables/deployable_combiner.csv (exp_pooled_rel minus fixed_pooled_rel over the 75 cells: +7.24 for D1_concat_mlp long_form h=5, +7.08 for C5_qwen3 long_form h=20 with in_primary_grid False, +6.68 for C2_finbert_s1 long_form h=5; next largest is +3.61)

Two limits on the reading. The figure does not compare text models: the text
forecasts are identical across the two schemes and only the weight-estimation
rule changes, so a cell's movement is a property of its combiner and not of
its representation. And the expanding-scheme statistic is not a frozen-weight
comparison of two forecasts: it compares two forecasting \emph{methods}
inclusive of their recursive estimation schemes, because the sixteen
deployable blocks are concatenated into one pseudo-out-of-sample path. The
figure offers no explanation of why the prompted arms are stable under the
refit; that lies outside what these two tables can show, and the report does
not supply one.

\begin{figure}[tbp]
\centering
\includegraphics[width=\textwidth,height=0.83\textheight,keepaspectratio]{figures/stability_weight_window_sensitivity__a.pdf}
% caption numbers src: results/tables/freeze_window_sensitivity.csv and .md (40 cells x 3 fit windows, day-clustered DM, scored against the recalibrated HAR); results/tables/valwindow_sensitivity.csv and .md (69-cell grid x calm/COVID/committed, scored against the fitted five-model price pool on the archived single-seed basis)
\caption[How much of the increment depends on the weight-fitting window]{
How
much of the increment depends on the window the combination weights are fitted
on. Both experiments hold the trained text forecasts and the test rows fixed
and move only the rows the weights are estimated on. 
Panel~(a) moves the
freeze origin across the 40 cells that already carry a confirmed increment.
Those 40 are the \emph{archived single-seed} 38 genuine cells plus the six
prompted-arm cells, deduplicated --- not the seed-ensemble 38 of the primary
basis, which under the identical construction gives 39 and a different set;
the two 38-cell sets overlap in 30. The window experiment was run and frozen
against the archived screen, so its cell list is that screen's. Each thin line is one cell and the heavy line the median.
Removing the acute pandemic months (January--June 2020) shifts the typical
cell by a paired $+0.10$ percentage points and costs a net eight of the 35 --- eleven stop clearing and three start
significant cells; the unpaired median is $+1.48$ either way. Refitting on the
2018--19 tail of
the training era drives that median to $-15.96$ and leaves 3 of 32
significant, eight further cells having too few rows there to fit. % src: results/tables/freeze_window_sensitivity.csv (median rel_impr_pct and count of clustered DM < 0 with p < .05 by fit_window; median paired val_ex_h1 minus val_full difference = +0.1008pp, positive in 26/40; the unpaired column medians are 1.4794 and 1.4847, both printed as +1.48 on the panel)
Bases differ: panel~(a) is scored against the recalibrated HAR with
day-clustered inference; panel~(b), Figure~\ref{fig:app-weight-window-b},
against the fitted five-model price pool on the archived single-seed basis.
Neither carries a Holm family: panel~(a)'s significant counts are uncorrected
cell by cell, and neither revises a survivor count.
}
\label{fig:app-weight-window}
\end{figure}

\begin{figure}[tbp]
\centering
\includegraphics[width=\textwidth,height=0.83\textheight,keepaspectratio]{figures/stability_weight_window_sensitivity__b.pdf}
% caption numbers src: results/tables/freeze_window_sensitivity.csv and .md (40 cells x 3 fit windows, day-clustered DM, scored against the recalibrated HAR); results/tables/valwindow_sensitivity.csv and .md (69-cell grid x calm/COVID/committed, scored against the fitted five-model price pool on the archived single-seed basis)
\caption[How much of the increment depends on the weight-fitting window (b)]{Continued from Figure~\ref{fig:app-weight-window}. How
much of the increment depends on the window the combination weights are fitted
on. Both experiments hold the trained text forecasts and the test rows fixed
and move only the rows the weights are estimated on. 
Panel~(b) turns to which half of the validation block the weights see, over the
69-cell grid, as paired differences. The calm half against the
pandemic half averages $+0.43$ percentage points, 37 of 69 cells favour the
calm fit, paired $t = +1.12$, $p = 0.268$, and the mean increment is negative
under both halves ($-0.09$ and $-0.52$). % src: results/tables/valwindow_sensitivity.md (mean +0.433pp, median +0.064pp, calm higher in 37/69, paired t = +1.12, p = 0.268, calm -0.087pp and COVID -0.519pp)
Bases differ: panel~(a), Figure~\ref{fig:app-weight-window}, is scored against
the recalibrated HAR with day-clustered inference; panel~(b) against the fitted
five-model price pool on the archived single-seed basis. Neither carries a Holm
family: panel~(a)'s significant counts are uncorrected cell by cell, and neither
revises a survivor count.
}
\label{fig:app-weight-window-b}
\end{figure}

\begin{figure}[tbp]
\centering
\includegraphics[width=\textwidth,height=0.83\textheight,keepaspectratio]{figures/stability_weight_window_sensitivity__c.pdf}
% caption numbers src: results/tables/freeze_window_sensitivity.csv and .md (40 cells x 3 fit windows, day-clustered DM, scored against the recalibrated HAR); results/tables/valwindow_sensitivity.csv and .md (69-cell grid x calm/COVID/committed, scored against the fitted five-model price pool on the archived single-seed basis)
\caption[How much of the increment depends on the weight-fitting window (c)]{Continued from Figure~\ref{fig:app-weight-window}. How
much of the increment depends on the window the combination weights are fitted
on. Both experiments hold the trained text forecasts and the test rows fixed
and move only the rows the weights are estimated on. 
Panel~(c) is the shared scoreboard of sign survival.
Bases differ: panel~(a), Figure~\ref{fig:app-weight-window}, is scored against
the recalibrated HAR with day-clustered inference; panel~(b),
Figure~\ref{fig:app-weight-window-b}, against the fitted five-model price pool
on the archived single-seed basis. Neither carries a Holm family: panel~(a)'s
significant counts are uncorrected cell by cell, and neither revises a survivor
count.
}
\label{fig:app-weight-window-c}
\end{figure}

Figures~\ref{fig:app-weight-window}, \ref{fig:app-weight-window-b} and~\ref{fig:app-weight-window-c} close an objection the project raised
against itself: that combination weights fitted on a validation block
containing the 2020 market disruption might have credited text with a regime
the test period does not repeat. Panel~(b) is the direct test, and its answer
is a qualified no. The calm-half fit does credit text slightly more, which is
the direction the objection predicts, but the difference is under half a
percentage point on average, favours the calm half in barely more than half
the cells, and does not approach significance. The mean increment is
negative under both halves, so neither fit produces the positive the objection
worries about.

Panel~(a) sits beside it as a boundary rather than as a second test of the
same thing, and the distinction matters enough to state twice. Removing the
pandemic months from the validation window is a legitimate perturbation, and
it barely moves the increment. Refitting on the tail of the \emph{training}
era is not a legitimate perturbation of the same kind. Those rows were used to
train every text model, so a combiner fitted there is fitting memorised
predictions rather than out-of-sample forecasts, and the resulting collapse
measures memorisation and not regime dependence. The source table says so
explicitly, this appendix repeats it, and the number is reported rather than
omitted, because dropping an adverse-looking arm on the grounds that it is
uninterpretable is exactly the practice this report's validation chapter
argues against.

Two things the figure does not show. The panels use different denominators and
different references, so their heights are not comparable. Panel~(a) is
conditioned on cells that already carry an increment and is scored against the
recalibrated HAR, whereas panel~(b) is the full 69-cell grid (most of it
null before any window is moved), scored against the fitted five-model
pool. And neither panel re-runs a model or claims significance for panel~(b),
which reports point estimates and a single paired test. The alternative
2018--19 window fits in only 57 of 69 cells in that arm and is reported for
completeness rather than as a comparison.

\FloatBarrier
\section{The Label-Proxy Relabelling, Cell by Cell}
\label{app:res-rangebased-cells}

Table~\ref{tab:app-rangebased} gives the range-based relabelling as four
summary panels; the cells that occupy its conjunction are named in the
cell-by-cell tables below rather than in it.
Tables~\ref{tab:app-range-ed} and~\ref{tab:app-range-lf} give the same
exercise cell by cell, which is the resolution at which the perturbation
can be told apart from noise: a reader can see that the fall from 38
genuine cells to 19 is not concentrated in the cells with the largest
increments, that several increments \emph{grow} under the new proxy, and
that the minimum detectable effect rises in most cells. Part of the
fall is therefore power and part is signal, cell by cell rather than on average.

Two patterns are worth naming. The long-form B2 TF--IDF cells, which hold
three of the grid's five largest primary-rung increments, and the long-form C6
prompted cells, which do not, mostly survive the
relabelling with their pool verdicts intact.
% src: results/tables/rangebased_cascade.csv (old_rel ranked over the 38 genuine cells: long_form B2 at ranks 1, 3 and 5; long_form C6 at 14, 18 and 36)
The fusion arms swing
hardest: long-form D4 at $h{=}20$ moves from $+0.60$ to $-0.89$ per cent
and its DM statistic from $-3.58$ to $+2.95$. And the event-driven channel's
minimum detectable effects rise proportionally further than the long-form
channel's, though that does not show in the genuine counts: under Parkinson
both channels lose exactly half (event-driven 8 to 4, long-form 30 to 15), and
under Garman--Klass the long-form channel loses the larger share (8 to 6
against 30 to 15).
% src: results/tables/rangebased_cascade.csv (genuine by channel under old/pk/gk)
Event-driven B2 at $h{=}20$ goes from 1.36 to
1.60 per cent against an increment that falls from $+1.84$ to $+1.23$.

\begin{table}[p]
\centering
\caption[The range-based relabelling, event-driven cells]{The range-based
relabelling cell by cell, \textbf{event-driven channel}, all 24 cells.
``old'' is the committed close-to-close label, ``PK'' the Parkinson
relabelling and ``GK'' the Garman--Klass restatement. Text forecasts are
frozen throughout and only the labels, the price references, the
recalibration and the combiner weights change, an asymmetry conservative
for text. \emph{Positive rel.\% and negative DM mean text improves.} The
three detection columns give the primary, firm-identity and maximal-pool
verdicts of that cell as a triple, Y for a Holm survivor and a dot
otherwise, so a cell that clears all three is ``YYY''. ``plac.'' is the
Parkinson label-shuffle placebo, pass band $|\mathrm{DM}|<2$; MDE is the
per-cell minimum detectable effect at 80\% power. Row supports are the full
event-driven panel, unchanged by the relabelling to within a single
observation.}
\label{tab:app-range-ed}
\singlespacing\footnotesize
\setlength{\tabcolsep}{3pt}
% src: results/tables/rangebased_cascade.csv (old_*, pk_*, gk_* columns, all 24 event-driven rows)
\begin{tabular}{@{}lrrrrrrlllrrr@{}}
\toprule
& & \multicolumn{3}{c}{Primary rel.\%} & \multicolumn{2}{c}{Primary DM} & \multicolumn{3}{c}{Detections (prim./firm/pool)} & & \multicolumn{2}{c}{MDE \%} \\
\cmidrule(lr){3-5}\cmidrule(lr){6-7}\cmidrule(lr){8-10}\cmidrule(lr){12-13}
Model & $h$ & old & PK & GK & old & PK & old & PK & GK & plac. & old & PK \\
\midrule
B1 BoW ridge & 5 & $+1.33$ & $+1.44$ & $+1.49$ & $-3.35$ & $-3.32$ & YY$\cdot$ & YY$\cdot$ & YY$\cdot$ & $+1.13$ & $0.80$ & $0.83$ \\
B1 BoW ridge & 10 & $+1.23$ & $+0.83$ & $+0.85$ & $-3.25$ & $-1.39$ & YY$\cdot$ & $\cdot$Y$\cdot$ & $\cdot$Y$\cdot$ & $+0.02$ & $0.82$ & $1.00$ \\
B1 BoW ridge & 20 & $+1.53$ & $+0.95$ & $+0.97$ & $-3.10$ & $-1.19$ & Y$\cdot$$\cdot$ & $\cdot$$\cdot$$\cdot$ & $\cdot$$\cdot$$\cdot$ & $+0.87$ & $1.22$ & $1.33$ \\
\addlinespace
B2 TF--IDF ridge & 5 & $+1.21$ & $+1.35$ & $+1.44$ & $-2.76$ & $-2.04$ & $\cdot$$\cdot$$\cdot$ & $\cdot$Y$\cdot$ & $\cdot$Y$\cdot$ & $+0.98$ & $0.76$ & $0.95$ \\
B2 TF--IDF ridge & 10 & $+1.35$ & $+1.11$ & $+1.12$ & $-2.97$ & $-1.46$ & $\cdot$Y$\cdot$ & $\cdot$Y$\cdot$ & $\cdot$Y$\cdot$ & $+0.40$ & $0.87$ & $1.14$ \\
B2 TF--IDF ridge & 20 & $+1.84$ & $+1.23$ & $+1.25$ & $-3.11$ & $-1.07$ & Y$\cdot$$\cdot$ & $\cdot$$\cdot$$\cdot$ & $\cdot$$\cdot$$\cdot$ & $+0.36$ & $1.36$ & $1.60$ \\
\addlinespace
B3 L--M linear & 5 & $+0.25$ & $+0.43$ & $+0.45$ & $-2.43$ & $-1.83$ & $\cdot$$\cdot$$\cdot$ & $\cdot$$\cdot$$\cdot$ & $\cdot$$\cdot$$\cdot$ & $-0.22$ & $0.20$ & $0.34$ \\
B3 L--M linear & 10 & $+0.20$ & $+0.23$ & $+0.23$ & $-1.23$ & $-0.60$ & $\cdot$$\cdot$$\cdot$ & $\cdot$$\cdot$$\cdot$ & $\cdot$$\cdot$$\cdot$ & $+0.43$ & $0.36$ & $0.60$ \\
B3 L--M linear & 20 & $+0.26$ & $+0.20$ & $+0.12$ & $-1.12$ & $-0.24$ & $\cdot$$\cdot$$\cdot$ & $\cdot$$\cdot$$\cdot$ & $\cdot$$\cdot$$\cdot$ & $+0.79$ & $0.55$ & $0.91$ \\
\addlinespace
B4 L--M features & 5 & $+0.18$ & $+0.32$ & $+0.36$ & $-0.99$ & $-2.86$ & $\cdot$$\cdot$$\cdot$ & $\cdot$Y$\cdot$ & YYY & $-0.03$ & $0.18$ & $0.19$ \\
B4 L--M features & 10 & $+0.08$ & $+0.10$ & $+0.11$ & $-2.10$ & $-2.99$ & $\cdot$$\cdot$$\cdot$ & $\cdot$$\cdot$$\cdot$ & YYY & $+0.03$ & $0.08$ & $0.08$ \\
B4 L--M features & 20 & $+0.25$ & $+0.27$ & $+0.25$ & $-2.01$ & $-1.36$ & $\cdot$$\cdot$$\cdot$ & $\cdot$$\cdot$$\cdot$ & $\cdot$$\cdot$$\cdot$ & $-0.55$ & $0.37$ & $0.56$ \\
\addlinespace
C2 FinBERT (S1) & 5 & $+2.57$ & $+2.51$ & $+2.59$ & $-5.92$ & $-4.73$ & Y$\cdot$$\cdot$ & YYY & YY$\cdot$ & $+0.29$ & $0.86$ & $1.04$ \\
C2 FinBERT (S1) & 10 & $+2.42$ & $+2.29$ & $+2.37$ & $-5.68$ & $-3.60$ & Y$\cdot$$\cdot$ & Y$\cdot$$\cdot$ & Y$\cdot$$\cdot$ & $+0.54$ & $0.93$ & $1.38$ \\
C2 FinBERT (S1) & 20 & $+1.58$ & $-1.45$ & $-1.90$ & $-1.94$ & $+0.80$ & $\cdot$$\cdot$$\cdot$ & $\cdot$$\cdot$$\cdot$ & $\cdot$$\cdot$$\cdot$ & $-0.12$ & $2.63$ & $3.26$ \\
\addlinespace
C6 prompted LLM & 5 & $+1.21$ & $+1.41$ & $+1.40$ & $-5.04$ & $-3.67$ & YY$\cdot$ & YY$\cdot$ & YY$\cdot$ & $-0.47$ & $0.50$ & $0.69$ \\
C6 prompted LLM & 10 & $+1.00$ & $+0.84$ & $+0.81$ & $-3.76$ & $-1.54$ & YY$\cdot$ & $\cdot$$\cdot$$\cdot$ & $\cdot$$\cdot$$\cdot$ & $+2.02$ & $0.56$ & $0.81$ \\
C6 prompted LLM & 20 & $+0.66$ & $+0.49$ & $+0.46$ & $-1.98$ & $-0.45$ & $\cdot$Y$\cdot$ & $\cdot$$\cdot$$\cdot$ & $\cdot$$\cdot$$\cdot$ & $+0.60$ & $0.70$ & $1.04$ \\
\addlinespace
D2 Gated fusion & 5 & $+0.56$ & $-0.06$ & $+0.63$ & $-1.70$ & $+0.86$ & $\cdot$$\cdot$$\cdot$ & $\cdot$$\cdot$$\cdot$ & $\cdot$$\cdot$$\cdot$ & $+0.74$ & $0.37$ & $0.50$ \\
D2 Gated fusion & 10 & $+0.18$ & $-1.09$ & $-0.68$ & $-0.42$ & $+4.22$ & $\cdot$$\cdot$$\cdot$ & $\cdot$$\cdot$$\cdot$ & $\cdot$$\cdot$$\cdot$ & $+0.49$ & $0.33$ & $0.72$ \\
D2 Gated fusion & 20 & $-2.76$ & $-0.17$ & $-0.11$ & $+3.55$ & $+1.84$ & $\cdot$$\cdot$$\cdot$ & $\cdot$$\cdot$$\cdot$ & $\cdot$$\cdot$$\cdot$ & $-0.26$ & $2.25$ & $0.27$ \\
\addlinespace
D4 price+C6 & 5 & $-0.01$ & $-0.35$ & $-0.05$ & $+3.37$ & $+2.62$ & $\cdot$$\cdot$$\cdot$ & $\cdot$$\cdot$$\cdot$ & $\cdot$$\cdot$$\cdot$ & $+0.49$ & $0.01$ & $0.44$ \\
D4 price+C6 & 10 & $-0.01$ & $-0.31$ & $-0.02$ & $+0.66$ & $+2.64$ & $\cdot$$\cdot$$\cdot$ & $\cdot$$\cdot$$\cdot$ & $\cdot$$\cdot$$\cdot$ & $+0.14$ & $0.04$ & $0.38$ \\
D4 price+C6 & 20 & $-0.35$ & $-0.09$ & $+0.08$ & $+4.69$ & $+1.54$ & $\cdot$$\cdot$$\cdot$ & $\cdot$$\cdot$$\cdot$ & $\cdot$$\cdot$$\cdot$ & $-0.32$ & $0.20$ & $0.22$ \\
\bottomrule
\end{tabular}
\end{table}

\begin{table}[p]
\centering
\caption[The range-based relabelling, long-form cells]{The range-based
relabelling cell by cell, \textbf{long-form channel}, all 45 cells.
Columns, proxies, sign rules and conventions exactly as in
Table~\ref{tab:app-range-ed}.}
\label{tab:app-range-lf}
\singlespacing\footnotesize
\setlength{\tabcolsep}{3pt}
% src: results/tables/rangebased_cascade.csv (old_*, pk_*, gk_* columns, all 45 long-form rows)
\begin{tabular}{@{}lrrrrrrlllrrr@{}}
\toprule
& & \multicolumn{3}{c}{Primary rel.\%} & \multicolumn{2}{c}{Primary DM} & \multicolumn{3}{c}{Detections (prim./firm/pool)} & & \multicolumn{2}{c}{MDE \%} \\
\cmidrule(lr){3-5}\cmidrule(lr){6-7}\cmidrule(lr){8-10}\cmidrule(lr){12-13}
Model & $h$ & old & PK & GK & old & PK & old & PK & GK & plac. & old & PK \\
\midrule
B1 BoW ridge & 5 & $+1.65$ & $+1.85$ & $+1.99$ & $-3.83$ & $-3.11$ & Y$\cdot$$\cdot$ & $\cdot$$\cdot$$\cdot$ & Y$\cdot$$\cdot$ & $-0.18$ & $1.11$ & $1.19$ \\
B1 BoW ridge & 10 & $+1.44$ & $+1.17$ & $+1.33$ & $-4.15$ & $-2.70$ & Y$\cdot$$\cdot$ & $\cdot$$\cdot$$\cdot$ & $\cdot$$\cdot$$\cdot$ & $+1.25$ & $0.81$ & $0.79$ \\
B1 BoW ridge & 20 & $+2.99$ & $+2.82$ & $+3.02$ & $-5.45$ & $-3.10$ & Y$\cdot$$\cdot$ & $\cdot$$\cdot$$\cdot$ & $\cdot$$\cdot$$\cdot$ & $-0.07$ & $1.38$ & $1.72$ \\
\addlinespace
B2 TF--IDF ridge & 5 & $+3.33$ & $+4.10$ & $+4.49$ & $-5.39$ & $-5.34$ & Y$\cdot$Y & Y$\cdot$Y & Y$\cdot$Y & $-0.84$ & $1.67$ & $2.10$ \\
B2 TF--IDF ridge & 10 & $+3.48$ & $+3.23$ & $+3.61$ & $-8.89$ & $-6.19$ & Y$\cdot$Y & Y$\cdot$Y & Y$\cdot$Y & $+2.27$ & $1.25$ & $1.44$ \\
B2 TF--IDF ridge & 20 & $+5.92$ & $+5.13$ & $+5.47$ & $-9.04$ & $-4.95$ & Y$\cdot$Y & Y$\cdot$Y & Y$\cdot$Y & $+0.86$ & $1.86$ & $2.35$ \\
\addlinespace
B3 L--M linear & 5 & $+0.49$ & $+0.96$ & $+0.94$ & $-2.58$ & $-1.96$ & $\cdot$$\cdot$$\cdot$ & $\cdot$$\cdot$$\cdot$ & $\cdot$$\cdot$$\cdot$ & $+1.59$ & $1.70$ & $1.43$ \\
B3 L--M linear & 10 & $+1.79$ & $+1.76$ & $+1.74$ & $-4.62$ & $-2.84$ & YY$\cdot$ & $\cdot$$\cdot$$\cdot$ & $\cdot$$\cdot$$\cdot$ & $+0.79$ & $1.70$ & $1.38$ \\
B3 L--M linear & 20 & $+3.48$ & $+3.07$ & $+3.06$ & $-4.69$ & $-2.09$ & Y$\cdot$$\cdot$ & $\cdot$$\cdot$$\cdot$ & $\cdot$$\cdot$$\cdot$ & $-1.35$ & $1.87$ & $1.82$ \\
\addlinespace
B4 L--M features & 5 & $+0.11$ & $+0.39$ & $+0.42$ & $+1.02$ & $+0.71$ & $\cdot$$\cdot$$\cdot$ & $\cdot$$\cdot$$\cdot$ & $\cdot$$\cdot$$\cdot$ & $+1.44$ & $0.21$ & $0.34$ \\
B4 L--M features & 10 & $-0.92$ & $-1.73$ & $-1.72$ & $+3.32$ & $+2.88$ & $\cdot$$\cdot$$\cdot$ & $\cdot$$\cdot$$\cdot$ & $\cdot$$\cdot$$\cdot$ & $-0.07$ & $0.57$ & $0.82$ \\
B4 L--M features & 20 & $-1.92$ & $-2.57$ & $-2.43$ & $+3.38$ & $+1.76$ & $\cdot$$\cdot$$\cdot$ & $\cdot$$\cdot$$\cdot$ & $\cdot$$\cdot$$\cdot$ & $-3.58$ & $0.70$ & $1.05$ \\
\addlinespace
C1 BERT (S1) & 5 & $+3.42$ & $+3.87$ & $+4.09$ & $-3.25$ & $-2.73$ & Y$\cdot$$\cdot$ & $\cdot$$\cdot$$\cdot$ & $\cdot$$\cdot$$\cdot$ & $+1.73$ & $2.99$ & $4.06$ \\
C1 BERT (S1) & 10 & $-0.85$ & $-0.11$ & $-0.03$ & $+3.06$ & $+0.63$ & $\cdot$$\cdot$$\cdot$ & $\cdot$$\cdot$$\cdot$ & $\cdot$$\cdot$$\cdot$ & $+1.14$ & $0.84$ & $0.91$ \\
C1 BERT (S1) & 20 & $+2.70$ & $+2.24$ & $+2.31$ & $-6.96$ & $-3.67$ & Y$\cdot$$\cdot$ & Y$\cdot$Y & Y$\cdot$Y & $+0.97$ & $1.27$ & $1.58$ \\
\addlinespace
C2 FinBERT (S1) & 5 & $+1.90$ & $+1.33$ & $+1.47$ & $-4.53$ & $-2.06$ & Y$\cdot$$\cdot$ & $\cdot$$\cdot$$\cdot$ & $\cdot$$\cdot$$\cdot$ & $-0.12$ & $0.95$ & $1.26$ \\
C2 FinBERT (S1) & 10 & $+2.62$ & $+2.67$ & $+2.97$ & $-6.67$ & $-4.23$ & Y$\cdot$$\cdot$ & Y$\cdot$Y & Y$\cdot$Y & $-0.73$ & $1.27$ & $1.71$ \\
C2 FinBERT (S1) & 20 & $-0.59$ & $-3.92$ & $-4.49$ & $-0.40$ & $+1.77$ & $\cdot$$\cdot$$\cdot$ & $\cdot$$\cdot$$\cdot$ & $\cdot$$\cdot$$\cdot$ & $+1.29$ & $2.82$ & $4.02$ \\
\addlinespace
C2 FinBERT (S2) & 5 & $+1.21$ & $+1.31$ & $+1.63$ & $-5.57$ & $-6.05$ & Y$\cdot$Y & Y$\cdot$Y & Y$\cdot$Y & $+1.38$ & $0.59$ & $0.68$ \\
C2 FinBERT (S2) & 10 & $+0.48$ & $-0.18$ & $-0.16$ & $-7.59$ & $+6.12$ & Y$\cdot$$\cdot$ & $\cdot$$\cdot$$\cdot$ & $\cdot$$\cdot$$\cdot$ & $+1.02$ & $0.26$ & $0.13$ \\
C2 FinBERT (S2) & 20 & $+1.68$ & $+0.47$ & $+0.45$ & $-4.89$ & $-0.97$ & Y$\cdot$$\cdot$ & $\cdot$$\cdot$$\cdot$ & $\cdot$$\cdot$$\cdot$ & $+0.00$ & $0.93$ & $0.89$ \\
\addlinespace
C2 FinBERT (S3) & 5 & $+2.90$ & $+2.02$ & $+2.38$ & $-5.36$ & $-4.13$ & Y$\cdot$$\cdot$ & Y$\cdot$Y & Y$\cdot$Y & $+0.94$ & $1.35$ & $1.71$ \\
C2 FinBERT (S3) & 10 & $+2.31$ & $+1.15$ & $+1.36$ & $-5.26$ & $-3.47$ & Y$\cdot$$\cdot$ & Y$\cdot$Y & Y$\cdot$Y & $+0.38$ & $2.36$ & $3.00$ \\
C2 FinBERT (S3) & 20 & $-3.86$ & $-8.19$ & $-9.21$ & $+1.68$ & $+3.71$ & $\cdot$$\cdot$$\cdot$ & $\cdot$$\cdot$$\cdot$ & $\cdot$$\cdot$$\cdot$ & $-1.83$ & $3.65$ & $4.31$ \\
\addlinespace
C2 FinBERT (S4) & 5 & $+1.46$ & $+1.66$ & $+1.72$ & $-4.78$ & $-6.53$ & Y$\cdot$$\cdot$ & Y$\cdot$Y & Y$\cdot$Y & $+0.44$ & $1.07$ & $0.89$ \\
C2 FinBERT (S4) & 10 & $+0.36$ & $+0.09$ & $+0.02$ & $-6.82$ & $-4.76$ & Y$\cdot$$\cdot$ & Y$\cdot$$\cdot$ & Y$\cdot$$\cdot$ & $-0.29$ & $0.22$ & $0.08$ \\
C2 FinBERT (S4) & 20 & $+3.08$ & $+2.64$ & $+2.82$ & $-8.34$ & $-5.41$ & Y$\cdot$Y & Y$\cdot$Y & Y$\cdot$Y & $+0.90$ & $1.31$ & $1.63$ \\
\addlinespace
C3 RoBERTa (S1) & 5 & $+0.30$ & $+0.24$ & $+0.38$ & $-5.13$ & $-3.64$ & Y$\cdot$$\cdot$ & Y$\cdot$$\cdot$ & Y$\cdot$$\cdot$ & $+0.07$ & $0.14$ & $0.19$ \\
C3 RoBERTa (S1) & 10 & $+1.84$ & $+1.56$ & $+1.77$ & $-5.06$ & $-3.34$ & Y$\cdot$$\cdot$ & Y$\cdot$$\cdot$ & Y$\cdot$Y & $-0.73$ & $1.11$ & $1.23$ \\
C3 RoBERTa (S1) & 20 & $+0.02$ & $+0.06$ & $+0.09$ & $-3.90$ & $-1.24$ & Y$\cdot$$\cdot$ & $\cdot$$\cdot$$\cdot$ & $\cdot$$\cdot$$\cdot$ & $+0.21$ & $0.02$ & $0.18$ \\
\addlinespace
C4 Longformer (S5) & 5 & $+1.47$ & $+1.09$ & $+1.30$ & $-6.18$ & $-3.96$ & Y$\cdot$Y & Y$\cdot$Y & Y$\cdot$$\cdot$ & $-0.49$ & $0.69$ & $0.89$ \\
C4 Longformer (S5) & 10 & $-2.95$ & $-2.25$ & $-2.20$ & $+10.32$ & $+7.45$ & $\cdot$$\cdot$$\cdot$ & $\cdot$$\cdot$$\cdot$ & $\cdot$$\cdot$$\cdot$ & $-0.15$ & $0.96$ & $1.06$ \\
C4 Longformer (S5) & 20 & $+0.91$ & $+0.13$ & $+0.05$ & $-4.38$ & $-0.51$ & Y$\cdot$$\cdot$ & $\cdot$$\cdot$$\cdot$ & $\cdot$$\cdot$$\cdot$ & $+0.84$ & $0.63$ & $0.77$ \\
\addlinespace
C6 prompted LLM & 5 & $+1.79$ & $+1.48$ & $+1.39$ & $-6.31$ & $-5.03$ & Y$\cdot$Y & Y$\cdot$Y & Y$\cdot$Y & $+1.95$ & $1.11$ & $0.87$ \\
C6 prompted LLM & 10 & $+2.25$ & $+1.46$ & $+1.42$ & $-7.92$ & $-5.24$ & Y$\cdot$Y & Y$\cdot$Y & Y$\cdot$Y & $-2.50$ & $1.07$ & $0.84$ \\
C6 prompted LLM & 20 & $+0.27$ & $+0.04$ & $-0.02$ & $-3.23$ & $-3.33$ & YY$\cdot$ & Y$\cdot$$\cdot$ & $\cdot$$\cdot$$\cdot$ & $-0.66$ & $0.37$ & $0.05$ \\
\addlinespace
D1 Concat-MLP & 5 & $-1.04$ & $+0.01$ & $+0.14$ & $+2.66$ & $-1.86$ & $\cdot$$\cdot$$\cdot$ & $\cdot$$\cdot$$\cdot$ & $\cdot$$\cdot$$\cdot$ & $+1.47$ & $1.21$ & $0.10$ \\
D1 Concat-MLP & 10 & $-0.46$ & $-0.66$ & $-1.31$ & $+3.54$ & $+3.74$ & $\cdot$$\cdot$$\cdot$ & $\cdot$$\cdot$Y & $\cdot$$\cdot$$\cdot$ & $-1.06$ & $0.62$ & $0.84$ \\
D1 Concat-MLP & 20 & $+0.12$ & $-1.39$ & $-1.74$ & $-5.57$ & $+3.66$ & Y$\cdot$Y & $\cdot$$\cdot$$\cdot$ & $\cdot$$\cdot$$\cdot$ & $-0.63$ & $0.09$ & $2.84$ \\
\addlinespace
D2 Gated fusion & 5 & $+0.19$ & $-0.37$ & $-0.25$ & $-2.11$ & $+2.45$ & $\cdot$$\cdot$$\cdot$ & $\cdot$$\cdot$$\cdot$ & $\cdot$$\cdot$$\cdot$ & $+1.20$ & $0.29$ & $0.89$ \\
D2 Gated fusion & 10 & $-0.02$ & $-2.04$ & $-3.15$ & $+2.02$ & $+6.56$ & $\cdot$$\cdot$$\cdot$ & $\cdot$$\cdot$$\cdot$ & $\cdot$$\cdot$$\cdot$ & $-0.95$ & $0.03$ & $1.40$ \\
D2 Gated fusion & 20 & $-2.27$ & $-0.22$ & $-1.15$ & $+5.56$ & $+0.78$ & $\cdot$$\cdot$$\cdot$ & $\cdot$$\cdot$$\cdot$ & $\cdot$$\cdot$$\cdot$ & $-0.13$ & $3.04$ & $0.37$ \\
\addlinespace
D4 price+C6 & 5 & $+0.15$ & $-0.80$ & $-1.24$ & $-0.75$ & $+3.57$ & $\cdot$$\cdot$$\cdot$ & $\cdot$$\cdot$$\cdot$ & $\cdot$$\cdot$$\cdot$ & $+0.95$ & $0.32$ & $1.35$ \\
D4 price+C6 & 10 & $-0.15$ & $-1.09$ & $-1.35$ & $-0.43$ & $+3.51$ & $\cdot$$\cdot$$\cdot$ & $\cdot$$\cdot$$\cdot$ & $\cdot$$\cdot$$\cdot$ & $+0.12$ & $0.84$ & $1.94$ \\
D4 price+C6 & 20 & $+0.60$ & $-0.89$ & $-1.34$ & $-3.58$ & $+2.95$ & Y$\cdot$$\cdot$ & $\cdot$$\cdot$$\cdot$ & $\cdot$$\cdot$$\cdot$ & $-0.45$ & $0.73$ & $2.46$ \\
\bottomrule
\end{tabular}
\end{table}

\clearpage
\FloatBarrier
\section{The Licence-Free Public Variant}
\label{app:res-public}

The benchmark's labels, its recalibrated reference and every price baseline
are built from licensed daily returns, so a reader without that licence
cannot rebuild the study exactly. Section~\ref{sec:data-public} prices the
coverage of the public alternative; Tables~\ref{tab:app-public}
and~\ref{tab:app-public-cells} price what the substitution does to the
verdicts, which is the question that decides whether the licence-free
variant is a usable artefact or a gesture.

The exercise separates the two things a public rebuild changes at once. On
panel~B the rows are restricted to those a public source covers, but the
labels stay licensed, isolating survivorship. On panel~C the same rows
carry public labels as well, adding label measurement error. The standalone
census is completely unmoved under the variance-unit QLIKE convention: 7 Holm winners of 180 on all three panels,
the same seven price arms, 0 of 153 for text and fusion, and \emph{zero}
per-comparison verdict flips out of 180 on either panel. The 69-cell
cascade moves a little, and not in one direction: the primary rung falls
from 38 to 35 and 36, the identity rung rises from 8 to 9, and the pool
rung rises from 9 to 11 and 13. The conjunction moves from 0 to 1. The
single cell that occupies it on both panels is the same one: long-form
B3 Loughran--McDonald linear at $h{=}10$, at $+2.09$ and $+2.04$ per cent
over the recalibrated HAR, $+0.61$ and $+0.60$ over the identity reference and
$+0.91$ over both against the price pool. % src: results/tables/public_variant_cascade.csv (long_form B3_lm_linear h=10: b_firm_rel 0.6064, c_firm_rel 0.6048, b_pool_rel 0.9125, c_pool_rel 0.9087)

That flip is the finding rather than a defect, and its location is what
makes it interpretable. The coverage the public source loses is not random:
active firms are covered on 97.3 per cent of rows and firms that left the
index on 31.5 per cent. Panel B is a mildly survivorship-thinned
panel, and the conjunction cell appears on the thinned panel and not on the
full one. A reader rebuilding the benchmark without the licence would therefore
reach the same standalone verdict comparison for comparison, and a
one-cell-weaker version of the central negative. Both facts belong in the
same sentence, which is why they are tabulated together.

\begin{table}[htbp]
\centering
\caption[Verdict objects under the licence-free rebuild]{What a
licence-free rebuild recovers. Panel A is the committed anchor; panel B
restricts the rows to those the public price source covers while keeping
the licensed labels, isolating survivorship; panel C additionally replaces
the labels and the price features with public ones. Text forecasts are
frozen throughout; the combiner, the recalibration and the firm-mean term
are refitted on each panel's own rows. All counts are Holm-corrected within
their own family. The standalone block re-runs the 180-comparison census
\textbf{under the variance-unit QLIKE convention of the third row of
Table~\ref{tab:app-census}, which is off the squared-error primary and is
tagged as such}. It is the only convention under which any challenger
beats A2, and it is why panel A here reads 7 and 153 where the squared-error
census of Table~\ref{tab:app-dm180} reads 0 and 155 on the same 180
comparisons. \emph{The A-to-B difference is survivorship and
the B-to-C difference is label measurement error}; they must not be added.}
\label{tab:app-public}
\singlespacing\footnotesize
\setlength{\tabcolsep}{3pt}
% src: results/tables/public_variant_cascade.csv (a_*, b_*, c_* detection and MDE columns); results/tables/public_variant_leaderboard.csv (panel-wise standalone counts and verdict flips; the loss column is dm_qlike_var_clu, i.e. variance-unit QLIKE, giving better_holm 7 and worse_holm 153 on panels A, B and C alike); results/tables/dm_pairwise_clustered.csv (the squared-error primary on the same 180 comparisons: better_clust 0, sig_worse_clust 155); public_variant_cascade.md coverage block
\begin{tabular}{@{}lrrrrr@{}}
\toprule
Verdict object & A & B & C & A$-$B & B$-$C \\
\midrule
\multicolumn{6}{@{}l}{\emph{(a) The 69-cell cascade, Holm survivors}} \\
Primary rung: text over recalibrated HAR & 38 & 35 & 36 & $+3$ & $-1$ \\
\quad after the placebo gate & 38 & 32 & 31 & $+6$ & $+1$ \\
Firm-identity-augmented reference & 8 & 9 & 9 & $-1$ & $0$ \\
Maximal five-model price pool & 9 & 11 & 13 & $-2$ & $-2$ \\
\textbf{Full conjunction} & \textbf{0} & \textbf{1} & \textbf{1} & $-1$ & $0$ \\
Median per-cell MDE, rel.\ QLIKE \% & 0.823 & 0.762 & 0.782 & --- & --- \\
\midrule
\multicolumn{6}{@{}l}{\emph{(b) The 180-comparison standalone census}} \\
Challengers better than A2 (Holm) & 7 & 7 & 7 & $0$ & $0$ \\
\quad of which text or fusion, of 153 & 0 & 0 & 0 & $0$ & $0$ \\
Challengers significantly worse (Holm) & 153 & 153 & 153 & $0$ & $0$ \\
Per-comparison verdict flips against A & --- & 0 of 180 & 0 of 180 & --- & --- \\
\midrule
\multicolumn{6}{@{}l}{\emph{(c) What the public source covers, of modelled rows}} \\
Training / validation / test split rows & \multicolumn{5}{l}{72.6\% / 89.3\% / 95.2\%} \\
Benchmark rows, clean coverage & \multicolumn{5}{l}{79.83\%} \\
Rows of firms in the index / of firms that left & \multicolumn{5}{l}{97.3\% / 31.5\%} \\
\bottomrule
\end{tabular}
\end{table}

\begin{table}[htbp]
\centering
\caption[Cells whose ladder verdict moves under the public rebuild]{Every
cell of the 69 whose detection triple changes on either public panel:
twelve of 69. The other 57 return the identical primary, identity and pool
verdicts on all three panels. Detections are written as a triple
(primary/identity/pool), Y for a Holm survivor and a dot otherwise, so the
conjunction cell is the one reading ``YYY''. Increments are relative QLIKE
improvements over the recalibrated HAR in volatility units. The row counts
show how much of each move is a change of sample: panel B and panel C share
their rows almost exactly, so a cell that moves between them moves on the
labels alone.}
\label{tab:app-public-cells}
\singlespacing\footnotesize
\setlength{\tabcolsep}{1.5pt}
% src: results/tables/public_variant_cascade.csv (n_test, primary rel% and detection triples on panels a, b and c)
\begin{tabular}{@{}llrrrrrrrlll@{}}
\toprule
& & & \multicolumn{3}{c}{Rows} & \multicolumn{3}{c}{Primary rel.\%} & \multicolumn{3}{c}{Detections} \\
\cmidrule(lr){4-6}\cmidrule(lr){7-9}\cmidrule(lr){10-12}
Ch. & Model & $h$ & A & B & C & A & B & C & A & B & C \\
\midrule
ED & B2 TF--IDF ridge & 20 & 24{,}732 & 23{,}551 & 23{,}547 & $+1.84$ & $+1.79$ & $+1.79$ & Y$\cdot$$\cdot$ & $\cdot$$\cdot$$\cdot$ & $\cdot$$\cdot$$\cdot$ \\
\addlinespace
LF & B2 TF--IDF ridge & 5 & 7{,}951 & 7{,}565 & 7{,}564 & $+3.33$ & $+3.24$ & $+3.20$ & Y$\cdot$Y & Y$\cdot$$\cdot$ & Y$\cdot$Y \\
LF & B3 L--M linear & 10 & 7{,}933 & 7{,}548 & 7{,}547 & $+1.79$ & $+2.09$ & $+2.04$ & YY$\cdot$ & YYY & YYY \\
LF & B3 L--M linear & 20 & 7{,}902 & 7{,}526 & 7{,}525 & $+3.48$ & $+3.91$ & $+3.87$ & Y$\cdot$$\cdot$ & YY$\cdot$ & YY$\cdot$ \\
LF & C2 FinBERT (S1) & 5 & 7{,}951 & 7{,}565 & 7{,}564 & $+1.90$ & $+1.88$ & $+1.85$ & Y$\cdot$$\cdot$ & Y$\cdot$Y & Y$\cdot$Y \\
LF & C2 FinBERT (S1) & 10 & 7{,}933 & 7{,}548 & 7{,}547 & $+2.62$ & $+2.76$ & $+2.74$ & Y$\cdot$$\cdot$ & Y$\cdot$Y & Y$\cdot$Y \\
LF & C2 FinBERT (S2) & 5 & 7{,}951 & 7{,}565 & 7{,}564 & $+1.21$ & $+0.95$ & $+0.94$ & Y$\cdot$Y & Y$\cdot$$\cdot$ & Y$\cdot$$\cdot$ \\
LF & C2 FinBERT (S4) & 5 & 7{,}951 & 7{,}565 & 7{,}564 & $+1.46$ & $+1.41$ & $+1.41$ & Y$\cdot$$\cdot$ & Y$\cdot$Y & Y$\cdot$Y \\
LF & C3 RoBERTa (S1) & 10 & 7{,}933 & 7{,}548 & 7{,}547 & $+1.84$ & $+1.90$ & $+1.91$ & Y$\cdot$$\cdot$ & Y$\cdot$Y & Y$\cdot$Y \\
LF & C3 RoBERTa (S1) & 20 & 7{,}902 & 7{,}526 & 7{,}525 & $+0.02$ & $-0.28$ & $-0.29$ & Y$\cdot$$\cdot$ & $\cdot$$\cdot$$\cdot$ & $\cdot$$\cdot$$\cdot$ \\
LF & D1 Concat-MLP & 20 & 7{,}902 & 7{,}526 & 7{,}525 & $+0.12$ & $-0.46$ & $+0.34$ & Y$\cdot$Y & $\cdot$$\cdot$$\cdot$ & Y$\cdot$$\cdot$ \\
LF & D4 price+C6 & 20 & 7{,}902 & 7{,}526 & 7{,}525 & $+0.60$ & $+0.07$ & $+0.18$ & Y$\cdot$$\cdot$ & Y$\cdot$$\cdot$ & Y$\cdot$Y \\
\bottomrule
\end{tabular}
\end{table}

Figures~\ref{fig:app-public-labels}, \ref{fig:app-public-labels-b} and~\ref{fig:app-public-labels-c} draw the
coverage gradient behind these counts, the label agreement on the rows the public
source supplies, and the verdict objects side by side.


\clearpage

\begin{figure}[tbp]
\centering
\includegraphics[width=\textwidth,height=0.83\textheight,keepaspectratio]{figures/AR4_public_label_variant__a.pdf}
% caption numbers src: results/tables/label_parity.csv (coverage, coverage_by_exit_year, ticker_status, parity_corr, verdict_agreement, verdict_combo); results/tables/public_variant_leaderboard.csv (360 comparisons over the two perturbed panels, verdict_flip_vs_A)
\caption[What a licence-free rebuild of the price layer would recover]{
What a
licence-free rebuild of the price layer would recover. Labels, recalibrated reference and price
baselines rest on licensed daily returns, so the labels were recomputed from a
public adjusted-close snapshot on the benchmark's own label windows and every
verdict object re-run on three panels: the full test panel with licensed labels, the public-coverage
intersection with licensed labels, and that intersection with public labels.
Panel~(a) gives coverage by the year a firm left the index, cohort sizes in
thousands of rows beneath each bar. 
}
\label{fig:app-public-labels}
\end{figure}

\begin{figure}[tbp]
\centering
\includegraphics[width=\textwidth,height=0.83\textheight,keepaspectratio]{figures/AR4_public_label_variant__b.pdf}
% caption numbers src: results/tables/label_parity.csv (coverage, coverage_by_exit_year, ticker_status, parity_corr, verdict_agreement, verdict_combo); results/tables/public_variant_leaderboard.csv (360 comparisons over the two perturbed panels, verdict_flip_vs_A)
\caption[What a licence-free rebuild of the price layer would recover (b)]{Continued from Figure~\ref{fig:app-public-labels}. What a
licence-free rebuild of the price layer would recover. Labels, recalibrated reference and price
baselines rest on licensed daily returns, so the labels were recomputed from a
public adjusted-close snapshot on the benchmark's own label windows and every
verdict object re-run on three panels: the full test panel with licensed labels, the public-coverage
intersection with licensed labels, and that intersection with public labels.
Panel~(b) accounts for all 431,245
rows: 80.2 per cent covered, 15.9 with no public data, 3.1 with an incomplete
window, 0.8 whose ticker was reused. 
}
\label{fig:app-public-labels-b}
\end{figure}

\begin{figure}[tbp]
\centering
\includegraphics[width=\textwidth,height=0.83\textheight,keepaspectratio]{figures/AR4_public_label_variant__c.pdf}
% caption numbers src: results/tables/label_parity.csv (coverage, coverage_by_exit_year, ticker_status, parity_corr, verdict_agreement, verdict_combo); results/tables/public_variant_leaderboard.csv (360 comparisons over the two perturbed panels, verdict_flip_vs_A)
\caption[What a licence-free rebuild of the price layer would recover (c, d)]{Continued from Figure~\ref{fig:app-public-labels}. What a
licence-free rebuild of the price layer would recover. Labels, recalibrated reference and price
baselines rest on licensed daily returns, so the labels were recomputed from a
public adjusted-close snapshot on the benchmark's own label windows and every
verdict object re-run on three panels: the full test panel with licensed labels, the public-coverage
intersection with licensed labels, and that intersection with public labels.
This part opens with the four-swatch key for the stacked bar of the previous
part, which the crop carried across the boundary with it.
Panel~(c) gives the correlation of public and
licensed log realised volatility by split and horizon; 
panel~(d) counts
verdict objects surviving on the public-coverage intersection, its columns being
that intersection under licensed and under public labels. % src: results/tables/label_parity.csv (coverage section: 345,804 / 68,497 / 13,564 / 3,380 of 431,245 rows)
The failure mode is coverage, not measurement. On covered rows the labels are
near-duplicates (0.9981 overall, 1.0000 on the test years, against 0.9714
before eighteen recycled ticker symbols are screened out). None of the 180
standalone comparisons flips on either panel; all eighteen day-clustered signs
and Holm verdicts agree; five of six per-panel QLIKE rankings are identical;
and genuine combination increments move only from 8 of 12 to 7 of 12. % src: results/tables/label_parity.csv (parity_corr, verdict_agreement, verdict_combo); results/tables/public_variant_leaderboard.csv (verdict_flip_vs_A false in all 360 rows)
Basis: the benchmark's own label windows and its committed verdict code paths,
re-run unchanged on each panel. Family F13 covers the standalone and
combination panels the public variant replicates, so it governs the
leaderboard and combination rows of panel~(d) and not the agreement counts
beside them, which compare two label sources and enter no family. % src: results/tables/holm_families.md (F13 "Public-variant panels", size "as F1/F3", carried by public_variant_cascade.csv)
}
\label{fig:app-public-labels-c}
\end{figure}

Figures~\ref{fig:app-public-labels}, \ref{fig:app-public-labels-b} and~\ref{fig:app-public-labels-c} price the project's own deliverable. The
report releases a licence-free core of the benchmark precisely so that readers
without a subscription to the licensed price data can use it, and the honest
question is what such a reader would and would not recover. Panels~(c)
and~(d) answer the optimistic half. On rows a public source covers, the labels
are the same labels to four decimal places on the test years, and every class
of verdict the report reports survives the substitution. The one movement,
from 8 genuine combination increments to 7, comes through a placebo term
rather than through a Diebold--Mariano statistic. That is a stronger result
than the project expected, and it is the reason the licence-free core is
shipped at all.

Panel~(a) answers the pessimistic half, and it governs how the release must be
described. Coverage of firms still in the index is 97.7 per cent; coverage of
firms that left it runs from 3.7 per cent in the worst exit cohort to
58.7 per cent in the best, with the three earliest cohorts all at or below
10 per cent, pooling to 31.9 per cent. % src: results/tables/label_parity.csv (coverage_by_exit_year: 0.9768 for the active cohort, 0.037 to 0.587 across exit years, 31.9 per cent pooled over exiting firms)
Since the entire point of a point-in-time universe is to keep exiting firms in
the panel, a licence-free rebuild reintroduces exactly the survivorship bias
the benchmark was constructed to remove --- unevenly, and worst in the
earliest years. A licence-free variant is therefore a survivorship-tilted
subsample and would have to be labelled as one, the more so because the screen
that removes recycled ticker symbols itself needs the licensed data to run.
The correlation of 0.9714 before that screen and 0.9981 after it shows the
screen doing real work that a public-only rebuild cannot reproduce.

Two limits close the figure. It shows no rerun of the reference ladder on
public labels: every count in it is a standalone or combination verdict,
and the identity and pool rungs are untouched. It bounds what a
licence-free reader recovers of the report's verdicts and not of its ladder.
And it settles no distribution question: the snapshot used here permits
research use but not redistribution, and its coverage percentages come from a
single fetch of a source whose composition drifts. The percentages above are
the committed fetch. The parallel recomputation recorded in the source file was run on a later
refetch of the same source.
It reads 79.83 per cent clean coverage, 97.3 per cent of active-firm rows and
31.5 per cent of exited-firm rows, a drift of 0.35 percentage points that the
source file itself records. % src: results/tables/label_parity.csv (coverage section, the committed fetch: coverage_clean 0.8019 on 345,804 of 431,245 rows, 0.9768 active, 0.319 pooled exited); results/tables/public_variant_cascade.md line 138 (committed 2026-07-09 fetch 80.19 per cent against this refetch 79.83 per cent, drift -0.353 pp, inside the 0.5 pp pre-registered threshold)
The stable object is the mechanism, not the third decimal place.

\FloatBarrier
\section{Where the Event-Driven Residual Lives}
\label{app:res-itemgeography}

The surviving positive in this study is a bounded increment of the prompted
arm over the firm-identity reference on the 8-K channel. A third of that
channel is earnings announcements, and the obvious deflationary reading is
that the model is reading a results number rather than a document.
Section~\ref{sec:res-residual} closes that reading with an
indicator control (adding an Item 2.02 flag to the reference leaves
$+0.21/+0.18/+0.20$ per cent standing), but an indicator answers a
different question from a partition. Table~\ref{tab:app-item-qwen} splits
the residual itself into six disjoint item groups and asks where it lives.

The grouping concedes the objection by construction: a filing carrying Item
2.02 is assigned to the earnings group first, so the five remaining groups
contain no results announcement at all and any increment inside them cannot
be number-parroting. On the primary family the answer is mixed and is
reported as mixed. Item 2.02 carries 33 per cent of the filings and 54 per
cent of the pooled residual summed over horizons, so more than half of it
is the earnings stream. The earnings-free pooled remainder is nonetheless
favourable and significant at all three horizons, at $+0.22/+0.22/+0.17$
per cent, and two of the fifteen narrative partition cells survive Holm on
their own: Item 7.01 Reg-FD disclosures at $h{=}10$ and $h{=}20$.
Table~\ref{tab:app-item-llama} repeats the decomposition on the
cross-family replication, where the concentration is sharper. Item 2.02
carries 69 per cent of that family's residual, the earnings-free pooled
remainder is significant at one horizon of three, and no narrative
partition cell survives Holm. Partly event reading, partly the earnings
number, and more of the latter in the larger model.

Two basis warnings attach to both tables. They sit on the \emph{full}
event-driven panel of 25{,}109,\ 25{,}001 and 24{,}732 rows, not on the
merged-grid rows the ladder is scored on (23{,}855,\ 22{,}785 and
22{,}318). The pooled row of Table~\ref{tab:app-item-qwen} therefore reads
$+0.45/+0.25/+0.20$ per cent where Table~\ref{tab:app-grid-ed}'s
identity column reads $+0.52/+0.24/+0.21$: the same arm and reference on a
wider row set, not a discrepancy. And the replication family is scored
single-pass here, at $+0.83/+0.64/+0.39$, against the three-seed ensemble's
$+0.84/+0.64/+0.38$ in Table~\ref{tab:app-crossfamily}.

\begin{table}[p]
\centering
\caption[The 8-K residual by item code, primary family]{The event-driven
residual of the \textbf{primary prompted family} (Qwen3-32B) decomposed by
8-K item code. The reference is the firm-identity-augmented recalibrated
HAR, with the firm-mean map and both regressions fitted on the full
validation split and frozen on test; test residuals are then
\emph{partitioned} by item code and nothing is refitted per stratum.
Filings are labelled by the first code present in the order 2.02, 5.02,
7.01, 8.01, 5.07, other, so the five non-2.02 groups contain no results
announcement. Volatility-unit QLIKE, day-clustered DM, \emph{negative DM
and positive rel.\% mean text improves}; Holm within the 18 disjoint
partition cells of this family, with the two pooled rows carrying raw $p$
only and excluded from Holm to avoid double counting. An asterisk marks
Holm significance. ``share'' is the signed share of the pooled absolute
QLIKE reduction carried by that group. \textbf{Support: the full
event-driven panel}, wider than the merged grid of
Table~\ref{tab:app-grid-ed}. The percentages are not interchangeable across
the two supports.}
\label{tab:app-item-qwen}
\singlespacing\footnotesize
\setlength{\tabcolsep}{3.5pt}
% src: results/tables/row11_item_stratified.csv (family = qwen3_32b; all 24 rows)
\begin{tabular}{@{}rlrrrrrrrrr@{}}
\toprule
$h$ & Item group & $n$ & days & \% filings & rel.\% & DM & raw $p$ & Holm $p$ & share \% & vs.\ HAR \% \\
\midrule
5 & \emph{All 8-K filings (pooled)} & 25{,}109 & 996 & 100.0 & $+0.45$\phantom{*} & $-5.26$ & $<$.0001 & --- & $+100.0$ & $+1.21$ \\
5 & Item 2.02 results & 8{,}348 & 821 & 33.2 & $+0.99$* & $-5.48$ & $<$.0001 & $<$.0001 & $+64.7$ & $+2.89$ \\
5 & Item 5.02 leadership & 4{,}706 & 978 & 18.7 & $+0.18$\phantom{*} & $-1.33$ & .1837 & 1.0000 & $+8.2$ & $+0.37$ \\
5 & Item 7.01 Reg-FD & 3{,}343 & 932 & 13.3 & $+0.38$\phantom{*} & $-1.48$ & .1384 & .9689 & $+11.2$ & $+0.71$ \\
5 & Item 8.01 other events & 4{,}951 & 989 & 19.7 & $+0.30$\phantom{*} & $-1.74$ & .0821 & .6564 & $+13.9$ & $+0.86$ \\
5 & Item 5.07 shareholder vote & 1{,}426 & 406 & 5.7 & $-0.02$\phantom{*} & $+0.52$ & .6019 & 1.0000 & $-0.3$ & $+0.20$ \\
5 & Other narrative items & 2{,}335 & 884 & 9.3 & $+0.11$\phantom{*} & $-0.29$ & .7686 & 1.0000 & $+2.3$ & $+0.06$ \\
5 & \emph{All non-2.02 items (pooled)} & 16{,}761 & 996 & 66.8 & $+0.22$\phantom{*} & $-2.96$ & .0031 & --- & $+35.3$ & $+0.52$ \\
\addlinespace
10 & \emph{All 8-K filings (pooled)} & 25{,}001 & 991 & 100.0 & $+0.25$\phantom{*} & $-5.16$ & $<$.0001 & --- & $+100.0$ & $+1.00$ \\
10 & Item 2.02 results & 8{,}331 & 817 & 33.3 & $+0.33$\phantom{*} & $-2.82$ & .0050 & .0646 & $+35.3$ & $+1.63$ \\
10 & Item 5.02 leadership & 4{,}681 & 973 & 18.7 & $+0.13$\phantom{*} & $-1.87$ & .0613 & .5518 & $+11.0$ & $+0.41$ \\
10 & Item 7.01 Reg-FD & 3{,}319 & 927 & 13.3 & $+0.32$* & $-3.56$ & .0004 & .0063 & $+18.1$ & $+0.84$ \\
10 & Item 8.01 other events & 4{,}920 & 984 & 19.7 & $+0.28$\phantom{*} & $-2.89$ & .0040 & .0557 & $+22.6$ & $+1.16$ \\
10 & Item 5.07 shareholder vote & 1{,}424 & 405 & 5.7 & $+0.22$\phantom{*} & $-1.93$ & .0546 & .5457 & $+5.1$ & $+1.05$ \\
10 & Other narrative items & 2{,}326 & 882 & 9.3 & $+0.19$\phantom{*} & $-1.32$ & .1875 & 1.0000 & $+8.0$ & $+0.46$ \\
10 & \emph{All non-2.02 items (pooled)} & 16{,}670 & 991 & 66.7 & $+0.22$\phantom{*} & $-4.54$ & $<$.0001 & --- & $+64.7$ & $+0.77$ \\
\addlinespace
20 & \emph{All 8-K filings (pooled)} & 24{,}732 & 981 & 100.0 & $+0.20$\phantom{*} & $-3.79$ & .0002 & --- & $+100.0$ & $+0.66$ \\
20 & Item 2.02 results & 8{,}305 & 811 & 33.6 & $+0.25$* & $-3.70$ & .0002 & .0039 & $+36.3$ & $+1.40$ \\
20 & Item 5.02 leadership & 4{,}610 & 963 & 18.6 & $+0.07$\phantom{*} & $-1.11$ & .2655 & 1.0000 & $+6.9$ & $-0.13$ \\
20 & Item 7.01 Reg-FD & 3{,}269 & 917 & 13.2 & $+0.28$* & $-3.33$ & .0009 & .0134 & $+20.2$ & $+0.69$ \\
20 & Item 8.01 other events & 4{,}842 & 974 & 19.6 & $+0.24$\phantom{*} & $-2.72$ & .0066 & .0731 & $+24.9$ & $+1.07$ \\
20 & Item 5.07 shareholder vote & 1{,}415 & 403 & 5.7 & $-0.18$\phantom{*} & $+1.00$ & .3192 & 1.0000 & $-5.0$ & $-1.98$ \\
20 & Other narrative items & 2{,}291 & 872 & 9.3 & $+0.30$\phantom{*} & $-2.75$ & .0060 & .0722 & $+16.6$ & $+0.86$ \\
20 & \emph{All non-2.02 items (pooled)} & 16{,}427 & 981 & 66.4 & $+0.17$\phantom{*} & $-3.61$ & .0003 & --- & $+63.7$ & $+0.38$ \\
\bottomrule
\end{tabular}
\end{table}

\begin{table}[p]
\centering
\caption[The 8-K residual by item code, replication family]{The same
decomposition for the \textbf{cross-family replication} (Llama-3.1-70B,
AWQ-INT4). Columns, grouping rule, reference, inference and Holm family
exactly as in Table~\ref{tab:app-item-qwen}, with Holm applied within this
family's own 18 partition cells and never pooled across the two families.
\textbf{Single-pass basis} for this family, so its pooled row reads
$+0.83/+0.64/+0.39$ against the three-seed ensemble's
$+0.84/+0.64/+0.38$ in Table~\ref{tab:app-crossfamily}.}
\label{tab:app-item-llama}
\singlespacing\footnotesize
\setlength{\tabcolsep}{3.5pt}
% src: results/tables/row11_item_stratified.csv (family = llama70_awq; all 24 rows)
\begin{tabular}{@{}rlrrrrrrrrr@{}}
\toprule
$h$ & Item group & $n$ & days & \% filings & rel.\% & DM & raw $p$ & Holm $p$ & share \% & vs.\ HAR \% \\
\midrule
5 & \emph{All 8-K filings (pooled)} & 25{,}109 & 996 & 100.0 & $+0.83$\phantom{*} & $-2.58$ & .0100 & --- & $+100.0$ & $+1.39$ \\
5 & Item 2.02 results & 8{,}348 & 821 & 33.2 & $+2.82$* & $-5.17$ & $<$.0001 & $<$.0001 & $+99.3$ & $+4.84$ \\
5 & Item 5.02 leadership & 4{,}706 & 978 & 18.7 & $+0.13$\phantom{*} & $-0.25$ & .8056 & 1.0000 & $+3.3$ & $-0.18$ \\
5 & Item 7.01 Reg-FD & 3{,}343 & 932 & 13.3 & $+0.29$\phantom{*} & $-0.79$ & .4287 & 1.0000 & $+4.7$ & $+0.29$ \\
5 & Item 8.01 other events & 4{,}951 & 989 & 19.7 & $+0.01$\phantom{*} & $+0.10$ & .9223 & 1.0000 & $+0.2$ & $+0.45$ \\
5 & Item 5.07 shareholder vote & 1{,}426 & 406 & 5.7 & $-0.68$\phantom{*} & $+1.08$ & .2814 & 1.0000 & $-5.1$ & $-1.06$ \\
5 & Other narrative items & 2{,}335 & 884 & 9.3 & $-0.21$\phantom{*} & $+0.66$ & .5066 & 1.0000 & $-2.4$ & $-0.55$ \\
5 & \emph{All non-2.02 items (pooled)} & 16{,}761 & 996 & 66.8 & $+0.01$\phantom{*} & $-0.19$ & .8460 & --- & $+0.7$ & $-0.03$ \\
\addlinespace
10 & \emph{All 8-K filings (pooled)} & 25{,}001 & 991 & 100.0 & $+0.64$\phantom{*} & $-2.00$ & .0460 & --- & $+100.0$ & $+1.17$ \\
10 & Item 2.02 results & 8{,}331 & 817 & 33.3 & $+1.30$\phantom{*} & $-1.60$ & .1110 & 1.0000 & $+55.6$ & $+3.10$ \\
10 & Item 5.02 leadership & 4{,}681 & 973 & 18.7 & $+0.11$\phantom{*} & $-0.14$ & .8889 & 1.0000 & $+3.6$ & $-0.30$ \\
10 & Item 7.01 Reg-FD & 3{,}319 & 927 & 13.3 & $+0.85$\phantom{*} & $-2.52$ & .0120 & .2040 & $+18.9$ & $+1.17$ \\
10 & Item 8.01 other events & 4{,}920 & 984 & 19.7 & $+0.31$\phantom{*} & $-0.84$ & .4029 & 1.0000 & $+10.1$ & $+0.82$ \\
10 & Item 5.07 shareholder vote & 1{,}424 & 405 & 5.7 & $+1.05$\phantom{*} & $-0.10$ & .9215 & 1.0000 & $+9.4$ & $+1.23$ \\
10 & Other narrative items & 2{,}326 & 882 & 9.3 & $+0.14$\phantom{*} & $-0.14$ & .8885 & 1.0000 & $+2.4$ & $-0.06$ \\
10 & \emph{All non-2.02 items (pooled)} & 16{,}670 & 991 & 66.7 & $+0.39$\phantom{*} & $-1.21$ & .2252 & --- & $+44.4$ & $+0.46$ \\
\addlinespace
20 & \emph{All 8-K filings (pooled)} & 24{,}732 & 981 & 100.0 & $+0.39$\phantom{*} & $-1.34$ & .1812 & --- & $+100.0$ & $+0.70$ \\
20 & Item 2.02 results & 8{,}305 & 811 & 33.6 & $-0.35$\phantom{*} & $+1.00$ & .3180 & 1.0000 & $-25.6$ & $+0.38$ \\
20 & Item 5.02 leadership & 4{,}610 & 963 & 18.6 & $+0.33$\phantom{*} & $-0.74$ & .4581 & 1.0000 & $+17.7$ & $-0.06$ \\
20 & Item 7.01 Reg-FD & 3{,}269 & 917 & 13.2 & $+0.69$\phantom{*} & $-2.02$ & .0436 & .6984 & $+25.3$ & $+1.28$ \\
20 & Item 8.01 other events & 4{,}842 & 974 & 19.6 & $+0.51$\phantom{*} & $-1.61$ & .1087 & 1.0000 & $+26.8$ & $+0.53$ \\
20 & Item 5.07 shareholder vote & 1{,}415 & 403 & 5.7 & $+2.99$\phantom{*} & $-2.00$ & .0464 & .6984 & $+42.1$ & $+4.94$ \\
20 & Other narrative items & 2{,}291 & 872 & 9.3 & $+0.49$\phantom{*} & $-0.52$ & .6061 & 1.0000 & $+13.7$ & $+0.42$ \\
20 & \emph{All non-2.02 items (pooled)} & 16{,}427 & 981 & 66.4 & $+0.68$\phantom{*} & $-2.47$ & .0137 & --- & $+125.6$ & $+0.83$ \\
\bottomrule
\end{tabular}
\end{table}

Figures~\ref{fig:app-item-geography}, \ref{fig:app-item-geography-b} and~\ref{fig:app-item-geography-c} draw the
decomposition, including the earnings-free residual as a panel of its own.

Figures~\ref{fig:app-firm-identity-rung} and~\ref{fig:app-firm-identity-rung-b} draw
the identity rung itself, and what a zero-text firm mean achieves without any text
at all, on both of the row bases that quantity is reported on.

\clearpage

\begin{figure}[tbp]
\centering
\includegraphics[width=\textwidth,height=0.83\textheight,keepaspectratio]{figures/F6_firm_identity_rung__a.pdf}
% caption numbers src: results/tables/firm_identity_ensemble.csv (panels a and b: all 69 cells, rel_impr_pct_firm, adds_holm, hurts_holm, firm coverage, and the merged-grid zero-text firm-mean comparisons); results/tables/firm_identity_control_zerotext.csv (panel c: the same six comparisons on every reference row)
\caption[The firm-identity rung, and what a zero-text firm mean achieves alone]{
The
firm-identity rung, and what a zero-text firm mean achieves on its own.
Panel~(a) plots all 69 cells against the firm-identity-augmented reference,
sorted by the relative QLIKE improvement text delivers over it.
That reference adds one zero-text scalar per firm to the recalibrated
forecast (the firm's mean realised volatility over its validation rows), so
signal that merely re-identifies the filer is absorbed. Eight cells add at Holm, 35 hurt and 26 carry no verdict. The channel means are $-2.25$ per cent long-form and $-0.27$ per
cent event-driven; 38 of the 45 long-form cells turn negative. % src: results/tables/firm_identity_ensemble.csv (adds_holm 8, hurts_holm 35, remainder 26; channel means of rel_impr_pct_firm; 38 of 45 long-form cells below zero)
Basis: seed-ensemble forecasts, volatility-unit QLIKE, day-clustered
inference with Holm inside the 69-cell family, which governs this panel alone.
}
\label{fig:app-firm-identity-rung}
\end{figure}

\begin{figure}[tbp]
\centering
\includegraphics[width=\textwidth,height=0.83\textheight,keepaspectratio]{figures/F6_firm_identity_rung__b.pdf}
% caption numbers src: results/tables/firm_identity_ensemble.csv (panels a and b: all 69 cells, rel_impr_pct_firm, adds_holm, hurts_holm, firm coverage, and the merged-grid zero-text firm-mean comparisons); results/tables/firm_identity_control_zerotext.csv (panel c: the same six comparisons on every reference row)
\caption[The firm-identity rung, and what a zero-text firm mean achieves alone (b, c)]{Continued from Figure~\ref{fig:app-firm-identity-rung}. The
firm-identity rung, and what a zero-text firm mean achieves on its own.
Panels~(b) and~(c) remove text entirely and ask whether the
firm-mean term alone improves on the plain recalibrated HAR, on the merged grid
rows and on every reference row respectively; they sit outside the fifteen
pre-declared families and carry no Holm correction, every printed $p$ being
unadjusted, and Holm over their six comparisons leaves three standing. Basis:
seed-ensemble forecasts, volatility-unit QLIKE, day-clustered inference; the
69-cell Holm family governs panel~(a) of
Figure~\ref{fig:app-firm-identity-rung} alone, not these two panels. % src: results/tables/firm_identity_ensemble.csv (p_firmMeanOnly_vs_fR over the six comparisons: 0.000050, 0.000493, 0.001787, 0.024108, 0.082580, 0.863043; step-down Holm over those six puts long-form h=5 at 0.072); results/tables/firm_identity_control_zerotext.csv (p_holm: 0.462281, 0.000054, 0.010419, 0.009461, 0.183015, 0.937142, three below .05)
}
\label{fig:app-firm-identity-rung-b}
\end{figure}

The identity rung of Figures~\ref{fig:app-firm-identity-rung} and~\ref{fig:app-firm-identity-rung-b} is the visual form of a finding
Chapter~\ref{ch:results} states as a count, and the mass of the distribution
in panel~(a) is what a count cannot convey. Once a reference already knows the
firm's typical volatility level, adding long-form text makes the forecast
\emph{worse} more often than better, and the eight survivors are drawn at full
strength beside thirty-five losers rather than netted against them. The report
does not read the thirty-five as evidence that text is harmful. A negative
increment against a reference handed a free firm-level term is what one
expects when the text forecast was carrying that same firm level and is now
being asked to add to it.

Panels~(b) and~(c) exist to keep one number honest. The unit there is the
channel-by-horizon comparison rather than the cell, because a zero-text
forecast takes exactly one distinct value per channel and horizon. Six is
the honest denominator; the 53-cell expansion of the same flag is
deliberately kept off every axis. The firm mean beats the recalibrated HAR in
four of six comparisons on the merged grid rows on which the rung is scored,
and in three of six when the identical test is run on every reference row. The
count moves with the row basis and both counts are printed, the comparison
that moves being long-form $h=5$, which passes at $p=0.024$ on the merged grid
rows and fails at $p=0.231$ on every reference row. % src: results/tables/firm_identity_ensemble.csv (firm_beats_fR true in 4 of 6 merged-grid comparisons; long_form h=5 p_firmMeanOnly_vs_fR = 0.0241, firm_beats_fR True); results/tables/firm_identity_control_zerotext.csv (3 of 6 on every reference row; long_form h=5 p_clustered = 0.2311, no longer a win; event_driven h=10 is negative on both bases, p = 0.0826 and p = 0.0610)
One inversion is shown rather than smoothed. At long-form $h=20$ the
observation-weighted improvement is $-1.39$ per cent against a day-clustered
statistic of $-3.13$, so the comparison counts as a win under the primary
inference while the observation-weighted number sits beside it. % src: results/tables/firm_identity_ensemble.csv (long_form h=20: rel_impr_firmMeanOnly_vs_fR = -1.391, dm_firmMeanOnly_vs_fR = -3.135, p = 0.0018)
The two averages disagree because one weights filings and the other the 765
trading days carrying the comparison, and the day-weighted test is primary.

The figure does not decompose the identity term into identity strings versus
firm-stable content. That separation is the work of the anonymisation and
matched-swap instruments of Chapter~\ref{ch:results}, and nothing here
distinguishes a model that recognises a company name from one that has learned
a company's typical behaviour. It reports no minimum detectable effect, so a
cell with no verdict is not evidence that text is worth nothing there. The
Holm family is the 69-cell grid, so these verdicts are not comparable with the
smaller residual family used for the prompted-arm increment. Panel~(a)
fixes one specification of the firm statistic, the validation-window mean, the
other three specifications of the registered battery being reported in
Appendix~\ref{app:fullresults}.

\begin{figure}[tbp]
\centering
\includegraphics[width=\textwidth,height=0.83\textheight,keepaspectratio]{figures/AF1_item_geography_8k__a.pdf}
% caption numbers src: results/tables/row11_item_stratified.csv (all 48 rows: rel_firm_pct, dm_firm, p_firm, p_firm_holm, share_of_pooled_residual_pct, n_test); cross-checked against results/tables/itemcode_control.md and results/tables/crossfamily_llama70.csv
\caption[The item geography of the surviving 8-K residual]{
The item geography
of the surviving 8-K residual. 
Panel~(a) decomposes the pooled event-driven
increment over the firm-identity-augmented reference into six disjoint item
groups, for the prompted Qwen3-32B arm and for its Llama-3.1-70B replication,
at each horizon. A filled marker denotes a cell surviving Holm within its
family's eighteen group-by-horizon tests. Four cells survive for Qwen3-32B
(Item 2.02 at $h=5$ and $h=20$, Item 7.01 at $h=10$ and $h=20$) and one for
the 70B. % src: results/tables/row11_item_stratified.csv (p_firm_holm < .05 with dm_firm < 0 among each family's 18 partition cells)
}
\label{fig:app-item-geography}
\end{figure}

\begin{figure}[tbp]
\centering
\includegraphics[width=\textwidth,height=0.83\textheight,keepaspectratio]{figures/AF1_item_geography_8k__b.pdf}
% caption numbers src: results/tables/row11_item_stratified.csv (all 48 rows: rel_firm_pct, dm_firm, p_firm, p_firm_holm, share_of_pooled_residual_pct, n_test); cross-checked against results/tables/itemcode_control.md and results/tables/crossfamily_llama70.csv
\caption[The item geography of the surviving 8-K residual (b)]{Continued from Figure~\ref{fig:app-item-geography}. The item geography
of the surviving 8-K residual. 
Panel~(b) sets Item 2.02's 33.4 per cent share of test rows against its share
of the horizon-pooled QLIKE reduction, 53.6 per cent for Qwen3-32B and 69.3
for the 70B, the per-horizon shares beneath being far from stable. % src: results/tables/row11_item_stratified.csv (share_of_pooled_residual_pct by horizon; abs_reduction_firm pooled; n_test shares)
}
\label{fig:app-item-geography-b}
\end{figure}

\begin{figure}[tbp]
\centering
\includegraphics[width=\textwidth,height=0.83\textheight,keepaspectratio]{figures/AF1_item_geography_8k__c.pdf}
% caption numbers src: results/tables/row11_item_stratified.csv (all 48 rows: rel_firm_pct, dm_firm, p_firm, p_firm_holm, share_of_pooled_residual_pct, n_test); cross-checked against results/tables/itemcode_control.md and results/tables/crossfamily_llama70.csv
\caption[The item geography of the surviving 8-K residual (c)]{Continued from Figure~\ref{fig:app-item-geography}. The item geography
of the surviving 8-K residual. 
Panel~(c) pools the five non-2.02 groups into an earnings-free residual,
significant at all three horizons for Qwen3-32B ($+0.22/+0.22/+0.17$ per cent,
day-clustered DM $-2.96/-4.54/-3.61$) and at one of three for the 70B. Basis:
the event-driven panel only, on the full $25{,}109/25{,}001/24{,}732$ test
rows over $996/991/981$ trading days; volatility-unit QLIKE; Holm inside each
18-cell family, the pooled rows carrying raw $p$ only; the grouping concedes
earnings first, so no non-2.02 group contains a results release.
}
\label{fig:app-item-geography-c}
\end{figure}

Figures~\ref{fig:app-item-geography}, \ref{fig:app-item-geography-b} and~\ref{fig:app-item-geography-c} carry the decomposition
Chapter~\ref{ch:results} does not. That chapter controls for earnings
windows by adding an Item 2.02 \emph{indicator} to the reference and reports
that the residual survives it. How the residual is distributed across item
types is a different question, and it has a two-sided answer. Earnings
releases are a third of the rows and a little over half of the pooled QLIKE
reduction for the anchor arm, so they are over-represented but do not account
for the residual. The earnings-free pool of the other five groups is
positive and significant at every horizon for that arm. The split differs by
family, however, and by a lot: for the replication the earnings group carries
69.3 per cent of the pooled reduction ($+99.3$, $+55.6$ and $-25.6$ per
cent by horizon). Its earnings-free pool is significant at one horizon
of three. % src: results/tables/row11_item_stratified.csv (llama70_awq 2.02_earnings share_of_pooled_residual_pct by horizon; narrative_ALL p_firm 0.846 / 0.225 / 0.014)
The honest reading is that the residual is partly event reading and partly the
earnings number, and that the balance between the two is not a stable property
of the channel.

Four things the figure does not show bound that reading. It does not show
survival at the maximal-pool rung or at the full conjunction, which is empty
over the whole 69-cell grid, so a filled marker here is survival at one rung
and not a surviving cell of the study. Its percentages sit on the full
event-driven panel and are therefore not interchangeable with the ladder
counts computed on $23{,}855/22{,}785/22{,}318$ merged-grid rows. The 70B rows
are drawn on the single-pass basis, whose pooled triple is
$+0.83/+0.64/+0.39$ per cent, rather than on the three-seed ensemble the
report uses for that arm elsewhere. % src: results/tables/row11_item_stratified.csv (llama70_awq ALL rows); results/tables/crossfamily_llama70_ens.csv (the ensemble triple +0.84/+0.64/+0.38 used in Chapter 4)
And the item groups partition the test residual with nothing refitted inside a
group, so a group's number is a property of that subset of test rows and not
of a model trained on it. The earnings-conceding assignment rule (a filing
takes the first code present in a fixed order beginning with 2.02) makes
the non-2.02 groups a conservative construction. The earnings-free residual
is a lower bound on what non-earnings disclosure carries and not an estimate
of it.

\FloatBarrier
\section{Contamination, Elicitation and Input Parity}
\label{app:res-probes}

Three threats attach to a finding read from a prompted language model, and
each received an executed control: the model may have memorised the
outcomes, the increment may be an artefact of one prompt, and the curated
excerpts rather than the prompting may be doing the work.
Section~\ref{app:res-cutoff} tabulates the memorisation control's temporal
stratification. This section tabulates the other three instruments the
report cites without showing: the zero-content prompt arms, the same-model
identity reference, the elicitation sweep and the input-parity arm.

\subsection{Zero-content prompts and the same-model identity reference}

Table~\ref{tab:app-contam-arms} runs the full combination protocol on three
prompts that differ only in what they contain. A date-only prompt carries
no increment anywhere (zero of six cells Holm-significantly positive,
and significantly \emph{negative} at two long-form horizons). Nothing
in the arm's gain comes from prompting as such or from era information in
the date. A date-plus-ticker prompt, still carrying no filing text, is a
different matter: it reproduces between 31 and 77 per cent of the
full-text increment in the three cells where that increment is large enough
for the ratio to be well identified --- an increment above one per cent and
Holm-significant, which the event-driven ten-day cell misses on the underlying
value ($0.998$) while printing as $+1.00$.
% src: results/tables/llm_contamination.csv (variant fulltext rel_pct 1.7942, 2.2516, 0.2748 long-form and 1.2141, 0.9983, 0.6582 event-driven, with p_holm)
At long-form $h{=}20$ the identity
prompt alone ($+1.11$ per cent) \emph{exceeds} the full-text increment
($+0.27$ per cent), a ratio of 405 per cent that is disclosed here rather
than trimmed. The uncontrolled prompted-arm percentage is therefore not a
text effect and must never be quoted as one.

Table~\ref{tab:app-contam-joint} is the control that follows from that. The
reference is augmented with the same model's own date-and-ticker forecast,
so everything the instrument can produce from the firm's name is already
inside what text must beat. Full text still adds in six cells of six after
Holm correction, from $+0.17$ to $+1.95$ per cent. Both statements are
true at once, and the pair is the honest form of the prompted-arm result: the
name alone reproduces a large share of the raw percentage, and what survives
the same-model identity control is still 63 to 96 per cent of it --- small in
absolute terms, positive and repeatedly significant. The two shares do not
partition the increment and are not meant to; they are two instruments on the
same cell, and the section above says they sum to more than the whole. % src: results/tables/llm_contamination.csv (joint beyond_identity rel_pct over the matching variant fulltext rel_pct: 83.5, 86.5, 62.6 long-form and 63.1, 68.4, 95.7 event-driven at h = 5/10/20)

Table~\ref{tab:app-70b-probe} repeats the same construction inside the
replication family, which the report reports without a probe of its own.
The zero-content date-and-ticker prompt run through the 70B recovers 5.2
and 7.3 per cent of that model's full-text edge over the recalibrated HAR
at $h{=}5$ and $h{=}10$, and 0.10 and 0.49 per cent of its edge over the
firm-identity reference. At $h{=}20$ the denominator is too small for the
ratio to be identified and the source file flags it as such, which is why
the 31 per cent figure in that row is printed with the flag rather than
quoted. With the probe inside the reference, the 70B's full text still adds
at two horizons of three after Holm. Content, not the name: for the
replication family as well as for the anchor.

\begin{table}[htbp]
\centering
\caption[Zero-content prompt arms against the recalibrated reference]{Three
prompts through the identical combination protocol: the committed full-text
arm, a prompt carrying the filing date and the firm's ticker but no
document text, and a prompt carrying only the date. Volatility-unit QLIKE,
combiner weights fitted on validation and frozen on test, day-clustered DM,
Holm within this 18-cell block. \emph{Positive rel.\% means the arm lowers
loss.} ``repro.'' is the arm's increment as a fraction of the full-text
increment in the same cell (a ratio, not a percentage), and is
meaningful only where the full-text increment is itself large and
significant. The $+4.05$ at long-form $h{=}20$ is such a case failing, the
identity prompt there being larger than the full-text arm it is divided by,
and it is shown rather than suppressed. Single-pass arms, seed 2026.}
\label{tab:app-contam-arms}
\singlespacing\small
\setlength{\tabcolsep}{4pt}
% src: results/tables/llm_contamination.csv (block = variant; all 18 rows)
\begin{tabular}{@{}llrrrrrr@{}}
\toprule
Channel & $h$ & Prompt contents & QLIKE $f_U$ & rel.\% & DM & Holm $p$ & repro. \\
\midrule
Long-form & 5 & Full text & $0.1187$ & $+1.79$ & $-6.31$ & $<$.0001 & --- \\
Long-form & 5 & Date $+$ ticker, no text & $0.1202$ & $+0.55$ & $-4.52$ & $<$.0001 & +0.31 \\
Long-form & 5 & Date only, no text & $0.1209$ & $-0.00$ & $-0.93$ & 1.0000 & -0.00 \\
\addlinespace
Long-form & 10 & Full text & $0.0853$ & $+2.25$ & $-7.92$ & $<$.0001 & --- \\
Long-form & 10 & Date $+$ ticker, no text & $0.0866$ & $+0.83$ & $-5.72$ & $<$.0001 & +0.37 \\
Long-form & 10 & Date only, no text & $0.0877$ & $-0.49$ & $+3.46$ & .0056 & -0.22 \\
\addlinespace
Long-form & 20 & Full text & $0.0699$ & $+0.27$ & $-3.23$ & .0115 & --- \\
Long-form & 20 & Date $+$ ticker, no text & $0.0693$ & $+1.11$ & $-5.82$ & $<$.0001 & +4.05 \\
Long-form & 20 & Date only, no text & $0.0712$ & $-1.65$ & $+4.34$ & .0002 & -5.99 \\
\addlinespace
Event-driven & 5 & Full text & $0.1250$ & $+1.21$ & $-5.04$ & $<$.0001 & --- \\
Event-driven & 5 & Date $+$ ticker, no text & $0.1253$ & $+0.94$ & $-2.78$ & .0439 & +0.77 \\
Event-driven & 5 & Date only, no text & $0.1265$ & $-0.00$ & $+0.95$ & 1.0000 & -0.00 \\
\addlinespace
Event-driven & 10 & Full text & $0.0874$ & $+1.00$ & $-3.76$ & .0020 & --- \\
Event-driven & 10 & Date $+$ ticker, no text & $0.0876$ & $+0.71$ & $-1.24$ & 1.0000 & +0.72 \\
Event-driven & 10 & Date only, no text & $0.0886$ & $-0.39$ & $+1.64$ & .6025 & -0.39 \\
\addlinespace
Event-driven & 20 & Full text & $0.0641$ & $+0.66$ & $-1.98$ & .3340 & --- \\
Event-driven & 20 & Date $+$ ticker, no text & $0.0645$ & $+0.01$ & $+0.12$ & 1.0000 & +0.01 \\
Event-driven & 20 & Date only, no text & $0.0645$ & $+0.00$ & $-0.69$ & 1.0000 & +0.01 \\
\bottomrule
\end{tabular}
\end{table}

\begin{table}[htbp]
\centering
\caption[Full text beyond the model's own identity forecast]{The same-model
identity control. The reference is
$\exp[a + b\log f_{\mathrm{HAR}} + c\log f_{\mathrm{date+ticker}}]$, which
already contains everything the same instrument produces from the firm's
name and the date. The test asks what the filing text adds beyond
it. Volatility-unit QLIKE, day-clustered DM, Holm within this six-cell
block, weights validation-fitted and test-frozen. $g$ is the fitted
coefficient on the text forecast. All six cells are Holm-significant
improvements.}
\label{tab:app-contam-joint}
\singlespacing\small
\setlength{\tabcolsep}{4pt}
% src: results/tables/llm_contamination.csv (block = joint; all 6 rows)
\begin{tabular}{@{}llrrrrrr@{}}
\toprule
Channel & $h$ & QLIKE $f_{R'}$ & QLIKE $f_{U'}$ & rel.\% & $g$ & DM & Holm $p$ \\
\midrule
Long-form & 5 & $0.1202$ & $0.1184$ & $+1.50$ & $+0.241$ & $-5.83$ & $<$.0001 \\
Long-form & 10 & $0.0866$ & $0.0849$ & $+1.95$ & $+0.321$ & $-7.13$ & $<$.0001 \\
Long-form & 20 & $0.0693$ & $0.0692$ & $+0.17$ & $+0.070$ & $-2.51$ & .0247 \\
Event-driven & 5 & $0.1253$ & $0.1244$ & $+0.77$ & $+0.233$ & $-3.79$ & .0006 \\
Event-driven & 10 & $0.0876$ & $0.0870$ & $+0.68$ & $+0.245$ & $-2.85$ & .0132 \\
Event-driven & 20 & $0.0645$ & $0.0641$ & $+0.63$ & $+0.214$ & $-2.01$ & .0446 \\
\bottomrule
\end{tabular}
\end{table}

\begin{table}[htbp]
\centering
\caption[The replication family's own zero-content probe]{The zero-content
identity probe inside the \textbf{cross-family replication}
(Llama-3.1-70B). Panel (a) runs the date-and-ticker prompt through the
combination protocol against both references. Panel (b) expresses that
probe's increment as a share of the same family's full-text increment
against each reference. The $h{=}20$ share is flagged in the source as
having a denominator that is not well identified (the full-text edge
there is small), and is printed with the flag rather than quoted as a
result. Panel (c) puts the probe inside the reference and asks what the
70B's full text still adds. Volatility-unit QLIKE, day-clustered DM, Holm
within each panel's own family; event-driven channel throughout, the only
channel this family ran.}
\label{tab:app-70b-probe}
\singlespacing\small
\setlength{\tabcolsep}{4pt}
% src: results/tables/crossfamily_llama70_probe.csv (blocks probe_m1, probe_share and beyond_identity)

\emph{(a) The zero-content date-and-ticker prompt, as an arm}\\[0.4ex]
\begin{tabular}{@{}lrrrrrr@{}}
\toprule
$h$ & vs.\ HAR \% & DM & Holm $p$ & vs.\ identity \% & DM & Holm $p$ \\
\midrule
 5 & $+0.073$ & $-1.62$ & .4665 & $+0.0008$ & $+0.42$ & 1.0000 \\
10 & $+0.085$ & $-1.55$ & .4665 & $+0.0031$ & $+0.36$ & 1.0000 \\
20 & $+0.218$ & $-3.54$ & .0025 & $+0.0218$ & $-1.68$ & .4665 \\
\bottomrule
\end{tabular}

\vspace{1.2ex}
\emph{(b) How much of the family's own full-text edge the probe recovers}\\[0.4ex]
\begin{tabular}{@{}lrrrrrl@{}}
\toprule
$h$ & probe \% & full text \% & share & probe \% & full text \% & share \\
 & \multicolumn{3}{c}{vs.\ recalibrated HAR} & \multicolumn{3}{c}{vs.\ firm identity} \\
\midrule
 5 & $+0.073$ & $+1.395$ & 5.2\% & $+0.0008$ & $+0.835$ & 0.10\% \\
10 & $+0.085$ & $+1.163$ & 7.3\% & $+0.0031$ & $+0.637$ & 0.49\% \\
20 & $+0.218$ & $+0.692$ & \emph{31.5\%, denominator not identified} & $+0.0218$ & $+0.383$ & \emph{5.7\%, ditto} \\
\bottomrule
\end{tabular}

\vspace{1.2ex}
\emph{(c) Full text with the probe inside the reference}\\[0.4ex]
\begin{tabular}{@{}lrrrrrl@{}}
\toprule
$h$ & QLIKE $f_{R'}$ & QLIKE $f_{U'}$ & rel.\% & DM & Holm $p$ & Verdict \\
\midrule
 5 & 0.1264 & 0.1247 & $+1.374$ & $-3.17$ & .0046 & adds \\
10 & 0.0882 & 0.0872 & $+1.136$ & $-2.27$ & .0465 & adds \\
20 & 0.0644 & 0.0640 & $+0.653$ & $-0.75$ & .4512 & --- \\
\bottomrule
\end{tabular}
\end{table}

\subsection{Elicitation sensitivity}

Table~\ref{tab:app-elicit} answers the objection that a near-null is
simply a badly written prompt, by running five elicitations of the same
model on the same deterministic 4{,}000-filing subsample: the committed
template decoded twice, two semantically equivalent paraphrases, and the
model's thinking mode. Its support is the subsample and not the test
panels, so its day counts are 286--290 on long-form and 787--800 on
event-driven and none of its percentages is comparable with those of
Table~\ref{tab:app-grid-ed}. What is comparable is the pattern across arms,
because every arm sees the identical rows.

Three findings follow. Decoding is near-deterministic but not bit-exact:
94 to 97 per cent of forecasts are identical across repeats, rank agreement
above 0.92. A repeat decode is therefore a weak perturbation and the two base
rows are expected to agree, as they do. Rewriting the prompt is a
\emph{strong} perturbation: individual forecasts correlate with the base
template at Spearman 0.39 to 0.58, meaning a paraphrase changes what the
model says about a particular filing far more than it changes the base
template's own noise. And the two channels respond to that perturbation
completely differently. On the event-driven channel every one of the five
arms returns a positive increment at $h{=}5$, from $+0.63$ to $+1.13$ per
cent, all five raw-significant. On the long-form channel the increment
flips sign between the two paraphrases ($+2.32$ against $-0.44$ at
$h{=}5$, $+5.95$ against $-2.16$ at $h{=}10$), and thinking mode does not
rescue it. The event-driven residual is elicitation-robust; the long-form
apparent gain is prompt-fragile, which is consistent with no long-form
prompted cell clearing the firm-identity and maximal-pool rungs together under
Holm in Table~\ref{tab:app-grid-lf}, each horizon clearing one and failing the
other. % src: results/tables/control_intersection_ensemble.csv (long_form C6_llmtext: maximal_holm True/True/False and firm_holm False/False/True at h=5/10/20; AND_maximal_firm_holm False in all three)

\begin{table}[htbp]
\centering
\caption[Elicitation sensitivity of the prompted arm]{Five elicitations of
the same model on the same rows. Panel (a): repeat-decode determinism under
an identical configuration. Panel (b): Spearman rank agreement of each
arm's individual forecasts with the committed base template at $h{=}10$.
Panel (c): each arm's combination increment over the recalibrated HAR
reference, volatility-unit QLIKE, weights validation-fitted and
test-frozen, day-clustered DM. \textbf{Support: the registered
4{,}000-filing stratified subsample} --- 286 to 290 long-form and 787 to 800
event-driven trading days --- not the full test panels, which carry 809/803/794
and 996/991/981. No percentage
here is comparable with the grid tables, and the comparison the panel
supports is arm against arm on identical rows. \textbf{The committed source
carries raw $p$ only for panel (c); no multiplicity correction is applied
and an asterisk marks raw significance at 5\%.}}
\label{tab:app-elicit}
\singlespacing\footnotesize
\setlength{\tabcolsep}{3.5pt}
% src: results/tables/elicitation_sensitivity.csv (sections repeat, agreement and m1)

\emph{(a) Repeat-decode determinism, $n = 4{,}000$}\\[0.4ex]
\begin{tabular}{@{}lrrr@{}}
\toprule
Horizon & Exactly equal & Spearman & Mean $|$rel.\ diff.$|$ \\
\midrule
$h=5$ & 0.970 & 0.9639 & 0.59\% \\
$h=10$ & 0.945 & 0.9538 & 0.68\% \\
$h=20$ & 0.936 & 0.9244 & 1.08\% \\
\bottomrule
\end{tabular}

\vspace{1.2ex}
\emph{(b) Forecast agreement with the base template, $h = 10$, $n = 4{,}000$}\\[0.4ex]
\begin{tabular}{@{}lr@{}}
\toprule
Arm & Spearman vs.\ base \\
\midrule
Paraphrase 1 & 0.578 \\
Paraphrase 2 & 0.413 \\
Thinking mode & 0.391 \\
\bottomrule
\end{tabular}

\vspace{1.2ex}
\emph{(c) Combination increment over the recalibrated HAR, per arm}\\[0.4ex]
\begin{tabular}{@{}llrlrlrl@{}}
\toprule
& & \multicolumn{2}{c}{$h{=}5$} & \multicolumn{2}{c}{$h{=}10$} & \multicolumn{2}{c}{$h{=}20$} \\
\cmidrule(lr){3-4}\cmidrule(lr){5-6}\cmidrule(lr){7-8}
Channel & Arm & rel.\% & raw $p$ & rel.\% & raw $p$ & rel.\% & raw $p$ \\
\midrule
Event-driven & Base template, decode 1 \emph{(the committed arm)} & $+0.82$* & .0019 & $+0.45$\phantom{*} & .1627 & $+0.10$\phantom{*} & .6660 \\
Event-driven & Base template, decode 2 & $+0.89$* & .0015 & $+0.62$* & .0205 & $+0.17$\phantom{*} & .9959 \\
Event-driven & Paraphrase 1 & $+1.13$* & $<$.0001 & $+1.58$* & .0001 & $+1.58$* & .0018 \\
Event-driven & Paraphrase 2 & $+0.81$* & .0034 & $+0.44$\phantom{*} & .9752 & $-0.16$\phantom{*} & .3517 \\
Event-driven & Thinking mode & $+0.63$* & .0010 & $+0.39$\phantom{*} & .0642 & $-0.59$\phantom{*} & .5877 \\
\addlinespace
Long-form & Base template, decode 1 \emph{(the committed arm)} & $+0.20$\phantom{*} & .0667 & $+4.20$* & .0142 & $+0.44$\phantom{*} & .1684 \\
Long-form & Base template, decode 2 & $+0.21$* & .0352 & $+4.21$* & .0166 & $+0.58$\phantom{*} & .1864 \\
Long-form & Paraphrase 1 & $+2.32$* & .0064 & $+5.95$\phantom{*} & .0758 & $+3.14$* & .0018 \\
Long-form & Paraphrase 2 & $-0.44$\phantom{*} & .0704 & $-2.16$\phantom{*} & .0585 & $-3.77$* & $<$.0001 \\
Long-form & Thinking mode & $-0.15$\phantom{*} & .8498 & $+1.55$\phantom{*} & .2466 & $+1.32$\phantom{*} & .1408 \\
\bottomrule
\end{tabular}
\end{table}

\subsection{Input parity: prompting against pooling}

Section~\ref{sec:design-models} registers C5x as a designed control: the
same lineage's 8B embedding model applied to the \emph{byte-identical}
excerpts the prompted arm reads, so that any gap between them is
attributable to the interface rather than to which text was selected.
Table~\ref{tab:app-c5x} is that control's result. Prompting adds
$+0.27$ to $+2.25$ per cent and is genuine at all three long-form
horizons. The same lineage pooling representations of the same bytes loses
at every horizon, $-4.12$ to $-0.32$ per cent, and is genuine at none. The
standard-input embedding arm, on its natural full-document input, loses
further still at $h{=}5$ and $h{=}20$ and marginally less at $h{=}10$ ($-3.10$
on the primary ensemble and $-3.13$ on seed 2026, against $-3.74$).
% src: results/tables/c5x_input_parity.csv (rel_impr_pct: C5_qwen3 ens -1.024/-3.101/-6.389, s26 -1.035/-3.135/-6.647; C5x_qwen3exc -0.325/-3.740/-4.121)
The dissociation is between eliciting a forecast and
pooling a representation, not between curated and uncurated text.

Two limits belong with that reading and are printed on the table. The two
arms share a lineage, not a parameter count: a 32B decoder against an 8B
embedder with a ridge head. The comparison therefore isolates the interface and
not scale, and the claim must be written ``same-lineage'', never
``same-size''. And the control ran on the long-form channel only, so it
licenses nothing about the event-driven residual. The diagnostic columns
are worth reading beside the verdicts. The prompted arm's forecasts take
only 8 to 16 distinct values after rounding to two decimals against the
embedding arms' 15 to 38, and its standalone QLIKE is the worst of the
four. The arm that adds the most in combination is also the coarsest and
the least accurate on its own, which is what an incremental-value
instrument is designed to be able to detect.

\begin{table}[htbp]
\centering
\caption[Input parity: prompting against same-lineage pooling]{The
input-parity control on the long-form channel. C5x applies
Qwen3-Embedding-8B to the byte-identical excerpts the prompted C6 arm reads
(the risk and MD\&A items, else head truncation, with a ridge head on the
log target and Duan smearing \cite{duan1983smearing}); C5 is the same embedding model on its
standard full-document input, shown on both the primary ensemble basis and
seed 2026; C6 is the prompted 32B decoder. Combiner weights are
validation-fitted and test-frozen throughout, QLIKE is in volatility units
for the increment columns, DM is day-clustered and Holm runs within this
one 12-cell family. \emph{Genuine} requires DM $<0$, Holm $<.05$ and a
label-shuffle placebo inside $|\mathrm{DM}|<2$. ``uniq.'' counts the
distinct forecast values after rounding to two decimals, a coarseness
diagnostic. \textbf{The comparison is at parity of input, not of parameter
count} (a 32B decoder against an 8B embedder with a linear head), and
it ran on the long-form channel only.}
\label{tab:app-c5x}
\singlespacing\footnotesize
\setlength{\tabcolsep}{3.5pt}
% src: results/tables/c5x_input_parity.csv (all 12 rows and all listed columns)
\begin{tabular}{@{}lrrrrrrrrc@{}}
\toprule
Arm & $h$ & QLIKE alone & $R^2$ & uniq. & rel.\% & DM & Holm $p$ & plac. & Gen. \\
\midrule
C5 Qwen3-Emb., ensemble \emph{(primary)} & 5 & $0.1540$ & $-0.013$ & 15 & $-1.02$ & $+2.64$ & .0160 & $+0.47$ & -- \\
C5 Qwen3-Emb., ensemble \emph{(primary)} & 10 & $0.1207$ & $-0.007$ & 15 & $-3.10$ & $+3.62$ & .0012 & $+0.45$ & -- \\
C5 Qwen3-Emb., ensemble \emph{(primary)} & 20 & $0.0937$ & $-0.004$ & 15 & $-6.39$ & $+5.38$ & $<$.0001 & $+0.42$ & -- \\
\addlinespace
\quad the same arm, seed 2026 & 5 & $0.1548$ & $-0.021$ & 16 & $-1.03$ & $+2.66$ & .0160 & $+0.47$ & -- \\
\quad the same arm, seed 2026 & 10 & $0.1241$ & $-0.032$ & 17 & $-3.13$ & $+3.70$ & .0011 & $+0.45$ & -- \\
\quad the same arm, seed 2026 & 20 & $0.0943$ & $-0.011$ & 16 & $-6.65$ & $+5.45$ & $<$.0001 & $+0.42$ & -- \\
\addlinespace
C5x Qwen3-Emb.\ on C6's excerpts & 5 & $0.1584$ & $-0.007$ & 32 & $-0.32$ & $+5.66$ & $<$.0001 & $+0.32$ & -- \\
C5x Qwen3-Emb.\ on C6's excerpts & 10 & $0.1288$ & $-0.020$ & 29 & $-3.74$ & $+8.60$ & $<$.0001 & $-0.58$ & -- \\
C5x Qwen3-Emb.\ on C6's excerpts & 20 & $0.0972$ & $+0.036$ & 38 & $-4.12$ & $+9.31$ & $<$.0001 & $+1.58$ & -- \\
\addlinespace
C6 prompted Qwen3-32B & 5 & $0.2072$ & $-0.174$ & 8 & $+1.79$ & $-6.31$ & $<$.0001 & $+0.83$ & \textbf{Y} \\
C6 prompted Qwen3-32B & 10 & $0.1778$ & $-0.201$ & 13 & $+2.25$ & $-7.92$ & $<$.0001 & $-1.45$ & \textbf{Y} \\
C6 prompted Qwen3-32B & 20 & $0.1601$ & $-0.345$ & 16 & $+0.27$ & $-3.23$ & .0038 & $-1.44$ & \textbf{Y} \\
\bottomrule
\end{tabular}
\end{table}

Figures~\ref{fig:app-elicitation}, \ref{fig:app-elicitation-b} and~\ref{fig:app-elicitation-c} carry the three
input manipulations of this section one to a part (input parity, the zero-content
prompts, and the same construction inside the replication family) with their
three separate multiplicity families kept apart.

\subsection{What the surviving instrument emits}

The coarseness diagnostic of Table~\ref{tab:app-c5x} counts distinct forecast
values on the long-form channel. Table~\ref{tab:app-instrument-readout} reads
the same instrument on the cell that carries the surviving residual, and
records what its near-binary output separates. The arm emits ten distinct
values on that cell, of which 0.22 (49.3 per cent of forecasts) and 0.18 (43.2
per cent) carry 92.6 per cent between them. % src: results/tables/instrument_readout.md
Splitting the test rows on that flag, mean realised volatility is higher in the
0.22 group in all ten deciles of the HAR-RV forecast. The difference is not
strictly increasing across the deciles: the generator prints the monotonicity
flag as false, and rows 3 to 4 are the visible reversal. % src: results/tables/instrument_readout.md
The stratification is by decile because the flag is itself correlated with the
price forecast, the 0.22 group rising from 1{,}067 rows in the lowest decile to
1{,}410 in the highest while the 0.18 group falls from 1{,}324 to 644. Nothing
here is a test. There is no clustering, no Holm correction and no placebo, the
table enters none of the fifteen pre-declared families of
Table~\ref{tab:app-holm-families}, and it adds no cell to any survivor count.
The one quantity it shares with Figures~\ref{fig:app-item-geography}, \ref{fig:app-item-geography-b} and~\ref{fig:app-item-geography-c} is the
panel: those figures cut the same event-driven residual by item code, this table
cuts it by the instrument's own output.

\begin{table}[htbp]
\centering
\caption[What the surviving instrument emits]{What the surviving prompted arm
emits, and what its modal split separates. C6 on the event-driven channel at
$h=5$, all 25{,}109 test rows, forecasts frozen. Rows are the deciles of the
recalibrated HAR-RV forecast; \emph{high} is the group whose text forecast is
the modal value 0.22 and \emph{low} the group at 0.18, the two values carrying
92.6 per cent of its forecasts. Realised volatility is in
volatility units. \textbf{This is a descriptive stratification of frozen
forecasts, not a test}: no clustering, no Holm and no placebo, no pre-declared
family, and no survivor claim. The differences are positive in all ten deciles
but are not strictly increasing.}
\label{tab:app-instrument-readout}
\singlespacing\footnotesize
\setlength{\tabcolsep}{6pt}
% src: results/tables/instrument_readout.md (all ten decile rows and both group counts); results/tables/instrument_readout.csv
\begin{tabular}{@{}rrrrrr@{}}
\toprule
HAR-RV decile & $n$ (high) & $n$ (low) & mean RV, high & mean RV, low & difference \\
\midrule
1 & 1{,}067 & 1{,}324 & 0.1919 & 0.1849 & $+0.0071$ \\
2 & 1{,}144 & 1{,}259 & 0.2092 & 0.1986 & $+0.0106$ \\
3 & 1{,}188 & 1{,}197 & 0.2240 & 0.2077 & $+0.0163$ \\
4 & 1{,}173 & 1{,}214 & 0.2416 & 0.2256 & $+0.0160$ \\
5 & 1{,}225 & 1{,}157 & 0.2634 & 0.2437 & $+0.0197$ \\
6 & 1{,}219 & 1{,}139 & 0.2680 & 0.2471 & $+0.0209$ \\
7 & 1{,}337 & 991 & 0.2959 & 0.2713 & $+0.0246$ \\
8 & 1{,}333 & 976 & 0.3205 & 0.2915 & $+0.0289$ \\
9 & 1{,}293 & 953 & 0.3694 & 0.3220 & $+0.0474$ \\
10 & 1{,}410 & 644 & 0.4534 & 0.3790 & $+0.0744$ \\
\bottomrule
\end{tabular}
\end{table}

\clearpage

\begin{figure}[tbp]
\centering
\includegraphics[width=\textwidth,height=0.83\textheight,keepaspectratio]{figures/F13_elicitation_not_curation__a.pdf}
% caption numbers src: results/tables/c5x_input_parity.csv (panel a); results/tables/llm_contamination.csv and .md (panel b); results/tables/crossfamily_llama70_probe.csv (panel c)
\caption[Elicitation, not curation, and not the name]{
Elicitation, not
curation, and not the name. Three input manipulations isolate what the
prompted arm uses. All three panels report volatility-unit relative QLIKE
improvement over the recalibrated HAR, and all three carry \emph{separate}
multiplicity families, so no count may be pooled across them. 
Panel~(a) is the
input-parity comparison, long-form only, that being the only channel with a
parity cell. The prompted 32B arm adds $+1.79/+2.25/+0.27$ per
cent and is genuine in 3 of 3 cells, a same-lineage 8B embedder on the
byte-identical excerpts loses at every horizon
($-0.32/-3.74/-4.12$, genuine in 0 of 3), and so does the standard-input
seed-ensemble embedding arm ($-1.02/-3.10/-6.39$, also 0 of 3). % src: results/tables/c5x_input_parity.csv (rel_impr_pct and genuine for C6_llmtext, C5x_qwen3exc and the C5_qwen3 ensemble)
Basis: volatility-unit QLIKE, day-clustered inference, each panel
Holm-corrected inside its own block; the denominator here is 12 cells. The
zero-content constructions of panels~(b) and~(c) are
Figures~\ref{fig:app-elicitation-b} and~\ref{fig:app-elicitation-c}.
}
\label{fig:app-elicitation}
\end{figure}

\begin{figure}[tbp]
\centering
\includegraphics[width=\textwidth,height=0.83\textheight,keepaspectratio]{figures/F13_elicitation_not_curation__b.pdf}
% caption numbers src: results/tables/c5x_input_parity.csv (panel a); results/tables/llm_contamination.csv and .md (panel b); results/tables/crossfamily_llama70_probe.csv (panel c)
\caption[Elicitation, not curation, and not the name (b)]{Continued from Figure~\ref{fig:app-elicitation}. Elicitation, not
curation, and not the name. Three input manipulations isolate what the
prompted arm uses. All three panels report volatility-unit relative QLIKE
improvement over the recalibrated HAR, and all three carry \emph{separate}
multiplicity families, so no count may be pooled across them. 
Panel~(b) removes the document text and runs the zero-content construction on
the anchor family: a date-only prompt and a date-plus-ticker prompt beside the
full-text arm, on both channels at all three horizons. Its denominator is 6
cells and its Holm block is its own, separate from those of panels~(a)
and~(c) (Figures~\ref{fig:app-elicitation} and~\ref{fig:app-elicitation-c}).
}
\label{fig:app-elicitation-b}
\end{figure}

\begin{figure}[tbp]
\centering
\includegraphics[width=\textwidth,height=0.83\textheight,keepaspectratio]{figures/F13_elicitation_not_curation__c.pdf}
% caption numbers src: results/tables/c5x_input_parity.csv (panel a); results/tables/llm_contamination.csv and .md (panel b); results/tables/crossfamily_llama70_probe.csv (panel c)
\caption[Elicitation, not curation, and not the name (c)]{Continued from Figure~\ref{fig:app-elicitation}. Elicitation, not
curation, and not the name. Three input manipulations isolate what the
prompted arm uses. All three panels report volatility-unit relative QLIKE
improvement over the recalibrated HAR, and all three carry \emph{separate}
multiplicity families, so no count may be pooled across them. 
Panel~(c) repeats that zero-content construction inside the larger
replication family, on the 8-K channel only, the percentage printed by each
probe bar being its share of the full-text bar. The panel carries no
pre-registered branch and is descriptive; its denominator is 3 cells, its Holm
block separate from those of panels~(a) and~(b)
(Figures~\ref{fig:app-elicitation} and~\ref{fig:app-elicitation-b}).
}
\label{fig:app-elicitation-c}
\end{figure}

Figures~\ref{fig:app-elicitation}, \ref{fig:app-elicitation-b} and~\ref{fig:app-elicitation-c} draw the control Chapter~\ref{ch:data}
designs and Section~\ref{sec:res-residual} reports: a same-lineage embedding
arm reading byte-identical excerpts. The chapter gives the three triples; what
is added here is the cell-by-cell view and the two companion panels. They answer the most natural deflation of the prompted
arm's result: that the excerpts handed to the prompted model were
themselves the informative object, so any model would have done as well on
them. Neither embedding arm did. Both lose at every horizon --- the parity arm
on exactly those excerpts, the seed-ensemble arm on its standard full-document
input --- while the prompted arm adds at all three, so the difference is
the interface and not the input. The comparison is at parity of \emph{input}
and not of parameter count (a 32B decoder against an 8B embedder with a
ridge head), so the licensed phrase is same-lineage, not same-size. In panel~(b) a date-only prompt is significantly positive in 0 of 6 cells, while a date-and-ticker prompt reproduces 31 to 77 per cent of the full-text increment in the three cells whose denominator is large enough to interpret.
% src: results/tables/llm_contamination.csv (date-only and date+ticker arms)

Two uncomfortable cells in that family of parts are drawn rather than cropped, and
both are stated here because a reader who noticed only the favourable panel
would be misreading the artefact. In panel~(b), at long-form $h=20$ the
full-text increment is only $+0.27$ per cent while the date-and-ticker prompt
alone returns $+1.11$ per cent, so the reproduction ratio prints as a nominal
405 per cent. % src: results/tables/llm_contamination.md (long_form h=20 datefirm row: rel +1.11 per cent, repro_frac 4.05, against a full-text increment of +0.27 per cent)
The ratio is a small-denominator artefact, drawn in the attention colour
and explicitly not quotable as a contamination estimate.
In panel~(c) the 70B's zero-content probe recovers 5.2 and 7.3 per cent of
that family's full-text increment at the two shorter horizons and a flagged
31.5 per cent at $h=20$, whose denominator the source records as not well
identified. % src: results/tables/crossfamily_llama70_probe.csv (share_har_pct 5.229 / 7.309 / 31.480; denominator_well_identified false at h=20)
The complement is printed in full. With the same-model zero-content forecast
inside the reference, the full-text arm still adds $+1.37$ per cent at $h=5$
(Holm .0046) and $+1.14$ per cent at $h=10$ (Holm .0465) and is not
distinguishable at $h=20$ ($+0.65$ per cent, Holm .4512). It retains 94 to 98
per cent of its uncontrolled increment in point estimate. % src: results/tables/crossfamily_llama70_probe.csv (beyond_identity block: rel_pct 1.374 / 1.136 / 0.653 at p_holm 0.0046 / 0.0465 / 0.4512, against fulltext_ens_rel_har 1.395 / 1.163 / 0.692)

The figure does not show what these arms do against the firm-identity
reference, which is a separate rung and a smaller residual. It does not
establish absence of contamination, since a zero-content prompt bounds what
firm identity alone can reproduce rather than proving it absent. It does not
show the era-split block of the same contamination file, which is a third
multiplicity family. And panel~(a) licenses no claim that prompting beats
pooling in general: only that on these excerpts, at these three horizons,
on long-form filings, the prompted arm adds where both embedding arms lose.

\FloatBarrier
\section{Within-Firm Text Representations}
\label{app:res-withinfirm}

Chapter~\ref{ch:background} raises the change-based reading of the
disclosure literature: that a periodic filing mostly repeats its
predecessor and the \emph{changes} carry what predictive content there is
\cite{cohen2020lazy}. The obvious objection to this study's null is
that it was run on levels. Three arms were executed to answer it.
Section~\ref{sec:res-standalone} gives the first two a sentence each; the
third is reported only here. The first replaces the TF--IDF
features with their within-sequence differences, excluding each firm-form
sequence's first filing. The second leaves the features alone and makes the
\emph{target} within-firm, training FinBERT on log realised volatility
minus the firm's training-window mean and restoring the predictions to
level units. Both are within-firm by construction, which is exactly the
property the firm-identity rung exists to test for.
The third builds the change features themselves, in the construction of
Cohen et al.~\cite{cohen2020lazy}: a ridge on the cosine and Jaccard
distance between each long-form filing and the same filer's own previous
filing of the same form, the log ratio of their lengths and a 10-K dummy,
which makes its change features within-firm in the same sense wherever a
predecessor was found, as one was for 94.9 per cent of unique long-form
filings. % src: results/tables/textchange_probe.md (prev-filing match 29,991/31,601 unique long-form docs; row-level imputed fraction 4.7 per cent)
Added to a validation-frozen combiner that already carries the
recalibrated HAR and the level TF--IDF forecast, the change features do
not add to the level features; they subtract from them. The joint
combiner's volatility-unit QLIKE is 2.26, 2.92 and 6.14 per cent
\emph{higher} than the level-only combiner's at $h{=}5$, 10 and 20, with
Diebold--Mariano statistics of $+5.23$, $+4.06$ and $+4.99$ against it,
positive meaning the joint model is the worse of the two, and raw
$p<.001$ at every horizon. % src: results/tables/textchange_probe.md (block 3, joint vs level-only combiner)
This arm is long-form only and a seed-2026 run off the seed-ensemble
primary basis; its Diebold--Mariano statistics are observation-level with
HAC lag $h-1$ rather than day-clustered and sit in no multiplicity family,
so it is read for direction and magnitude only.

Tables~\ref{tab:app-wf-har} and~\ref{tab:app-wf-firm} report the first two
arms against the two references in turn, and that pair is the answer. Against the
recalibrated HAR the total count is unchanged (6 of 9 cells genuine,
against 6 of 9 for the level counterparts on the matching seed-2026 basis),
but the composition inverts. Change features \emph{lose} where levels gain: the
difference arm returns $-1.46$ to $-0.16$ per cent on the long-form cells
where the level TF--IDF arm returns $+3.32$ to $+6.02$ per cent on the
identical rows. It is Holm-significantly worse at all three horizons. The
demeaned-target arm gains in all six of its cells. Against the
firm-identity reference, however, both tabulated arms return zero of nine,
exactly as
their level counterparts do. One cell clears Holm there (long-form
demeaned FinBERT at $h{=}20$, at $+2.52$ per cent) and fails the
label-shuffle placebo at $-2.27$, meaning permuted text improves that
reference too. The gain is therefore a property of the demeaned forecast's
marginal distribution rather than of the filing. Building the model
within-firm does not recover a filing-specific increment; it relocates
which cells look favourable against the weaker reference and changes
nothing at all against the stronger one.

\begin{table}[htbp]
\centering
\caption[Within-firm text arms against the recalibrated HAR]{The two
within-firm text arms against the recalibrated HAR reference. B2d replaces
the TF--IDF features with within-sequence differences, excluding each
firm-form sequence's first filing. That is why its row support is about
100 observations below the level arm's, exactly 103 at every horizon here and
98 to 100 on the identity rung of Table~\ref{tab:app-wf-firm}. C2dm trains FinBERT on a
firm-demeaned log target and restores level units;
\textbf{it is a seed-2026 reduced-form run, off the seed-ensemble primary
basis}. Volatility-unit QLIKE, day-clustered DM, Holm within this
nine-cell family, \emph{genuine} requiring DM $<0$, Holm $<.05$ and a
placebo inside the band. The last two columns re-run the corresponding
\emph{level} arm on the same reduced row set, so that the row exclusion
cannot explain the difference between the two.}
\label{tab:app-wf-har}
\singlespacing\footnotesize
\setlength{\tabcolsep}{2pt}
% src: results/tables/row2_delta_demeaned_m1.csv (har_* columns, genuine_har, sr_orig_har_* same-row level counterparts; har_n_test 7,848/7,830/7,799 against orig_har_n_test 7,951/7,933/7,902, a gap of 103 at each horizon, and firm_n_test 7,450/7,068/6,999 against orig_firm_n_test 7,550/7,167/7,097, a gap of 100/99/98)
\begin{tabular}{@{}llrrrrrrcrr@{}}
\toprule
& & & & \multicolumn{5}{c}{Within-firm arm} & \multicolumn{2}{c}{Level arm, same rows} \\
\cmidrule(lr){5-9}\cmidrule(lr){10-11}
Channel & Arm & $h$ & $n$ & rel.\% & DM & Holm $p$ & plac. & Gen. & rel.\% & DM \\
\midrule
Long-form & B2d $\Delta$TF--IDF & 5 & 7{,}848 & $-0.16$ & $+2.11$ & .0357 & $-1.04$ & -- & $+3.32$ & $-5.26$ \\
Long-form & B2d $\Delta$TF--IDF & 10 & 7{,}830 & $-0.22$ & $+2.52$ & .0357 & $+0.77$ & -- & $+3.55$ & $-9.13$ \\
Long-form & B2d $\Delta$TF--IDF & 20 & 7{,}799 & $-1.46$ & $+2.39$ & .0357 & $+0.05$ & -- & $+6.02$ & $-8.95$ \\
\addlinespace
Long-form & C2dm demeaned FinBERT & 5 & 7{,}951 & $+1.14$ & $-4.54$ & $<$.0001 & $-0.48$ & Y & $+0.56$ & $-0.84$ \\
Long-form & C2dm demeaned FinBERT & 10 & 7{,}933 & $+0.54$ & $-3.95$ & .0004 & $+0.86$ & Y & $+4.56$ & $-6.46$ \\
Long-form & C2dm demeaned FinBERT & 20 & 7{,}902 & $+9.62$ & $-5.70$ & $<$.0001 & $-1.75$ & Y & $-0.08$ & $+0.88$ \\
\addlinespace
Event-driven & C2dm demeaned FinBERT & 5 & 25{,}109 & $+1.12$ & $-3.39$ & .0029 & $+1.01$ & Y & $+2.14$ & $-4.95$ \\
Event-driven & C2dm demeaned FinBERT & 10 & 25{,}001 & $+1.46$ & $-4.56$ & $<$.0001 & $+0.57$ & Y & $+2.10$ & $-5.52$ \\
Event-driven & C2dm demeaned FinBERT & 20 & 24{,}732 & $+1.86$ & $-4.67$ & $<$.0001 & $+0.72$ & Y & $+0.92$ & $-0.50$ \\
\bottomrule
\end{tabular}
\end{table}

\begin{table}[htbp]
\centering
\caption[Within-firm text arms against the firm-identity reference]{The
same two arms against the \textbf{firm-identity-augmented} reference: a
reference that knows which firm is speaking and reads no text at all. A
model built to be within-firm by construction that still cannot beat this
reference carries no filing-specific increment. Columns, conventions and
bases exactly as in Table~\ref{tab:app-wf-har}; the row supports are the
smaller five-price-model intersection on which the identity rung is always
scored. Zero of nine cells are genuine. The one cell clearing Holm,
long-form C2dm at $h{=}20$, fails the label-shuffle placebo at $-2.27$
and is reported as an artefact rather than removed.}
\label{tab:app-wf-firm}
\singlespacing\footnotesize
\setlength{\tabcolsep}{2pt}
% src: results/tables/row2_delta_demeaned_m1.csv (firm_* columns, genuine_firm, sr_orig_firm_* same-row level counterparts)
\begin{tabular}{@{}llrrrrrrcrr@{}}
\toprule
& & & & \multicolumn{5}{c}{Within-firm arm} & \multicolumn{2}{c}{Level arm, same rows} \\
\cmidrule(lr){5-9}\cmidrule(lr){10-11}
Channel & Arm & $h$ & $n$ & rel.\% & DM & Holm $p$ & plac. & Gen. & rel.\% & DM \\
\midrule
Long-form & B2d $\Delta$TF--IDF & 5 & 7{,}450 & $+0.21$ & $+0.97$ & .4404 & $+0.04$ & -- & $-0.71$ & $+4.04$ \\
Long-form & B2d $\Delta$TF--IDF & 10 & 7{,}068 & $+0.11$ & $+1.23$ & .4404 & $+0.21$ & -- & $-3.99$ & $+8.05$ \\
Long-form & B2d $\Delta$TF--IDF & 20 & 6{,}999 & $-0.60$ & $+1.61$ & .3241 & $+0.86$ & -- & $-8.23$ & $+8.01$ \\
\addlinespace
Long-form & C2dm demeaned FinBERT & 5 & 7{,}550 & $-0.69$ & $+3.72$ & .0013 & $-0.22$ & -- & $-0.18$ & $+1.89$ \\
Long-form & C2dm demeaned FinBERT & 10 & 7{,}167 & $-2.21$ & $+2.36$ & .0738 & $+0.08$ & -- & $-0.60$ & $+4.09$ \\
Long-form & C2dm demeaned FinBERT & 20 & 7{,}097 & $+2.52$ & $-3.88$ & .0008 & $-2.27$ & -- & $-0.92$ & $+0.17$ \\
\addlinespace
Event-driven & C2dm demeaned FinBERT & 5 & 23{,}855 & $-0.49$ & $+3.65$ & .0014 & $+0.69$ & -- & $-0.28$ & $+4.88$ \\
Event-driven & C2dm demeaned FinBERT & 10 & 22{,}785 & $-1.32$ & $+5.25$ & $<$.0001 & $+0.05$ & -- & $-1.22$ & $+6.33$ \\
Event-driven & C2dm demeaned FinBERT & 20 & 22{,}318 & $-2.33$ & $+5.00$ & $<$.0001 & $+0.84$ & -- & $+0.58$ & $-0.81$ \\
\bottomrule
\end{tabular}
\end{table}

\clearpage

\begin{figure}[tbp]
\centering
\includegraphics[width=\textwidth,height=0.83\textheight,keepaspectratio]{figures/AF2_longform_section_geography__a.pdf}
% caption numbers src: results/tables/section_ablation.csv (all 18 rows: m1_rel_improve_pct, qlike_vs_B2_pct, frac_nonempty_all, frac_nonempty_10Q); results/tables/section_ablation.md (method disclosures, table-of-contents skip re-extraction, sanity re-run)
\caption[Where inside a long-form filing the apparent increment lives]{
Where
inside a long-form filing the apparent increment lives. Four
TF-IDF-and-ridge variants, each trained on a single named span of every annual
and quarterly report or on the complement of the three named spans, replicate
the archived bag-of-words recipe and are combined with the same recalibrated
HAR reference as the primary grid. 
Panel~(a): the three text-bearing spans add and none
recovers the whole filing, while the market-risk item subtracts at every
horizon ($-1.42/-2.11/-3.63$), its text column being empty in every
quarterly report. The management discussion is the strongest
single carrier at $+2.26/+2.04/+3.21$ per cent against full text's
$+3.33/+3.48/+5.92$, risk factors add $+1.37/+1.20/+1.81$, and the remainder
of the document still adds $+2.08/+1.51/+2.72$. % src: results/tables/section_ablation.csv (m1_rel_improve_pct by section and horizon)
}
\label{fig:app-section-geography}
\end{figure}

\begin{figure}[tbp]
\centering
\includegraphics[width=\textwidth,height=0.83\textheight,keepaspectratio]{figures/AF2_longform_section_geography__b.pdf}
% caption numbers src: results/tables/section_ablation.csv (all 18 rows: m1_rel_improve_pct, qlike_vs_B2_pct, frac_nonempty_all, frac_nonempty_10Q); results/tables/section_ablation.md (method disclosures, table-of-contents skip re-extraction, sanity re-run)
\caption[Where inside a long-form filing the apparent increment lives (b, c)]{Continued from Figure~\ref{fig:app-section-geography}. Where
inside a long-form filing the apparent increment lives. Four
TF-IDF-and-ridge variants, each trained on a single named span of every annual
and quarterly report or on the complement of the three named spans, replicate
the archived bag-of-words recipe and are combined with the same recalibrated
HAR reference as the primary grid. 
Panel~(b): the three disjoint spans' increments sum to $5.71/4.75/7.73$ per
cent against the full-text $3.33/3.48/5.92$, so they are substantially
redundant rather than additive. 
Panel~(c): each named span is nonetheless a
\emph{worse} standalone forecaster than the full text, variance-unit test
QLIKE 3.6 to 14.2 per cent higher, the sole exception being the remainder,
which is better. % src: results/tables/section_ablation.csv (qlike_vs_B2_pct: +3.58 to +14.16 for the three named spans, -1.91 to -4.29 for the remainder)
Basis: the long-form panel only, $7{,}951/7{,}933/7{,}902$ test rows at
$h=5/10/20$; single seed 2026; observation-level Diebold--Mariano rather than
the report's day-clustered primary, and no Holm family: an off-basis
diagnostic. The reference is the recalibrated HAR alone, the ladder's first
rung, so every bar is an \emph{apparent} increment.
}
\label{fig:app-section-geography-b}
\end{figure}

The finding that makes Figures~\ref{fig:app-section-geography} and~\ref{fig:app-section-geography-b} worth a page is
panel~(b), and it is a negative one. If the three disjoint spans carried
separable information their increments would sum to roughly the full-text
increment; instead they sum to substantially more, at every horizon. Whatever
the bag-of-words pipeline extracts is therefore largely the same information
wherever in the document it reads. That is consistent with the report's wider
finding that the surviving signal is cross-sectional (about which firm is
filing) rather than about the content of a particular section. Panel~(c)
records a second dissociation: each named span is a worse standalone
forecaster than the full document, and two of the three still add
incrementally while Item 7A subtracts at all three horizons; the
remainder is a better standalone forecaster and adds slightly less.
% src: results/tables/section_ablation.csv (m1_rel_improve_pct: item1a +1.372/+1.199/+1.808, item7 +2.262/+2.038/+3.205, item7a -1.415/-2.112/-3.631, rest +2.078/+1.509/+2.719)
Standalone
skill and incremental value are simply different quantities, which is the
methodological point Chapter~\ref{ch:results} makes at the level of the whole
model spectrum, reproduced here at the level of spans within one document.

The figure does not license a section-level attribution of the increment:
panel~(b) is precisely the demonstration that no such decomposition is
available. It does not test any span against the firm-identity or
maximal-pool rungs, so every bar is of exactly the kind those rungs go on to
dissolve. Because the whole experiment sits at the first rung, on a single
seed, with observation-level inference and no multiplicity correction, no bar
in it is a survivor and none should be quoted as one. One negative result is
an artefact and is labelled as such on the figure. The market-risk item's bars
run from $-1.42$ to $-3.63$ per cent, and its text column is non-empty for
24.6 per cent of the panel and for none of the quarterly reports. % src: results/tables/section_ablation.csv (B2sec_item7a: m1_rel_improve_pct -1.415 / -2.112 / -3.631; frac_nonempty_all 0.2461; frac_nonempty_10Q 0.0)
The
negative is a statement about an empty column rather than about that item's
content.
Against that, the re-implemented recipe reproduces the archived full-text run
to within 0.01 per cent of test QLIKE at every horizon, so the four variants
are comparable with the committed pipeline and not merely with each other.

\FloatBarrier
\section{The Earnings-Call Panel, Cell by Cell}
\label{app:res-maec}

Section~\ref{sec:val-external} reports the earnings-call audit as a single
count: nothing detectable in 24 cells. Table~\ref{tab:app-maec} is the
24 cells, because a count of zero is the claim in this report most easily
dismissed as underpowered, and the only defence against that dismissal is
to print the minimum detectable effect beside every one of them.

The construction repeats the SEC ladder on a published benchmark. A text
arm is combined with a recalibrated past-volatility reference; the
reference is then augmented with a same-ticker expanding mean, the
identity term; and the residual increment of text over that augmented
reference is the object adjudicated. Under the benchmark's own published
convention the text arms do reproduce a large apparent gain: the
TF--IDF arm cuts pooled loss by 13.6 to 53.1 per cent against the raw
past-volatility baseline the literature discards. Recalibrating that
baseline and adding the identity term absorbs it. Of the 24
identity-controlled residual cells, none is a Holm-significant improvement
and none is a Holm-significant deterioration. Every point estimate for the
two fitted arms is positive, between $+1.94$ and $+4.61$ per cent, and not
one is distinguishable from zero. The identity share (the fraction of
each arm's combination increment that the zero-text same-ticker mean
reproduces) runs from 89 to 133 per cent for the two fitted arms, so the identity term
accounts for the whole of the increment in the cells at or above 100 per cent
and for very nearly all of it in the six that fall below. % src: results/tables/maec_audit.md (share d4/d3 over the fitted-arm cells; six read 89, 90, 90, 91, 94 and 90 per cent)

The power statement is what makes that a finding. The panel is small: 672
test observations over 143 call dates and 461 entities per horizon, against
the SEC panel's tens of thousands of filings. The per-cell minimum
detectable effect runs from 2.97 to 12.80 per cent, so this table can rule
out only effects of that size, and it says so in every row. What licenses
the ``fully absorbed'' reading is that the published-convention gain being
repriced is \emph{larger} than the minimum detectable effect in every cell
where such a gain exists (17 to 120 per cent of the reference loss).
The instrument could therefore have seen the published gain survive and did not.
One qualification belongs with that statement, and it is not a small one. The
published scorer was run on two of the three arms; the embedding arm has no
published-convention gain of its own, and its cells are compared against the
TF-IDF arm's. The audit records that proxy as biased upward, so for that arm
the ``powered'' verdict is the less conservative reading rather than the more,
and its tightest cell --- thirty days against the autoregressive reference, a
10.07 per cent detectable effect against a 17.4 per cent proxy gain --- would
fall to underpowered if the embedding arm's own gain were below about 58 per
cent of the TF-IDF arm's. % src: results/tables/maec_audit.md sections 5 and 7-5 (published scorer run for tfidf and prompted_qwen only; qwen_emb G_conv uses the tfidf Delta_pub as a PROXY, biased upward; tightest cell h30 r_ar MDE(ent) 10.07 per cent against G_conv 17.4 per cent, flipping below ~58 per cent)
Where an arm that \emph{was} scored has no published-convention gain to absorb,
as the prompted arm does not at three of its four horizons, the share column
carries a dash rather than a claim of absorption. The dash is set on a related
but distinct condition --- the source writes it whenever the combination
increment itself is non-positive, so the division is undefined --- and the two
coincide on those cells rather than being the same test. % src: scripts/analysis/maec_audit_table.py (identity_share is NaN iff d3 <= 0); results/tables/maec_audit.md section 2

\begin{table}[p]
\centering
\caption[The earnings-call identity-controlled residual, all 24 cells]{The
earnings-call audit: the identity-controlled text residual, all 24 cells of
the primary alignment. The reference is a recalibrated past-volatility
forecast (autoregressive or HAR, both shown) augmented with the
same-ticker expanding pre-call volatility mean. The row tests what the text
arm adds beyond it. Loss is squared error on the volatility label,
inference is call-date-clustered DM with Holm inside each arm's eight-cell
family, and the intervals are date-block bootstrap 95\% intervals. ``G4''
is the label-shuffle placebo and ``G4b'' the within-date swap diagnostic,
whose dirty cells are excluded from prose claims but shown. ``share'' is
the fraction of that cell's combination increment reproduced by the
zero-text identity term, so a share at or above 100\% means the identity
term accounts for the whole increment. ``MDE'' is the per-cell minimum
detectable effect at 80\% power in the same units as the increment.
Support: 672 test observations over 143 call dates and 461 entities per
horizon. \textbf{No cell is a Holm-significant improvement and none is a
Holm-significant deterioration.} Nothing in this table shares a loss, a
panel or a denominator with any QLIKE percentage elsewhere in this report.}
\label{tab:app-maec}
\singlespacing\footnotesize
\setlength{\tabcolsep}{1.8pt}
% src: results/tables/maec_audit.csv (stage = row5, alignment = primary; three headline arms, both references, four horizons)
\begin{tabular}{@{}llrrrrrlllrr@{}}
\toprule
Text arm & Ref. & $h$ & rel.\% & DM & raw $p$ & Holm $p$ & 95\% CI & G4 & G4b & share & MDE \\
\midrule
TF--IDF & AR & 3 & $+3.24$ & $-0.75$ & .4536 & 1.0000 & $[-3.65,\,8.27]$ & pass & dirty & 114\% & 8.44\% \\
TF--IDF & AR & 7 & $+4.61$ & $-1.72$ & .0883 & .7066 & $[-0.08,\,8.27]$ & pass & dirty & 119\% & 5.85\% \\
TF--IDF & AR & 15 & $+4.13$ & $-1.22$ & .2227 & 1.0000 & $[-0.80,\,8.03]$ & pass & clean & 112\% & 6.77\% \\
TF--IDF & AR & 30 & $+2.90$ & $-1.05$ & .2975 & 1.0000 & $[-0.66,\,5.75]$ & pass & clean & 119\% & 5.00\% \\
\addlinespace
TF--IDF & HAR & 3 & $+2.59$ & $-1.22$ & .2263 & 1.0000 & $[-1.67,\,6.45]$ & pass & clean & 90\% & 6.06\% \\
TF--IDF & HAR & 7 & $+3.09$ & $-1.69$ & .0924 & .7066 & $[-0.20,\,5.79]$ & pass & clean & 91\% & 4.18\% \\
TF--IDF & HAR & 15 & $+2.77$ & $-1.14$ & .2544 & 1.0000 & $[-0.76,\,5.29]$ & pass & clean & 90\% & 4.66\% \\
TF--IDF & HAR & 30 & $+2.49$ & $-1.42$ & .1565 & .9392 & $[-0.16,\,4.53]$ & pass & clean & 89\% & 3.19\% \\
\addlinespace
Qwen3-Emb.\ 8B & AR & 3 & $+4.27$ & $-1.95$ & .0536 & .3750 & $[0.16,\,6.86]$ & pass & dirty & 112\% & 5.10\% \\
Qwen3-Emb.\ 8B & AR & 7 & $+2.45$ & $-0.24$ & .8107 & 1.0000 & $[-4.83,\,6.16]$ & pass & dirty & 133\% & 7.65\% \\
Qwen3-Emb.\ 8B & AR & 15 & $+3.75$ & $-0.30$ & .7658 & 1.0000 & $[-5.35,\,7.79]$ & pass & clean & 111\% & 9.97\% \\
Qwen3-Emb.\ 8B & AR & 30 & $+2.92$ & $-0.23$ & .8204 & 1.0000 & $[-4.58,\,6.01]$ & pass & clean & 117\% & 10.07\% \\
\addlinespace
Qwen3-Emb.\ 8B & HAR & 3 & $+2.24$ & $-2.23$ & .0275 & .2197 & $[0.34,\,4.13]$ & pass & clean & 113\% & 2.97\% \\
Qwen3-Emb.\ 8B & HAR & 7 & $+1.94$ & $-0.34$ & .7315 & 1.0000 & $[-3.21,\,4.61]$ & pass & clean & 123\% & 5.15\% \\
Qwen3-Emb.\ 8B & HAR & 15 & $+2.68$ & $-0.19$ & .8519 & 1.0000 & $[-4.10,\,5.03]$ & pass & clean & 94\% & 7.15\% \\
Qwen3-Emb.\ 8B & HAR & 30 & $+2.63$ & $-0.27$ & .7859 & 1.0000 & $[-3.32,\,4.89]$ & pass & clean & 90\% & 7.48\% \\
\addlinespace
Prompted 32B & AR & 3 & $-0.37$ & $+0.35$ & .7234 & 1.0000 & $[-8.67,\,5.02]$ & pass & clean & 7890\% & 10.31\% \\
Prompted 32B & AR & 7 & $-2.21$ & $+0.49$ & .6266 & 1.0000 & $[-11.86,\,6.91]$ & pass & clean & --- & 12.80\% \\
Prompted 32B & AR & 15 & $-3.73$ & $+0.79$ & .4330 & 1.0000 & $[-9.34,\,3.26]$ & pass & clean & --- & 9.61\% \\
Prompted 32B & AR & 30 & $-0.79$ & $-0.94$ & .3482 & 1.0000 & $[-2.38,\,8.30]$ & pass & clean & 2807\% & 5.90\% \\
\addlinespace
Prompted 32B & HAR & 3 & $-1.17$ & $+0.77$ & .4445 & 1.0000 & $[-5.30,\,1.44]$ & pass & clean & --- & 5.25\% \\
Prompted 32B & HAR & 7 & $-2.22$ & $+1.06$ & .2896 & 1.0000 & $[-7.45,\,2.07]$ & pass & clean & --- & 6.52\% \\
Prompted 32B & HAR & 15 & $-3.37$ & $+1.47$ & .1434 & 1.0000 & $[-9.37,\,0.78]$ & pass & clean & --- & 7.07\% \\
Prompted 32B & HAR & 30 & $+0.02$ & $-0.73$ & .4658 & 1.0000 & $[-2.31,\,5.81]$ & pass & clean & 3955\% & 3.81\% \\
\bottomrule
\end{tabular}
\end{table}

A second alignment of the call transcripts to their outcome windows was
pre-registered as a sensitivity, and it moves two cells: the prompted
arm at $h{=}15$ against both references becomes a Holm-significant
% src: results/tables/maec_audit.csv (stage = row5, alignment = shifted)
\emph{deterioration}, at $-5.32$ and $-3.87$ per cent. No cell becomes a
significant improvement under either alignment.
The audit therefore travels to this panel with the same verdict it reaches
at home and with its power stated: on an external benchmark whose published
convention credits text substantially, recalibrating the discarded baseline
and adding the entity's own history absorbs the gain, and the residual is
zero of 24 at a minimum detectable effect this table prints cell by cell.

\FloatBarrier
\section{The Registered Mechanism Analysis That Failed}
\label{app:mechanism}

One pre-registered analysis returned a falsification, and it is set out here
rather than only summarised in Chapter~\ref{ch:conclusions}, because a
registered analysis that fails is evidence in the same sense as one that does
not. The question was whether the size of the identity control's effect can be
predicted from how strongly a panel's own reference forecasts already encode
entity identity. The registered prediction was a \emph{negative} rank
correlation: the more identity a price reference already carries, the less an
explicit identity control should be able to remove.

The predictor $x$ is the $R^{2}$ of an ordinary least-squares fit of the
reference's log-space \emph{validation} predictions on entity dummies, computed
as the one-way between-group variance share within each panel, horizon and
reference. The outcome $y$ is the identity control's relative gain in each
panel's own loss, read from that panel's frozen artefact: relative QLIKE on the
six SEC cells, and relative squared error on the two Yelp and the eight MAEC
cells, those two panels being scored on squared error throughout this report.
Sixteen cells were available. No test row is read anywhere in the construction.
% src: results/tables/yelp_cascade.csv (mse and delta_rel_pct columns);
% results/second_domain/maec/protocol_tfidf_primary.json (horizons[h][ref] deltas computed from the mse block);
% results/tables/firm_identity_ensemble.csv (the six SEC cells, volatility-unit QLIKE)
% src: results/tables/mechanism_identity.md (the 16 points, x and y per cell);
% configs/prereg_mechanism_and_labels.md section C, tag prereg-cd-v1.0

The observed Spearman correlation is $\rho = +0.585$, the opposite sign to the
one registered. A one-sided permutation test in the registered direction, at
10{,}000 draws, returns $p = 0.991$ pooling the panels and $p = 0.026$
permuting within them; the conservative reading, and the one the analysis record
takes, is the pooled figure. A median split on $x$ gives a one-sided Fisher
exact $p = 1.000$ in the predicted direction. The falsification branch fired:
the mechanism claim does not hold, no wording anywhere in this report escalates
the identity control's behaviour to a mechanism, and the descriptive language of
Chapter~\ref{ch:results} stands unchanged.
% src: results/tables/mechanism_identity.md (pre-registered statistics block:
% rho +0.5853, permutation p global 0.9908 and within-panel 0.0259, Fisher exact
% one-sided 1.0000; verdict block "fired branch: FALSIFIED")

Three properties of the design limit what the failure itself licenses. The
outcome is not one quantity: ten of the sixteen cells are relative squared error
and six are relative QLIKE, so the rank correlation is taken across two loss
conventions and its magnitude is not interpretable as a single elasticity. The
sixteen cells come from three panels and are not independent draws: cells within
a panel share references, entities and overlapping outcome windows, which is why
the within-panel permutation and the pooled one disagree so widely. And the
predictor is mechanically panel-size dependent, because entities with few
validation rows fit their own dummies almost perfectly, so $x$ is high wherever
rows per entity are few. The registered statistic is rank-based, which blunts
but does not remove that dependence. The analysis therefore falsifies a
prediction; it does not establish the reverse.
% src: results/tables/mechanism_identity.md (Disclosures: non-independence,
% 3 panels with 6/2/8 cells; mechanical caveat on x, n_val and n_entities per cell)

\FloatBarrier
\section{The Founding Design, Replicated and Dissolved}
\label{app:res:kogan}

Section~\ref{sec:val-kogan} runs the founding design in both directions and
quotes its counts in prose; this section carries them cell by cell, because
the convention this report keeps is that a headline count names the table
behind it. Two tables follow. The first is the ladder on this panel: the
ported recipe at full strength in sample, then the same recipe repaired one
element at a time until the apparent gain reverses. The second is the
off-panel cascade on the original corpus years, where the same repairs are
applied to a reading that first has to clear a reproduction gate against the
published figures.

\begin{table}[p]
\centering
\caption[The founding design replicated, then repaired rung by rung]{The
founding design of \cite{kogan2009predicting} ported to this panel and
repaired one element at a time. The ported object is that study's
\emph{evaluation recipe} --- TF--IDF text features beside a single lagged
log-volatility control, log-volatility mean squared error, year-by-year
out-of-sample scoring and naive observation-level inference --- not its
corpus, period or horizon; the single text arm is the archived B2 TF--IDF
ridge recipe throughout, so each rung's movement is attributable to one
repair. The rungs are L0, the ported recipe scored per year over 2020--25;
L1, the declared chronological test split only; L2, the recalibrated-HAR
log-combiner in place of the single weak control, scored in volatility-unit
QLIKE; L3, day-clustered rather than observation-level inference; L4, a
firm-identity-augmented reference; L5, the maximal five-model price pool on
top of that firm control. Panel~(a) is the \emph{in-sample} arm and is a
deliberate, labelled look-ahead: it fits and evaluates on the same rows and
is the defect being replicated, never a forecasting claim. L0's year-by-year
arm carries a second labelled defect --- it scores 2020--21, years this
benchmark reserves for validation --- and L1 is the rung that repairs it.
\emph{Sign
conventions differ by column and are opposite to one another}: for the naive
$t$ a \emph{positive} value means text is better, while for both DM columns a
\emph{negative} value means text is better; the gain column is positive when
text is better throughout. That gain column changes unit at L2 by design:
L0 and L1 are log-volatility MSE reductions of the text-augmented arm over
the single control, L2 onwards volatility-unit QLIKE reductions of the
text-augmented combination over its own reference, and a percentage in one
convention may not be set beside a percentage in the other. Holm is applied inside exactly one pre-declared
family, the three L5 day-clustered $p$-values; rungs L0--L4 report raw $p$ by
design, each reporting what its own progressively less broken protocol would
have reported. Panel~(c) carries three quantities the main text does not use
--- the clustering-day count, the share of test rows whose firm has a
validation-window mean, and the QLIKE levels behind each percentage --- and
records that L5's inner join over five price models shrinks the row support.
The $h{=}20$ cell that reverses at L5 is a single-arm ladder cell on that
shrunken support, not a conjunction survivor: the same channel--model cell
stands at $-8.09$ per cent on the committed firm-identity rung. % src: results/tables/firm_identity_control.csv (long_form, B2_tfidf_ridge, h=20, rel_impr_pct_firm -8.089490)
}
\label{tab:app-kogan-ladder}
\singlespacing\footnotesize
\setlength{\tabcolsep}{4pt}
% src: results/tables/kogan_dissolve.md and results/tables/kogan_dissolve.csv (in-sample arm; THE DISSOLVE LADDER; L4/L5 detail)

\emph{(a) The in-sample arm: the published-style positive at full strength (labelled look-ahead)}\\[0.4ex]
\begin{tabular}{@{}rrrrrr@{}}
\toprule
$h$ & $n$ & $R^2$ (control only) & $R^2$ ($+$ text) & MSE reduction \% & naive obs $t$ \\
\midrule
 5 & 31{,}466 & 0.250 & 0.381 & $+17.42$ & $+42.9$ \\
10 & 31{,}426 & 0.312 & 0.455 & $+20.87$ & $+46.8$ \\
20 & 31{,}345 & 0.336 & 0.495 & $+23.92$ & $+47.4$ \\
\bottomrule
\end{tabular}

\vspace{1.2ex}
\emph{(b) The dissolve ladder, rung by rung}\\[0.4ex]
\begin{tabular}{@{}llrrrlrll@{}}
\toprule
Rung & Repair added & $h$ & $n$ & gain \% & Statistic & Value & $p$ (raw) & Holm $p$ \\
\midrule
L0 & ported recipe, per-year OOS &  5 & 11{,}907 & $-0.68$ & naive obs $t$ & $-1.14$ & $.2523$ & --- \\
L0 & & 10 & 11{,}883 & $-4.85$ & naive obs $t$ & $-7.90$ & $3.14\times10^{-15}$ & --- \\
L0 & & 20 & 11{,}845 & $-7.65$ & naive obs $t$ & $-12.76$ & $5.11\times10^{-37}$ & --- \\
\addlinespace
L1 & $+$ declared test split only &  5 & 7{,}951 & $+0.16$ & naive obs $t$ & $+0.22$ & $.8281$ & --- \\
L1 & & 10 & 7{,}933 & $-3.16$ & naive obs $t$ & $-3.78$ & $1.58\times10^{-4}$ & --- \\
L1 & & 20 & 7{,}902 & $-6.91$ & naive obs $t$ & $-7.15$ & $9.73\times10^{-13}$ & --- \\
\addlinespace
L2 & $+$ recalibrated-HAR ref. &  5 & 7{,}951 & $+3.33$ & obs-level DM & $-11.34$ & $1.43\times10^{-29}$ & --- \\
L2 & & 10 & 7{,}933 & $+3.48$ & obs-level DM & $-15.30$ & $4.27\times10^{-52}$ & --- \\
L2 & & 20 & 7{,}902 & $+5.92$ & obs-level DM & $-19.62$ & $1.07\times10^{-83}$ & --- \\
\addlinespace
L3 & $+$ day-clustered inference &  5 & 7{,}951 & $+3.33$ & day-clust.\ DM & $-5.39$ & $9.21\times10^{-8}$ & --- \\
L3 & & 10 & 7{,}933 & $+3.48$ & day-clust.\ DM & $-8.89$ & $3.90\times10^{-18}$ & --- \\
L3 & & 20 & 7{,}902 & $+5.92$ & day-clust.\ DM & $-9.04$ & $1.17\times10^{-18}$ & --- \\
\addlinespace
L4 & $+$ firm-identity reference &  5 & 7{,}951 & $-1.04$ & day-clust.\ DM & $+5.04$ & $5.77\times10^{-7}$ & --- \\
L4 & & 10 & 7{,}933 & $-4.12$ & day-clust.\ DM & $+7.78$ & $2.30\times10^{-14}$ & --- \\
L4 & & 20 & 7{,}902 & $-6.71$ & day-clust.\ DM & $+7.29$ & $7.45\times10^{-13}$ & --- \\
\addlinespace
L5 & $+$ maximal price pool &  5 & 7{,}550 & $-0.33$ & day-clust.\ DM & $+3.59$ & $3.56\times10^{-4}$ & $7.12\times10^{-4}$ \\
L5 & & 10 & 7{,}167 & $-1.45$ & day-clust.\ DM & $+5.02$ & $6.49\times10^{-7}$ & $1.95\times10^{-6}$ \\
L5 & & 20 & 7{,}097 & $+0.20$ & day-clust.\ DM & $-2.85$ & $.0045$ & $.0045$ \\
\bottomrule
\end{tabular}

\vspace{1.2ex}
\emph{(c) L4 and L5 in detail: support, coverage and the QLIKE levels behind the percentages}\\[0.4ex]
\begin{tabular}{@{}lrrrrrrr@{}}
\toprule
Rung & $h$ & $n$ & Days & Firm-val.\ cov. & QLIKE ($f_R$) & QLIKE ($f_U$) & gain \% \\
\midrule
L4 &  5 & 7{,}951 & 809 & 0.92 & $0.1220$ & $0.1232$ & $-1.04$ \\
L4 & 10 & 7{,}933 & 803 & 0.92 & $0.0859$ & $0.0895$ & $-4.12$ \\
L4 & 20 & 7{,}902 & 794 & 0.92 & $0.0722$ & $0.0770$ & $-6.71$ \\
\addlinespace
L5 &  5 & 7{,}550 & 792 & 0.91 & $0.1242$ & $0.1246$ & $-0.33$ \\
L5 & 10 & 7{,}167 & 766 & 0.91 & $0.0900$ & $0.0914$ & $-1.45$ \\
L5 & 20 & 7{,}097 & 765 & 0.91 & $0.0859$ & $0.0857$ & $+0.20$ \\
\bottomrule
\end{tabular}
\end{table}

\begin{table}[p]
\centering
\caption[The off-panel cascade on the founding corpus, year by year]{The
same cascade run off panel, on the founding study's own tokenised 10-K
management-discussion corpus, 1996--2006, at its 12-month horizon. The
corpus is integrity-checked before use: 30{,}474 filings from 7{,}207
registrants, with every year's document count matching the published census
exactly. % src: results/tables/kogan_corpus_audit.md (G-K0 row/key-space consistency; G-K3 CIK coverage)
Panel~(a) is the pre-registered reproduction gate, computed under the
published convention (five-year rolling training window, test years
2001--06, count-weighted micro-average) because that is the only arm
commensurable with the published table; the gate operationalises ``same
order of magnitude'' as a factor of ten, and at $7.91\times$ the reproduction
sits inside that bound but not inside a stricter same-power-of-ten reading,
under which the falsification branch would have fired. Panel~(b) is the
binding arm, whose rungs are L0, the ported convention on the pre-registered
expanding split; L1, a recalibrated control; L2, a same-registrant identity
term added to the reference, read leave-one-out as the declared primary of
two pre-registered readings; L3, filing-date-clustered inference; L4, Holm
within the family of per-year $p$-values; L5, the conjunction $\mathrm{L}1
\wedge \mathrm{L}2 \wedge \mathrm{L}4$. Each test year contributes between
202 and 227 distinct filing dates and a Newey--West lag of 4, the more
permissive of the two pre-declared lag rules and so the one most favourable
to text. \emph{A negative DM means text is better}; the gain columns are
positive when text is better, and are identical from L2 onwards because
those rungs share one text arm and one reference. The placebo permutes the
text rows over five seeds and refits, and the gate is $|\mathrm{DM}| < 2$;
because the ported convention makes the text arm a single joint model rather
than a combiner over the reference, shuffled text retains a structural edge,
so the gate is informative but not null-calibrated. The gain columns are
reductions in log-volatility mean squared error, share no denominator with
any QLIKE percentage elsewhere in this appendix, and must not be set beside
one.}
\label{tab:app-kogan-corpus}
\singlespacing\footnotesize
\setlength{\tabcolsep}{3pt}
% src: results/tables/kogan_corpus_audit.md and results/tables/kogan_corpus_audit.csv (G-K1 gate and per-year comparison; THE LADDER -- L0_prereg, binding; Placebo -- label shuffle)

\emph{(a) The reproduction gate: this reading against the published one, under the published convention}\\[0.4ex]
\begin{tabular}{@{}lrrrrl@{}}
\toprule
Test year & MSE (pub.\ ref.) & MSE (pub.\ $+$ text) & published gain \% & this reading \% & sign agrees \\
\midrule
2001 & $0.1747$ & $0.1919$ & $-9.85$ & $-3.01$ & yes \\
2002 & $0.1600$ & $0.1618$ & $-1.12$ & $+3.92$ & no \\
2003 & $0.1873$ & $0.1965$ & $-4.91$ & $+13.72$ & no \\
2004 & $0.1442$ & $0.1246$ & $+13.59$ & $+20.22$ & yes \\
2005 & $0.1365$ & $0.1276$ & $+6.52$ & $+10.21$ & yes \\
2006 & $0.1463$ & $0.1403$ & $+4.10$ & $+8.74$ & yes \\
\midrule
Micro-average & $0.1576$ & $0.1557$ & $\mathbf{+1.21}$ & $\mathbf{+9.54}$ & 4 of 6 \\
\bottomrule
\end{tabular}

\vspace{0.4ex}
\begin{tabular}{@{}p{0.95\textwidth}@{}}
\emph{Gate: sign agrees and the ratio is $7.91\times$, inside the pre-registered factor-of-ten bound --- \textbf{PASS}. This reading's micro-averaged reference MSE reproduces the published $0.1576$ exactly; the text arm is the more favourable of the two, not the less.} \\
\end{tabular}

\vspace{1.2ex}
\emph{(b) The binding cascade, per test year, with the label-shuffle placebo}\\[0.4ex]
\begin{tabular}{@{}lrrrrrllrrll@{}}
\toprule
Test & & \multicolumn{3}{c}{gain \%} & \multicolumn{3}{c}{L3 / L4 inference} & \multicolumn{2}{c}{placebo DM} & & \\
\cmidrule(lr){3-5}\cmidrule(lr){6-8}\cmidrule(lr){9-10}
year & $n$ & L0 & L1 & L2$+$ & DM & $p$ (raw) & Holm $p$ & mean & max $|\cdot|$ & Gate & L5 \\
\midrule
1997 & 2{,}260 & $+3.10$ & $+1.14$ & $+11.40$ & $-2.50$ & $.0131$ & $.0523$ & $-2.08$ & $2.87$ & fail & --- \\
1998 & 2{,}462 & $+7.73$ & $+4.65$ & $+4.40$ & $-4.05$ & $7.17\times10^{-5}$ & $5.74\times10^{-4}$ & $+1.12$ & $1.64$ & \textbf{pass} & \textbf{Y} \\
1999 & 2{,}524 & $+6.07$ & $+0.88$ & $+0.22$ & $+0.99$ & $.3241$ & $.6482$ & $+1.87$ & $2.70$ & fail & --- \\
2000 & 2{,}425 & $+9.69$ & $+6.41$ & $+5.60$ & $-2.80$ & $.0056$ & $.0337$ & $+0.44$ & $0.98$ & \textbf{pass} & \textbf{Y} \\
2001 & 2{,}596 & $-3.01$ & $+6.38$ & $-0.61$ & $-0.09$ & $.9248$ & $.9248$ & $+5.19$ & $6.16$ & fail & --- \\
2002 & 2{,}846 & $+4.72$ & $+5.63$ & $+6.44$ & $-3.00$ & $.0030$ & $.0208$ & $-1.39$ & $3.05$ & fail & (Y) \\
2003 & 3{,}612 & $+14.70$ & $+19.25$ & $+19.12$ & $-8.10$ & $3.38\times10^{-14}$ & $3.38\times10^{-13}$ & $-0.37$ & $1.89$ & \textbf{pass} & \textbf{Y} \\
2004 & 3{,}559 & $+20.74$ & $+16.89$ & $+23.82$ & $-6.91$ & $5.34\times10^{-11}$ & $4.80\times10^{-10}$ & $-2.82$ & $3.91$ & fail & (Y) \\
2005 & 3{,}474 & $+11.69$ & $+7.89$ & $+12.96$ & $-2.72$ & $.0071$ & $.0355$ & $-2.38$ & $2.99$ & fail & (Y) \\
2006 & 3{,}308 & $+8.68$ & $+6.28$ & $+10.04$ & $-2.17$ & $.0315$ & $.0944$ & $-0.59$ & $1.48$ & \textbf{pass} & --- \\
\midrule
\multicolumn{12}{@{}p{0.97\textwidth}@{}}{\emph{``L5'' marks survival of the conjunction $\mathrm{L}1 \wedge \mathrm{L}2 \wedge \mathrm{L}4$: \textbf{6 of 10} test years before the placebo gate, of which the three in parentheses fail it. The placebo-gated survivor count, the one the branch is decided on, is \textbf{3 of 10} --- 1998, 2000 and 2003. The gate passes in only 4 of 10 splits, and \textbf{13 of the 50} individual placebo draws show the shuffled-text arm beating the reference at $\mathrm{DM} < -2$.}} \\
\bottomrule
\end{tabular}
\end{table}
